% 本文件由 LaTeX 排版 Agent 根据有效 translation.json 自动生成。
% author=Arnobius; work_id=0631; title=驳异教徒

\chapter{驳异教徒}

% structure: p0001 source=h1 target=omitted reason=page-title-already-rendered-by-parent

\Emph{第一卷}\newline{}\Emph{第二卷}\newline{}\Emph{第三卷}\newline{}\Emph{第四卷}\newline{}\Emph{第五卷}\newline{}\Emph{第六卷}\newline{}\Emph{第七卷}\par

\SourceRecord{Translated by Hamilton Bryce and Hugh Campbell.；Ante-Nicene Fathers；Vol. 6.；Edited by Alexander Roberts, James Donaldson, and A. Cleveland Coxe.；Buffalo, NY: Christian Literature Publishing Co.,；1886.；<http://www.newadvent.org/fathers/0631.htm>.}

% structure: p0001 source=h1 target=section reason=page-title

\section{驳外教人（卷一）}

1. 我既见有些人自以为智慧极高，行事仿佛受了灵感，并以神谕的全权宣告：自从基督徒存在于世以来，宇宙便趋于毁灭，人类遭受了种种灾祸，甚至神明们也舍弃了惯常的执掌——往日他们正是藉此执掌关心人事的——已被逐出地上的领域。我遂决意，只要我才力与微薄的语言能力允许，便要对抗公众的偏见，并驳斥那些诽谤性的指控；以免一方面这些人以为他们是在宣告什么重大之事，实际却只是在传播流俗的谣言；另一方面，如果我们回避这样的争辩，他们便会认为他们是赢得了一场因其自身固有缺陷而失败、而非因辩护者的缄默而放弃的诉讼。因为我不会否认，那指控是极其严重的，而且若显得我们是使宇宙偏离其法则、神明远遁、世代人类遭受如此众多苦难的原因，那么我们便完全应当承受公敌的恨恶。\par

2. 因此，让我们仔细检视那种意见的真实意义，以及那指控的性质；并摒除一切争辩的欲念——这种欲念常会模糊\Emph{甚至}阻断对事物的冷静观察——让我们公正地权衡双方的考虑，以验证那被指控的是否真实。因为必会有一系列令人信服的论证证明，并非我们被发现有更多的不虔，而是那些自称为神明敬拜者、古旧迷信的沉迷者，他们自己被定了那罪。首先，我们以友善平静的语言向他们提出这一问题：自从基督教之名在地上开始使用时，所谓的\Quote{\InnerQuote{万物的本性}}可曾感受到或遭受过何种先前未见、先前未闻的现象？可曾有什么事件违悖了太初所立的定律？原初的元素——众所周知，万物皆由其所组成——可曾被转变为性质相反的元素？这宇宙的机体与整体——我们皆被其覆盖，并被其包容其中——可有任何部分松弛了，或是崩解了？我们习以为常的球体自转，可有脱离其原初运动的速率，开始运行得太慢，或是急遽地加速旋转？星辰可曾开始从西方升起，而星座的沉落发生于东方？太阳本身，那众天体的首领，以其光明覆被万物，以其热量化生万物，可曾以加增的猛烈喷射光芒？他可曾变得没那么温暖，并且他那惯常作用于大地的调和的温度，可曾转劣成相反的状况？月亮可曾停止重新塑造自身，并停止倚靠新相位的不断更替而变回前态？冬日的严寒、夏日的酷热、春秋的和暖，可曾因不相配的时节的混杂而被改变？冬天可曾开始拥有长昼？夜晚可曾开始唤回夏日极迟的暮光？诸风可曾完全耗尽它们的暴烈？天空可曾因诸风力道失尽而不再聚云？田地被甘霖润湿后可曾不繁茂？大地可曾拒绝接受托付给它的种子，或是树木可曾不肯披上绿叶？佳果的滋味可曾改变？葡萄的汁液可曾变化？橄榄果中可曾压榨出污浊的血液，而熄灭的灯盏可曾不再得到油的供应？陆地与海中的动物可曾没有了性欲，不再受孕怀胎？它们可曾不按各自的习性本能，保卫腹中所生的后代？最后，人类本身——一股活跃的力量已凭其原初的冲动把他们散布于可居住的大地——可曾不按合宜的礼规结成婚姻？他们可曾不生养亲爱的子女？他们可曾不照料公共的、个人的和家庭的事务？他们可曾不如各人所愿，将其才能运用于各种职业、不同类型的学问？他们可曾不从勤勉的投入中收获果实？那些被分配如此的人，可曾不执掌王权或军权？人们岂不是天天在荣誉的职位上、在权力的官位上得到擢升？他们岂不是在法庭的辩论中主持？他们岂不解释法典？他们岂不阐明公义的原则？人类生命所围绕、所赖以为基的一切其他事物，难道不是所有人在各自的族群中，按照其国家风俗的既定秩序，都在实行着吗？\par

3. 既然如此，既然没有任何奇异的影响突然出现，以打断事件之持续进程并中断其连续，那么，所谓一场瘟疫在\OriginalTerm{基督宗教}{Christian religion}传入世界之后降临大地，并且在它揭示了隐藏真理的奥秘之后，这一指控有何根据？然而，我的对手说，\Quote{瘟疫、干旱、战争、饥荒、蝗虫、老鼠、冰雹，以及其他侵害人类财产的有害事物，是神明们因你们的恶行和过犯而震怒，降在我们身上的。}如果说在这些已经清楚、无需辩护的事上纠缠不休还不算愚蠢，那么我确可以通过展开过去时代的历史来表明，你们所说的那些灾祸并非不为人知，它们的降临也非突如其来；而且，瘟疫并非在我们这一教派赢得此称号之荣耀时才突然爆发，人间事务并非从那时起才开始受到各种危险的攻击。因为如果我们当受责备，如果这些瘟疫是因为我们的罪而设计出来的，那么古人何以得知这些灾祸之名？他们何以给战争命名？凭借什么观念，他们能指出瘟疫和冰雹，或者如何能在他们的言语中引入这些词汇，使得言语清楚明白？因为如果这些灾祸是全新的，如果它们源于近期的过犯，古人又怎能创造出这些事物的词汇——他们既知道自己从未经历过，又未曾听闻在祖先时代发生过？我的对手说，\Quote{农产品短缺、谷物供应不足，压得我们更重。}因为，\Emph{我请问}，前几代人，甚至最古老的人，在哪个时期完全免于这种不可避免的灾祸？这些灾祸所借以指称的词语本身，难道不正是证据，并大声宣告，任何世人都未曾全然幸免于它们？但如果此事难以置信，我们可以根据作家们的见证，力陈有多少伟大的民族，哪些个别的民族，以及这些民族多少次经历了可怕的饥荒，并因接连不断的荒芜而灭亡。许多冰雹降下，袭击万物。我们不是发现古代文献中确凿记载，甚至石雨常常毁坏整个地区吗？暴雨使庄稼枯萎，给各地带来饥荒：——古代人难道真的幸免了这些灾祸吗？而我们已知，甚至大河曾干涸，河床淤泥龟裂。瘟疫的传染性毁灭人类：——遍查各种语言的历史记录，你将会发现，所有地区都屡屡被荒废，丧失了居民。各种庄稼被蝗虫和老鼠吞噬、吃尽：——查阅你们自己的年代记，你们将从这些灾害中得知，先前的时代如何频繁地遭受它们，如何频繁地被逼至贫困潦倒。城市被强烈地震撼动，摇摇欲坠：——什么！昔日难道没有见证过城市连同其居民被大地的巨大裂口吞没吗？或者它们曾享有免遭此类灾难的状态？\par

4. 人类何时被洪水毁灭？岂不是在我们之前吗？世界何时被火烧，化为焦炭与灰烬？岂不是在我们之前吗？最大的城市何时被海浪吞没？岂不是在我们之前吗？与野兽的战争何时发生，与狮子的战斗何时进行？岂不是在我们之前吗？何时整个社群被毒蛇带来毁灭？岂不是在我们之前吗？因为，你们惯于将频繁的战争、城市的荒废、日耳曼人与斯基泰人的入侵归咎于我们，请允许我说——你们急于诽谤我们，却未觉悟所指责之事的真正性质。\par

5. 难道是我们造成，一万年前，如\Person{柏拉图}所告，大批人众从名为\Place{亚特兰蒂斯}的\Person{尼普顿}之岛涌出，彻底毁灭并抹除了无数部落吗？难道这成了我们的不利证据：古代\Person{尼努斯}与\Person{琐罗亚斯德}领导下的亚述人与巴克特里亚人之间，不仅凭借刀剑与武力，而且借助术士及迦勒底人的神秘学问维系了一场斗争吗？难道\Person{海伦}在神明们的指引和教唆下被掳走，并成了她自己的时代及后世之可怕命运，这也要归咎于我们的宗教吗？难道是因为我们之名，那位狂妄的\Person{薛西斯}才让海洋涌入陆地，并且步行越过大海吗？难道是我们制造并引发了那些原因，使得一个青年从\Place{马其顿}出发，将东方诸王国与民族虏为奴役与捆绑吗？难道，饶是我们，我们催促神明们发狂，以致罗马人近来像一股暴涨的洪流，推翻所有民族，将他们卷入波涛之下吗？但如果没有人敢将那些久远以前发生的事归因于我们的时代，当一切并非新事，而是旧事，古人无一不知，我们又如何能成为当前不幸之原因呢？\par

虽然你声称你所提及的那些战争是因憎恨我们的宗教而起，但要证明，自基督之名闻于世界之后，它们非但没有增加，反而因狂怒之情的遏制而在很大程度上减少了，这并不困难。因为我们这人数众多的团体，从祂的教导与律法中得知，不应以恶报恶\BibleRef{玛 5:39}，宁可受屈，也不加害于人，宁可流自己的血，也不以他人之血玷污我们的双手与良心，一个不知感恩的世界，如今长时期享受着由基督而来的益处，因着祂，野蛮的凶暴之怒已得以缓和，并开始收回那加诸同类之血的敌意之手。然而倘若所有无一例外地，自觉其为人并非在于肉体之形状，而在于理性之能力的人，愿稍微倾听祂那有益且和平的规诫，且不因启蒙之骄矜自负，信赖自己的感官胜于祂的劝诫，那么，整个世界，将钢铁之用途转向更和平的劳作，现今必生活于至为宁静的平安中，并在融洽的和谐中团结一致，持守条约之神圣不可侵犯。\par

然而，我的对手们会说，如果你们没有对人类事务造成任何损害，那么那些如今压迫并压倒可怜凡人的灾祸又从何而来？你要求我给出一个明确的陈述，而这对本案绝无必要。因为我并未着手就此进行即时且准备的讨论，以表明或证明每一事件因何原因、缘何理由而发生；而是为要证明，如此严重的指控所含的谴责，距我们之门相去甚远。而若我证明了这一点，若藉着事例与有力的论证使事实真相昭然若揭，我便不在乎这些灾祸从何而来，或出自何种源头与初始之因。\par

然而，为免我似乎对此类问题全无见解，或于被问及时显得无话可说，我可以说：倘若那弥漫于宇宙四元素中的原始质料，其自身结构中便蕴含一切痛苦之因，又将如何？倘若天体的运行在特定的星象、地区、季节与区域中产生这些灾祸，并将种种危险之必然加诸其下的事物，又将如何？倘若宇宙间在定期的间隔内发生变迁，犹如海潮，繁荣一时涌来，一时退去，灾祸与之交替，又将如何？倘若我们践踏于足下的质料之污浊，被赋予此等条件，以致散发最为有害的瘴气，藉此败坏了我们的空气，将疫病加诸我们身体，削弱人类，又将如何？倘若——而这似乎最接近真相——凡我们所视为逆厄者，实际上对世界本身并非祸患，又将如何？倘若我们以自己的利益衡量一切发生之事，因错误的判断而责怪自然的结果，又将如何？哲学家们崇高的领袖和柱石\Person{柏拉图}，在其著作中宣称，那些残酷的洪水与世界的大火，乃是大地之净化；那位智者也不忌惮称人类之倾覆、毁灭、败亡与死亡为万有之更新，并断言，一种青春活力，仿佛藉此更新之力而得以保全。\par

我的对手说，天不降雨，我们正因谷物收成的异常匮乏而受苦。那么，莫非你要求元素成为你需求之奴仆？为了你能更安逸、更精致地生活，难道顺遂之季节就该为你的便利效力？倘若如此，一位志在航海的人抱怨说，长久以来全无风起，天上的风已永远平息，又将如何？难道因此便说宇宙的宁静是有害的，只因它妨碍了商人的愿望？倘若一个习惯于晒太阳以求得身体干燥的人，同样抱怨云彩夺去了晴朗天气的乐趣，又将如何？那么，难道云彩就该被称为以有害的帷幕笼罩，只因无聊的欲望不被容许在灼热中烤炙，并设计饮酒之借口？一切发生之事，在这宇宙巨体之下成就与发生之事，均不应被视为为着我们微末之利而遣发，而应被视为与自然自身的计划与安排相符。\par

倘若有何事发生，未能以喜乐之成功滋养我们自身或我们的事务，便不可立即视之为恶，或为有害之事。世界降雨或不降雨：为其自身而降雨或不降雨；而或许你对此茫然不知，它若非以燃烧般的干旱减少过度的湿气，便是藉着倾盆之雨缓解持续过久之干燥。它兴起瘟疫、疾病、饥荒，及其他种种灾害之祸端：你又怎能断定，它是否因此除去过剩之物，是否藉着事物自身之损失，为那易于滋蔓者设置界限？\par

你岂敢断言，在这宇宙中，此物或彼物为恶，而其根源与缘由你竟未能解释与分析？因其干涉了你合法、甚或不法之享乐，你便说它有害且逆性？那么，因为寒冷使你肢体不适，并易于冷冻你血液之温暖，冬季便因此不当存于世间？因为你无法忍受太阳最炽热的光芒，夏季便应从一年中删除，并以不同的法则设立自然之不同进程？\OriginalTerm{嚏根草}{Hellebore}于人乃毒药；它便因此不应生长？狼潜伏在羊圈旁；自然便全然有罪，因她产生了一种对羊群极危险的野兽？蛇以其噬咬夺走生命；确然，创造应受谴责，因为它为动物增添了如此残忍的怪物？\par

12. 当你甚至不是自己的主人，甚至当你是另一个人的财产时，对那些更强大的人发号施令，是相当放肆的；希望发生你所愿的事，而非那些你发现已由事物的\OriginalTerm{原初构造}{original constitution}所固定的事。因此，如果你希望你的抱怨有根据，你必须首先告知我们你从何而来，或你是谁；这个世界是为着你而受造并形成的，还是你作为寄居者从其他地方来到其中？既然你不能说出或解释你在这穹苍之下生活的目的，就当停止相信任何东西属于你；因为所发生的事并非为局部利益而发生，而是顾全整体的利益。\par

13. 我的对手们说，因着基督徒的缘故，诸神降下所有灾祸给我们，天上的神明使我们的庄稼遭到毁坏。请问，当你们说这些话时，难道没有看见你们是在用厚颜无耻、用明显且已被清楚证明的谎言指责我们吗？自从我们基督徒开始存在并被世人注意，差不多有三百年——或多或少——在这些年里，战争是连接不断的吗？庄稼是年复一年歉收的吗？地上没有和平吗？没有物价低廉、万物丰盛的季节吗？因为指责我们的人首先必须证明，这些灾祸是没完没了、连续不断的，人完全没有喘息的时间，没有任何松弛地经历了许多形式的危险。\par

14. 然而，我们岂不见，在这些过去的年代和季节里，从被征服的敌人那里取得了无数的胜利——帝国的疆界得以扩展，我们以前不曾听闻其名的民族被置于我们的权下——往往有了最丰收的谷物、物价低廉的季节，以及如此丰富的商品，以致所有的商业都因价格标准而瘫痪？因为，若非自然的生产力继续供应日用所需的一切，事务如何能进行，人类又如何能存留到如今？\par

15. 然而，有时也有短缺的季节，但它们被丰足的时期缓解了。再者，有一些战争是违背我们的意愿进行的，但它们后来由胜利和成功补偿了。那么，我们该说什么呢？——诸神一时记得我们的恶行，另一时又忘掉了？如果每逢饥荒，诸神就被说成是对我们发怒，那么在丰足时，他们就是没有发怒，也难以平息的；这样，事情就成了：他们以嬉戏的奇想收起又发出怒气，总是因想起冒犯的因由而重新燃起愤怒。\par

16. 然而，人无法通过任何合理的推理过程发现这些说法的含义。如果诸神愿意\OriginalTerm{阿勒曼尼人}{Alemanni}和\OriginalTerm{波斯人}{Persians}被战胜是因为基督徒住在他们的部落中，那么当基督徒也住在\OriginalTerm{罗马人}{Romans}当中时，他们如何将胜利赐予\OriginalTerm{罗马人}{Romans}呢？如果他们愿意在\Place{亚细亚}和\Place{叙利亚}有极多的老鼠和蝗虫成群涌出，也是因为基督徒住在他们当中，为什么在\Place{西班牙}和\Place{高卢}同时却没有这种现象，尽管无数的基督徒也住在那些行省呢？如果在\OriginalTerm{格图里人}{Gaetuli}和\OriginalTerm{廷吉塔尼人}{Tinguitani}中，诸神因此使庄稼干旱枯萎，为什么在同一年他们给\OriginalTerm{摩尔人}{Moors}和\OriginalTerm{诺马德人}{Nomads}极其丰富的收成，而相似的宗教也住在这些地区呢？如果在任何一个国家，诸神因厌恶我们的名字而使许多人饿死，为什么在同一国家，他们使那些不属于我们团体的人，甚至基督徒自己，因谷物高价而变得更富裕，甚至极其富有呢？因此，要么如果我们是一切不幸的原因，就不该有任何祝福，因为我们在所有民族中；要么当你们看见祝福与不幸混合时，就不要再把有损你们利益的事归于我们，因为我们丝毫不干涉你们的祝福和繁荣。因为，如果我使你们遭殃，为什么我不阻止你们得福呢？如果我的名字是极重饥荒的因由，为什么我无力阻止极丰盛的生产呢？如果人们说我带来了战争中负伤的\Emph{厄运}，那么当敌人被击杀时，我为何不是凶兆；为何我不因这坏兆头的\Emph{厄运}而被设置来对抗好希望呢？\par

17. 然而，啊，你们这些诸神的伟大崇拜者和祭司们，既然你们断言那些至圣的诸神对基督徒团体发怒，为什么你们同样没有觉察到，没有看到你们将何等卑鄙的情感、何等不雅的疯狂归给了你们的神明？因为，发怒，除了是疯狂、是咆哮、是被逼到复仇的欲望中，并且由于野蛮性情的疯狂而在他人忧伤的苦恼中狂欢之外，还能是什么呢？那么，你们的伟大诸神便知道，遭受并感受到那些只合于野兽、只合于怪异禽兽，以及致命植物\OriginalTerm{纳特里克斯毒草}{natrix}在它有毒的根中所包含的东西了。那远胜于他者、且建立在不可动摇之德行的坚固根基上的本性，正如你们所声称的，竟经历了属于人的不稳定性，属于地上动物的过失。那么必然得出的结论岂不是：他们的眼睛射出火花，火焰迸发，气喘的胸膛从口中吐出急促的呼吸，并且由于他们焦灼的言辞，焦干的嘴唇变得苍白？\par

18. 但是，如果你们所说的是真的——如果这已经过验证并充分证实，神明们确实怒火中烧，而且这类冲动使神明们激动不安，那么一方面他们并非不死者，另一方面他们也绝不能被视作分有神性。因为正如哲学家们所认为的，凡有冲动之处，必有情欲存在。有情欲之处，自然会有精神上的激动。有精神激动之处，便有忧伤及痛苦存在。有忧伤及痛苦之处，便已有衰弱与朽坏的余地；而若这两者折磨他们，灭绝便近在咫尺，即死亡，它终结万物，并从一切有感觉之物那里夺走生命。\par

19. 再者，你们以这种方式将他们描绘为不仅反复无常且易于激动，而且——这是所有人都同意与神性的特征相去甚远的——处事不公，如同作恶者，并且最终，确实连一丝最起码的公正都不具备。因为还有什么比迁怒于某些人、伤害其他人、抱怨人类、毁坏无害的五谷庄稼、憎恨基督徒之名、并以各种损失来毁灭基督的敬拜者，更为不义的呢？\par

20. 他们是否为此也向你们发泄怒气，好让你们因自身的创伤而被激起，起来为他们复仇？看来，神明们竟寻求凡人的帮助；倘若没有你们热切的护卫，他们自身便不能抵御并报复加诸他们的侮辱。不但如此，倘若他们确然怒火中烧，那就给他们自卫的机会，让他们施展并尝试其固有的力量，以报复他们受冒犯的尊严。通过酷热、通过有害的严寒、通过毒风、通过最隐秘的疾病，他们能杀死我们，能消灭我们，并能将我们完全逐出一切人际交往；或者，若以暴力攻击我们是不智的，就让他们发出某些愤怒的迹象，以此向众人显明，我们生活在天之下，仍处于他们强烈的愤恨之中。\par

21. 让他们赐予你们健康，却给我们疾病，甚至是最严重的疾病。让他们以应时的甘霖浇灌你们的田庄；却从我们的小块田地赶走一切柔和的降雨。让他们使你们的羊群因众多后裔而繁殖增多；却给我们的羊群带来不幸的不育。从你们的橄榄树和葡萄园，让他们带来满载的收成；却让我们这里，一滴汁液也不能从任何嫩枝压榨出来。最后，也是他们最恶毒的一招，让他们下令，使土地所产在你们口中保持其天然品质；相反，在我们口中，蜂蜜变苦，流淌的油腐败变质，啜饮的美酒刚一沾唇，便骤然化作令人失望的醋。\par

22. 既然事实本身证明这种情况从未发生，并且显然我们所得的世间福泽并不比你们少，你们所得的也并不比我们多，那么，何苦如此固执地断言神明对基督徒怀有敌意甚至仇恨呢？这些基督徒无论在极大的逆境中，还是顺境中，都与你们毫无二致。若你们允许真相直言相告，毫无保留，那么，这些指控不过是言辞——我再说一遍，不过是言辞；不，更是凭着诬蔑之词深信不疑之事，未经任何确凿证据证实。\par

23. 可是，真正的神明，以及那些配得并配戴此名尊严者，既不动怒，也不怀怨，更不会设阴谋诡计以伤害另一方。因为，实在说来，相信那智慧且至福的本性，若有人谦卑朝拜、俯伏于其前，便心中得意；若此朝拜不被献上，便自觉受轻蔑，并视自己从荣耀的顶峰跌落——这是亵渎的，且超越一切亵圣之举。这是幼稚、软弱而小气的，对于那些有学识者的经验久已称为半神和英雄者而言，几乎是不相称的，他们竟然不通晓天界之事，反而离弃自身本有的境地，忙于地上粗鄙之事。\par

24. 这些便是你们的想法，这些便是你们的意见，它们被不虔地构想出来，又被更加不虔地相信。更坦率地说，是那些占卜官、解梦者、占卜师、先知和小祭司们，他们素来虚妄，编造了这些故事；因为他们唯恐自己的伎俩归于乌有，唯恐自己只能从如今稀少而零落的信徒那里榨取微薄的捐献，所以每当发觉你们愿意让他们的行当蒙受恶名，便大声疾呼：\Quote{神明被轻忽了，神殿里如今人迹罕至。往日的仪式沦为笑柄，一度神圣的制度中历史悠久的礼仪，也已沉没于新宗教的迷信之下。人类遭受如此众多紧迫的灾祸，实属应当；被如此众多艰苦的磨难所折磨，实属应当。}而人们——这无知的族类——由于天生的盲目，甚至连摆在明处之事都无法看见，竟在狂乱中斗胆断言，你们在清醒状态下却深信不疑而毫不脸红之事。\par

25. 为避免有人以为我们因对自身答复缺乏信心，而将诸神塑造成性情平和，赋予他们无怨恨、远离一切激动的心智，既然你们喜欢，就让我们假定，他们向我们发泄怒气，渴望我们的鲜血，且长久以来急切地要将我们从人类世代中剪除。然而，若你们不嫌麻烦，若不觉冒犯，若不偏不倚、本着真理探讨这论证的各点乃是我们共同的义务，我们要求听你们解释，这究竟为何，原因何在；为何一方面诸神只对我们施以残酷，而另一方面众人又对我们愤恨相向。反对我们的人说，你们信奉亵渎的宗教体系，奉行举世未闻的仪式。你们这些被赋予理性的人啊，胆敢断言什么？胆敢妄谈什么？在口不择言的鲁莽中，你们又试图提出什么？敬拜天主为至高存在，为万有之主，为一切崇高者中居最高位者，在困苦中恭敬顺服地祈求祂，用我们全部感官——可以这么说——依恋祂，爱祂，以信德仰望祂——这竟是一种可憎且不虔的宗教，充满不敬与亵圣，以其新奇迷信玷污了自古设立的礼仪吗？\par

26. 我请问，这就是那胆大妄为、罪大恶极的行径，竟使天上的大能者对我们磨砺愤怒的刺芒，也使你们自己每逢兽性发作时，便抢夺我们的财物，将我们逐出祖辈的家园，处以极刑，施以酷刑，残毁我们，焚烧我们，最终将我们投给野兽，让怪物撕碎吗？凡是谴责我们这种信仰，或认为这应作为指控加诸我们的人，他配称作人吗，哪怕他自以为如此？或者，他该被认作神，尽管他借千百先知之口自称为神？是 \Person{特洛福尼乌斯}，还是 \Place{多多纳}的\Person{朱庇特}，宣告我们是邪恶的？而那向侍奉\OriginalTerm{无上君王}{King Supreme}的人指控不虔，或因无上君王的尊威与敬拜超越他自己而受嫉妒折磨的人，他自身还能被称为神，被列入神祇之列吗？\par

\Person{阿波罗}，不论被称为得利安的、克拉里安的、狄底玛的、菲莱西亚的，还是皮提亚的，他既不知\OriginalTerm{无上统治者}{Supreme Ruler}，也不晓得我们每日在祈祷中向祂祈求，难道还能被视为神明吗？即便他不晓得我们心中的秘密，未能发觉我们最深处的意念，他也总该凭耳闻，或凭我们祈祷时所用声调，感知到我们呼求\OriginalTerm{无上天主}{God Supreme}，并从祂那里祈求我们所需要的一切吧。\par

27. 此处不宜细究所有诽谤我们的人，他们是谁，从何而来，有何能力，有何知识，为何听到基督之名便颤抖，为何视他的门徒为仇敌、为可憎之人；而是要针对那些愿运用普遍理性的人，以适用于我们所有人的措辞，明确陈述——\Emph{就我们自身而言}，我们基督徒不过是无上君王与元首的敬拜者，在我们的主基督之下。你若仔细考察，便会发现那宗教别无他义。这是我们一切作为的总纲，这是我们神圣职责指定的目标与界限。我们众人都按习俗俯伏在祂面前；我们以共同的祈祷朝拜祂；我们向祂祈求公正、尊荣且值得祂垂听的事。并非祂需要我们的祈求，或喜爱目睹万千民众俯伏在祂脚下的崇敬。这乃是我们的益处，关乎我们的利益。因为我们易于犯错，由于天生的软弱，容易屈从于各种私欲与嗜好，所以祂容许自己常被我们思念，好让我们在恳求祂、努力赚取祂的恩赐之际，能获得对纯洁的渴望，并通过除去我们所有缺失，洗净一切污点。\par

28. 何言，神圣律法的诠释者啊，你们认为那些敬拜\OriginalTerm{拉瑞斯·格伦杜勒斯}{Lares Grundules}、\OriginalTerm{艾伊·洛库提}{Aii Locutii}和\OriginalTerm{利门提尼}{Limentini}的人，比我们这些钦崇天地万物的天主圣父的人，更趋向美善的缘由吗？你们看来，那些敬畏\OriginalTerm{法乌尼}{Fauni}、\OriginalTerm{法图厄}{Fatuae}与\OriginalTerm{国家的守护神}{genii of states}，敬拜\OriginalTerm{帕乌西}{Pausi}和\OriginalTerm{贝洛娜}{Bellonae}的人，谨慎、明智、洞彻事理，无可指摘；而我们全心事奉天主，万物皆因祂的旨意和喜悦而存在，并藉祂永恒的命定稳固不移，却被断为迟钝、愚拙、疯癫、懵懂、失智。是你们提出此论？是你们订立此法？是你们颁布此令：凡朝拜你们的奴仆者，当受最高荣冠；而向你们为主者祈祷者，当受十字架极刑？在至大的邦国和最强大的民族中，竟以公共名义为那些在古时靠不洁赚取报酬、向所有人的情欲卖身的妓女举行神圣仪式；\Emph{然而即便如此}，神明们却毫无愤慨之情。人们为猫、甲虫和小母牛建造了高耸的庙宇：神明们虽受此等侮辱，其力量却默不作声；他们看到这些卑贱动物的标志竟敢与他们分庭抗礼，亦不生出丝毫嫉妒之心。难道神明们单与我们为敌？只因我们钦崇他们的创造者——倘若他们存在，他们正是由祂而始有，并由此而获得其权能与尊威之实体，可以说，他们从祂那里领受了神性，遂悟己之存在，悟己被列于万有之中，凭着祂的旨意和命令，他们既能灭亡消散，亦能不消散不灭亡——便对我们如此绝情？倘若我们皆承认唯有一至高伟者，在漫长岁月中无有任何先于祂者，则必然结论为：万物由祂始有而产生，且各从其类，迸入存在。若此理确凿无疑，你们势必被迫承认，一方面神明乃是受造；另一方面，他们实赖万有之伟大本源而获其存在之泉源。若他们受造且被生出，无疑亦会湮没并遭遇危险；然而他们却被信为不朽、永存、不受灭尽。他们虽本性易逝易灭，却有特权历无数世代而保持不变，这亦来自他们的创造者天主的恩赐。\par

29. 但愿我得以将这番论证，在仿佛全宇宙汇成一场大会时，陈于全人类之耳！因此，我们在你们面前被指控信奉不虔敬的宗教吗？我们因以钦崇的侍奉接近宇宙的元首与支柱，就要被你们视为——用你们责骂我们的字眼——当被远避的人，无神者吗？谁比那认识、寻求或信奉任何异于我等所信天主之神明的人，更配蒙受这些名号之耻辱呢？我们岂不首先由祂而获此恩：我们存在，我们被称为人，我们无论是由祂遣发，或自祂堕落，我们被幽禁于此躯壳之黑暗中？我们行走、呼吸与生活，岂非从祂而来？祂岂非以生命之大能，使我们存在，并以生灵之活动运行？岂非从祂那里流出众多缘由，使我们藉各种逸乐之丰盈供给以维持康健？你们所生活的世界属谁？谁授权你们持守其产物与拥有？谁赐予那普照之光，使我们能明辨其下万物，触摸之，审察之？谁命定太阳之火存在，以滋长万物，免得含蕴生机之元素因沉沦于倦怠僵滞而麻木？当你们信太阳为神时，你们岂不追问其创造者、塑造者为谁？你们既以月亮为女神，难道同样关心其创始与制作者为谁吗？\par

30. 你们岂不曾想到反思并省察：你们生活在谁的领域？你们占有谁的财产？你们所耕种之大地属谁？你们所吸入并呼出之空气属谁？你们丰享之泉源属谁？水属谁？谁调理风之吹拂？谁筹谋含雨之云？谁以殊异特性甄别种子之生殖力？是\OriginalTerm{阿波罗}{Apollo}给你们雨水吗？是\OriginalTerm{墨丘利}{Mercury}从天上赐你们水吗？还是\OriginalTerm{埃斯库拉庇乌斯}{Aesculapius}、\OriginalTerm{赫丘利}{Hercules}或\OriginalTerm{狄安娜}{Diana}策划了阵雨与风暴的方案？当你们声称他们生于地上，并在某一时期获得生命知觉时，这又如何可能？因为若世界在漫长岁月中先于他们，若在他们出生前自然早已经历雨水与风暴，那么后来出生的他们便无权司雨，亦不能插手那些他们发现早已在此运行、并源自更伟大创造者的方法。\par

31. 噢，至大、至高的不可见之物的创造者！噢，祢自己本是不可见、不可领悟的！祢堪当，祢实在堪当——只要可朽的唇舌能谈及祢——一切有气息、有理智的受造，都应永不停止感受并感谢祢；毕生屈膝，献上永不停息的祈祷恳求。因为祢是元始之因；受造之物在祢内存有，祢是一切万物根基安息之所在。祢无边无界，无生无始，不死不灭，永存无尽，唯祢是独一的\OriginalTerm{天主}{God}，无任何形体可代表，无任何轮廓可描绘；美德无可言喻，伟大无可界定；不受地域、运动、状态之限，关于祢，人类言语的意义无法清晰地表达任何事物。要理解祢，我们必须保持静默；那迷误的揣测要在朦胧云雾中追寻祢，便不可发出一词。至高的君王，请宽恕那些迫害祢仆役的人；因祢慈祥的本性，赦免那些逃避敬拜祢圣名、守奉祢宗教的人。祢不为人知，不足为奇；祢竟能被清晰领悟，才更令人惊异。\par

然而，或许有人敢于—— 因为这正是狂肆之举—— 迟疑不定，并表明怀疑：那天主是否存在；祂是凭可靠证据的证实而被相信，还是凭虚妄传闻的想象。因为那些献身哲学的人，我们听说过，有的否认任何神性力量的存在，有的日日寻问其是否存在；还有的以偶然事故和随机碰撞构建整个宇宙框架，以不同形状的原子聚合将之塑造；此时我们绝无意与这种人展开讨论其乖戾的信念。因为明智者说，辩驳明显愚妄之事，是更大愚妄的标志。\par

32. 我们的讨论针对的，是那些承认有神明族类存在的人，他们怀疑那些更高等级和大能的神明，却承认存在更低微、更卑下的神祇。那么呢？我们难道要努力辛劳，通过论证来取得这样的结果吗？愿这种疯狂远离我们；如俗语所言，愿这愚妄，我说，从我辈身上被驱除。因为试图以论证证明天主是至高者，其危险等同于试图以这种理性去探索祂是否存在。你否认祂的存在，或确认并承认祂，这都无关紧要；因为不论是如此主张，还是那不信之对手的否认，同样是有罪的。\par

33. 岂有任何人在其生命之初日，未曾怀有那位伟大元首的意念吗？在谁身上，这意念未曾被天性植入？在谁身上，这意念未曾被铭刻，耶，甚至几近在他母胎之中就被烙印？在谁心中，没有一种本然的直觉，认出祂是君王、是主，是万有之主宰？总之，倘若哑口的动物也能结结巴巴地吐露其思想，倘若它们能使用我们的语言；非但如此，倘若树木，倘若土块，倘若石头藉由活力的知觉能生出有声之音，道出清晰之言，那么在此情形下，它们岂不是要以天性为向导与教师，怀着未腐之纯真，既感受到有一位天主，又宣告唯有祂是万有之主吗？\par

34. 但有人说，你们以无根无据、诽谤性的指控来攻击我们，仿佛我们否认有更高种类的神明，这都是徒劳的，因为\Person{朱庇特}被我们称为并尊为至善至大者；并且我们已经为他奉献了最为神圣的住所，竖立了宏伟的\Place{卡匹托尔}神庙。你正试图将不相类同的事物硬拉在一起，归为一类，\Emph{从而}造成混乱。因为据所有人的一致判断，并据全人类共同的认可，全能的天主被视为从未出生，从未被带至新的光明之中，未曾于任何时间或世代开始存在。因为祂自己是万有的源头，是世代与季节之父。因为它们并非自存，而是从祂永恒的恒久中，在无间断、无止境的流动中运行。然而，你们所声称的\Person{朱庇特}，确然有父有母，有祖父祖母，还有兄弟：新近在母亲胎中受孕，在十个月中完全成形并发育完善，他带着生机之感，迸入他先前所未知的光明中。那么，若果真如此，\Person{朱庇特}怎能是\Emph{至高的}神呢？既然明显可见，那一位是永恒的，而前者据你们所描述，是有了诞辰，并因对陌异景象的恐惧而发出了哀号？\par

35. 但假设他们是同一的，如你所愿，在神格与尊威的任何能力上并无不同，那么你就因此以不应受的仇恨迫害我们吗？如果我们同样敬拜你所敬拜的神明，为什么你一提及我们的名字，就犹如闻听最不祥之兆那般惊惧战栗？或者说，你为什么声称众神对你们友好，却对我们怀有敌意，而且是最大的敌意，尽管我们与他们之间的关系是相同的？因为若我们与你们有共同的宗教，神明的忿怒便会平息；但如果他们唯独对我们有敌意，那么显然，你们和他们都不认识天主。而那天主并非\Person{朱庇特}，此事由那些神明的怒气本身便可得证。\par

36. 但，我的对手说，那些神祇对你们并非怀有敌意，因为你们崇拜全能的天主；而是因为你们宣称，有一位像人一样出生、又被钉死在十字架上——即便对无价值的人来说也是一种耻辱的刑罚——竟就是天主，并且你们相信祂仍然活着，每日在祈求中崇拜祂。如果你们乐意，朋友们，请明确说出，那些相信我们对基督的崇拜有损于他们的神祇，究竟是哪些？是\Person{雅努斯}，\Place{雅尼库鲁姆}的创建者，和\Person{萨图恩}，萨图恩王国的创立者吗？是\Person{法乌努斯}的妻子\Person{法乌娜·法图阿}，被称为好女神，但她在饮酒方面却更善良、更值得称赞吗？是那些\Emph{\OriginalTerm{本土神}{Indigetes}}，在河中游泳、生活在\Place{努弥基乌斯}河道中，与青蛙和小鱼为伴吗？是\Person{埃斯库拉庇乌斯}和父\Person{巴克斯}吗？前者由\Person{科洛尼斯}所生，后者被闪电从母腹中击出。是\Person{迈亚}之子\Person{墨丘利}，以及更为神圣的、美丽的\Person{迈亚}吗？是持弓的神祇\Person{狄安娜}和\Person{阿波罗}吗？他们陪伴母亲漂泊，在浮岛上几乎无法安身。是\Person{狄俄涅}之女\Person{维纳斯}吗？她是特洛伊家族一个男子的情人，出卖她隐秘的魅力。是生于\Place{西西里}地区的\Person{刻瑞斯}，和采花时被掠走的\Person{普罗塞庇娜}吗？是\Place{底比斯}的\Person{赫拉克勒斯}还是腓尼基的\Person{赫拉克勒斯}？——后者葬于\Place{西班牙}地区，前者在\Place{奥塔山}上被火焚化。是\Person{廷达瑞俄斯}之子、兄弟\Person{卡斯托耳}和\Person{波鲁克斯}吗？——一个惯于驯马，另一个是优秀的拳击手，带着未鞣制的拳套所向无敌。是\Person{提坦}和\Place{摩尔人}的\Person{波科雷斯}，还有那些卵生的\Place{叙利亚}神祇吗？是生于\Place{伯罗奔尼撒}，在\Place{埃及}被称为\Person{塞拉庇斯}的\Person{阿庇斯}吗？是被\Place{埃塞俄比亚}的日头晒黑的\Person{伊西斯}吗？她哀悼失去的儿子和四肢被撕裂的丈夫。我们略过不提\Person{俄普斯}的王族后裔，你们的作家在他们的书中已向你们阐明，告诉你们他们是谁，品性如何。那么，这些神祇难道会因听到基督受崇拜、被我们接纳并视为一位神性者，就感到被冒犯了吗？他们忘记了自己不久前的地位和状态，难道不愿意与他人分享已赐给他们自己的东西吗？这就是天上神祇的公正吗？这就是神祇们的正义判断吗？这难道不是一种恶意和贪婪吗？这难道不是一种卑劣的嫉妒吗——只愿自己的命运上升，他人的命运被贬低、被践踏于轻蔑的卑微之中？\par

37. 我们崇拜一位生而为人者。那又怎样？难道你们不崇拜任何生而为人者吗？你们不是崇拜一个又一个，是的，无数神祇吗？不，你们岂不是从凡人中挑选了所有如今在你们庙宇中供奉的那些；把他们安置在天上，列于星座之间吗？因为，如果你们或许忽略了，他们也曾经历人类的命运，与众人有相同的状况，那就去查考最古老的文献，遍览那些生活最接近远古时代的人们的著作，他们以不加掩饰的真理、不作谄媚地记述一切：你们就会详细了解到，他们各自出自怎样的父、怎样的母，生于哪个地区、哪个部族；他们造就了什么，做了什么，忍受了什么，如何度过一生，在履行职能时经历了何种逆顺的际遇。然而，如果你们明知他们生于母腹，以地上出产为生，却仍然指责我们崇拜一位像我们一样出生的人，那你们就极其不公：你们以自己惯常所为的事，在我们身上却视为该受谴责；或者，你们允许自己合法的事，却不愿同样合法地允许他人。\par

38. 然而，且依从你们的想法，暂时承认\Person{基督}是我们当中的一员——在心智、灵魂、肉身、软弱和处境上与我们相似；难道祂不配因祂众多无量的恩惠，被我们称为并尊为天主吗？因为，你们既然将\Person{利伯尔}列入群神之列，是由于他发现了酒的用途；将\Person{刻瑞斯}列入，是由于她发现了面包的用途；将\Person{埃斯库拉庇乌斯}列入，是由于他发现了药草的用途；将\Person{密涅瓦}列入，是由于她出产了橄榄；将\Person{特里普托勒摩斯}列入，是由于他发明了犁；将\Person{赫丘利}列入，是由于他制服并捆绑了野兽、强盗和多头的水蛇——那么，祂以多大的殊荣应受我们尊敬呢？是祂将真理灌输到我们心中，使我们脱离重大的谬误；当我们在各处漂泊，如同盲人且无人引导时，祂将我们从险峻偏邪的路径上拉回，使我们的脚踏上更为平坦的地方；祂指出了对于人类最有益处、最能带来救恩的事；祂向我们展示天主是谁，祂有怎样的属性，何等伟大而又何等美善；祂容许并教导我们，按我们有限的能力去构想并理解祂那深不可测、无法言述的奥秘；祂以莫大仁慈，使我们得以知晓，这世界是由哪一位创造者、哪一位造物主所建立和造成的；祂解释了世界的本原及其基本实体，这是先前从未有人类思想可以想像的；为何太阳的光芒获得了赋予生命的热力；为何月亮在其运行中始终不受损害，人们相信是由于理智的原因而使光明与幽暗交替；动物的起源是什么，怎样的规则支配着种子；是谁设计并塑造了人本身，或者说，祂是用何种材料构建了我们的身体结构；知觉是什么；灵魂是什么，是它自己飞临到我们身上，还是与我们的肉体一同被产生并进入存在；它是寄居在我们内，分担死亡，还是被赐予了无终的永生；当我们脱离了松弛于死亡的肉体后，等待着我们的将是怎样的状况；我们是会保留知觉，还是将不再记得先前的感受与记忆；祂遏制了我们的傲慢，使我们那因骄傲而扬起的颈项承认了自身的软弱；祂指出我们乃是受造得不完美的存在，我们信赖空洞的期望，我们无法透彻理解任何事物，我们一无所知，我们对摆在眼前的东西视而不见；祂引领我们从虚假的迷信走向真实的宗教——这份恩赐超越并凌驾于祂所有的其他恩赐之上；祂将我们的思想，从以最卑贱泥土所塑成的粗野偶像那里，提升至天上，并使我们得以在感恩与祈祷中，与宇宙的主交心密谈。\par

39. 哦，可悲的盲目啊！不久之前，我还叩拜那些出自火炉的形像，那些在铁砧上和锤子下造成的神明，那些象骨，绘画，悬挂在老树上的花环；我若瞥见一块涂了油的石头、一块抹了橄榄油的石头，便觉得其中似有神力寓居，于是我朝它下拜，向它祈祷，向无灵的木头求福。而且，正是这些我深信存在的神明，我却曾以极其不敬的态度对待——当时我认为它们就是木头、石块和骨头，或想像它们居住在那些物件的质料内。如今，既经如此伟大的一位师者引领，踏上了真理的道路，我才明白了这一切究竟是什么；对堪受尊敬者，我心存敬意；对任何神的名号，我不再加以侮辱；并将各所应得的敬礼，无论大小，按着分明的等差，依据不同的权威，分别授予。那么，\Person{基督}难道不应被我们视作天主吗？在一切其他方面，祂都堪称为至尊至大者，难道竟不应以神圣的敬拜来尊崇祂？我们从祂那儿，在有生之日已领受了如此浩大的恩赐，并且，待那日子来临，我们还期望领受更大的恩赐。\par

40. 但祂被钉死在十字架上。这与论证何干？因为，死亡的方式与其羞辱，既不能更改祂的言行，也不会使祂教导的份量显得略有减轻；这是因为祂并非经由自然的分离，而是因遭受的暴力而离去，从肉身的桎梏中释放了自己。\Place{萨摩斯}的\Person{毕达哥拉斯}，因涉嫌图谋王权这一不公的嫌疑，而在一座庙宇中被活活烧死。难道他的学说便因此失去了特有的影响，只因他并非出于自愿，而是因残酷的袭击而呼出最后一息吗？同样，\Person{苏格拉底}因同城公民的判决而遭受极刑：难道他关于道德、德性与责任的论述，便因此化为虚无，只因他不公平地被夺去了生命吗？还有不可胜数的人，以其声名、功绩和公众地位而显赫，却经历了最残忍的死亡方式，如\Person{阿奎利乌斯}、\Person{特雷博尼乌斯}和\Person{雷古卢斯}：难道他们死后就被判定为卑鄙吗？只因为他们并非遵照命运的一般规律，而是在遭受了最惨烈的酷刑折磨之后才死去的吗？任何无辜之人，死于卑污的方式，并不因此蒙羞；任何遭受严酷刑罚的人，只要这刑罚不是源于他自身的恶行，而是由于迫害者的残暴，那么任何卑劣的烙印就绝不会玷污他。\par

41. 然而，你们这些因我们敬拜一位死于耻辱之死的人而嘲笑者啊，你们不也是通过向他献祠而尊崇\Person{利伯尔}父吗？他可是被\Person{提坦}们撕成了碎块。你们不是也在他受罚并被闪电劈死后，将医药的发现者\Person{埃斯库拉庇乌斯}奉为健康、力量与安全的守护者和保护者吗？你们不是也用祭品、牺牲和点燃的乳香呼求伟大的\Person{赫丘利}本身吗？而你们自己声称他在受罚后被活活烧死，在致命的火葬柴堆上被焚化。你们不是也在\Place{高卢}人一致的赞同下，在\Person{大母神}的庙宇中，将那位被肢解并丧失男根的弗里吉亚的\Person{阿提斯}当作慈祥而神圣的神来呼求吗？至于\Person{罗穆卢斯}父本身，他被一百名元老院议员的手撕碎，你们不是称他为\Person{奎里努斯·马尔提乌斯}，并用祭司和华丽的长椅尊崇他，在最宽敞的庙宇中敬拜他吗？而且，除此之外，你们不是还断言他已升天了吗？所以，要么你们也应受嘲笑，因为你们将死于最残酷折磨的人视为神明；要么，如果你们觉得这样做确有可靠根据，那么就请允许我们也确信我们这样做是出于何种原因、基于何种理由。\par

42. \Emph{我的对手说：}\Quote{你们敬拜的只是\Emph{仅仅}一个凡人。}即使果真如此，如前文所述，但基于祂赐给我们的许多慷慨恩赐，祂也理应被称为天主、被称呼为天主。\Emph{然而}，祂既在事实上是真实无伪、毫无疑问的天主，难道你们认为我们会否认我们以全副热忱敬拜祂，并奉祂为我们身体的守护者吗？\Quote{那么，你的那位\Person{基督}是天主吗？}某个狂怒、愤激的人会这样说。我们回答：\Quote{是天主，而且是内在权能的那位天主；}而---这可能会让不信者更加痛苦地受折磨---祂是由\Emph{至尊君王}为了一个极其重大的目的而差遣给我们的。我的对手变得更加疯狂、更加狂乱，或许会问这事能否如我们声称的那样被证明。没有比祂所行之事迹的可信度、比祂\Emph{所显示}的非凡异能之卓绝、比所有那些致命条例的征服与废除---这一切都是各民族各支派在光天化日之下看见的，毫无反对之声---更大的证明了；并且，\Emph{甚至那些}祂显明其古老法律或邦国法律充满了虚空和最愚昧迷信的人，（甚至他们）也不敢宣称这些事情是假的。\par

43. 我的对手或许会用许多其他常见的毁谤和幼稚的指控来反驳我。\Person{耶稣}是个术士；祂用秘术行这一切。祂从\Place{埃及}的神祠里偷窃了大能天使的名字和遥远国度的宗教体系。愚昧人啊，你们为何谈论那些未经查考、你们所不知晓的事，凭借着冒失舌头的唠叨胡说呢？那么，那些所行的事，难道是魔鬼的把戏和法术的诡计吗？你能给我具体指出，在过往任何时代存在过的所有术士中，有哪一位做过哪怕是千分之一类似于\Person{基督}所做的呢？有谁能在不借助任何咒语之力、不用任何草药汁液、不必焦虑守候祭献、奠酒或特定时辰的情况下，做成这些事呢？因为我们并不追究，探究他们自称能做什么，亦不探究他们全部的学识与经验通常包含在何种行动中。谁不知道这些人要么致力于预知那即将发生、无论他们愿不愿意都必然按预定而来的事？要么随心所欲地将致命且耗损的疾病加于某人；要么离间亲属的情感；要么不用钥匙打开锁闭之处；要么用沉默封住人的口；要么在战车比赛中削弱、催赶或阻碍马匹；要么在妻子们和外人的子女---无论男女---心中点燃非法爱情的火焰和疯狂欲望？或者，即便他们似乎尝试有益的事，也不是靠自己的能力，而是凭借他们所呼求的那些神祇的威力才能做成。\par

44. 然而，众所周知，\Person{基督}所行的一切奇迹，是在没有任何外物辅助、不守任何仪式、不循任何固定程序的情况下完成的，而\Emph{单单}靠着祂权柄的内在本力；并且，正如那位\Emph{真}天主的本分所当为的，与祂的本性相合，与祂的身份相称，凭着祂丰沛大能的慷慨，祂所赐予的绝非伤害或损害，而\Emph{只}是那有益、有裨益且充满对人类美善祝福的事物。\par

45. 你们又说什么呢，你们啊---？那么，祂是人，祂是我们中的一员吗？在祂的命令下、在祂的声音---以可闻可解的言语发出---之下，虚弱、疾病、热症及身体的其他病痛就都逃遁了；祂是我们中的一员吗？那些附在人身上的魔鬼族类，一见祂的面、一见祂的容，便无法承受，被陌生的能力所惊吓，逃遁而去；祂是我们中的一员吗？祂一命令，令人厌恶的麻风病立时止息，顺从复原，让原先斑驳的身体恢复了一致的肤色；祂是我们中的一员吗？祂轻轻一触，血漏就止住了，停止了过多的流淌；祂是我们中的一员吗？昏睡性水肿的水液从祂手前逃遁，流窜的体液远远避开；肿胀的身体变得健康干燥，得到了舒缓；祂是我们中的一员吗？祂吩咐瘸子行走。难道也是祂的工作，让残废的人伸直了手，甚至先天僵硬的关节也得以放松；让瘫子起身站立，那些刚才还被人抬在肩上的人，如今自己拿着床回家；让瞎子复明，生来没有眼睛的人如今看见了天和日光？\par

46. 我问你们：祂可是我们中的一位？我是说，那凭一次干预就立刻治好了一百个或更多患各种病症和疾病的人；那只需一句话，汪洋怒海便平息，旋风风暴便止息；那脚不沾湿越过最深的水潭，踏在深渊的脊背，波浪都惊愕，自然俯首为奴；那用五块饼饱饫了跟随祂的五千人，并且，为免不信和心硬的人以为这只是幻象，还用剩下的碎屑装满了十二只大篮。祂可是我们中的一位？那命令已离去的生命气息返回躯体，吩咐埋葬的人从坟中走出、并在三日后解开殡葬者的裹布。祂可是我们中的一位？那看透沉默者心中每个人所忖度、每个人隐秘的意念。祂可是我们中的一位？那说出一句话，远方言语不通的民族就以为听见了熟稔的声音和各人特有的语言。祂可是我们中的一位？那将不可辩驳的宗教义务教导门徒时，突然充满了整个世界，藉着彰显祂无边无际的权威，显明祂有多伟大、祂是谁。祂可是我们中的一位？那在身体被安放坟墓之后，光天化日下向无数人显现；和他们说话，听他们说话；教导他们，责备劝诫他们；为免他们以为被虚渺幻象所欺，一次、两次，甚至多次，在日常交谈中显现；甚至现在仍向爱祂、心灵无玷的义人显现，不是在虚幻的梦中，而是以全然单纯的形象；祂的名字一经听见，就驱逐恶魔，使占卜者缄口，阻止人求问鸟卜师，使狂妄术士的图谋挫败；并非如你们所声称的因祂的名令人畏惧，而是凭更大权能的自由运行？\par

47. 我们在圣所中陈明的这些事实，并非基于这样的假设：行事者的伟大只在这些异能中得见。因为尽管这些事多么伟大，若一旦显明祂从怎样的领域而来、是哪一位天主的仆役，它们就会显得何其琐碎微不足道。然而，祂所行的事，绝非为夸耀自己陷入空洞的炫耀，而是为使心硬不信的人确信所宣称的并非虚妄，并且现在能从祂工作的恩惠中学习想象，何为真天主。同时，我们也希望让人知道，正如前面所说，在扼要列举了祂的事迹之后，基督不仅能做祂所做的这些事，甚至也能战胜命运的定命。因为，若如显然可见且众人公认，疾病和身体的痛苦、耳聋、畸形、哑巴、筋络萎缩和眼瞎临到我们，是由命运的安排所带来，而唯有基督纠正了这些，恢复了并治愈了人，那么，比太阳本身更清楚的是：祂比命运更强大，因为祂解开了并制服了那些被永恒的结所捆绑、被不可更改的必然所固定的事物。\par

48. 然而，有人会说：你们为基督要求如此之多也是徒然，因为我们现今知道，过去也知道有其他神祇曾给许多病者施予疗法，并治好了许多人的疾病和虚弱。我并不追问，也不要求知道，是哪一位神这样做了，或是什么时候；他解救了谁，或使怎样的残躯恢复了康健：我只渴望听到这一点，是否他不用添加任何物品——即任何药物施敷——仅凭触摸就命令疾病从人身上逃去；是否他命令并迫使病因被根除，使虚弱者的身体恢复自然的力量。因为众所周知，基督或是以手触摸患处，或是仅凭声音的命令，就开了聋子的耳朵，从眼睛驱除了瞎盲，叫哑巴说话，松解僵硬的关节，使萎缩的人能行走——祂惯常以一言和一命，医好麻风病、疟疾、水肿以及各种其他病症，就是某种邪恶权势愿意人身所忍受的。所有这些神祇曾做过怎样类似的事？你们声称他们曾给患病和受困者带来帮助。因为如果他们何时曾命令——如所传闻——给某人施用药剂或特定饮食，或饮下药水，或将草木汁液敷在不适之处，或\Emph{曾命令}人们行走、静养、戒绝某些有害之物——只要你们留心省察，就可以显而易见，这并不算什么大事，也不值得大加赞叹——医生们也采用类似疗法，医生是出自尘土的受造物，并非依靠真知，而是建立在猜测体系上，在估量可能性时摇摆不定。用药剂消除侵扰人的病痛，并没有\Emph{特殊的}功德：疗效属于药物，而非施药者内在的德能；尽管知道何种医药、何种方法适合治疗人是值得称赞的，这种称赞仍可归功于人，而非神明。因为\Emph{至少}，借助外物改善人健康并不羞耻：一位神若自己不能成就，而\Emph{只能}依靠外物才赐下健全和安康，那才是神明的耻辱。\par

49. 既然你们将基督与其他神祇在赐予健康的福佑上进行比较，请问你们希望我们向你们展示多少千名病弱者；多少身患消耗性疾病的人，纵然他们走遍所有庙宇，在众神明面前俯伏，用嘴唇扫过门槛——只要一息尚存，他们便以祈祷使其厌倦，以最可怜的誓愿纠缠那位他们称之为健康赐予者的\OriginalTerm{埃斯库拉庇乌斯}{Aesculapius}本人——却没有任何方法能使他们复原？我们难道不知道有些人死于自己的疾病？有些人因疾病的折磨痛苦而衰老？有些人在日复一日、夜复一夜不断的祈祷和期待慈悲之后，反而开始过起更加放纵的生活？那么，指出一两个可能得到医治的例子又有何用呢？当时有成千上万的人未获援助，庙宇里充满了所有悲惨不幸的人。除非你们或许会说，神明帮助善人，而恶人的苦难则被忽视。然而，基督却同等地扶助善人和恶人；凡在逆境中寻求帮助以对抗暴力和厄运的人，没有一个被祂拒绝。因为，这就是真神和君王权能的标志：不拒绝施恩于任何人，也不考虑谁配得谁不配；因为，使人成为罪人的，是本性软弱，而不是他欲望的选择或清醒的判断。再者，说神明在苦难中援助配得的人，这会使你们的断言悬而未决且值得怀疑：以至于那得痊愈的人看似是偶然得救，而那不得愈的人则看似未能祛除疾病，并非由于他的缺失，而是由于一种上天所降的软弱。\par

50. 此外，祂凭自己的权能不仅行了我们简要记述的——而非按事情的重要性所要求的——那些奇迹；更崇高的是，祂还允许许多其他人尝试，并藉着祂的名去施行这些奇迹。因祂预知你们将成为祂作为及神圣工作的诽谤者，为免任何潜伏的怀疑存留，以为祂是藉法术赐下这些恩惠与福泽，于是从广大的、以仰慕惊奇之情争取祂恩惠的群众中，祂拣选了渔夫、工匠、农夫，以及类似的没有学识的人，好使他们被差遣到万国中，毫无欺诈、毫无任何物质辅助地行这一切奇迹。祂凭一句话就缓和了疼痛肢体的剧痛；他们也凭一句话止息了令人发狂的痛苦折磨。祂凭一道命令从人身上驱逐魔鬼，使那些失去知觉的人恢复神智；他们也以毫无二致的命令，使那些遭受这些\Emph{魔鬼}折磨的人恢复了健康与心智的清明。祂藉按手除去癞病的痕迹；他们也以相似的触摸使身体恢复本来的皮肤。祂命水肿肿胀的肉体恢复天然的干爽；祂的仆人同样止住了流窜的体液，命其循自身的管道流动，避免伤害躯体。对那不肯愈合的巨型疮口，祂仅凭插入一句话，便阻止了它们继续侵蚀肉体；他们也同样通过限制其肆虐，迫使顽固而无情的毒疮局限于一处瘢痕。祂赐予跛子行走的能力，赐予盲眼光明，使死者复生；他们也同样确实地松弛紧绷的筋脉，使已失明的眼睛充满光明，并命令死者从坟墓中返回，逆转葬礼的仪式。凡祂所行、引起众人惊奇赞叹的一切，祂无不自由地准许那些卑微的乡下人去施行，并置其于他们权下。\par

51. 啊，不信、顽固、刚硬的心灵！你们有何话说？你们那位伟大的\OriginalTerm{卡皮托利山的朱庇特}{Jupiter Capitolinus}，可曾将这种权能赋予过任何人？他可曾将此权授予任何一位\OriginalTerm{库里亚}{curia}的祭司、\OriginalTerm{大祭司长}{Pontifex Maximus}，不，甚至是那位其名号\Emph{显示}他为生命之神的\OriginalTerm{狄亚利斯}{Dialis}？我且不说，\Emph{他是否曾赋予}能力去复活死者、赐盲者光明、使虚弱瘫痪者肢体复原，但\Emph{他是否甚至使任何人}能够凭口里的话或手的触摸，治好一个脓包、一处倒刺、一颗粉刺？那么，这是一种人天性固有的能力吗？或者说，由一位以地上粗俗食粮喂养长大的人，能赐予这样的权利、这样的许可吗？这难道不是一种神圣且属神的恩赐吗？或者说，若此事可容许任何夸张，它岂不超越神圣属神之上吗？因为，如果你做你所能做的，做与你力量和能力相称的事，就没有表达惊奇的理由；因为你已做了你能做的事，你的能力必该完成的事，为要使行为与行事者完全相符。能够将自己的能力转移给人，与最脆弱的存在分享唯独你能做到的事，这乃是证明那统御万有、掌管一切成因与方法手段之自然法则的至高权能。\par

52. 那么，让某位祆教祭司\Person{琐罗亚斯德}从遥远的地方来，穿越炎热的地区，如果我们相信权威\Person{赫密帕斯}的话。也让这些人与他一同——那位巴克特里亚人，其事迹\Person{克特西亚斯}记载于他的\WorkTitle{历史}第一卷中；\Person{霍斯坦尼斯}的孙子、那位亚美尼亚人；以及\Person{潘菲鲁斯}，\Person{居鲁士}的密友；\Person{阿波罗尼乌斯}、\Person{达米格罗}和\Person{达达努斯}；\Person{维鲁斯}、\Person{尤利安努斯}和\Person{贝布鲁斯}；若有其他任何被认为在这种魔术技艺中拥有特殊能力和名声的人。让他们将能力赐给某个人，使哑者之口适于言语，使聋者之耳得以通畅，使天生目盲者获得视觉的自然能力，并使久已冰冷而亡的尸体恢复感觉与生命。或者，如果这\Emph{太}困难，并且他们无法将行这等事的能力传授给他人，就让他们自己施行，用他们自己的仪式。无论大地从怀中生出何种有害的草药，无论那些喃喃的咒语和伴随的符咒含有何种力量——这些让他们添加，我们不嫉妒；\Emph{那些}让他们收集，我们并不禁止。我们想试验并发现，他们能否在自己神明的帮助下，成就那些无学问的基督徒常仅凭一句话就已完成的事。\par

53. 停止你们无知地用辱骂对待如此伟大的事迹吧，这丝毫不会伤害那行事者，却会给你们自己带来危险——危险，我说，危险绝非微小，而是关乎重大的事，甚至是最为重大的事，因为毫无疑问，灵魂是宝贵的，对人来说，没有什么比他自己更珍贵。在\Person{基督}身上，并没有你们所设想的魔术，没有人为的、欺骗的或诡诈的东西；在祂内没有隐藏的奸计，尽管你们一如惯常地嘲弄微笑，纵声狂笑。祂是高天之上的天主，是至深本性中的天主，是来自不可知领域的天主，由万有的主宰派遣而来，作为救主天主；无论是太阳自身，还是任何星辰，若它们有感知，无论是世上的统治者和君王，还是那些大神，或是那些自封为王、恐吓全人类的，都无从知晓或猜测祂是从哪里来的，祂是谁——这也很自然。但是，当祂脱离了身体——那身体不过是祂携带的一小部分——祂让自己被看见，并\Emph{让世人知道}祂是何等伟大时，宇宙的诸元素被这奇异事件所震惊，陷入混乱。地震摇动了世界，海水从深处涌起，天空被黑暗笼罩，太阳的炽热火焰被遏制，其热度变得温和；当祂被揭示为天主，而之前祂被认作我们中的一员时，还能发生别的什么呢？\par

54. 但你们不信这些事；然而那些目睹了这些事发生、亲见它们在自己眼前成就的人——他们是最好的见证人、最可信的权威——自己相信了，并传给了我们这些追随者，要我们以不小的信心相信。这些人是谁？你们也许会问。是部落、民族、国家，以及那个不信的人类种族；但如果这事并不明显，正如俗话所说，比白昼本身还清楚，他们就决不会如此轻信地同意这类事件。但我们能说那时的人如此不可信、虚伪、愚蠢和粗鄙，以至于他们假装看见从未见过的事，以虚假的证据提出从未发生的事，或用孩子气的断言声称从未发生的事吗？而且当他们本可以与你们和谐相处、保持友好关系时，竟然肆意招致仇恨，并被人憎恶？\par

55. 但如果这些事件的记载如你们所说都是虚假的，那怎么会在如此短的时间内，全世界充满了这样的宗教呢？或者，那些相隔遥远、被气候和天穹的弧度所分隔的民族，怎能归于同一个信念呢？\Emph{我的对手会说}，他们只是被单纯的宣称所说服，被引向虚无的希望；在他们不顾一切的疯狂中，自愿选择面对死亡的危险，尽管他们迄今未见任何能以其奇异独特性质诱使他们接受这种敬拜方式的迹象。不，因为他们看见这一切都是\Emph{\Person{基督}}亲自和祂的宗徒们所做的，宗徒们被派遣到全世界，带去圣父的祝福，将祝福分施以造福人的心灵和身体；他们被真理本身的力量所征服，既献身于天主，又把将身体交给你们、让肉体被撕裂视为小小的牺牲。\par

56. \Emph{有人会说}，我们的作者是厚颜无耻地提出这些叙述的；他们把小事情吹嘘得过高，以最自夸的傲慢夸大了微不足道的事。但愿一切都能被记载成书——无论是祂亲自做的，还是祂的宗徒们以同等的权威和能力成就的。然而，这样一堆奇迹会使你们更加不信；也许你们能发现一些段落，似乎很有可能在事实之外有所添加，在著作和注释中插入了虚假的东西。但在那些作者所不知的民族中，在那些自己不懂文字使用的民族中，所成就的一切不可能被收录于记载，甚至无法传到所有人的耳中；或者，如果某些事被付诸于书面连贯的叙述，魔鬼和类似之人出于恶意会做些插入和添加，他们的心思就是阻碍这真理的传播：会有一些字词和音节的改动和残缺，既败坏谨慎者的信德，又损害这些事迹的道德效果。但他们绝不会因\Emph{仅仅}从书面证据中得知\Person{基督}是谁、是怎样的一位而获益；因为祂的事业建立在这样的基础上，即如果我们所说的是真实的，那么众人的承认就证明祂是天主。\par

57. 你们不信我们的著作，我们也不信你们的著作。\Emph{你们说}我们捏造关于基督的谎言；而你们提出关于你们神明的毫无根据和虚假的说法：因为没有神明从天上降下，或以其自身及生活勾勒出你们的体系，或以同样的方式贬损我们的体系和礼仪。这些是人写的；那些也是人写的——以人的言语表述；无论你们想怎样评论我们的作者，请记住，你们也会发现对你们的作者同样有力的说法。你们希望你们著作中所包含的被当作真实；同样，我们书中证实的那些事，你们也必须承认是真实的。你们指责我们的体系虚妄；我们也指责你们的体系虚妄。但你们说，我们的更古老，因此最可信可靠；仿佛古老确实不是错误最肥沃的源头，且未曾亲自产生那些在可耻的寓言中给神明带来莫大丑闻的事。因为谎言难道不能在一万年前被说出并相信吗？难道更靠近我们时代的事不是比起相隔漫长岁月的事更可信吗？因为我们这些是根据见证人的证词提出的，你们那些是基于意见；而在近期发生的事中，编造应远少于那些深埋在远古黑暗中的事，这更自然得多。\par

58. 但它们是由无学问和无知的人写的，因此不应轻信。要注意这反而可能是更强有力的理由，让人相信它们未被任何虚假陈述所玷污，而是由心地单纯的人提出，他们不懂得用浮华的装饰来粉饰故事。但语言是粗俗低下的。因为真理从不寻求虚假的润饰，对于已经确凿无疑的事，也不会沉溺于过度冗长。三段论、省略推理、定义，以及人们用来建立主张的所有那些修饰，能帮助摸索真理的人，却不能清晰标示其重要特征。但真正了解讨论主题的人，既不定义，也不推理，也不追求其它言语技巧，使听众常被欺骗，被诱骗而勉强同意某个命题。\par

59. 我的对手说，你们的叙述充满了野蛮用语和语法错误，并被极严重的错误所歪曲。这种指责确实显示出孩子气和狭隘的精神；因为，如果我们承认它是合理的，那么就让我们停止食用某些种类的果实，因为它们长有刺，以及其他不可食用的植物，一方面它们不能滋养我们，另一方面却不妨碍我们享用那些特别优越的、自然设计来对我们最有益的东西。请问，语言表达得流畅还是粗糙，究竟如何妨碍或延迟对\Emph{陈述的理解}呢？是该用重音的地方用了锐音，还是该用锐音的地方用了重音？或者，如果在数、格、介词、分词或连词上犯了错误，陈述的真实性又如何会减少呢？让那种浮夸的风格和严格规范的文字留给公共集会、诉讼、广场和法庭吧，无论如何，把它交给那些追求愉悦感觉的抚慰效果、全力关注语言华丽的人。\Emph{但是}，当我们讨论远离炫耀的事情时，我们应该考虑所说的是什么，而不是说得多么迷人，也不是如何取悦耳朵，而是它给听者带来什么益处，特别是因为我们知道，有些甚至投身于哲学的人，不仅漠视文风的精致，而且当他们本可以用更优美和丰富的语言时，却故意采用粗俗朴实的表达，以免削弱言语的严肃庄重，而去沉溺于\OriginalTerm{智术师}{Sophists}的虚饰炫耀。因为，在重要的事情上寻求享乐，确实证明了一颗毫无价值的心；而当你必须与患病者打交道时，向他们的耳朵灌输甜美的声音，却不把药物敷在他们的伤口上。然而，如果你考虑事情的真相，没有任何语言是自然完美的，同样，也没有任何语言是有缺陷的。因为，有什么自然理由，或有什么写在宇宙构造中的法则，使\OriginalTerm{paries}{paries}应称为\OriginalTerm{hic}{hic}，而\OriginalTerm{sella}{sella} \OriginalTerm{hoec}{hoec}呢？——因为它们既没有以男性和女性来区分的性别，最博学的人也不能告诉我\OriginalTerm{hic}{hic}和\OriginalTerm{hoec}{hoec}是什么，或者为什么其中一个指阳性，而另一个用于阴性。这些约定俗成是人为的，并且对所有人形成语言并非不可或缺；因为如果起初就约定这样称呼，并且如果后代在日常交谈中保持这种用法，那么\OriginalTerm{paries}{paries}也许本可以称为\OriginalTerm{hoec}{hoec}，而\OriginalTerm{sella}{sella}称为\OriginalTerm{hic}{hic}，而不会发现任何错误。然而，你们这些指责我们著作有可耻缺陷的人啊，在你们那些最完美、奇妙的书中，难道就没有这些语法错误吗？你们中难道没有人把\OriginalTerm{uter}{uter}的复数形式说成\OriginalTerm{utria}{utria}吗？另一个人说成\OriginalTerm{utres}{utres}吗？你们不也说\OriginalTerm{Coelus}{Coelus}和\OriginalTerm{coelum}{coelum}，\OriginalTerm{filus}{filus}和\OriginalTerm{filum}{filum}，\OriginalTerm{crocus}{crocus}和\OriginalTerm{crocum}{crocum}，\OriginalTerm{fretus}{fretus}和\OriginalTerm{fretum}{fretum}吗？还有\OriginalTerm{hoc pane}{hoc pane}和\OriginalTerm{hic panis}{hic panis}，\OriginalTerm{hic sanguis}{hic sanguis}和\OriginalTerm{hoc sanguen}{hoc sanguen}吗？\OriginalTerm{candelabrum}{candelabrum}和\OriginalTerm{jugulum}{jugulum}不也以同样方式写作\OriginalTerm{jugulus}{jugulus}和\OriginalTerm{candelaber}{candelaber}吗？因为，如果每个名词不能有多个性别，如果同一个词不能既是这个性又是那个性，因为一个性不能转为另一个性，那么按照阴性规则发出阳性性的词，或者将阳性的冠词用于阴性性的词的人，所犯的错误就同样严重。然而，我们看到你们不加区分地将阳性用如阴性，将阴性用如阳性，将你们称为中性的词时而这样用，时而那样用。因此，要么随意使用这些词并非错误，\Emph{在这种情况下}你们说我们的著作被极严重的语法错误所歪曲就是徒劳的；要么，如果每个词的用法是固定不变的，你们自己也会陷入类似的错误，尽管有所有那些\Person{埃皮卡狄}、\Person{凯塞利}、\Person{维里}、\Person{司考里}和\Person{尼西}站在你们一边。\par

60. 但是，我的对手们说，如果基督是天主，祂为什么以人的形象出现，又为什么按照人的方式被死亡所终结？那无形无体的能力，若不取一个更坚实的形体作为外壳，以承受目光，并使最粗心的观察者的视线也能集中于其上，又如何能来到世上，适应世界，并与人类社会相融合呢？因为，如果祂决定以其本原的性质，以其本身的特质和天主性来到世上，有哪个凡人能看见祂，能认出祂呢？因此，祂取了人的形体；借着我们族类的伪装，祂约束了祂的能力，以便能被看见和仔细观察，能说话和教导，并在不侵犯至高君王的主权和治理的情况下，完成祂来到世上的所有目的。\par

61. 那么，\Emph{我的对手}说，难道至高者就不能不伪装成人，而完成祂预定要在世上做的事吗？如果必须像你说的那样做，祂或许会那样做；但因为非必须，祂采取了别的方式。祂为何选择这种方式而不选择那种方式的原因，无人知晓，笼罩在极大的幽暗中，几乎无人能理解；然而，在为你解释你所寻求了解和倾听的事之前，如果你没有预备好不去理解，没有硬着心走向不信，你或许早已理解这些了。\par

62. 但是，\Emph{你会说}，祂像人一样被死亡剪除了。并非\Emph{基督}本身；因为死亡不可能临于\OriginalTerm{天主性}{divine}之物，也不可能使那\OriginalTerm{本体}{substance}为一、非复合、也不由各部分组合而成者，在死亡中消没。那么，\Emph{你会问}，是谁被看见挂在十字架上？谁死了？\Emph{我回答}，是祂所取、所随身背负的人性形体。\Emph{你会说}，这是一个难以置信的故事，笼罩在暗昧模糊之中；如果你愿意，它并不暗昧，而且是一个极其贴切的类比所确立的。如果\Person{西彼拉}在倾吐她的预言和神谕时，如你所说，被\Person{阿波罗}的神力充满，而那时不虔敬的强盗把她砍杀了，能说\Person{阿波罗}在她里面被杀吗？如果\Person{巴基斯}、\Person{赫勒努斯}、\Person{马基乌斯}和其他占卜者，在发狂受感时同样被夺去生命与光明，有谁会说那借着他们的口向求问者宣告当行之事的人，是按着人生命的条件灭亡了呢？你所说的死亡，乃是祂所取的人体的死亡，不是祂自己的——是那被背负者的，不是背负者的；而且，若不是为处理如此重大的事，并将\OriginalTerm{命运}{fate}那不可测度的计划在隐藏的奥迹中揭示出来，祂甚至不会屈尊忍受这样的死亡。\par

63. \Emph{你会说}，这些隐藏而不可见的奥秘是什么呢？世人不能知晓，就连那些称为世界的诸神者，也不能凭幻想和猜测以任何方式触及；\Emph{这奥秘}，除了那些蒙\Emph{基督}亲自乐意赐予如此大知识的福分、并被领进\Emph{智慧}内库的隐秘深处的人以外，\Emph{无人能发掘}。那么，你难道看不出，如果祂曾决意不让人对祂施暴，祂本应竭尽全力使敌人远离祂，甚至将自己的能力指向他们？\Emph{那么}，那位恢复了瞎子视力的主，若有必要，难道不能使\Emph{祂的敌人}眼瞎吗？那位\Emph{曾赐}力量给软弱者，要使这些人衰弱，对祂来说，是困难或麻烦的事吗？那位吩咐瘸子行走的，难道不晓得如何取去他们一切活动肢体的能力，使他们的筋腱僵硬吗？那位从坟墓中拉出死人的，难道将死亡施加在祂所愿意的人身上，会对祂困难吗？但是，因为理性要求那些已定的计划必须也在世界这里完成，并且以所完成的那种方式、不以别的方式完成，祂就以超越理解和信仰的温柔，将人们对祂的恶行视为不过幼稚的玩笑，忍受了野蛮和极其顽梗之强盗的暴力；祂也不认为值得计较他们的大胆所企图的是什么，只要祂向自己的\Emph{门徒}显明了他们有责任从祂那里期盼什么。因为，当祂就灵魂的许多危险、就种种灾祸……另一方面，这位引进者、师傅和教师制定了祂的法律和规条，使它们在合宜的职责中达到目的；祂岂没有摧毁骄傲者的狂妄吗？祂岂没有熄灭贪欲之火吗？祂岂没有遏制贪婪的欲求吗？祂岂没有从他们手中夺下武器，并从他们撕下每种\Emph{形式的}腐败的一切根源吗？总之，祂本身岂不正是温柔的、和平的、容易接近的、交谈时友善的吗？祂岂不因人们的苦难而悲伤，以祂无可比拟的仁慈怜悯凡在种种患难和身体疾病中受苦的人，把他们领回并恢复健康吗？\par

64. 那么，是什么强迫你们，是什么刺激你们去辱骂、去诋毁、去势不两立地仇恨那位没有人能指控有任何罪行的那一位呢？那些暴君和你们的君王，他们抛开对众神的\Emph{一切}畏惧，掠夺劫掠庙宇的财库；他们通过公敌宣告、放逐和屠杀，剥光了国家的贵族；——\Emph{这些人}你们却称之为\OriginalTerm{本土神}{indigites}和\OriginalTerm{神君}{divi}，并且用躺椅、祭坛、庙宇和其他仪式来敬拜他们，通过庆祝他们的庆典和生辰，而他们正是你们本应以最激烈的仇恨攻击的人。还有\Emph{所有那些}人，他们著书以各种形式用尖刻的指责攻击公共道德；他们谴责、烙印、撕碎你们奢靡的习惯和生活；他们以传世之作将自身的时代的恶评传给后代；他们试图说服人，婚姻的权利应当共有；他们与少年同寝，漂亮、淫荡、赤身露体；他们宣称你们是野兽、逃亡者、流放犯、和最无价值品性的疯狂暴怒的奴隶；——\Emph{对所有这些}，你们惊叹赞美，将他们高举至天堂的星辰，置于图书馆的神龛中，赠以战车和雕像，并且尽你们所能，以不朽称号所作的见证，仿佛赋予他们某种不死不灭。唯有\Person{基督}，你们要撕裂，要扯碎，如果你们能这样\Emph{对待一位神的话}；不，\Emph{唯有他}，如果允许，你们会用血淋淋的嘴啃咬，打碎他的骨头，像田野的野兽吞吃他。因为他做了什么，告诉我，我请求你们，因为什么罪行？他做了什么让正义的进程偏离，并激起你们被令人疯狂的折磨煽动起的凶猛仇恨呢？\Emph{难道}是因为他声称自己是由唯一真君王差遣，作你们灵魂的守护者，并带给你们那你们相信、基于少数人的断言\Emph{已经}拥有的不死不灭吗？但\Emph{即使}你们确信他说谎，他提出的希望毫无根据，即使如此，我也看不出\Emph{任何}理由你们应当仇恨他，\Emph{并}以尖刻的指责定罪他。非但如此，如果你们心灵宽厚温柔，你们甚至应当仅为此而敬重他：他向你们许诺了那些你们原可愿望并期盼的事物；他是好消息的传达者；他的信息乃是这样的：不扰乱任何人的心智，反而使\Emph{所有}人充满更少焦虑的期待。\par

65. 哦，忘恩负义而不敬虔的时代，以其异常的顽固预备了自身的毁灭！如果有一位医生从遥远、你们以前未知的地方来到你们这里，带来一种药物，能完全防止你们遭受各种疾病和不适，你们不都会急切地赶去\Emph{他那里}吗？你们不都会以各种奉承和尊敬接待他进你们的家，并善待他吗？你们不希望那种药物确实是\Emph{确实可靠、}货真价实的吗？它承诺即使在生命的最远端，你们都能免于如此无数的身体痛苦。即使这事尚属可疑，你们仍会将\Emph{自己托付}给他；你们也不会犹豫去喝下那未知的药剂，受摆在你们面前的健康希望和求安之心所驱使。\Person{基督}闪耀并显现，要告诉我们最重要的消息，带来一个繁荣的预兆，并对相信的人传达了安全的讯息。请问，这残酷、这种野蛮，更确切地说，是傲慢的骄傲，是什么意思？不仅用辱骂的言语骚扰如此大礼的使者和传递者，甚至以猛烈的敌意，以一切能倾泻于他的武器，以\Emph{一切毁灭的方式}来攻击他？他的话不中听吗？你们听见时被触怒了吗？只把它们看作\Emph{不过}占卜者的空谈吧。他说得极愚蠢，并许诺了愚拙的礼物吗？像智慧人一样嗤笑，\Emph{把他留在}他的愚拙中，任凭他在他的错谬间颠簸吧。这种凶猛，——不止一次地重复说过的话——，这种如此凶残的激情是什么意思？对一位未曾行过任何值得你们如此对待之事的人宣示势不两立的敌意；如果允许，就渴望将他肢解，而这一位不仅没有伤害任何人，反而以一致的仁慈告诉他的敌人，有什么救恩正从\OriginalTerm{至高天主}{God Supreme}那里带给他们，必须做什么才能逃脱毁灭并获得他们不知道的不死不灭？当那些所展示的奇怪且前所未闻的事，使听见的人心智动摇，令他们犹豫不信时，\Emph{他虽}是万有的主宰，死亡本身的毁灭者，却容许自己的人形被杀，好叫他们从结果中知道，那关于灵魂救恩的、他们长久怀抱的希望是稳妥的，且没有别的方法能避免死亡的危险。\par

\SourceRecord{Translated by Hamilton Bryce and Hugh Campbell.；Ante-Nicene Fathers；Vol. 6.；Edited by Alexander Roberts, James Donaldson, and A. Cleveland Coxe.；Buffalo, NY: Christian Literature Publishing Co.,；1886.；<http://www.newadvent.org/fathers/06311.htm>.}

% structure: p0001 source=h1 target=section reason=page-title

\section{驳异教徒（卷二）}

1. 在此，若有任何办法可寻，我愿稍微偏离原先设定的辩护，如此与所有憎恨基督之名的人对话：—— 若你认为在被问及问题时作答并非耻辱，那么请向我们解释并说明，是什么样的原因、什么样的理由，使你以如此刻骨的敌意追迫基督？或者你记得祂做过何种冒犯，以致一提及祂的名，你便激起阵阵狂怒与残忍的暴戾？祂可曾为自己要求王者的权柄，让最凶猛的士兵队伍遍布全世界，并将原本和平的各族中，有些人灭绝除尽，\Emph{又}强迫其他人臣服于祂的轭下并侍奉祂？祂可曾为贪婪的欲念所激动，以权利为名将所有财富据为己有——那些财富正是人们竭力争夺以求富足的？祂可曾为情欲所驱使，用暴力摧毁贞洁的屏障，或暗中觊觎他人的妻子？祂可曾因傲慢自负而趾高气扬，任意妄为地施以伤害与侮辱，不分对象？(B) 祂不配让你们听从并相信\Emph{祂，然而}，即使因这缘故——祂向你们显示了关乎你们\OriginalTerm{救恩}{salvation}的事，为你们铺设了通往天堂的道路，以及你们所渴望的\OriginalTerm{永生}{immortality}——你们也不该轻看祂；尽管祂既未将生命之光普赐给所有人，也未将所有人从不学无术所招致的危险中解救出来。\par

2. 然而，\Emph{或许有人会说}，祂理应招致我们的仇恨，因为祂将宗教从世间驱逐，因为祂阻止人们寻求敬拜神明。那么，那位将\OriginalTerm{真宗教}{true religion}带入世界，为盲目、实际生活于不虔敬中的人敞开敬虔之门，并指示他们应向谁俯伏者，难道竟被谴责为宗教的摧毁者与不虔敬的煽动者？又或者，有比这更真实的宗教——\Emph{一种}更有效、更有力、更\Emph{正确}的宗教——胜过认识\OriginalTerm{至高天主}{supreme God}，懂得如何向至高天主祈祷的？惟独祂是一切美善的源头与泉源，是万有之存续的创造者、建立者与塑造者；藉着祂，地上的一切与天上的一切得以活跃，充满生命的搏动；若没有祂，便确实没有任何具名、\Emph{有任何}实体的东西存在。但也许你怀疑这位我们所谈论的主宰是否存在，而更\Emph{倾向于}相信阿波罗、狄安娜、墨丘利、玛尔斯的存在。请作出正确的判断；环顾我们所见的这一切事物，\Emph{任何人}宁愿怀疑其他所有神明的存在，也不会对那位我们\Emph{所有人}藉着本性所认识的天主迟疑不决——无论当我们呼求\Quote{天主啊！}，还是当我们以天主为恶\Emph{行}的见证，并仰面向天仿佛祂看见我们时。\par

3. 但是，祂不允许人们向那些较低级的神明祈求。那么，你知道这些较低级的神明是谁，或在何处吗？对他们或是对他们的谈论方式的不信任，可曾触动过你，以至于你理所当然地愤慨他们的崇拜被废止、被剥夺了一切的尊荣？但若希腊人所谓的骄傲与狂妄，没有阻挡并妨碍你，你或许早已能够明白祂禁止什么事，或为何禁止；祂要让真宗教居于怎样的界限之内；你认为的服从给你带来了怎样的危险？或者你若脱离那危险的迷误，将会避免怎样的恶事？\par

4. 但当我们进一步论述时，这一切都将被更清晰、更分明地指明。因为我们将表明，基督并没有教导万民不虔敬，而是从那些最邪恶地苛待他们的人手中，拯救了无知的、不幸的人。你说，我们不相信祂说的是真的。那么，怎么样呢？你对你所说为不真实的事情，竟毫不怀疑，尽管这些事因为还\Emph{只是}即将发生，尚未显明，无论如何也无法被反驳？但祂同样没有证实祂的许诺。的确如此；因为，如我所说，\Emph{对于尚在}未来的事，无法给出证明。既然未来的性质如此，无法通过任何预期加以把握与理解，那么，对于两件不确定、悬于疑惑之中的事，选择相信那\Emph{带来}某些希望的事，岂不比相信那\Emph{毫无}带来希望的事更为合理？因为在第一种情况下，即使那所谓将要临到的事证实为空虚且无根据，也不会有危险；而在另一种情况下，则会有极大的损失，甚至是救恩的损失，如果时候到了，显明\Emph{所宣告的}并无虚谎。\par

5. 无知的你们啊，我们真想为你们哭泣悲伤，你们还有什么可说的呢？你们如此毫无畏惧，难道这些被你们藐视嘲弄的事可能是真实的吗？难道你们在私下的思念中，至少不考虑一下，恐怕你们今日以执拗的顽固拒绝相信的事，将来会为时太晚地被证实为真，而无尽的悔恨将惩罚\Emph{你们}？难道至少这些证据还不能让你们产生信德去相信吗？\Emph{即：}在如此短暂的时间内，这支庞大军队的誓愿已传遍全球？即：如今已没有任何民族如此粗鲁凶悍，而不被祂的爱所改变，征服其凶悍，并以前所未有的宁静变得性情温和？即：那些具有如此大才能的人，演说家、批评家、修辞学家、律师和医生，还有那些探究哲学奥秘的人，都寻求学习这些事，鄙视他们片刻前还信赖的事吗？即：奴隶宁愿任由主人随意折磨，妻子宁愿被休，儿女宁愿被父母剥夺继承权，也不肯对基督不忠，背弃救恩战斗的誓愿吗？即：尽管你们对遵循这宗教规诫的人宣布了如此可怕的刑罚，它反而\Emph{甚至}更加增长，大批群众更勇敢地对抗一切想要阻止它的威胁与恐吓，并因这种阻碍的企图而被激发起热切的信德吗？你们当真认为这些事是徒然且随机的吗？认为这些情感是偶然相遇而采纳的吗？那么这难道不是神圣而属天的吗？抑或\Emph{你们相信}，没有天主的恩宠，他们的心灵竟能如此改变，以至于尽管如我们所说，致命的钩子和无数其它酷刑威胁着将要相信的人，他们却如同被某种魔力摄去(A)一般，出于对一切美德的热切渴望，接受他们所认识的信德的根基，并喜爱基督的友谊胜过世间一切？\par

6. 但或许在你们看来，那些如今正在全世界联合起来，以\Emph{你们所嘲笑的}那种乐从态度去同意的人是心志软弱和愚拙的。那么怎么样呢？难道唯独你们，充满了智慧与明达的真正能力，看到全然不同而深邃的事吗？难道唯独你们觉察到这一切都是琐屑小事吗？我们宣称将要从天主那里\Emph{确实}临到我们的事，难道在你们看来不过是空言和幼稚的荒谬吗？请问，从哪里给了你们如此多的智慧？从哪里给了你们如此多的精明和才智？或者，从什么科学训练中，你们能够获得如此多的智慧，得到如此多的远见？因为你们善于按格和时态来变格变位名词和动词，并且避免粗野的词语和表达；因为你们学会了或者用和谐、有序、恰当的言语来表达自己，或者知道何时言语粗鲁不文；因为你们铭记了\Person{卢西留斯}的\WorkTitle{福尔尼克斯}和\Person{庞波尼乌斯}的\WorkTitle{马耳叙阿斯}；因为\Emph{你们知道}诉讼中要提出的争点是什么，有多少种案件，有多少种辩护方式，什么是属，什么是种，以及用什么方法区分对立面与相反者——你们就因此认为你们知道什么是假的、什么是真的、什么能做、什么不能做，以及最低和最高者的本性吗？那众所周知的格言：\BibleQuote{人的智慧在天主前原是愚妄}（\BibleRef{格前 3:19}），难道从未在你们耳中回响过吗？\par

7. 首先，你们自己也清楚看见，每当你们讨论晦涩的问题，试图揭示自然的奥秘时，一方面你们并不知道你们所谈论的、所肯定的、常常特别热诚地坚持的那些事物本身；每个人都顽固地捍卫自己的臆测，仿佛它们是已被证实和确知的\Emph{真理}。因为我们怎可能靠我们自己知道我们是否领悟了真理，即使花费一切世代去寻求知识——\Emph{我们}这些被某种嫉妒的力量所生出，被造成如此无知和骄傲的人，以至于我们虽然一无所知，却仍自欺，且被骄傲和狂妄所抬举，竟自以为拥有知识？因为，暂且不提那些属神的事物，以及那些深陷于自然晦暗中的事物，有任何人能解释那位在\WorkTitle{斐德罗篇}中著名的\Person{苏格拉底}都无法理解的问题吗——人是什么，或他从何而来，这不确定、多变、诡诈、多重、种类繁多的人？他是为何目的而被造？是由谁的巧思所设计？他在世上做些什么？(C) 他为何遭受如此数不清的灾祸？他是否像虫和鼠一样从土中得生命，由于某种湿气的作用，大地受腐烂的影响而赋予他生命；还是他从某个创造者和塑造者的手中接受了这身体的轮廓，和\Emph{这种}面容的模样？我说，他能知道这些对所有人敞开、并可为\Emph{所有人的}共同感官所认知的事情吗——我们因什么原因陷入睡眠，又因什么醒来？梦境以何种方式产生，又以何种方式被看见？更有甚者——正如\Person{柏拉图}在\WorkTitle{泰阿泰德篇}中所疑惑的——我们究竟是否曾经清醒，还是那所谓的清醒状态本身就是一场不间断的沉睡的一部分？当我们说我们看见一个梦时，我们似乎在做些什么？我们是通过从眼睛射向物体的光线来看，还是物体的影像飞至并降落在我们的瞳孔上？味道是在\Emph{被尝的}事物中，还是产生于它们接触上颚时？头发由于什么原因脱去其天然的黑色，而并非一下子全都变灰，却是点点滴滴地增添？为什么所有液体混合时，会形成一个整体，而\Emph{油，相反}，却不容许其他液体注入其中，而总是清晰地聚拢为一团，保持自身那不可渗透的实体？最后，那被你们称为不朽和神圣的灵魂，又为何在\Emph{患病的人}身中患病，在儿童身中无觉，在昏聩、愚笨而疯癫的老年中衰耗？这些\Emph{理论的}软弱和可悲的无知，在此就更大了，因为虽则我们有时或许能说出某些真实的东西，但我们甚至不能确定这一点——我们究竟是否说出了任何真理。\par

8. 既然你们惯于嘲笑我们的信德，并且用滑稽的笑话也把我们的乐于相信撕得粉碎，那么，请说，哦，你们这些饱饮智慧纯酿的才智之士，人生中可有一种需要勤勉和行动的事务，是从事者在不相信\Emph{它能够完成}的情况下，就去承担、参与和尝试的吗？你们旅行、航海，难道不是相信在事务完成后会返回家中吗？你们用犁破开土地，将各类种子撒入其中，难道不是相信随着季节的更替，你们会收获果实吗？你们与婚姻中的伴侣结合，难道不是相信这会纯洁，且是对丈夫有益的联合吗？你们生育儿女，难道不是相信他们会平安地度过生命\Emph{不同的}阶段而达到老年的终点吗？你们将患病的身体交托在医生手中，难道不是相信疾病能因减缓其严重性而得以缓解吗？你们与敌人交战，难道不是相信你们会在战斗中获取胜利而赢得凯旋吗？你们敬拜和服事神明，难道不是相信他们存在，且垂听你们的祈祷而施恩吗？\par

9. 怎么，你们曾亲眼看见、亲手触摸过那些你们自己写下的、时常阅读的、关于超越人类知识之事物吗？难道不是人人都信任这个或那个作者吗？那被某人说服而相信是由另一人真实说出的东西，难道他不是用某种仿佛\Emph{类似于信仰的}赞同去捍卫吗？难道不是说火或水为万物本源的人，把他的信仰寄托在\Person{泰勒斯}或\Person{赫拉克利特}身上吗？把\Emph{万物的}原因置于\OriginalTerm{数字}{numbers}之中的，归信于\Place{撒摩}的\Person{毕达哥拉斯}和\Person{阿尔基塔斯}吗？把灵魂划分，并建立无形体之形式的，归信于\Person{苏格拉底}的门徒\Person{柏拉图}吗？把\OriginalTerm{第五元素}{fifth element}加到那些\OriginalTerm{第一原因}{primary causes}之上的，归信于逍遥学派的祖师\Person{亚里士多德}吗？用\Emph{火焚来}威胁世界，并说时机一到世界就将\Emph{被火}焚烧的，归信于\Person{帕奈提乌斯}、\Person{克吕西波}、\Person{芝诺}吗？总是从\OriginalTerm{原子}{atoms}构造诸世界、又将\Emph{它们}毁灭的，归信于\Person{伊壁鸠鲁}、\Person{德谟克利特}、\Person{梅特罗多鲁斯}吗？说人什么都不能理解、万事都包裹在浓重的晦暗之中的，归信于\Person{阿尔凯西劳斯}、\Person{卡尔内亚德}吗？——总之，归信于老学园和新学园的某位老师吗？\par

10. 最后，连那些已经提到的学派的领袖与创立者，不也是凭着对自己观念的相信而说出他们所说的那些话吗？因为，\Person{赫拉克利特}可曾\Emph{目睹}火的变化所产生的事物？\Person{泰勒斯}可曾\Emph{目睹}水凝结而生的事物？\Person{毕达哥拉斯}可曾\Emph{目睹}它们从数中生出？\Person{柏拉图}可曾\Emph{目睹}那\OriginalTerm{无形相}{bodiless forms}？\Person{德谟克利特}可曾\Emph{目睹}\OriginalTerm{原子}{atoms}的聚合？或者，那些断言人根本无法认识任何事物的人，是否知道他们所说的是真是假，从而明白他们所提出的命题本身就是真理的宣称？因此，既然你们并未发现、也未曾学到任何东西，却凭轻信而断言你们所写下的这一切，并写满了成千上万本书；请问，这是怎样一种判断？竟然如此不公，你们嘲笑我们身上的信德，却明明看到你们自己同样有着我们那种乐于相信的态度。但\Emph{你们却说}你们相信那些精通各种学问的智者！——那些实际上什么都不知道、所说的互不一致的人；他们为了自己的意见而与对手交战，总是怀着顽固的敌意激烈地争斗；他们互相推翻、反驳、摧毁对方的学说，使一切变得可疑，并且从他们彼此的冲突中表明，根本无法认识任何事物。\par

11. 但是，\Emph{假定}这些事丝毫不会阻碍或阻止你们在很大程度上必须相信并听从他们；那么，有什么\Emph{理由}让你们在这方面有更多\Emph{自由}，而我们\Emph{应有更少}呢？你们相信\Person{柏拉图}、\Person{克罗尼乌斯}、\Person{努梅尼乌斯}，或任何你们喜欢的人；我们则相信并信赖\Person{基督}。这多么不合理啊：当我们双方都遵守各自的老师，并且都共同拥有相信这一事情时，你们竟希望允许自己去接受他们所说的话，\Emph{却}不愿意倾听并观看由\Person{基督}所提出的东西！然而，如果我们选择将理由与理由相比较，我们更有能力指出我们在\Person{基督}内所随从的，胜过\Emph{你们指出}你们在哲学家们身上\Emph{所随从的}。的确，我们在祂内所随从的是这些——祂所彰显并展示在各种奇迹中的光荣事工和至大的德能，藉着这些奇迹，任何人都可被引领而\Emph{感受到}相信的必要性，并且\Emph{可以}满怀信心地断定，这些事并非可被视为人力所及的，而是\Emph{显示了}某种神圣而未知的大能。你们在哲学家们身上所随从的是什么德能，竟使你们比我们相信\Person{基督}更有理由\Emph{去相信}他们呢？他们当中可曾有人，仅凭一句话，或一个命令——我尚且不说是去平息、遏止大海的狂怒或风暴的肆虐；使盲者复明，或赐光明给生来失明的人；呼唤死者复活；结束多年的病痛——但，这是更容易得多的——仅凭一句斥责便治愈一个脓疮、一处疥癣，或一根扎入皮肤的刺呢？我们并非否认，他们因品行端正而值得称赞，也并非否认他们精通各种学问和知识；因为我们知道，他们能以极为优雅的语言说话，\Emph{其言辞}如经过雕琢的辞章般流畅；他们能以最敏锐的思维进行三段论推理；他们能恰当地安排推论；他们能藉定义来表达、划分、区分原理；他们论及数字\Emph{不同}种类，论及音乐；他们通过格言和训导也解决几何学的问题。但是，这些与事情何\Emph{干}呢？\OriginalTerm{省略三段论}{enthymemes}、\OriginalTerm{三段论}{syllogisms}及其他这类东西，就能使我们确信这些人知道真理吗？或者，他们因此就如此可信，以致对于极为隐晦的主题，我们必须给他们以相信么？对人物的比较，不能凭口才的犀利，而要凭\Emph{他们所成就的}事工的卓越来判定。不可称那清楚表达自己的人为好老师，而应称那以神圣的作为保证其许诺之人为好老师。\par

12. 你们提出反对我们的论证和臆测的诡辩，这些——我这样说，但愿不冒犯祂——倘若基督亲自在各民族聚集中使用，谁会赞同？谁会聆听？谁会说祂清楚明白地决断了什么？或者，即便有人鲁莽而全然轻信，在祂倾吐空洞无稽之言时，谁会跟随祂？祂的德能\Emph{已经}向你们显明，那掌管万有的前所未闻的能力，无论是祂公开施行的，还是那些传扬祂名的人普世所致力的，这能力已平伏情欲的烈火，并促使本性千差万别的种族、民族和国家同心一意，急赴领受同一信仰。因为那些在\Place{印度}、\Place{赛里斯}人、\Place{波斯人}和\Place{玛待人}中间，在\Place{阿拉伯}、\Place{埃及}、\Place{亚细亚}、\Place{叙利亚}，在\Place{迦拉达人}、\Place{帕提亚人}、\Place{弗里吉亚人}中间，在\Place{阿哈雅}、\Place{马其顿}、\Place{伊庇鲁斯}，在日出日落所照的一切岛屿和省区，最后在\Place{罗马}，即\Emph{世界}之都，在那里，尽管人们忙于\Person{努马}王所引进的礼俗和古代迷信的奉行，却仍急弃祖宗的生活方式，而皈依基督的真理。因为他们曾见术士\Person{西满}的战车和火轮，因\Person{伯多禄}的口而被吹散，并在呼唤基督之名时消失。我说，他们曾见\Emph{他}信赖假神，在恐惧中被假神遗弃，因自身重量倒栽下来，双腿折断，扑倒在地；\Emph{后来}，当他被抬往\Place{布伦达}时，因痛苦和羞耻而心力交瘁，又再次从一座极高房屋的屋顶纵身跳下。但这一切事迹你们既不知道，也不愿知道，从不认为这些对你们至关重要；你们只信赖自己的判断，将那份狂妄自恃称为智慧，便给欺骗者——就是那些有罪之人，我说，其利益在于贱辱基督徒之名——以兴云作雾、掩蔽如此重要真理的机会；他们夺取你们的信仰，代之以轻蔑，好使你们因觉得与他们相称的结局正在威胁他们，而在你们心中激起一种情感，令你们陷入危险，丧失天主的仁慈。\par

13. 然而，你们这些对博学之士的教义和哲学惊奇叹赏的人啊，你们岂不认为，讥笑嘲弄我们如同说愚妄无稽之言，是极不公义的吗？因为你们自己也被发现，当由我们说出和发表之时，你们所说的正是这些或几乎相同的话，你们反倒耻笑。我并非对那些分散在学派各条岔路、因意见分歧而形成各种微不足道之派系的人说话。你们，我正是对你们说话，你们热心跟随\Person{墨丘利}、\Person{柏拉图}和\Person{毕达哥拉斯}，以及你们其余志同道合、在同样教义道路上齐心同行的人。你们竟敢嘲笑我们，因为我们崇敬朝拜宇宙的造物主和主，并将我们的希望交托给祂？你们的\Person{柏拉图}在\WorkTitle{泰阿泰德篇}中，特予指出，究竟说了什么？他岂不勉励灵魂逃离尘世，并尽其所能，不断思念冥想于祂？你们竟敢嘲笑我们，因为我们说将有死者的复活？我们诚然承认我们这样说，但声明你们对它的理解与我们持有的大有出入。同一位\Person{柏拉图}在\WorkTitle{政治家篇}中说了什么？他岂不声称，当世界开始从西方升起并趋向东方时，人将再次从大地的腹中涌出，年迈、白发、垂老；而当更遥远的年代开始临近时，他们将经过现在成长为人所走的同样阶段，逐渐回归婴儿的摇篮？你们竟敢嘲笑我们，因为我们致力于我们灵魂的救恩？——也就是说，我们\Emph{关顾}我们自己：因为我们人是什么，不就是被禁锢在肉身内的灵魂吗？——而你们，确然不竭尽心力保灵魂之安宁，因为你们不戒绝一切恶习与情欲；对此你们所焦虑的，是紧紧依附\Emph{你们的}身体，仿佛被不可分割地束缚于它们。——那些奥秘仪式又是什么意思？你们在其中恳求某些\Emph{未知}的神力帮助你们，并在你们返回本然居所时，不设置任何阻碍来拦阻你们？\par

14. 当我们谈论地狱和不灭的火，就是我们得知灵魂被其仇敌投入其中的那个火时，你们竟敢嘲笑我们吗？怎么，难道你们的\Person{柏拉图}，在他所著的\WorkTitle{论灵魂不朽}一书中，不也提到了\OriginalTerm{阿刻戎}{Acheron}、\OriginalTerm{斯堤克斯}{Styx}、\OriginalTerm{科库托斯}{Cocytus}和\OriginalTerm{皮里佛勒革同}{Pyriphlegethon}这些河流，并断言灵魂在其中被翻滚、吞没并焚烧吗？然而，\Emph{虽然}他并非缺少智慧，且具有精确的判断力和洞察力，但他试图解决一个无法解决的问题；以至于他一方面说灵魂是不死、永存且没有形体的，另一方面却又说灵魂受惩罚，让它们受苦。但是，哪个人看不出那不死而单纯的东西不可能受任何痛苦；反之，那受痛苦的就不可能不死呢？尽管如此，他的观点离真理并不太远。因为这位温良仁慈的人虽然认为判灵魂死刑是残忍不人道的，但他却并非毫无道理地设想它们被投入火焰汹汹、深渊污浊的河流中。因为它们被投入其中，并被消灭，在永恒的毁灭中徒然消亡。因为从基督的教导中我们得知，它们处于一种中间状态；\Emph{它们}是这样的：如果它们不认识天主，就会灭亡；如果它们留意了祂的恐吓和\Emph{赐予的}恩惠，就能从死亡中得解救。为了显明那未知之事，这才是人真正的死亡，什么也不留下的死亡。因为肉眼所见\Emph{只是}灵魂与肉体的分离，而不是终局——灭亡：我说，这才是人真正的死亡，当那些不认识天主的灵魂将被怒火的长期折磨燃尽，某些极其残忍的\Emph{存在}者，\Emph{它们}在基督之前不为人知，唯藉祂的智慧才被显明，将把它们投入那火中。\par

15. 因此，没有理由让那话误导我们，让我们抱虚妄希望，这些话出自一些至今无名、自视过高的人之口，他们说灵魂是不死的，地位仅次于作为世界之主和统治者的天主，源自那位神圣、智慧、博学的父和祖先，且不受身体接触的影响。现在，因为这是真实而确定的，因为我们是由那位完美无瑕的主所创造的，所以我们生活得无可指摘，\Emph{我想}，因此无可指责；我们是良善、正义、正直的，丝毫没有堕落；没有情欲压倒我们，没有邪欲贬低我们；我们坚持不懈地操练一切德行。而且，因为我们所有的灵魂都有同一个根源，因此我们的想法完全一致；我们的品行没有差异，我们的信仰没有差异；我们都认识天主；世上的意见并不像世人的数目那样多，也不\Emph{分为}无限多样的种类。\par

16. 但是，\Emph{他们说}，当我们迅速向下移动趋向我们有死的肉身时，来自\OriginalTerm{世界诸圈}{world's circles}的原因追逐着我们，通过这些原因的作用，我们变坏了，唉，甚至变得极恶；我们欲火和怒火中烧，在可耻的行为中度日，且因出卖我们的身体而委身于一切贪欲。可是，物质怎能与非物质结合呢？或者，天主所造的，怎会被较弱的原因引导，借着恶习的实践而自甘堕落呢？人啊，你们声称天主是你们的父，并坚称自己像祂一样不死，你们愿意放下惯常的傲慢吗？你们愿意探问、审视、查究你们自己是什么，属于谁，你们\Emph{被认为}有何出身，你们在世做什么，你们如何出生，如何跃入生命吗？你们愿意放下\Emph{一切}偏心，在思想的静默中思考，我们或者是与其它受造物十分相似的受造物，或者并未被很大的差异分隔开吗？有什么能表明我们不像它们呢？或者我们有何优越之处，以至于我们不屑于被列入受造物之列呢？它们的身体由骨骸建构，以筋腱紧紧捆绑；我们的身体同样由骨骸建构，以筋腱紧紧捆绑。它们通过鼻孔吸气，又在呼吸时呼出；我们同样吸入空气，并以频繁的呼吸呼出。它们被分为雌雄两类；我们同样被我们的创造主塑造成同样的性别。它们的幼仔从子宫出生，经由两性结合而受孕；我们同样从性交合中出生，从母胎中诞生并送入生命。它们靠吃喝维持，并通过下体排泄残留的污物；我们同样靠吃喝维持，并以同样的方式处理自然所拒绝的东西。它们关心防范致死的饥荒，必然要留意食物。我们生活中如此紧迫的忙碌，除了寻找避免饥饿危险、消除焦虑忧虑的方法外，还能有什么目的呢？它们受疾病和饥饿的侵袭，最终因年老而失去力气。那么，怎么样呢？我们难道不受这些恶害的侵袭，难道我们不是同样被有害的疾病削弱，被消耗人的年老摧毁吗？但是，如果在更隐秘的奥秘中所说的也是真的，即恶人的灵魂在离开人的身体后，进入牲畜和其它动物内，那就\Emph{更加}清楚地表明我们与它们同类，没有很大的差距，因为正是基于同一理由，我们和它们才都被称为活物，并被如此对待。\par

17. 但是，\Emph{有人会说}，我们拥有理性，在悟性上超越所有哑口动物。我或可相信此说全然属实，但前提是所有人都能依循理性与智慧生活，从不偏离本分，戒绝禁行，远离卑污，且\Emph{若}无人因愚昧无知而妄求对自己有害且危险之事。然而，我愿知此理性究竟为何，竟能令我们超乎一切动物族群之上。\Emph{难道}是我们为己建造了房舍，藉此躲避冬寒与暑热？怎么！难道其他动物未曾在此方面展现深谋远虑？我们岂非眼见，有的在最便利处筑巢为居；有的藏身于岩石与高崖之间，以求\Emph{自身}得庇护稳妥；有的掘地造穴，在挖出的坑洞中为自己备下堡垒巢穴？然而，倘若赋予它们生命的本性，也肯赐予辅助的手，它们无疑同样会筑起高楼，创出新颖工艺。然而，即使在它们以喙爪所造之物中，我们见到诸多理性与智慧的迹象，无论我们如何思索，都无法仿效，尽管我们有双手，能灵巧胜任各类工作。\par

18. \Emph{有人会对我说}，它们未曾学会制作衣物、坐具、舟船、耕犁，简言之，即家庭生活所需的其他器具。这些并非知识的赠礼，乃是极度紧迫之需的提示；技艺也非随\Emph{人的}灵魂自至深之天界降临，而是在这地上，穷搜苦求，渐次显明，经长久时日与慎思方渐掌握。然而，倘若灵魂\Emph{自身}拥有那与\Emph{神圣而不朽}之族类相配的知识，那么所有人自始便当知晓一切；也就不会有任何时代对某门技艺一无所知，或缺乏实务知识。然而当下，一种匮乏且需求众多的生活，注意到某些事偶然发生而对自己有利，既模仿、尝试、试验，又失败、改铸、变更，从屡次失败中，为自己获取并铸就了对某些技艺的浅显认识，将许多世代的进展汇集于一点。\par

19. 然而，倘若人能透彻认识自己，或对天主有丝毫认识，就绝不会自诩拥有神圣而不朽的本性；也不会因\Emph{他们已制造了}烤架、盆、碗，制作了内衣、外衫、斗篷、披肩、官袍，以及刀、胸甲、剑、鹤嘴锄、手斧、耕犁，便自以为有何了不起。我敢说，他们绝不会被骄傲与狂妄所驱，而因他们发明了文法、音乐、雄辩术和几何学，便自视为最高品秩的神明，与那至高者平起平坐。因为我们看不出这些技艺有何等\Emph{如此}奇妙，竟使人相信，因着这些发明，灵魂便高于太阳及所有星辰，在宏伟和本质上超越整个宇宙——那原是其组成部分的宇宙。因为他们声称能宣告或教导的，无非是让我们学习知道名词的规则与差异，辨识\Emph{不同}音调的间隔，在诉讼中雄辩陈词，或丈量大地的边界罢了。若灵魂果真自天界携来这些技艺，且人不可能不知晓它们，那么普世之人早该孜孜于此，也就绝不会发现任何民族未能同样齐备地掌握这一切。然而当下，世上的乐师、逻辑学家和几何学家何其稀少！演说家、诗人、文评家又何其稀少！由此可见，正如一再强调的，这些事物乃是在时间与环境的催迫下被发现的，灵魂并非受神圣引导而降临于此，因为既非人人博学，亦非人人能学；且他们中许多人略欠机敏，甚至愚钝，\Emph{唯有}因畏惧鞭笞，才勉强致力于学问。然而，倘若我们所学之事不过是回忆——如古人学说中所主张的——那么，既出一源，我们便该同等地学习，同等地记忆，而非持有各异、繁多且矛盾的见解。而今，既然我们各执一词，便清晰显明：我们从天上一无所携，不过是熟识了此处兴起之事，并持守那在我们思想中扎下深根者。\par

20. 为使你们更清晰、分明地看出，你们相信与至高能力极为肖似的人，其价值几何，请设想此景；若能直接与之接触便可做到，且让我们就当实际接触一般来构想。那么，让我们设想地上掘出一处穴居之所，成形为室，屋顶四壁围护，冬无严寒，夏无酷暑，调适匀和，我们既不感觉寒冷，亦不受夏季之酷热。于此，任何声响呼号皆不可传入：无论鸟鸣、兽吼、风暴、人声——总之，任何噪声，乃至雷霆之可怖巨响。再让我们设想一种照明方式，不用引火，不见太阳，而是造出某种仿品，模仿日光，中间隔以黑暗。\Emph{但}不要只开一门，亦不可直接出入，须经迂曲盘旋方可抵达，若非万不得已，绝不开启。\par

21. 我们已经为我们的设想预备了地方，接着让我们迎接一位生来就住在那里的人，那里除了空无之外什么也没有——比如一位\Person{柏拉图}或\Person{毕达哥拉斯}一类的人，或是一位被认为具有超人才智者，或被诸神的神谕宣称为最具智慧的人。做完这些之后，就必须用适宜的食物来喂养和抚育他。因此让我们也找一位乳母，她总是裸体前来，永远沉默，一言不发，绝不开口动唇说话，只是在他吸乳之后，并做完其他必要之事，就让他熟睡，然后日夜守在紧闭的门前；因为乳母的照顾通常必须近在身旁，而且\Emph{她}必须留意他不停变化的活动。而当孩子开始需要更实在的食物支持时，仍由同一位乳母，依然不穿衣服，保持同样不破的沉默，将食物带进来。所带进的食物也必须始终完全一样，材料上毫无差别，也不通过不同调味重新烹煮；而应该是粟米粥，或是斯佩耳特小麦面包，或是效法古人，在热灰中烤熟的栗子，或从树上摘下的浆果。此外，绝不要让他学会喝酒，除了清纯的泉水之外，不用别的东西解渴，而\Emph{那}水如果可能，就用手心捧到嘴边来喝。因为习惯会变成\Emph{第二}本性，由习俗而变得熟习；他的欲望就不会扩展得更远，因为不知道还有别的什么可以追寻。\par

22. 那么，\Emph{你们问}，这些事用意何在？\Emph{我们提出这些}，是为了——既然人们相信人的\Emph{灵魂}是神圣的，因而是不死不灭的，并且它们带着一切知识进入人的身体——我们可以从这个我们设想以这种方式抚养的\Emph{儿童}身上，试验这种说法是否可信，抑或只是出于欺骗性的预想而轻率相信并视之为当然。那么，让我们设想他在一处与世隔绝的偏僻地方长大，度过任意多年，二十年或三十年——不，就让他活满四十年后才被带到人群当中；如果他的确是神圣本质的一部分，源自生命之泉而居住此世，那么在他认识任何事物，或熟悉人类的语言之前，先让他被人询问并回答：他是谁，或是由哪个父亲出生在哪个地区，如何或以何种方式被抚养长大；他前半生从事过什么工作或行业。那时他难道不会像哑巴一样站着，比任何野兽、木头、石头更缺乏才智和辨别力？当他接触到陌生而前所未见的事物时，岂不格外不认识自己？如果你问他，他能不能说出太阳是什么，大地、海洋、星辰、云、雾、阵雨、雷、雪、冰雹是什么？他能知道树木是什么吗，青草或杂草，一头公牛、一匹马或一只公羊，骆驼、大象或鸢是什么吗？\par

23. 如果你在他饥饿时给他一颗葡萄、一块葡萄汁饼、一个洋葱、一株蓟、一根黄瓜、一个无花果，他会知道所有这些都能解饿吗？或者知道每一样东西要怎样才\Emph{适宜食用}？如果你生起一堆大火，或用有毒的生物围绕他，他会不会穿过火焰、毒蛇、\OriginalTerm{狼蛛}{tarantulae}之中，而不知道它们有危险，甚至不知畏惧？再者，如果你把城居与乡居所需的衣服和器具摆在他面前，他真能分辨每样东西各适于什么用途吗？能说出它们适合完成什么服务吗？他能否说出\OriginalTerm{铺毯}{stragula}、头巾、腰带、束发带、坐垫、手帕、斗篷、面纱、餐巾、毛皮、鞋、凉鞋、靴子各有什么衣着上的用途？而且，如果你继续问什么是轮子、什么是雪橇、簸箕、罐子、盆、榨油机、犁铧、筛子、磨石、犁柄或轻锄；雕花座椅、针、\OriginalTerm{刮身板}{strigil}、水盆、无靠背座、长柄勺、大盘、烛台、高脚杯、扫帚、杯子、袋子；里拉琴、笛子、银、黄铜、金、书、杖、卷轴，以及人类生活所围绕和维持的其他一切器物呢？在如此情况下，他岂不像我们所说的那样，如一头牛或一头驴，一头猪，或其他更蠢的野兽，看着这些东西，观察它们各式各样的形状，却完全不知道它们是什么，也不明白它们为何被保存？如果以任何方式强迫他发出声音，他岂不会张着嘴，像哑巴通常那样，含混不清地喊出些什么？\par

24. 喔，柏拉图啊，你为何在\WorkTitle{美诺篇}中，向一位年轻的奴隶提出某些关于数的学说的问题，并力图通过他的回答来证明：我们所学习的东西其实并非学习，而\Emph{仅仅}是回忆起我们从前所知道的那些事物呢？如今，他若能正确回答你——因为我们不应不信你的话——他之所以\Emph{如此}，并非凭借真实的知识，而是凭借他的\OriginalTerm{智性}{intelligence}；并且，由于他每日使用数字而对其有所了解，这才在被问到时能领会\Emph{你的意思}，而乘法运算过程本身也总是提示他。然而，你若确然相信灵魂——\Emph{属人的}——是\OriginalTerm{不死不灭}{immortal}的，并且在\Emph{它们}飞至此世时赋有知识，那就停止向那位你明见其为无知且习惯人间行径的少年提问吧；召来那位四十岁的男子，并问他——\Emph{不是}问那些关于三角形、正方形的离奇或隐晦的问题，不问何为立方体或二次幂、九与八之比，甚或四与三之比；而是问他众人皆知之事：二的两倍是多少，或二的三倍是多少。我们愿见，我们愿知，当被问及这期望中的难题时，他会给出怎样的答案。在此情形下，他纵然双耳张开，是否能觉察你是在说些什么，问些什么，或要求他给出某种答复？他岂不将如木桩，或如俗语所说的\Place{马耳珀萨岩石}一般，哑然无言，呆立不动，甚至连这一点都不理解、不知晓——你是在和他交谈还是和别人交谈，是在跟另一人对话还是跟他对话；你口中所发出的，是可解的言语，抑或\Emph{仅仅}是毫无意义的叫喊，徒然拖长而无的放矢？\par

25. 你们这些将过多不属于自己的卓越归于自己的人哪，有何话可说？这就是你们所描述的那个有学问的灵魂吗？——不朽、完美、神圣，居于第四位，在天主、宇宙的主宰之下，也在那同类的神灵之下，从生命之泉涌出。这就是那珍贵\Emph{的存在}——人，赋有最高超的理性能力，据说是个\OriginalTerm{微观宇宙}{microcosm}，\Emph{是}按照整个\Emph{宇宙}的样式被创造和形成的；正如已被看到的，并不比任何禽兽优越，比木桩\Emph{或}石头更无感觉；因为他不认识人，总是生活着，在寂静的荒漠中闲荡，纵使富有，活了无数年岁，也从未挣脱身体的束缚。但是，\Emph{你们说}，当他进入学校，接受师长的教导时，他就变得智慧、有学问，并抛弃了先前一直缠身的无知。而一头驴，一头牛，若被长期强迫练习，也会学着犁地和磨谷；一匹马，学会听从轭具，奔跑时顺从缰绳；一匹骆驼，在装载或卸载时学会跪下；一只鸽子，放飞后会飞回主人的家；一条狗，发现猎物时会控制并压抑自己的吠声；一只鹦鹉，甚至能发音吐字；一只乌鸦，能说出名字。\par

26. 当我听到灵魂被说成是某种非凡之物，与天主亲密相近，\Emph{并且}来到此世时洞悉过往一切时，我愿它能教导人，而非学习；它不该如俗话所说，在有了高深知识后再倒退到初步，而应牢牢持守它进入尘世身体时所学的真理。因为若非如此，如何能分辨出灵魂是在回忆，还是\Emph{初次}学习所听到的？毕竟，相信它学习它所不知的，远比相信它忘记了自己前刻还知道的东西、且因身体的介入而丧失了回忆能力，要容易得多。而且，主张灵魂\Emph{是}无形无体，没有实体，这一学说又将如何？因为凡不与任何身体形式相连者，不会受另一对立者的阻碍，也无法使任何东西去摧毁那不可能为对立面所触及者。正如一种在身体中确立的比例，虽在千般事例中隐而不见，却仍保持不受影响、稳然安定；同样，灵魂若非如声称的那样是物质的，就必定无论被身体多么严密地封闭，都保持对过往的知识。再者，同样的推理不仅表明它们不是\OriginalTerm{非形体}{incorporeal}的，甚至剥夺了它们所有的\OriginalTerm{不死不灭}{immortality}，并将它们归于生命通常结束的界限之内。因为凡是因某些诱因而被迫变异改形，以至不能保持其自然状态者，必然从根本上被视为\OriginalTerm{被动的}{passive}。而那易于遭受、暴露于痛苦者，正因其忍受能力而被宣告为\OriginalTerm{可朽坏的}{corruptible}。\par

27. 因此，如果灵魂一与身体结合便丧失所有知识，它们必定经验了某种性质的事物，使其变得盲目健忘。因为，它们若非屈从于无论何种东西，便不可能在保持自然状态的同时抛开自己的知识，也不可能不经历自身的变化而转入另一状态。不，我们反倒认为，凡是一、不死不灭、单纯者，无论其为何物，都必然永远保持自己的本性，并且，如果它果真立志忍耐并常住于真正不死不灭的界限内，就不该也不能屈从于任何东西。因为一切痛苦都是死亡与毁灭的通道，一条通向坟墓、带来无法逃脱的生命终结的路；而灵魂要是容易受其影响，并屈服于其作用与侵袭，那么它们所享有的生命就只是暂借予当前用途，而非稳固的产业，尽管有些人得出另外的结论，并在如此重大的事情上信赖自己的论证。\par

28. 然而，为了我们离开你们时不至于像\Emph{从前那样}一无所知，请让我们听听你们是怎么说的：灵魂被包裹在尘世的身体中，对过去没有任何记忆；然而，一旦确实被安置在身体之内，并因与之结合而变得几近麻木之后，它却能牢牢地、忠实地记着许多年前——就算是八十年，甚至更久远——的所作所为、所受之苦、所言所闻。因为，如果由于受身体的羁绊，它不记得自己早先知道、在来到这世界之前就知道的事情，那么它更应该忘记自从被关进身体以来一次次所做的事情，而不是在进入身体之前、尚未与人联结时所做的事情。因为，那夺去进入其中的灵魂之记忆的同一具身体，本应使它内部所行的事也被全然忘记；因为同一个原因不能带来两种彼此对立的\Emph{结果}，致使一些事被遗忘，却让另一些事被行这事的人记住。然而，如果灵魂——按你们的称呼——因其\Emph{属肉体}的肢体而受阻，无法忆起先前的知识，它们又如何能记得正是在这些身体内所安排的事，并知道自己乃是\OriginalTerm{灵体}{spirits}、没有形体物质，且因其不朽者的身份而尊贵呢？\Emph{它们又如何知道}自己在宇宙中拥有的位阶，按什么秩序被\Emph{从其他存在中}分别出来？它们是如何来到这些宇宙最低的区域？它们在滑向这些区域时获得了什么属性，来自何种\OriginalTerm{天环}{circles}？我要问，它们又如何知道自己曾满腹学识，并因身体带来的障碍而丢失了知识？因为，对于这件事本身它们也该一无所知：即与身体的结合是否给它们带来了任何污点；因为知道你自己过去是什么，而今天不是，这并非你失去了记忆的征兆，而是记忆完全健全的证明与证据。\par

29. 既然事实如此，我恳请你们，不要再把琐细无谓之事估价过高。不要再把人置于上等的秩序，因为他是最低的；也不要再把人置于最高的等级，因为所考虑的只是他这个人，他一无所有，在居所住处上贫穷匮乏，\Emph{从未}有资格被宣称出身显赫。因为，身为义人和公义的持守者，你们本应制服骄傲与傲慢，这些恶习使我们众人高抬自己，充满空虚的虚荣；你们不仅认为这些恶习是自然产生的，而且——更糟糕的是——你们还在其中增添了使罪孽增长、邪恶无可救药的原因。因为，即便天性一向远避恶名与羞耻之事的人，当他听到极有智慧的人说灵魂是不朽的，且不受命运法命的辖制，又有谁不会全速投身于种种罪孽之中，\Emph{并且}毫无畏惧地参与和图谋不法之事呢？简言之，在不受拘束的欲望所需求的一切事上，\Emph{有谁}不会因确信自己安全且不受惩罚而变本加厉地去满足自己的私欲呢？因为，什么能阻止他这样做呢？是对在上权能和天主的审判的畏惧吗？而一个深信自己就如至高天主本身一样不朽，并且深信因为两者同样不朽，那不朽的存在不会受到另一个\Emph{仅是}它的等同者的困扰，所以天主无法向他宣布任何刑罚的人，又如何能被任何恐惧或害怕所压倒呢？\par

30. 但他\Emph{难道不会被}我们听说的\OriginalTerm{冥府}{Hades}中的惩罚所吓倒吗？这些惩罚，\Emph{如他们所设想的那样}，还伴有种种酷刑形式。谁又会如此愚昧，不明就里，竟相信对不朽的灵体而言，\OriginalTerm{塔尔塔罗斯}{Tartarus}的黑暗、火河、带深渊的泥沼、或疾旋于空中的轮子，能以任何方式造成伤害呢？因为，那不可触及、不受毁灭法则支配的东西，纵使被狂怒洪流的全部火焰包围、在泥泞中翻滚、被悬垂的岩石崩塌和巨山倾倒所掩埋，也必安然无恙，不受任何致命损伤。\par

此外，这种深信不仅因其\Emph{暗示的}犯罪自由而引人作恶，甚至也取消了哲学本身的依据，并声称从事哲学的研修是徒劳的，因为其工作困难，且毫无结果。因为，如果灵魂果真永无止境，并随一切世代不断前进，那么让自己沉溺于感官快乐——轻视、忽略那些美德，只因\Emph{顾及}这些美德，生活在享乐上就会更加拮据，且\Emph{变得}不那么诱人——并任由无节制的欲望急不可耐、毫无阻遏地驰骋于各种放荡淫逸之中，又会有什么危险呢？是\Emph{那种}被此类快乐耗损、被恶性娇气所腐化的危险吗？然而，不朽的、常存的、无受苦的，又怎会被腐化呢？是\Emph{那种}被污秽鄙贱的行为所污染的危险吗？然而，没有形体实质的，又怎会被玷污呢？或者，哪里能让腐化停留，连让这腐化本身打上印记的地方都没有呢？\par

但再则，如果灵魂如\OriginalTerm{伊壁鸠鲁}{Epicurus}的教义所断言的那样，会临近死亡之门，在这种情况下，同样没有充分的理由去探寻哲学，即使哲学确实能使灵魂从一切污秽中得到清洁与净化。因为，倘若它们全都死去，甚至在身体内那生命特有的感觉也消亡失落了，那么抑制天生的欲望，将生活方式限制在狭窄的界限内，不顺从自己的倾向，不去做我们情欲所要求与催促的事，就不单是一个极大的错误，而且\Emph{显示出}愚蠢的盲目，因为当死亡之日来临，你从身体的束缚中解脱时，如此巨大的辛劳并不会带给你任何赏报。\par

31. 那么，灵魂的某种\OriginalTerm{中性特质}{neutral character}，以及未决和疑惑的性质，为哲学创造了空间，并找到了追求它的理由：也就是说，一个人因其犯下的恶行而充满恐惧；另一个人则怀有极大的希望，如果他不作恶，并在责任与正义中度日的话。因此，在有学问的人，以及具备卓越才能的\Emph{人}当中，对于灵魂的本性存在争执，有些人说它屈服于死亡，且不能承当\OriginalTerm{神圣实体}{divine substance}；而另一些人则\Emph{主张}它是不朽的，不会沉沦于死亡的力量之下。然而，这是由灵魂\Emph{的}中性特质的法则引起的：因为，一方面，有一些论证呈现给一方，由此发现灵魂能够受苦，且可灭；另一方面，对于另一方的反对者，也不乏论证，显示灵魂是神圣且不朽的。\par

32. 既然这些事如此，且我们被最伟大的师尊教导，灵魂离死亡的敞开之口不远；然而，只要他们竭力认识那\OriginalTerm{至高主宰}{Supreme Ruler}，因祂的恩惠和仁慈，他们的生命仍可延长——因为认识祂是一种生命的酵母和粘合剂，将否则会飞散的事物联结在一起——那么就让他们放下野蛮残忍的本性，回归较温和的方式，以便能为将给予的做好准备。有什么理由你们认为我们粗野和愚昧，如果我们因这些恐惧而降服并把自己交给天主我们的解救者呢？我们常常寻求疗法以治伤口与毒蛇的咬伤，并用\Place{普绪利}人或\Place{马尔西}人及其他小贩和骗子所售的薄片来防护自己；并且为了不受寒冷或酷热之苦，我们用焦虑而细致的勤勉提供房屋的遮盖和衣物。\par

33. 既然死亡的恐惧，即我们灵魂的毁灭，威胁着我们，我们何尝不是如常人所行，出于对我们将有益处的事的感觉，因为我们紧紧抓住\Emph{祂}，祂向我们保证祂将是我们的解救者脱离此危险，拥抱祂，并将我们的灵魂托付给祂照管，只要这种交换是对的？你们将你们灵魂的救恩寄托在自己身上，并且深信单凭你们自己的努力就能成为神；而我们，相反地，不因自身的软弱而对已怀有任何希望，因为我们看到我们的本性毫无力量，在一切争竞中被自己的情欲所胜。你们认为，一旦你们离世，脱离肉体肢体的捆锁，你们将找到翅膀，借此飞升天堂，翱翔于星辰之间。我们却远避这种僭越，并不认为到达高处的居所是我们能力所及之事，因为我们对此甚至没有把握，我们是否配得生命并脱离死亡的律法。你们设想无需他人的帮助，\Emph{你们}就能归回主人的宫殿，好似归自己的家，无人拦阻；但我们相反，既无任何期待以为此事可能成就，除非因那\OriginalTerm{万有之主}{Lord of all}\Emph{的}旨意，也不认为如此大的权能和许可赐给了任何人。\par

34. 既然事实如此，这有什么不公平，以至于我们在那种\Emph{你们所嘲笑的}轻信上被你们视为愚拙，而我们却看到你们既有类似的信念，也怀有相同的希望？如果我们因为对自己怀有此种希望而被认为应受嘲笑，那么同样的嘲笑也等待着你们，你们为自己声称拥有不死的希望。如果你们持有并遵循理性的途径，也请让我们分享其中。如果\Person{柏拉图}在\WorkTitle{斐德若篇}中，或此\Emph{哲学家的}行列中的另一位，曾向我们应许这些喜乐——即一条逃脱死亡的途径，或者能够提供这条途径，并带\Emph{我们}达到他所应许的目的，那么我们寻求敬奉那位我们从他期待如此大恩赐和恩惠的人，是合宜的。然而现在，既然\Person{基督}不仅应许了它，也以祂的德能——\Emph{这些德能}如此伟大——证明了它是可以实现的，那么如果我们俯伏敬拜祂的名和尊威，我们期待从祂那里\Emph{领受}\Emph{这两样福分}，以便立即逃脱受苦的死亡，并得享永生，我们所行的是什么怪事，又凭什么理由被指控为愚昧呢？\par

35. 但是，我的\Emph{对手们}说，如果灵魂是可死的且具有中性特质，它们如何能从其中性属性成为不朽？如果我们说我们不知道此事，只因为由\Emph{一位}比\Emph{我们}更大能者所说而相信，那么当我们相信那全能的君王没有难成、没有困难的事，并且对我们不可能的事，对祂并随祂的旨意是可能的时，我们的轻信何时会显得错误呢？因为有\Emph{什么}能抵挡祂的旨意，或者祂所愿意的，岂非\Emph{必须}成就？难道我们应根据我们自己的区分来推断什么是能做的或不能做的；难道我们不该考虑我们的理智与我们自己一样是可死的，且在至高者面前毫无价值？然而，哦，你们这些不相信灵魂具有中性特质，并且被置于生与死之间的中界线上的人，难道所有幻想所认为存在的一切，诸神、天使、\OriginalTerm{邪神}{daemons}，或无论它们的名字是什么，它们自身不也是具有中性特质，并在其将来之不定中易于变化的吗？因为如果我们都同意有一位万有之父，祂\Emph{是}唯一不朽且\OriginalTerm{非受生}{unbegotten}的，并且\Emph{如果}在祂之先绝找不到任何可被命名的，那么结果就是：所有这些被人的想象信以为神的，要么是由祂所生，要么是奉祂的命令而产生的。它们是受造且受生的吗？它们也在次序和时间上是后起的：如果在次序和时间上是后起的，它们必然有一个起源，并有出生和生命的开始；但那具有\Emph{进入}生命并在其初始阶段拥有开始之物，必然也会有终结。\par

36. 但据说神是不死的。那么，不是出于本性，而是由于天主他们的父的善意和恩惠。因此，正如不死的恩赐是天主给予\Emph{那些确实被造者}的礼物，同样，祂也会惠然赐予灵魂永生，尽管残酷的死亡似乎能将它们割断，并将其彻底消灭，使其不复存在。神圣的\Person{柏拉图}，他的许多思想配得上天主，而非俗人的见解，在他那篇题为\WorkTitle{蒂迈欧篇}的讨论和论述中说，神和世界在本性上是可朽坏的，绝不可能免于死亡，但他们的存在靠着天主——\Emph{他们的}君王和主宰——的旨意得以永远维持；因为\Emph{即使}是被最牢固的纽带恰如其分地连接和捆绑在一起的，也\Emph{唯有}靠着天主的美善才得以保存；除了那位将\Emph{他们的}元素捆绑在一起者，再没有谁能在必要时将它们分解，并发出命令保存其存在。如果事情确实如此，而不宜作他想或相信别的，那你们为何惊奇我们说灵魂在特性上是\OriginalTerm{中性的}{neutral}，既然柏拉图说，就连神灵也是如此，但他们的生命靠着天主的恩宠得以维持，没有间断或终结？因为，如果你们出于偶然未曾知道，且因其新颖而先前不知，那么\Emph{现在}，虽然已晚，但仍可接受并学习那知晓并已将其显明的——\Person{基督}：灵魂并非至高主宰的儿女，也并非在由祂创造之后才开始具有自我意识，并以其特有的身份为人所言说；而是另有其父，在品秩和权能上远次于至高主宰，然而却属于祂的宫廷，并因其高贵而崇高的长子权而卓著。\par

37. 但是，如果灵魂如所说的那样，是主的儿女，由至高能力所生，那就什么也不缺，足以使它们完美，\Emph{因为它们原本会}出生时就具有最完美的优长：它们都会心意相同，\Emph{都}和谐一致；它们会常居王宫；不会绕过那些福乐之境——在那里它们已学习并铭记了最高尚的教导——而贸然寻求地上的这些区域，以便生活在幽暗的肉身中，处于痰和血之间，置身于这些污秽的囊袋和最令人作呕的尿器之中。然而，\Emph{有人会反驳说}，这些部分也需要住满，因此全能天主派遣灵魂到这里来，仿佛\Emph{建立}一些殖民地。那么，人类对世界有何用处，因何而必需，以免让人相信他们命定住在这里、成为地上肉身的租客而毫无目的？\Emph{我的对手说}，它们有份于成全这一庞大整体的完备，没有它们的加添，整个宇宙便不完整而有缺陷。那么又如何呢？如果没有人类，世界是否会停止运行其功能？星辰是否会不再运转变化？是否会没有夏季和冬季？狂风的呼啸是否会平息？从聚拢并\Emph{高悬}空中的云层，是否会不降下阵雨于大地以缓解旱情？而现在，万物必须按各自轨道运行，不会放弃遵循自然所定的安排，即使世上听不到人类的名字，这大地也静默如无人居住的荒原。那么，怎能声称必须给这些区域安排一位居民呢？因为很明显，人类对成全世界毫无\Emph{助益}，他的一切努力总是为了自己的私利，从未停止追求自己的利益。\par

38. 首先，且从重要之事说起：世间有最强大的君王，对世界何益之有？有暴君、贵族\Emph{和}其他无数显赫的权势，又有何益？有最富征战经验的将领，精通攻城略地；有稳固而全然不可战胜的士兵，在骑兵战中或徒步肉搏中无往不利——又有何益？有演说家、文法家、诗人、作家、逻辑学家、音乐家、芭蕾舞者、哑剧演员、演员、歌手、号手、长笛与芦笛演奏者，又有何益？有跑步者、拳击手、战车御者、跳跃者、踩高跷者、走绳索者、杂耍艺人，又有何益？有卖咸鱼的、腌鱼的、鱼贩、香料商、金匠、捕鸟者、编织簸箕和灯心草篮的织工，又有何益？有漂洗工、毛纺工、刺绣工、厨师、糖食师傅、卖骡子的、皮条客、屠夫、妓女，又有何益？有其他各类商贩，又有何益？\Emph{其他各类}教师和技艺——要一一列举，一生都嫌\Emph{太}短——对世界的计划与构成又贡献了什么呢，竟叫我们相信，没有人类，世界便不能建立，若缺少这悲惨而无用的生灵的努力，世界便无法达致完备？\par

39. 但也许，\Emph{有人会强调}，世界的主宰为了这个目的从祂自己遣发灵魂到此——如此说法对一个凡人来说实在是狂妄——让他们，那些曾与祂同享神性者，本不接触肉身与尘世局限，却应被埋葬于人的种子之中，从母胎中生出，爆发出并不断发出最愚蠢的哭啼，吮吸乳房，用自己的污秽弄脏并涂抹自身，然后因受惊乳母的摇晃与拨浪鼓的声音而安静下来。难道祂为此缘故遣发灵魂\Emph{到此}，让那些刚才还真诚坦荡、德行无瑕者，既成了人，竟学会伪装、掩饰、说谎、欺诈、欺骗、以谄媚者的卑贱去诱陷；心里藏着一套，脸上显出另一套；以诡计诱骗，蒙蔽无知者；借由恶念\Emph{启发}的无量计谋寻求毒害；并被塑造成以欺诈的多变来适应环境？难道祂为此遣发灵魂，让那些\Emph{直到那时}还生活在安宁、不受扰乱的平静中者，竟在肉身中找到理由，变得凶暴蛮横，怀藏仇恨与敌意，彼此争战，征服并倾覆邦国；给自己加上奴役之轭，且委身于其中；最终，被置于彼此权下，改变了他们出生时的处境？难道祂为此遣发灵魂，使他们被弄得不顾真理，忘记天主是何等，竟向不能移动的偶像祈求；将木、铜、石块当作超人的神明；用宰杀牲畜的血向他们求助；不提祂自己：更有甚者，让他们当中的一些人竟怀疑自己的存在，甚或完全否认任何事物存在？难道祂为此遣发灵魂，让那些在自己居所曾同心合意、在理智和知识上平等者，在穿上必死之形后，竟因意见分歧而分裂；对何为正义、有益、正当持不同见解；为所欲之事与所憎之物而争执；以不同方式界定至善与大恶；在寻求认识事物真理时，竟被其晦暗所阻碍；仿佛眼目被夺，一无所见清晰，偏离真理，被引至幻想的不定小径？\par

40. 难道祂为此遣发灵魂\Emph{到此}，让其他受造物由自然生长、无需播种的所喂养，且不为自己寻求房屋或衣物的庇护和遮盖，他们自己却在悲苦的必须中，以极大开销和无尽辛劳为自己建造房屋，为肢体预备遮盖，为日常需要制作各种\Emph{种类}的家具，向无言的受造物借取帮助以补软弱；对大地施以暴力，为使它不生长自己的草本，而是产出所需的果实；并且在耗尽全部力量征服大地之后，竟被迫因枯萎、冰雹、干旱而失去他们劳苦所寄的希望；最终为饥饿所迫，投奔人的尸体；及至释放，竟被消耗的疾病与人的形体分离？难道祂为此，让那些曾与祂同居、从未对财产有过任何贪恋者，竟变成极度贪婪，并因无法满足的渴求而燃起急切的占有欲；让他们挖掘高山，将大地未知的内部转变成物资和不同\Emph{的}用途；冒着生命危险强行通过远方的民族，在交换货物时总是争取高价\Emph{卖出}、低价\Emph{买进}，收取贪婪而过高的利息，增加他们不眠之夜，\Emph{耗费在}盘算从可怜人的生命之血中榨取来的成千上万；让他们不断扩展财产的范围，即使把整个省份变成一座庄园，还要为了一棵树、为了\Emph{一条}犁沟而用诉讼使法庭厌烦；让他们恶毒地恨自己的朋友和兄弟？\par

41. 难道祂为此遣发灵魂，让那些不久前还温文尔雅、不懂\Emph{何为}被烈怒激情所动者，竟为自己建造市场与圆形剧场，这流血的场所与公开邪恶之地；在一处，他们看见人被野兽吞噬和撕碎，且自己毫无理由地杀害他人，只为取悦观众，并在这些邪恶行径发生的日子里尽情享乐，以节庆的欢乐度假；而在另一处，他们又将可怜动物的肉撕开，像狗和秃鹫一样，有的抢一块，有的抢另一块，用牙齿将\Emph{之}磨碎，送入那永不餍足的口腹，并且，在如此凶猛而野蛮的面孔环绕下，那些因贫穷窘迫而被隔绝于此等盛宴之外的人，应哀叹自己的命运；他们的生命竟在如此野蛮的行径弄脏他们的口与脸时，还显得快乐而繁荣？难道祂为此遣发灵魂，让他们忘却自己\Emph{身为}神性的重要与尊严，竟以纯洁为代价去获取宝石、珍贵石头、珍珠；用它们缠绕颈项，刺穿耳尖，束住前额，寻求化妆品来妆饰身体，用散沫花染黑眼睛；并且，虽具人形，却不羞于用卷发钳卷曲头发，使身体皮肤光滑，露膝行走，并伴以其他各种\Emph{放荡}，既丢下男性的气力，又在柔弱中变得如同妇女的习惯与奢靡？\par

祂派遣灵魂，就是为了这个吗？为了让一些人侵扰大路与小径，另一些人诱骗不谨慎的人，伪造假的遗嘱，准备毒药；让他们夜间破门而入，收买\Emph{奴隶}，偷盗抢夺，不行正直，背信弃义，出卖\Emph{他们的信用}；让他们发明精致美味以悦口舌；让他们烹饪禽肉时，懂得接取滴下的油脂；让他们制作脆饼、香肠、肉馅、小块美食、\OriginalTerm{卢卡尼亚香肠}{Lucanian sausages}，以及母猪乳房和冰镇布丁？祂派遣灵魂，就是为了这个吗？为了让那神圣而庄严的族类在此练习歌唱和吹奏；让他们鼓起腮颊吹起笛子；让他们领头唱起污秽的歌曲，奏响嘈杂的响板声，由此引领另一群灵魂放纵地投身于笨拙的动作，舞之蹈之，围成圆圈，最后高抬臀胯，跟着腰肢的颤巍巍的扭动漂浮而过？\par

祂派遣灵魂，就是为了这个吗？让它们在男人身上变得不洁，在女人身上成为娼妓、演奏\OriginalTerm{三角琴}{triangle}与\OriginalTerm{索尔特里琴}{psaltery}的人；让它们为了金钱出卖身体，委身于所有人的情欲，等候在青楼之中，出没于花街柳巷，随时准备屈从任何事情，甚至预备对自己的嘴施暴？\par

你们这\OriginalTerm{至高神明}{Supreme Deity}的后裔与子孙，有何话说？这些灵魂，既然智慧，且由第一因所生，会熟悉如此种种卑鄙、罪恶与不良的情感吗？它们被命定居于此，披上人体的衣裳，是为了从事、为了实行这些邪恶的\Emph{行径}，并且屡屡为之吗？哪一个拥有理性的人会认为世界是因它们而建立的，而不是把它设为一种场所和居所，好让各式各样的邪恶日日上演，让一切恶行得以实施——阴谋、欺骗、诈伪、贪婪、抢劫、暴力、不虔、\Emph{所有的}放肆、淫秽、下流、可耻的事，\Emph{以及}世人在全地上怀着罪恶意念谋划，为彼此毁灭而设的一切罪恶？\par

然而，你们会说，它们是自行前来的，并非受其主的差遣。那么，全能造物主在哪里呢？祂那王权和崇高地位的权威在哪里，竟不阻止它们离开，也不禁止它们陷入危险的享乐？因为，如果祂知道它们会因更换处所而变得卑劣——而且，作为万有的安排者，祂必定知道——或者会有任何外来的事物触及它们，致使它们遗忘自身的伟大和道德的尊荣——我愿千万次恳求祂赦免\Emph{我的言辞}——这所有的原因无非是祂自己，因为祂允许它们有游离的自由，而祂预知它们不会持守其天真无罪的状态；因此，事情就弄成这样：无论它们是自行前来，还是服从祂的命令，都无关紧要了，因为祂没有阻止本应阻止的事，祂的不作为使罪咎归于自己，并且在\Emph{事情发生}之前，祂因疏于阻止其\Emph{行动}而容许了它。\par

但是，让这种怪异而不敬的空想\Emph{远离我们}：以为全能的天主，那造化与塑造者、伟大而不可见之万物的创造者，竟然被相信生就了如此反复无常、毫无庄重、坚定与稳固的灵魂，它们趋向恶习，偏向一切罪行；祂明知它们如此这般，却仍吩咐它们进入身体，囚禁其中，它们日复一日暴露在命运的暴风骤雨之下，时而做卑贱之事，时而遭受淫秽的对待；以致它们可能因沉船、事故、毁灭性的大火而灭亡；贫困压迫一些人，乞讨生活压迫另一些人；一些人被野兽撕碎，另一些人死于毒虫的螫咬；一些人走路跛足，另一些人失明，还有一些人关节蜷曲僵硬；总之，它们该当遭受一切疾病，这可怜而堪悲的人类，带着由各种苦难所生的极度痛苦，承受着这一切；然后，它们忘记自己拥有同一个本源、同一位父亲和首领，竟动摇和侵犯亲属关系的权谊，推翻自己的城池，像敌人一样蹂躏土地，奴役自由人，强暴少女和他人之妻，彼此憎恨，嫉妒他人的幸福与好运；更有甚者，人人诽谤、挑剔，以毒牙彼此撕扯碎骨。\par

然而，为了一再重申同样的话，让这种如此怪异而不敬的信念\Emph{远离我们}：以为那保存万有的天主，万德的根源与仁爱的至高者，且用人的赞美来颂扬，祂最智慧、最公义，使万有达致完美并恒久如此，竟然创造任何不完全、不尽然正确的事物，或成为任何存有之痛苦或危险的起因，或者安排、命令并指定，人的一生所度过和从事的那些行为，正是出于祂的安排而流淌出来。这些事与天主不相称，削弱祂伟大的力量；与其相信天主是这些事的制造者，倒不如说，凡想象人是出自天主的人，就犯了亵圣不敬的罪——这\Emph{人}，是何等悲惨、可怜的生灵，他哀叹自己的存在，憎恶、痛惜自身的状态，并且明白自己之所以产生，无非为了让邪恶拥有借以蔓延的渠道，好叫那常有的悲惨者，以其极度的痛苦，取悦于某种敌对人类的、不可见的残酷能力。\par

47. 但是，你说：\Quote{如果天主不是灵魂的\OriginalTerm{父和父亲}{parent and father of souls}，他们是由何父所生，又是如何被产生的呢？}如果你愿意听不加虚饰、不堆砌空洞华丽辞藻的陈述，我们也承认我们对此无知，并不知晓；并且我们认定，认识如此重大的事，不仅超越我们软弱脆弱的限度，\Emph{甚至也超越}世界内一切能力——就是那些篡夺人心中神祇之位的能力。但是，难道因为我们否认灵魂属于天主，我们就必须指出它们属于谁吗？这绝不必然得出；因为如果我们否认苍蝇、甲虫、臭虫、睡鼠、象鼻虫和蛾子是由\OriginalTerm{全能君王}{Almighty King}所造，我们并不因此而需要说出谁造了它们、形成了它们；因为不招致任何指责，我们完全可以不知道实际上谁赋予它们存在，却断言这些受造物并非由\OriginalTerm{至高天主}{Supreme Deity}所造——它们如此无用、如此多余、如此毫无目的，不仅如此，有时甚至有害，造成无法避免的伤害。\par

48. 同样，当我们否认灵魂是\OriginalTerm{至高天主}{God Supreme}的\OriginalTerm{后裔}{offspring of God Supreme}，并不必然意味着我们必须宣告他们由何父所出，以及由何原因产生。因为有谁禁止我们或者不知他们出自何处、来于何方，或者明知他们并非天主的后裔呢？你要问，以何种方法，何种方式？因为，正如我们常说的，最为真实而确定的是：至高者不会造成、制作、决定任何事，除非那事是合宜且应当成就的；除非是完整而齐全，在其纯全中彻底完美的。但是，进一步来看，我们看到人，也就是这些灵魂本身——因为人除了被束缚于肉体的灵魂，还能是什么？——他们自己无数次因堕入罪恶而歪曲，这表明他们绝非出身\OriginalTerm{贵族血统}{patrician race}，而是源于卑微的家族。因为我们看到有些人严苛、邪恶、放肆、鲁莽、不顾后果、盲目、虚伪、掩饰、说谎、骄傲、傲慢、贪婪、吝啬、好色、反复无常、软弱，且不能遵守自己的规诫；然而，如果他们有\OriginalTerm{原初的善}{original goodness}在维护，并且他们高贵的血统可追溯至宇宙之元首，他们必然不会如此。\par

49. 但是，你会说，世界上也有\OriginalTerm{善人}{good men}——明智、正直、品行无瑕且极其纯洁。我们并不质问是否曾经有过这样的人，在他们身上，所谈及的纯全毫无缺陷。即使他们是极其可敬的\Emph{人}，配得赞颂，已达完美的极顶，一生从未动摇沉沦於罪恶，我们仍要请你告诉我们，他们有多少，或者曾经有多少，以便我们从其数目判断，所作的比较是否公正均衡。一个、两个、三个、四个、十个、二十个、一百个，他们至少\Emph{确实}数目有限，或许屈指可数。但合乎情理的是，应当以全体来评定和衡量\OriginalTerm{人类}{human race}，而不应仅凭极少数善人，更要包括所有其余的人。因为部分是整体的部分，而非整体在部分之中；那为整体的，应吸引部分归向自己，而非整体被引向部分。试想，如果你说，一个人四肢全废，痛苦哀号，却因一个指甲不痛，就说他十分健康，那会怎样？或者说，大地是金子做的，因为一个小丘上有几粒细沙，溶解后能出产金子，聚合成块时就令人惊奇？整个质体才显示一种元素的特性，而非微细如气的粒子；同样，如果你将少许不太咸的水滴倒入海中，海水不会立刻变甜，因为那少量之物被其浩大的体积吞没；而且，不仅应视之为微不足道，\Emph{甚至}毫无作用，因为它散布全体，便在那浩大水体之无限广袤中消散、隔绝了。\par

50. 你说人类中有\OriginalTerm{善人}{good men}；或许，若将之与大恶者相比，我们可能容易被引导而相信有。请问，他们是谁？告诉我们。我想，是那些哲学家们，他们断言唯有自己最具智慧，并因此名之义而心高气傲——正是他们，日日与自己的偏情搏斗，奋力以较善品质的持续对抗，从心中驱逐、根除根深蒂固的偏情；他们为使自己在任何机会的暗示下不致陷于邪恶，遂躲避财富和遗产，好除去自己的绊脚石；但在这样做时，并对此忧心忡忡，他们极为清楚地表明，\Emph{他们的}灵魂因自身的软弱，易于倾向堕落。然而，依我们看来，\OriginalTerm{本性善者}{that which is good naturally}，无需纠正或责斥；不仅如此，若每一类本性能持守自身的纯全，它甚至不应知道恶为何物，因为两个相反者不能互相植入，平等不能存于不平等中，甘甜不能存于苦涩中。所以，那奋力纠正自己生来偏邪\OriginalTerm{堕落倾向}{inborn depravity}的人，极为清楚地表明他是不完美、应受谴责的，纵使他可能以全副热忱和坚定不移去努力。\par

51. 但是你们嘲笑我们的回答，因为，虽然我们否认灵魂有高贵的出身，却反过来并没有说它们是由什么原因和开端产生的。但是，对某事无知，或者坦率地承认你不知道你所无知的，这算什么罪行呢？或者，在你看来，是那种对自己不了解某个晦涩主题就不自以为有知识的人更值得嘲笑，还是那种认为自己对超越人类知识、且被笼罩在黑暗晦涩之中的事物最清楚明白的人更值得嘲笑呢？如果万物的本性被彻底思考，你们所处的立场也类似于你们所责难我们的。因为你们并没有说出任何\Emph{已被确证}、并在真理之光中明确确立的东西，因为你们说灵魂是从至高天主本身降落，并进入人的形体。因为你们是推测，并没有察觉到\Emph{这一点}；是臆断，并没有实际知道\Emph{它}；因为如果知道是指心中持守你自己亲眼见过或知道的事物，那么你们所肯定的那些事，没有一件你能说你曾见过——即灵魂从上面的居所和区域降落。因此你们是在运用推测，不是信赖明确的信息。但是，推测若不是对事物的可疑想象、以及将心思指向某种无法触及的东西，那又是什么呢？所以，推测者并不理解，也不行走在知识的光明中。然而，如果这在适当且极有智慧的判断者看来是真实且确定的，那么你们所信赖的推测，也必然应被视为\Emph{表明了你们的无知}。\par

52. 然而，为了不让你们以为只有你们自己才能运用推测和臆断，我们也能够提出它们，因为你们的问题对双方都适用。你们说，人是从哪里来的；这些人拥有怎样的灵魂，它们是从哪里来的？\Emph{我们反问}，大象、公牛、牡鹿、骡子、驴是从哪里来的？狮子、马、狗、狼、豹是从哪里来的？这些生灵的灵魂是什么，是从哪里来的？因为既不可信，那些灵魂是从柏拉图式的杯中出来的——那是\Person{蒂迈欧}所准备和混合的；也不应相信，\Emph{那些}蝗虫、老鼠、鼩鼱、蟑螂、青蛙、蜈蚣之所以被赋予生命并存活，是因为它们在元素本身中有产生的根源和来源，如果\Emph{在这些元素中}有秘密且极少为人所知的手段，产生居住于其中的生物。因为我们看到，一些智慧者说大地是人类的母亲，另一些则加上水，还有一些更添加空气之息，然而有些人说太阳是人类的塑造者，人类因他的光线而赋有生命，并被生命的搏动所充满。如果并非这些，而是有某种别的原因、别的方法、别的能力，总之，一种我们闻所未闻、以名称亦未知晓的，它也许塑造了人类，并将它与既存事物联系起来；难道人不可以是以这种方式出现的，其出生的原因并不追溯到至高天主吗？凭什么理由我们认为，伟大的\Person{柏拉图}——一个在智慧上虔敬而谨慎的\Emph{人}——会把塑造人类的工作从至高的天主那里撤回，并转移给某些次等的\Emph{神灵}，并且他不愿由他曾用来制造宇宙灵魂的那种纯粹混合来形成人类的灵魂，难道不是因为他认为塑造人类与天主不相称，塑造一个脆弱的存有不配于祂的伟大和卓绝吗？\par

53. 既然如此，我们相信人的灵魂属于中性的特质，因为它们是由次等的存有所产生的，隶属于死亡的法律，\Emph{并且}力量微弱，且是会朽坏的；\Emph{并且他们}被赋予不死不灭，只要他们将如此伟大恩赐的希望寄托于至高天主，只有祂有权力赐予这样的\Emph{福佑}，通过除去朽坏。但\Emph{你们说}，我们相信这一点是愚蠢的。那\Emph{与你们}有何相干？\Emph{我们如此相信，显得}极其荒谬、傻气。如果我们相信，当我们离开身体时，全能的天主会照顾我们，并从地狱的口中，正如所说的，拯救我们，那么我们在什么事上伤害了你们，或对你们行了什么错事或施加了什么损害呢？\par

54. 那么，有人会说，难道没有天主的旨意，任何事物能够被造成吗？我们必须仔细考虑，并以不小的努力去检验，以免我们认为通过这样的问题是在尊敬天主，却陷入了相反的罪过，损害了祂的至高威严。何以见得？因为，如果万物都是按照祂的旨意发生的，世界上的任何事物都不能违背祂的意愿而成功或失败，那么必然的结果就是，一切恶事也都是由祂的旨意而产生的。但是，反过来，如果我们选择说祂与恶无涉，也不产生恶，不把最邪恶的行为的原因追溯到祂那里，那么最坏的事情就会开始看起来要么是违背祂的旨意而做的，要么，一个可怕的说法是，是祂不知道的，\Emph{却}是无知且不觉察的。但是，更进一步，如果我们选择说没有恶存在，正如我们发现某些人所相信和主张的那样，那么所有种族和所有民族都会一起大声反对\Emph{我们}，向我们展示他们的苦难，以及人类每时每刻所遭受的各种各样的危险。然后他们会问我们：如果没有恶，为什么你们要避免某些行为和举动呢？为什么你们不做出急切渴望所要求或索要的任何事呢？最后，为什么你们要以可怕的法律为有罪者设定惩罚呢？因为还有什么比断言没有恶，而\Emph{同时}又杀害并谴责犯过错者，好像他们是邪恶的一样，更愚昧的怪异行为呢？\par

55. 然而，当我们被折服，承认这些事确实存在，并明确承认人间事务充满它们时，他们接下来会问：那么，为何全能的天主不除去这些恶，反而容其存在，且历世历代永无止息地延续？若我们已认识天主，至高的统治者，并已决意不徘徊于不虔与疯狂的臆测迷宫中，我们就必须回答：我们并不知晓这些事，也从未寻求、竭力去了解那非我们\Emph{所具}能力所能领悟的；我们甚至认为，宁可保持无知与匮乏，也比说\Quote{没有天主，无物造成}更好，以免使人以为，依祂的旨意，祂同时既是恶的源头，又是无数痛苦的缘由。那么，你或许要问，这一切恶从何而来？智者说：来自\OriginalTerm{元素}{elements}及其\OriginalTerm{相异性}{dissimilarity}；然而，无感觉、无判断之物怎能被视为邪恶或罪过？若要达成某种结果，而采用后来变得极坏而有害的东西，那施行者岂非更是邪恶、有罪？——这应由那些断言者去思量。那么，我们怎么说？从何而来？我们没有必要回答，因为无论我们是能说出\Emph{恶从何来}，或是力有不逮、无法言明，依我们看来，这都无关紧要；我们也不认为知道或不知此事有何重大意义，我们只满足于确立一事——从至高者天主那里不出任何有害与败坏。我们确信此点，我们明了此点，我们立足于这一知识与学问的真理——祂所造者无他，唯为万有之福祉，为可喜悦者，为充满爱、喜乐与欢欣者，为拥有无尽不朽之欢愉者；人人皆可于祈祷中求其临己，并视若无此，生命便为有害与致命。\par

56. 至于其他通常在探询与讨论中反复论及之事——它们源于何种源头，或由谁产生——我们既不力求知晓，也不留意探究或考察：我们将一切事物归诸其自身的原因，并不认为它们与我们祈望临己之事有何关联与牵涉。因为有才智的人，出于好辩之心，何事不敢推翻、破坏？即便他们试图否定的对象无可指摘、明显清晰，且显然带有真理的印记。再者，凭借看似合理的论据，他们又何事不能坚持？即便那些事非常明显不真，且是明白无误的虚妄？因为当人说服自己某事存在或不存在时，他喜欢肯定自己的所思，并显露比他人更精妙的心思，特别是当所讨论的主题逸出常轨，且本性深奥隐晦时。有的智者认为世界并非受造，也永不会消亡；有的认为世界虽说是受造而成的，却是不朽的；而第三派则选择说，世界既是受造而成的，也必将如其它事物一样消亡。然而，这三种意见中只有一种必为真实，他们却全都寻得论据，既可维护自己的学说，又能削弱、推翻他人的主张。有些人教导并宣称，这同一个\Emph{世界}是由四种\OriginalTerm{元素}{elements}构成，另一些人说是两种，第三派说是一种；有些人说\Emph{其构成}并非由这些，而是由原子形成，且为其原初之始。既然这些意见中只有一种真实，却没有一种是确定的，此处亦然，论据同样呈现于各方，使他们既能确证己说之真，又能指明他人意见中之谬误。同样，有些人彻底否认\OriginalTerm{神明}{gods}的存在；另有人说他们陷于疑惑，不知神明是否存在某处；然而，还有人说他们确实存在，但不干预人间事务；甚或有人认为，他们既参与人间之事，又引导世事的进程。\par

57. 情形既然如此，而所有这些意见中只能有一个为真，他们却全都各逞论辩，彼此争胜——且无一人缺乏看似有理之说，无论用以肯定己见，或反对他人意见。关于\OriginalTerm{灵魂}{soul}的状况，亦以完全相同的方式被讨论。因为有人认为灵魂既是\OriginalTerm{不朽}{immortal}的，并且在我们尘世生命终结后依然存留；那人则相信灵魂不存留，与肉身一同消亡：然而，又有一种意见认为，灵魂并非立刻遭受任何事，而是在人的\Emph{形式}被撇弃之后，允许它们再存活一段短暂时日，随后来到死亡的权下。所有这些意见既不能皆真，但所有\Emph{持此意见者}都竭力以强有力且极具分量的论据支持其立场，以致你无法找出任何看似为假之处，尽管你处处看见所言之事全然彼此冲突、互不相容；倘若人的好奇心能达致任何确定，或某一方认为已真正发现之事，竟能获得所有其他人的赞同佐证，那么上述现象断然不会发生。因此，提出某事如同你确知一般，或欲断言你知晓某事——即使它可能是真的，你却看见它可被驳斥——完全是徒劳、无益之举；或接受那并非真实、却如狂人呓语般被提出的东西为真，亦属虚妄。此理诚然，因为我们不以神圣的方法，而以人的方法去衡量、揣测神圣之事；正如我们设想某事应当如何造成，便断言其必然如此。\par

58. 那么，难道只有我们无知吗？难道只有我们不知道谁是创造者，谁是灵魂的形成者，是什么原因造了人，灾祸从何处爆发，或者为何至高主宰允许它们既存在又被施行，却不将其从世界驱逐？你们果真确实查明并知晓了这些事情中的任何一件吗？如果你们愿意放下大胆的臆测，你们是否能揭示并说明我们现在居住的这个世界是在某个时间被创造或建基的？如果它被建基和造成，那么请问是靠什么样的工程，或为了什么目的呢？你们能否提出并指明它为何不保持固定不动，却永远被带著作圆周运动的原因？它是凭自己的意愿和选择旋转，还是受某种力量的影响而转动？此外，它被安置并旋转于其中的位置和空间是怎样的：是无限的、有限的、空洞的，还是坚实的？它是被一根两端支在轴座上的轴所支撑，还是凭借自身的力量，并由其内在的精神自我维持？如果被问及，你们能否澄清，并极巧妙地表明，是什么使雪展开成羽片状？是什么原因导致黎明未从西方升起，并将其光明遮蔽于东方？太阳又是如何借同一种影响产生如此不同，甚至相反的结果？月亮是什么，星辰是什么？为什么一方面月亮不能保持同一形状，或者为什么这些火的微粒应该遍布整个世界是正确且必要的？为什么它们有的小，有的大而较大——这些发出暗淡的光，那些则闪耀着更明亮的光辉？\par

59. 如果我们所乐于知晓的是可达的，如果这样的知识对所有人开放，那么请向我们说明，告诉我们阵雨是如何产生的，以致水竟能悬于上方区域和半空中，尽管它本性易于滑落，且如此倾向于流动和向下流淌。我说，请解释并说明，是什么使冰雹\Emph{在空中}旋转，是什么使雨点一滴滴落下，是什么铺展了雨、羽片般的雪和片片闪电；风从何处升起，它是什么；为何制定了季节的变化，而本来可以规定只有一种季节、一种气候，使世界的完整性无所缺失？是什么原因，什么理由，使海水是咸的；或者使陆上的水有的甜，有的苦或冷？人体内部各部分是由何种材料形成并构建成坚固的？它们的骨骼是由什么变得坚实的？是什么使肠子和血管形状像管子，且容易通过？为什么，既然用几只眼睛给我们光明以\Emph{防范}失明的危险会更好，我们却被限制为两只？无限且不可胜数的怪物和蛇类，被形成或产生出来，是为了什么目的？猫头鹰、隼、鹰在世界中起什么作用？其他鸟类和有翼生物呢？各种\Emph{不同种类}的蚂蚁和蠕虫，以各种方式成为祸害，又有什么目的？跳蚤、纠缠不休的苍蝇、蜘蛛、鼩鼱及其他鼠类、水蛭、水蜘蛛呢？荆棘、蒺藜、野燕麦、稗子呢？各种药草或灌木的种子，有的闻起来芬芳，有的则气味难闻，又是为什么？不仅如此，如果你们认为任何事物能被知悉或领悟，那么请问小麦是什么——斯佩耳特小麦、大麦、小米、鹰嘴豆、蚕豆、小扁豆、甜瓜、莳萝、葱、韭葱、洋葱？因为\Emph{即使}它们对你们有用，并被列入不同种类的食物，要知道它们每一个是什么，为何以这样的形状形成，也并非轻易之事；\Emph{是否}有必要它们不应该具有与各自现有的味道、气味和颜色不同的其他味道、气味和颜色，或者它们是否也能接受其他的；还有，这些事物本身是什么——我的意思是味道等等；\Emph{以及}它们从哪些关系中得出其特性的差异。你们说，从元素，从事物的最初根源而来。那么，元素是苦的还是甜的呢？它们有任何气味或恶臭吗，使我们相信，通过它们的结合，特性被植入其产物中，从而产生甜味，或某种使感官不快的东西？\par

60. 那么，既然你们自己也不理解如此众多且如此重要的事物的起源、原因和理据，既无法说出也无法解释已造成的事物为何这样而非\Emph{别样}，你们为何要攻击和抨击我们的畏怯？我们坦言自己无法知悉那不可知之事，也无心去探寻和追问那些明明不可能理解的事物——即使人的猜测能扩展到千万颗心灵之中，也不可能。因此，\OriginalTerm{神圣的基督}{Christ the divine}——尽管你们不愿承认——我再说一遍，\OriginalTerm{神圣的基督}{Christ the divine}，（这话必须常被说出，好让不信者的耳朵震破撕裂），祂奉\OriginalTerm{至高天主}{Supreme God}的命令，以人的形象说话，因为祂知道人生来盲目，根本无法把握真理，也无法把摆在他们眼前的事视为确凿无疑，并且为了自己的猜测，毫不犹豫地提出和引发引起许多纷争的问题——祂吩咐我们放弃和漠视你们所说的这一切，不要把心思浪费在远非我们知识所能及的事物上，而要尽可能以全部心思和灵魂寻求\OriginalTerm{宇宙之主}{Lord of the universe}；要超脱这些事物，把我们尚在犹豫转向何处的心交给祂；要时常记念祂；虽然任何想象都无法把祂的真貌呈现出来，但仍要对祂形成某种模糊的概念。因为\Emph{基督说过}，在所有被含混地视为神圣和属神的事物中，唯有祂是无可置疑的，唯有祂是真实的，唯有祂是那种只有狂妄鲁钝的疯子才会怀疑的；尽管你除此之外一无所知，但认识祂就足够了；如果你通过认识而确实与天主，即世界之首，有了关联，你就获得了那真实而最重要的知识。\par

61. 祂说，这与你们何干？你们竟去细察和追问：谁制造了人？灵魂的起源是什么？是谁谋划了种种恶事的原因？太阳是否比地球大，还是只有一足之宽？月亮是以借来的光亮照耀，还是本有的光明——知道这些既无益处，不知道也无损失。把这些事留给天主吧，让祂知晓那是什么、为何以及从何而来；是否必须发生或不发生；某物是否一直存在，还是最初被造的；是否应被消灭或保存，被销毁、毁灭，或在新活力中恢复。你的理性不允许你卷入这般问题，徒劳地为那些远不可及的事费心劳神。你们的利益正岌岌可危——我指的是你们灵魂的救恩；除非你们致力于寻求认识\OriginalTerm{至高天主}{Supreme God}，否则当你们摆脱肉身的锁链时，残酷的死亡等候着你们，它不是带来突然的消灭，而是以严峻而漫长惩罚的痛苦来摧毁。\par

62. 不要被一些无知而又极其狂妄的冒充者所说的虚妄希望所欺骗或迷惑，他们说他们从天主所生，不受命运之律的支配；只要过有节制的生活，祂的宫殿就为他们敞开，死后作为人，他们便毫无阻碍地回归，如同回到父家；也不要\Emph{受}\OriginalTerm{玛哥}{Magi}人所声称的\Emph{此事}蒙骗，他们说自己拥有代祷的祈祷，借此赢得了某些神力，为那些奋力登天的人铺平道路；也不要\Emph{受}\Place{伊特鲁里亚}在\WorkTitle{阿刻戎之书}中所提出的\Emph{此事}蒙骗，即只要向某些神祇奉献某些动物的血，灵魂就能成为神圣，并从死亡的律法中获得解脱。这些都是空虚的迷惑，只会激发虚妄的渴望。惟有\OriginalTerm{全能者天主}{Almighty God}才能保存灵魂；除祂之外，无人能赐予长寿，并赋予永不死亡的精神，惟有那位独一不死、永恒常在、不受任何时限约束的祂才能。因为既然所有神祇——无论是真实的，还是仅凭传闻和猜测而被称为神的——都是因祂的善意和白白的恩赐而得以不死和永恒，那么，那些自身所有乃是得自另一更大能力的恩赐者，又怎能给予他人呢？让\Place{伊特鲁里亚}随意献祭吧，让智者禁绝生活的一切享乐吧，让\OriginalTerm{玛哥}{Magi}人安抚和软化所有\Emph{较小的}神力吧，\Emph{然而}，除非灵魂从\OriginalTerm{万有之主}{Lord of all things}那里领受了理性所要求的，并且是奉\Emph{祂的}命令而领受的，否则当它开始感受到死亡的临近时，必将为自己已沦为笑柄而深深悔恨。\par

63. 但是，我的对手们说，如果基督是天主为这个目的而派遣的，即解救不幸的灵魂免于毁灭和灭亡，那么在祂未降临以前就已处于必死状态而逝去的先前世代，又犯了什么罪呢？那么，你们能知道这些远古之人的灵魂后来怎样了吗？它们是否也以某种方式得到了援助、供给和照顾呢？我再说，你们能知道那些本可通过基督的教导而了解的事吗？即：自从人类开始在地上以来，世代是否有无限数目；灵魂最初是何时被缚于肉身的；是谁策划了这种束缚，不，是谁造了人自己；生活在我们之前的人的灵魂往何处去了；它们在世界的哪些部分或区域；它们是否可朽；如果没有基督在急需时作为他们的保全者前来，他们是否可能遭遇死亡的危险呢？放下这些挂虑，放弃这些你们找不到答案的问题吧。主的怜悯也已经向他们显明，属神的仁慈同样惠及众人；他们已被保全，已蒙解救，并已放下了必朽的命运和状况。我的对手们问：\Emph{什么样的、什么、何时？}如果你们不存预设、骄傲和自负，你们早就可从此师傅那里学到这一切了。\par

64. 但是，\Emph{我的对手们问}，\Quote{如果基督如你们所说作为人类的救主来临，为什么祂不以一致的慈爱毫无例外地解救所有人呢？}\Emph{我答复说}，\Quote{难道那位平等地邀请所有人的，不也平等地解救所有人吗？或者，祂曾拒绝或排斥任何人，不让他接近天主的仁慈吗？祂赐给所有人同等的能力来到祂面前：高贵者、最卑贱的奴仆、妇女、儿童？祂说，生命的泉源对所有人敞开着，没有人被阻挡或拦阻去饮用。}如果你如此挑剔，以至于唾弃祂仁慈提供的恩赐，甚至更甚，如果你的智慧如此之大，竟将基督提供的事物视为可笑和荒谬，那么为什么祂要继续邀请你呢？而祂唯一的义务就是使祂恩赐的享用取决于你自己的自由抉择。\Person{柏拉图}说，\Quote{天主并不使任何人选择自己生命中的命运；另一个人的选择也不能正确地归因于任何人，因为选择的自由已置于创造它的那位手中。}难道必须恳求你屈尊接受天主的救恩恩赐吗？难道天主的仁慈必须倾注到你心中，而你却轻蔑地拒绝、远远逃离呢？你选择接受所赐予的，并将它转化为自己的好处吗？\Emph{那样的话}，你就考虑了自己的利益。你轻蔑地拒绝、轻视并鄙弃它吗？\Emph{那样的话}，你就剥夺了自己从这恩赐中得益。天主不强迫任何人，也不用压倒性的恐惧恐吓任何人。因为我们的救恩对祂并非必需，以至于如果祂使我们成为神圣，或任由我们被腐朽所湮灭，祂会获得什么或遭受什么损失。\par

65. 不，\Emph{我的对手说}，\Quote{如果天主是全能的、仁慈的、愿意拯救我们，就让他改变我们的性情，强迫我们信赖他的许诺。}这便成了暴力，不是仁慈，也不是至高天主的恩赐，而是为求控制的一种幼稚而虚妄的斗争。因为有什么比强迫不情愿的人改变他们的倾向更不义呢？把自己不愿接受和躲避的东西强行印在他们心中；在施恩之前先伤害，并通过剥夺前者，把人带到另一种思想和感觉的方式中？你这希望自己被改变并遭受暴力的人，使你去做并被迫接受你所不愿的事物，为什么你不在被改变和转变后自愿选择你愿做的事呢？\Emph{他说}，\Quote{我不愿意，也没有愿望。}那么，你为什么责备天主，好像他亏负了你呢？你是否希望他来帮助你，而你不仅拒绝和躲避他的恩赐与恩惠，还称之为空话，并以诙谐的俏皮话攻击它呢？\Emph{我的对手说}，\Quote{除非我成为基督徒，否则我不能期望救恩。}正如你自己所说的。因为，为了带来救恩，并将那应当赐予并必须加诸灵魂的赋予它们，\Emph{唯独基督}从天主圣父那里接受了这一职责和托付，遥远的和更隐秘的原因已被如此安排。因为，像你们那里，某些神具有固定的职能、特权和能力，你不会向他们任何人要求不在其能力和许可范围内的事物，同样，唯独基督有权给予灵魂救恩并赐予他们永生。因为如果你相信父\Person{巴克斯}能赐予美酒，但不能解除疾病；如果你相信\Person{刻瑞斯}能赐予好收成，\Person{埃斯库拉庇乌斯}赐予健康，\Person{尼普顿}赐予一样东西，\Person{朱诺}赐予另一样，\Person{福耳图娜}、\Person{墨丘利}、\Person{伏尔甘}各自是某种固定和特定事物的赐予者——那么，你们也必须从我们这里接受这一点：灵魂不能从任何人那里获得生命和救恩，除了那至高统治者将这一职责和任务托付给他的那一位。世界的全能主宰已决定，这应是救恩的道路——可以说，这是生命之门；唯有通过他，才能进入光明：人们既不能从其他地方偷偷溜入或进入，所有其他道路都被不可逾越的屏障关闭和封锁了。\par

66. 因此，那么，即使你是洁净的，已洗清了一切恶习的污点，已争取并迷住了那些力量，使它们在你返回天堂时不给你关闭道路和阻碍你的通行，若不通过基督的恩赐，你领悟这真正的不朽究竟是什么，并被允许进入真实生命，你无论如何努力，也无法达到不朽的奖赏。因为关于你们习惯用来嘲讽我们的那一点：我们的宗教是新的，几乎才在不久之前兴起，并且你们不能放弃从祖先那里继承下来的古老信仰，而去接受野蛮而异邦的礼仪，这种指责完全没有道理。因为，如果我们以这种方式选择责备此前的、甚至最古老的时代，那会怎样呢？当他们发现如何种植庄稼时，他们鄙视橡实，并轻蔑地拒绝野草莓；在设计了更为有用和方便的布衣服之后，他们不再以树皮蔽体、身穿野兽皮革；或者，在建造了房屋和更舒适的住所之后，他们并不固守他们古老的棚屋，也不喜欢像田野的野兽那样留在岩石和洞穴下面？这是一种所有人都具有的倾向，从我们几乎在摇篮里时就被铭刻在我们身上：宁可选择好的而非坏的、有用的而非无用的，并更愉快地追求和寻求那些通常被认为格外珍贵的事物，并将我们对繁荣和顺利境遇的盼望寄托于此。\par

67. 因此，当你们指责我们背离了往昔\Emph{时代}的宗教时，你们理当审视我们为何这样做，而非究竟做了何事；不要只注意我们离弃了什么，却要特别观察我们遵从了什么。因为，若改变主张、从古老习俗转向新式状态与愿望是过失或罪行，那么这一指控同样适用于你们，你们也已如此频繁地改变了自己的习惯与生活方式，转而追随其他风俗与礼仪，以致你们也\Emph{与我们一样}被往昔时代所谴责。你们当真仍如你们祖先曾做的那样，将人民划分为五个等级吗？你们何曾经由人民投票选举官员？你们可知道军事民众会议、城市民众会议和平民民众会议\OriginalTerm{民众会议}{comitia}为何物？你们是否观测天象，是否因恶兆宣布而中止公共事务？当你们备战之时，可曾从卫城悬挂旗帜，或遵行战和事务祭司团\OriginalTerm{战和事务祭司团}{Fetiales}的仪轨，严正要求归还被劫掠之物？抑或面临战争危险时，你们是否也开始因矛尖所示的吉兆而怀有希望？就任官职时，你们仍遵守制约适当时机的律法吗？至于礼物与馈赠，\Emph{给辩护人的，你们可遵循}辛西亚法\OriginalTerm{辛西亚法}{Cincia lex}和禁奢法\OriginalTerm{禁奢法}{leges sumptuariae}来限制花费吗？你们可曾在幽暗的圣所中维持永燃之火？你们是否以安放盐碟及神像来祝圣餐桌？婚礼之时，你们可曾用托加\OriginalTerm{托加}{toga}铺陈床榻，并呼求丈夫们的守护神\OriginalTerm{守护神}{genii}？你们可曾用独身矛\OriginalTerm{独身矛}{hasta caelibaris}梳理新娘的头发？你们是否将少女的衣裳持往\Person{福尔图娜·维吉纳利斯}的神殿？你们的已婚妇女可在你们宅邸的大厅中劳作，公开显示其勤勉，并戒饮葡萄酒？她们的朋友与亲属获允亲吻她们，以此彰显她们的清醒与节制？\par

68. 在\Place{阿尔班山}，古时只许献祭雪白的公牛：你们岂非改变了这一习俗与宗教仪轨，并且\Emph{岂不是}通过元老院法令规定，红棕色的公牛亦可献上？在\Person{罗慕路斯}和\Person{庞皮留斯}统治期间，内脏被彻底煮熟烂软，之后在\Emph{献祭}给神明时焚烧；而在\Person{图利乌斯}王时，你们难道不是开始改为半生不熟、微温即献上，毫不顾及先前的惯例吗？在\Person{赫拉克勒斯}降临意大利之前，曾依\Person{阿波罗}的指示，用活人之首向父神\Person{迪斯}和\Person{萨图恩}祈求；你们不也以狡诈的欺瞒和模棱两可的名称，同样改变了这种风俗吗？因此，既然你们自己也时而遵循这些习俗，时而遵循不同的律法，并在看到自己的错误或遇见更好的事物时，摈弃和拒绝了许多东西，那么，我们选择了更伟大、更确定的事，不容自己受无理智的迷信所阻碍，又怎能算作有悖于常理和人之共有的辨别力呢？\par

69. 但有人对我们说，我们的名号是新的，我们所奉行的宗教是数日前才兴起的。暂且承认你们现在对我们的责难并非不实，那么我要问，在人间诸事中，有哪一件是由身体的努力和手工劳动所致，或由心智的学识与知识所获，又不是在某个时间起始，而后才进入普遍运用与实行呢？医学、哲学、音乐，以及其他一切构建并精炼了社会生活的技艺——这些是与人类同时诞生的吗？难道不反而是直到晚近——甚至可以说前不久——才开始为人所追求、理解和实践的吗？在伊特鲁里亚的\Person{塔格斯}出世以前，可有人知道或费心去得知和了解雷霆的堕落或祭牲内脏的纹理有何含义呢？星辰的运行或推算出生星位\OriginalTerm{出生星位}{nativitas}的技艺，是何时开始为人所知的？难道不是在埃及人\Person{修提斯}之后，或如某些人所说，在\Person{阿特拉斯}之后吗？——这位扛天者、支撑者、支柱和苍穹的承负者。\par

但我又何必\Emph{谈论}这些琐碎之事呢？那些不朽的神明自身，你们如今怀着敬畏进入其庙宇，跪拜其神性的神明，难道不也曾据你们的典籍和传说所载，在某些时候，因被赋予的名号和称号而开始存在、被认识、被呼求吗？因为，若说\Person{朱庇特}及其兄弟们是由\Person{萨图恩}与其妻所生，那么在\Person{奥普斯}成婚生子之前，\Person{朱庇特}——无论至高者和冥河之神——都还不存在，不，连海神、\Person{朱诺}也不存在，更有甚者，天界宝座中除了双亲外无人居住；反倒从他们的结合中，其余神明被孕育、出生，并呼吸生命的气息。这样看来，\Person{朱庇特}神在某个时候才开始存在，在某个时候才配受崇拜和\OriginalTerm{祭献}{sacrifices}，在某个时候才在权能上凌驾于其兄弟们之上。再说，若\Person{利伯尔}、\Person{维纳斯}、\Person{狄安娜}、\Person{墨丘利}、\Person{阿波罗}、\Person{赫拉克勒斯}、\Person{缪斯}、\Person{廷达瑞俄斯兄弟}，以及火之主宰\Person{伏尔甘}，都是由父\Person{朱庇特}所生，且由出自\Person{萨图恩}血统的母亲所出，那么在\Person{记忆女神}、\Person{阿尔克墨涅}、\Person{迈亚}、\Person{朱诺}、\Person{拉托娜}、\Person{勒达}、\Person{狄俄涅}和\Person{塞墨勒}为\Person{狄斯帕特}产子之前，这些神明同样在世界上、在宇宙的任何角落都不存在；而是因\Person{朱庇特}的拥抱，他们才被孕育、出生，并开始意识到自身的存在。如此看来，这些神明也是在某个时候开始存在，并开始被宣召参与\OriginalTerm{神圣仪式}{sacred rites}的神明行列。对\Person{密涅瓦}，我们也是如此说。因为，若如你们所言，她从未受生便从\Person{朱庇特}的头颅中跃出，那么在\Person{朱庇特}被孕育、在母腹中获得身体形廓之前，\Person{密涅瓦}显然未曾存在，也不被视为万物之一或实在的存在；而是从\Person{朱庇特}的头颅中出生，才开始有真实的存在。因此，她起初有源头，并在某个时候才被称为女神，被配享庙宇，并因宗教不可侵犯的规条而被\OriginalTerm{祝圣}{consecrated}。情况既如此，当你们谈论我们宗教的晚近时，难道你们自己的宗教就不进入你们的思绪，你们就不留心查考你们的众神是何时产生——他们有什么来源、什么起因，或出自怎样的族系而迸生涌现吗？然而，把你们自己也在做的事拿去责备别人，却拿来嘲笑和指控他人，而这些事又能转而回击到你们自己身上——这是多么可耻、多么厚颜无耻呢！\par

\Emph{有人声称，}我们的仪式是新的\OriginalTerm{仪式}{rites}；你们的则是古老的，且极其悠久。这于你们有何裨益，或对我们的事理和论证有何损害呢？我们持守的信仰是新的；有朝一日，它也会成为古老的：你们的信仰是古老的；然而当它兴起时，它也是全新、闻所未闻的。然而，宗教的可信不应由年代久远，而应由其神性来断定；你们所当思考的，不是你们何时开始崇拜，而是所崇拜的究竟是什么。\Emph{我的对手说，}四百年前，你们的宗教并不存在。而\Emph{我回答说，}两千年前，你们的神明并不存在。\Emph{你们会问，}凭哪种推算或哪种计算，能得出这推断？这并不困难，也不复杂，即便如俗语所说，任何人只要愿意着手，都能看见。谁生了\Person{朱庇特}及其兄弟们？据你们叙述，是\Person{萨图恩}与\Person{奥普斯}，他们出自\Person{刻卢斯}与\Person{赫卡忒}。谁生了\Person{皮库斯}——\Person{法乌努斯}之父、\Person{拉丁努斯}之祖父？据你们的书册和教师们所传述，仍是\Person{萨图恩}。因此，若情形如此，\Person{皮库斯}与\Person{朱庇特}便因亲缘纽带而相连，因为他们出自同一血脉与族系。可见，我们所说的属实。从\Person{朱庇特}与\Person{皮库斯}传至\Person{拉丁努斯}，共有多少代？如世系所示，有三代。你们愿假定\Person{法乌努斯}、\Person{拉丁努斯}、\Person{皮库斯}各享寿一百二十岁吗？因为人的寿命不能超越此限。这一估算基础稳固且清楚。那么，后三者之间共有三百六十年吗？正如计算所示。\Person{拉丁努斯}的岳父是谁？是\Person{埃涅阿斯}。他又是谁之父？他是\Place{阿尔巴}城创建者之父。在\Place{阿尔巴}，诸王统治了多少年？约四百二十年。据编年记载，\Place{罗马}城被证实有多大年岁？计有一千零五十年，或与之相差不远。如此，从\Person{朱庇特}——\Person{皮库斯}的兄弟、其余较低神明的父亲——直至今日，算来近乎，或略添时日，总共两千年。既然这无可辩驳，那么不但你们所信奉的宗教被显明为晚近才出现，而且\Emph{也显明}那些你们冒着引发疫病的风险堆上公牛和其他祭品的神明本身，都是年幼的娃娃，还当喂以母乳才对！\par

72. 但是你们的宗教比我们的宗教早许多年，因此，\Emph{你们说}，因为它得到古老权威的支持而更为真实。然而，即便它比\Emph{我们的}宗教早多少年随你们说罢，既然它始于某个特定时间，这又有什么意义呢？与无数个时代相比，两千年又算得了什么呢？不过，为了使我们不致因长久忽视而看似放弃\Emph{我们的}事业，请问，若不使你们不快的话，全能而至高的天主在你们看来是新近的吗？那些钦崇并敬拜祂的人，在\Emph{你们看来}，是在推行并引进一个闻所未闻、不为人知且突然兴起的宗教吗？有比祂更古老的吗？在实体、时间、名号上，能找到先于祂的存在吗？难道不是唯有祂是非受造的、不死不灭的、永恒的吗？万物之首与泉源是谁？不就是祂吗？永恒之名当归于谁？不归于祂吗？难道不正是因为祂是永恒的，世代才无尽地流转吗？这是毫无疑问且真实的：因此，\Emph{我们所追随的宗教}并非新的宗教，而是我们迟迟才认识到应当追随并尊崇的对象，应当将我们救恩的望德寄托于何处，并善用那赐予我们、为拯救我们\Emph{而施予的}援助。因为那要指导迷失\Emph{正道}者以路途，赐予那躺卧最深黑暗者以知识之光，并驱散其愚昧之盲目者，尚未显现。\par

73. 然而，难道只有我们处于这种境地吗？怎么！自从\Person{皮索}和\Person{伽比尼乌斯}担任执政官以来，你们不就已经把名为塞拉比斯和伊希斯的埃及神祇列入你们诸神的行列中了吗？怎么！当\Place{迦太基}人\Person{汉尼拔}蹂躏\Place{意大利}并觊觎世界霸权时，你们不曾开始认识、了解并以非凡的敬礼尊崇那\Place{弗里吉亚}母神吗——据说她最初是由\Person{弥达斯}或\Person{达耳达诺斯}立为女神的？那不久之前才被采纳的母亲刻瑞斯的圣礼，不正因为你们不认识它们而被称为\OriginalTerm{希腊的}{Graeca}，其名称便见证了它们的新奇吗？在博学之士的著述中，不是说过\Person{努马·庞皮利乌斯}的仪典中并未包含阿波罗之名吗？由此清楚显明，他也是你们所不认识的，但后来某个时候，他也开始为人所知了。因此，若有人问你们为何直到最近才开始敬拜我们方才提到的那些神祇，你们必定会回答说，要么因为我们 C 最近才不知道它们是神，要么因为我们现在得到了先知的警告，要么因为在极其艰难的处境中，我们蒙了它们的恩宠与援助而得保全。但若你们认为这样的回答出于你们之口是合理的，那么你们务必考虑，在我们这边，也已作出了类似的答复。我们的宗教乃是刚刚兴起；因为如今那一位已经到来，祂被派遣来向我们宣告它，引领\Emph{我们}归向它的真理，昭示天主为何者，将我们从纯然的臆测中召唤至对祂的敬拜。\par

74. 那么，\Emph{我的对手说}，为什么天主，那宇宙的统治者和主，竟决定在几小时前——据说如此——从高天之上派遣一位救主基督给你们呢？我们反过来也要问你们，是什么原因，什么理由，使得季节有时不在其相应的月份出现，而冬季、夏季和秋季却姗姗来迟？为什么在庄稼干枯、谷粒已死后，有时降下雨来，而这些雨本该在它们还未受损时落下，以应时之需？不仅如此，我们更要问，如果适宜于使\Person{赫拉克勒斯}、\Person{埃斯库拉庇乌斯}、\Person{墨丘利}、\Person{利柏尔}及其他一些神祇诞生，好让他们既能被列入诸神之列，又能为人类做些什么——那为什么\Person{朱庇特}如此迟延才产生他们，以至于只有较晚的时代才认识他们，而先前那些更早的时代却对他们一无所知呢？你们会说其中必有某种缘由。那么，我们种族的救主不是久远之前，而是今日才来临，这里也必有某种缘由。那么，你们会问，缘由何在？我们承认自己不知，并不否认。因为无人能够窥见天主的意念，或祂安排计划的途径。人，这盲目的受造物，甚至不认识自己，根本无法得知何事应当发生、何时发生或其性质为何：唯有圣父——万有的统治者和主——独自知晓。即使我未能向你们揭示为何某事以某种方式成就的缘由，也绝不意味着所成就之事便化为乌有，或一件已借如此大能与权能显明其为不容置疑之事，便成为不可置信的了。\par

75. 你们或许会反驳说：救主为何被差遣得如此之晚？\Emph{我们回答}，在无限、永恒的世代中，任何事都不可称为晚。因为那里既无终结也无开端，没有什么事是太早，也没有什么太晚。时间乃是由其起始与终结被人所感知，而绵延不绝、无始无终的世代序列则不可能有这些。试想，若需要救助之事本身要求某个适当的时间呢？试想，若远古的状况与后来的时代不同呢？若有必要以一种方式帮助古人，而以另一种方式眷顾他们的后代呢？难道你们没有听到你们自己的著作中念诵，曾经有世人\Emph{乃是}半神，以及体魄魁梧的英雄吗？难道你们没有读到，尚在母亲怀中的幼儿发出\Person{斯腾托尔}般的尖啸，而当他们的骨骸在地里各处被掘出时，发现者几乎怀疑它们是人类的肢体吗？因此，可能是全能的天主，唯一的天主，确是在人类变得\Emph{越来越}脆弱、软弱，开始成为我们如今这般模样之后，才差遣了基督。若今日所行之事能在数千年前成就，至高的统治者必已成就；又或者，若今日所成之事适宜再经过数千年后才实现，没有任何力量能迫使天主预先跳过那必要的时间间隔。祂的计划以确定的方式执行；而一旦定下的决断，绝无可能再行更改。\par

76. 那么，你们说，既然你们事奉全能的天主，且相信祂眷顾你们的安全与\OriginalTerm{救恩}{salvation}，为什么祂容许你们遭受如此迫害的风暴，经受各种惩罚与酷刑呢？让我们也反问一句：为什么你们崇拜如此众多又数不胜数的神祇，为他们兴建庙宇，铸造金像，宰杀成群的牲口祭献，\Emph{并且}将一箱箱的香堆在已然满载的祭台上，却仍每日生活在许多危险和\Emph{灾祸}的风暴之中，频遭种种致命的厄运呢？我说，为什么你们的神祇忽略不替你们免去诸般疾病与病症、海难、倒塌、火灾、瘟疫、不育、丧失子女、财产充公、不和、战争、仇恨、城邑被攻陷，以及失去生来自有者的权利而沦为奴隶呢？但是，\Emph{我的对手说}，在这样的灾祸中，我们同样丝毫得不到天主的帮助。原因显而易见。因为对于此生，我们没有蒙受任何希望的许诺，也没有为这肉身的躯壳获赐应许的帮助或规定的援助——不仅如此，我们还受教，对于命运的任何威胁，都要轻看淡视；倘若有什么极其沉痛的灾祸临到\Emph{我们}，就当在\Emph{那}不幸中视那必随之而来的结局为喜乐，而不畏惧或逃避它，以便我们能更容易地脱离肉身的桎梏，逃离我们的黑暗与盲目。\par

77. 所以，你们所说的那种迫害之苦，乃是我们的解救而非迫害，我们的受虐非但不会加害于我们，反将引领我们进入自由的光明。犹如一个冥顽不灵的家伙，以为若非对监狱本身发怒，砸碎其石头，烧毁其屋顶、墙壁、门窗，拆毁、推倒、摔碎其馀部分，就不算严酷残忍地惩罚了一个已入狱的人，却不知这样倒给他似乎正在伤害的人带来了光明，夺去了他可咒的黑暗：同样地，你们用火刑、流放、酷刑以及种种骇人听闻的手段撕裂扯碎我们的身体，并非夺去我们的生命，而是除去我们的皮囊，不知道你们越打击、越企图向我们的这些影子和形体泄愤，便越使我们脱离沉重压迫的锁链，斩断我们的捆锁，让我们飞升到光明之中。\par

78. 所以，人啊，不要用无谓的问题阻碍你们所盼望的；若有任何事与你们的想法不符，也不要宁可相信自己的意见，而不信那应受尊崇的。充满危险的时势催迫我们，致命的惩罚威胁着我们；让我们投奔天主我们的救主以获安全，不必要求解释这赐予的恩惠。当攸关我们灵魂的\OriginalTerm{救恩}{salvation}和我们自身的利益时，有些事即使没有理由也必须去做，正如\Person{阿里安}认可\Person{爱比克泰德}所说的话。我们怀疑，我们犹豫，猜疑所说之事的可信性；让我们把自己交托给天主，不要让我们的不信胜过祂圣名和能力的伟大，免得我们正在为自己寻找论据，借以使我们所不愿且拒绝承认为真的事显得虚假之时，末日偷偷临到我们，而我们被发现陷入我们仇敌——死亡——的口中。\par

\SourceRecord{Translated by Hamilton Bryce and Hugh Campbell.；Ante-Nicene Fathers；Vol. 6.；Edited by Alexander Roberts, James Donaldson, and A. Cleveland Coxe.；Buffalo, NY: Christian Literature Publishing Co.,；1886.；<http://www.newadvent.org/fathers/06312.htm>.}

% structure: p0001 source=h1 target=section reason=page-title

\section{\WorkTitle{反异教徒}（卷三）}

1. 所有这些指控，更确切地说应该称之为辱骂，早已由在这方面杰出且配得认识真理的人士给予了充分而周详的答复；没有任何探究之点被遗漏，而不以千百种方式、凭最坚实的理由加以论述。因此，我们无须在这一部分案情上再多耽搁。因为基督宗教即使没有辩护者，也并非不能自立；即使有众多人赞同并获得份量，也不因此而证明为真。它自有其充分的力量，奠基于自身真理的根基，即使无人为之辩护，不，即使所有声音都攻击反对它，并以共同的敌意联合起来摧毁对它的全部信仰，它也不会丧失其能力。\par

2. 让我们现在回到我们方才不得不岔开的顺序，以免我们的辩护因中断过久，而被认为给诋毁者提供了确定其指控的胜利理由。因为他们提出这些问题：如果你们对宗教是认真的，为什么不与我们一同事奉敬拜其他的神明，或与你们的同伴分享你们的神圣仪式，并将\Emph{不同的}宗教的仪式置于同等地位？我们姑且可以这样说：在试图接近神明时，至高天主为我们足够了——我说的是那位至高者，宇宙的创造者和主，祂安排并掌管万物：在祂内，我们事奉一切需要我们事奉的；\Emph{在祂内}，我们敬拜一切当受敬拜的——尊敬那要求我们至诚敬礼的。因为当我们把握住神明本身的源头，即一切神明所有神性的由来时，我们以为逐一亲近每个神明是徒劳的，因为我们既不知道他们是谁，也不知他们被称呼的名字；更无法得知、发现并确定他们的数目。\par

3. 正如在世间王国里，我们绝无必要明确地对王室家族成员与对君主本人一样行礼致敬，但凡属于他们的尊荣，都默然隐含在对国王本人的敬礼之中；同样，这些你们保证其存在的神明，无论他们是谁，如果他们是王室族裔，出自至高统治者，即使我们不明确向他们致敬，仍然感到他们与他们的主一同受尊崇，并在对祂的敬礼中有份。但\Emph{必须记得}，我们作出这一陈述，仅仅基于一个清晰且无可否认的假设，即除了这位统治者与主本身之外，还有其他的存在者，当他们按等级序列排列安置时，宛如构成一种平民群体。但不要试图向我们指看你们庙宇中那些替代神明的画像，以及你们竖立的偶像，因为你们自己也知道，只是不愿且拒绝承认，它们是用最无价值的泥土做成，出自工匠之手的幼稚人形。当我们与你们谈论宗教时，我们要求你们证明这一点：除了那位独一的至高天主之外，还有其他在性体、能力、名称方面是神明的——不是像我们所看见的显现于偶像之中，而是存在于那种理应与如此伟大尊荣的圆满相称的实体中。\par

4. 但我们不打算在这部分论题上再拖延，以免显得我们有意激起最激烈的纷争并卷入扰攘不休的争论。\par

就算如你们所肯定的，有那一大群神明，有数不清的神明家族；我们同意，赞同，\Emph{并且}不过分\Emph{仔细}地审查，也不在任何论点上攻击你们所持的犹疑不定的立场。然而，我们要求你们告诉我们，你们从何处发现，或如何得知，究竟有没有这些你们相信在天上、你们所事奉的神明，或是某些其他名号未闻、不知其名的神明？因为可能存在那些你们不认为其存在者；而你们确信其存在的那些，在宇宙中却无处可寻。因为你们从未被提升至天上的星辰，\Emph{从未}见过每一神明的面貌形状；\Emph{然后}才在此地建立对同样神明的敬拜，你们记得它们在那里，是你们所认识、\Emph{被你们}见过的。但这一点，我们也要再向你们学习：这些你们用来称呼他们的名字，是他们固有的，还是在洁净之日自取的？如果这些是属天的圣名，是谁传给你们的？但另一方面，如果这些名字是你们加于他们的，你们怎能给你们从未见过、其品德或境遇你们一无所知的神明命名呢？\par

5. 但\Emph{就假设}如你们所愿望、所相信、所确信的那样，有这些神明；也让他们被那通俗之人以为称呼那些次级神明所用的名字所称呼。然而，你们从何处得知，这些名号之下的神明名单是由谁组成的？有没有任何神明被他人所熟悉认识\Emph{于他人}，而他们的名字却是你们不知道的？因为很难知道，他们那众多的群体，在\Emph{数目上}是确定不变的；还是说，他们的众多，无法被任何计算所统计和限定。让我们假定你们敬奉一千位，甚至五千位神明；但在宇宙中，或许有十万位；可能还不止于此——不，正如我们稍前所说，也许神明的数目无法计算，无法以确定数字限定。那么，要么你们自己是不虔敬的，因为你们只事奉少数神明，却忽视了对其他神明应尽的义务；要么，如果你们声称对其它神明的无知应当得到宽恕，那么你们也同样会为我们求得类似的宽恕，如果我们以完全相同的方式拒绝敬拜那些我们对其存在完全无知的神明。\par

6. 然而，莫让人以为我们执拗地不愿顺服其他神祇，不管他们是谁！因为我们\Emph{举起}虔诚的心灵，伸手祈祷，并不拒绝前往你们召唤我们的任何地方；只要我们了解你们强加给我们的那些神灵究竟是谁，以及可以和谁分享我们对那统御万有的君王与元首的崇敬。\Emph{我的对手说}，那是\Person{萨图恩}、\Person{雅努斯}、\Person{密涅瓦}、\Person{朱诺}、\Person{阿波罗}、\Person{维纳斯}、\Person{特里普托勒摩斯}、\Person{赫拉克勒斯}、\Person{埃斯库拉庇俄斯}，以及所有其他神祇，古人的崇敬几乎在每座城市都为他们建造了宏伟的庙宇。你们也许本来能够吸引我们去崇拜你们提到的这些神明，若非你们自己率先以粗鄙不当的幻想，编造出如此的故事，不仅玷污了他们的名誉，更因赋予他们的本性，证明了他们根本不存在。因为，首先，我们无法被引导相信这一点——那不死不灭、至高无上的本性竟被区分为男女两性，有些是男神，有些是女神。然而，这一点确实早已被拉丁和希腊的热诚天才们充分探讨过；罗马人中最雄辩的\Person{西塞罗}，不惜担起不敬神的责难，尤其以更大的虔诚——大胆、坚定而坦率地——宣布了他对这种想法的看法；如果你们愿意从他那里接受凭真实见识写下的观点，而不只是\Emph{仅仅}精彩的语句，这案子本该早已了结；也不需要由我们这些弱者来作所谓的第二轮辩护。\par

7. 然而，我何必说人们追求他精妙的表达与华丽的辞藻呢？我知道有许多人回避并逃离他关于这个话题的书籍，不愿听取他的观点而被推翻成见；我还听到另一些人愤怒地咕哝，说\Quote{元老院应下令销毁这些著作，因为它们维护了基督宗教，压过了古老传统的分量}。但是，事实上，如果你们确信自己关于神明的任何说法都毫无疑问，那么请指出\Person{西塞罗}的错误，驳斥、反驳他轻率而不敬的话语，\Emph{并且}证明\Emph{它们确实如此}。因为当你们要带走著作，查禁一本公之于众的书时，你们不是在保卫神明，而是畏惧真理的证据。\par

8. 然而，为了避免任何轻率之人对我们虚加指控，仿佛我们相信我们所崇拜的天主是男性——因为我们谈论祂时使用阳性词汇——当让他明白，这表达的并非性别，而是祂的名号，以及按习惯这名的含义，和我们惯常使用词语的方式。因为神性并非男性，而是祂的名在语法上属阳性：但在你们的礼仪中，你们却不能这样说；因为你们在祈祷时习惯于说\Quote{\Emph{无论你是男神还是女神}}，而这种不确定的描述，即便从它们彼此对立来看，也表明你们将性别归于神明。因此，我们无法被说服相信神圣是有形体的；因为形体必然以性别差异来区分，若是男女有别的话。因为，不论一个人的能力如何低下，有谁不知道，不同性别的两性是由地上万物的创造者所命定和形成的，只是为了通过身体的交合与结合，使那短暂易逝的万物能因不断更新与维持而持久不衰呢？\par

9. 那么，我们该怎么说呢？神明生育，也为神明所生？因此他们领受了生育器官，从而能够繁衍后代，并且，随着每一新种族的兴起，一种定期发生的替代，应当弥补先前时代所扫除的一切？那么，若果真如此——也就是说，如果上面的神明生育\Emph{其他神明}，且受这些性别条件的制约，却又是不死不灭的，不因岁月寒凉而衰败——其必然结果是，世界应充满神明，无数的天穹也容纳不下他们那众多的数量，因为他们既自身不断生育，而他们无数的后代，始终在增多，又经由自己的子孙而不断增长；或者，若如适宜的那样，神明并不因受制于性欲冲动而被贬低，那么，又有什么原因或理由能指出，他们竟以那些肢体为特征，而两性正是借助这些肢体，在各自欲望的驱使下彼此辨认呢？因为，他们不大可能毫无目的地拥有这些，或自然愿在他们身上捉弄自身的疏忽，为他们提供了毫无用处的肢体。因为，正如手、脚、眼，以及构成我们身体的其他肢体，均被安排用于特定用途，各为其目的，我们同样有充分理由相信，这些肢体的赐予，乃是为履行其职能；否则，便必须承认，在神明的身体中有某些无目的之物，它们被徒劳而无益地造成了。\par

10. 你们这些宗教的圣洁无玷的守护者啊！那么，神明真的有性别吗？他们是否被那些部位所丑化，那些部位的名字端庄的嘴唇连提都不能提呢？那么，现在剩下的只是相信他们像不洁的野兽一样，被强烈的情欲驱动，带着疯狂的欲望投入彼此的怀抱，最终身体破碎毁坏，因纵欲而衰弱。并且因为某些事情是女性所特有的，我们必须相信女神们也在适当时期屈从于这些状况，怀着厌恶受孕怀胎，流产，足月，有时早产。哦，纯洁、圣善、毫无任何可耻污点的神性啊！心灵渴望并燃烧着要看见，在天庭的宏大殿堂中，神和女神们裸露着身体，丰乳的\Person{刻瑞斯}哺育\Person{伊阿科斯}，正如\Person{卢克莱修}的缪斯所唱，\Place{赫勒斯滂}的\Person{普里阿普斯}在处女和已婚女神之间炫耀着他那永远准备交媾的器官。我说，它渴望看到怀孕的女神，拖着身孕，随着肚子天天变大，脚步蹒跚，因为所负的重担而苦恼；另一些，经过长久迟延，到了分娩，寻求接生婆的帮助；还有一些，遭受剧烈阵痛和剧痛时尖叫，受尽折磨，并在这一切之下祈求\Person{朱诺·卢西娜}的援助。与其假装虔诚，却以如此可耻的迷信观念看待神明，岂不远远不如斥骂、诽谤或以其他方式侮辱神明吗？\par

11. 你们竟敢指控我们冒犯神明，尽管经审查发现，冒犯的缘由极其明显地在于我们自己，而不是由于你们所认为的侮辱。因为如果神明确如你们所说，会被愤怒感动，心中燃起盛怒，我们为什么不应认为，他们极其恼怒于你们将性别归于他们，如同狗和猪被造那样，并且由于这是你们的信仰，他们被如此描绘，以一种可耻的方式公开暴露。那么，既然如此，你们是一切灾难的缘由——你们引领神明，你们刺激他们用各种邪恶骚扰大地，每天策划各种新的不幸，以便他们能为自己复仇，因被你们引起的如此多的错待和侮辱所激怒。我要说，因你们的侮辱和冒犯，部分是在卑鄙的故事中，部分是在可耻的信条中，这些是你们的神学家、你们的诗人，还有你们自己，在丢脸的仪式中所庆贺的，你们会发现人类事务已被毁坏，神明已丢弃了舵柄，如果人类的命运确实是由他们照管和安排的。因为对我们，他们确实没有理由发怒，他们看到并察觉我们既不像所说的那样嘲笑他们，也不敬拜他们，并且对于他们名号的尊荣，我们的思考和信仰远比你们更相称。\par

12. 关于性别就谈到这里。现在让我们来看看你们相信天上神明被赋予的外貌和形状，你们确实用这些外形来塑造他们，并把他们安置在你们最辉煌的居所——你们的神庙中。在此，谁也不要拿犹太人的寓言和撒杜塞人的寓言来反驳我们，好像我们也把形象归于天主；因为据说这是他们的经书中所教导的，并且仿佛以确信和权威所断言的。因为这些故事要么与我们无关，与我们毫无共同之处，要么如果如你们所信，\Emph{我们也}有份，你们必须寻求更有智慧的教师，通过他们你们或许可以学会如何最好地克服那些经文的幽暗深奥之说。我们对此事的看法如下：——整个天主性体，既不在任何时候开始存在，也永不会结束生命，是没有任何形貌特征的，也不具有任何像各个肢体末端通常构成部位集合的那种形式。因为凡是具有这种特性的，我们认为都是有死有坏的；我们也不相信任何被不可避免的终结所封闭的东西能永远持续，纵然包围它的界限是最遥远的。\par

13. 但你们以形式局限神明还不够：——你们甚至将他们限制在人形上，并且更加不体面地将他们封闭在尘世的身体内。那么，我们该怎么说呢？神明有一个造型完美对称的头，由筋腱牢固地连接在背和胸上，并且为了允许颈部必要的弯曲，它由\OriginalTerm{脊椎骨}{vertebrae}的组合和骨质基础支撑着？但如果我们相信这是真的，那么他们也有耳朵，被弯曲的耳道所穿透；转动的眼球，被眼眉边缘所遮蔽；一个鼻子，安置为通道，废物和气流可轻易通过；三种牙齿，用来咀嚼食物，适应三种功用；双手去做工作，通过关节、手指和灵活的手肘轻松活动；双脚支撑身体，调节步态，启动行走的最初动作。但若神明\Emph{具有}这些可见的部分，那么他们也应该具有那些皮肤掩盖在肋骨构架之下，以及包裹内脏的膜里的部分；气管、胃、脾、肺、膀胱、肝、缠绕长长的肠子，以及紫血的血管，与气管结合，流遍全部内脏。\par

14. 那么，这些神明的身体就没有这些畸形吗？既然他们不吃人的食物，我们是否应该相信，他们像孩子一样没有牙齿，没有内脏，就像充气的膀胱，由于肿胀的身体内部空虚而没有力量？此外，若是如此，你必须看看神明们是完全一样，还是在形体轮廓上有所不同。因为如果每一个都有一模一样的外形，那么相信他们会彼此认错、受骗，也就不足为奇了。但另一方面，如果他们通过面容来区分，那么我们应该由此推断，这些差异被赋予，无非是为了让他们能够通过不同的标记特点各自认出自己。因此我们可以说，一些神有大头、凸额、宽眉、厚唇；另一些有长下巴、痣和高鼻子；这些有扩张的鼻孔，那些是塌鼻子；有些因颚部肿胀或脸颊丰满而圆胖，矮的，高的，中等身材，瘦的，光润的，胖的；有些有卷曲的头发，另一些剃光了，头秃而光滑。现在你们的工坊显示并指出，我们的看法并非虚妄，因为当你们塑造和制作神明时，你们把一些描绘成长发，另一些光滑赤裸，有年老的，有青年的，有少年的，有肤黑的，灰眼的，黄肤的，半裸的，赤裸的；或者，为了让他们不受寒冷，披着飘逸的衣服。\par

15. 凡有判断力的人，谁会相信神明的身体上长着毛发和茸毛？相信他们之间还有年龄区分？而且他们穿着各种样式的衣袍，防暑御寒？但如果有人相信这些，他也必须接受如下真实：有些神明是漂洗工，有些是理发师；前者清洗圣衣，后者当头发长得浓密缠结时为其打薄。这难道不是极度的贬低、最为不敬和侮辱，竟将脆弱易朽动物的特征归于神明？赋予他们那些，任何端庄的人都不敢细述、描述，或在想象中描摹的肢体，而不致因极度的猥亵而战栗？这就是你们对我们的蔑视——这就是你们自以为傲的智慧，借此你们视我们为无知，以为一切宗教知识都归你们所有？你们嘲笑埃及人的奥迹，因为他们将哑巴动物的形象嫁接于神明的原因上，且用许多香火及这类仪式中所用的一切来崇拜这些形象：你们自己却崇拜人的形象，仿佛他们是强大的神明，并且不以赋予这些神明世俗受造物的面容为耻，指责他人错误愚昧，自己却被发现犯了同样可耻的错误。\par

16. 但是你们或许会说，神明确实另有其形，而你们赋予他们人的外表，\Emph{仅仅}是出于敬意，并且为了形式，而这比因无知而陷入任何错误更加侮辱神明。因为如果你们承认，你们是将自己所设想和相信的东西归于了神的形体，那么你们源于偏见的错误，还不至于那么可责。但现在，当你们相信是一回事，而塑造的是另一回事，你们既侮辱了那些你们承认其不属本性却被你们添加了那外形者，又显示你们的不敬，在于崇拜你们所塑造的，而非你们认为真实存在、确实如是的。如果驴、狗、猪具有人类的智慧和设计技巧，并想通过某种崇拜来尊敬我们，并通过为我们立像来表示敬意，若它们决定我们的像应具有并呈现它们自己身体的形状，这会在我们内心激起怎样的狂怒、引发怎样的情感风暴啊！我再说一遍，如果罗马的建立者\Person{罗慕路斯}被塑造成驴脸，可敬的\Person{庞皮利乌斯}被塑造成狗脸，如果在猪的形象下写着\Person{加图}或\Person{马尔库斯·西塞罗}的名字，他们会如何让我们充满愤怒、激起我们的激情呢？那么，难道你们以为，你们的神明若真有笑，不会嘲笑你们的愚蠢吗？或者，既然你们相信他们可能暴怒，难道你们以为他们不会被如此深重的错待和侮辱所激怒，被气疯，而不想为之复仇，并将由懊恼所支配、由刻骨仇恨所设计的惩罚倾泻在你们身上吗？还不如赋予他们大象、豹或虎、公牛和马的形象呢！因为人身上有什么美呢——我请问，有什么可钦佩的，或俊美的——除了，如某位诗人所坚称的，他与猿猴所共有的那个？\par

17. 但是，他们说，如果你不满意我们的意见，你们自己指出，告诉我们，天主的形式是什么。如果你们想听真理，要么天主没有形式；或者若祂确具某形体，我们实在不知道那是什么。再者，我们不以对从未见过的事物无知为耻；我们也不会因此便无法反驳他人的意见，因为我们自己在此事上并无己见可陈。因为，正如若有人说地球是由玻璃、银、铁制成的，或是由脆土聚集成形，我们不会犹豫断言此为不实，尽管我们不知它是由什么做的；同样，当讨论天主的形式时，我们表明它并非你们所主张的，即使我们\Emph{仍然}更不知该如何解释它是什么。\par

18. 那么，有人会说：\Quote{难道\OriginalTerm{天主}{the Deity}不听吗？难道祂不说话吗？难道祂看不见呈于祂面前的事物吗？难道祂没有视觉吗？祂可以按祂自己的方式，但不是按我们的方式。} 但在如此重大的事上，我们根本无从认识真理，也不能藉思辨抵达；因为很明显，这些思辨在我们这里是无根基的、欺骗性的，如同虚妄的梦。因为如果我们说祂以与我们相同的方式观看，就会推论出祂有眼睑覆盖瞳孔，祂会闭眼、眨眼，藉光线或影像观看，或者就像所有眼睛那样，若没有其他光的存在，就什么也看不见。因此，对于听觉、言语的形式，以及话语的发出，我们也必须同样推断。如果祂凭耳听闻，那么，\Emph{我们必须说}，祂也有被曲折的通道穿透的耳朵，声音携带着语义得以潜入；或者，如果祂的话语是从口中发出，那么祂就有唇齿，通过舌头的触碰和不同运动而清晰地发出声音。\par

19. 如果你们愿意听取我们的结论，\Emph{那么请知道}，我们并不将形体归于\OriginalTerm{天主}{the Deity}，甚至不敢将心灵的恩惠以及极少数人难能可贵地彰显的卓越品质归于如此伟大的存在。因为谁会称天主勇敢、坚定、良善、智慧呢？谁会\Emph{说}祂正直、节制，甚至说祂具有知识、理解力、先见呢？谁会说祂使自己所决定的行动朝向固定的道德目标呢？这些品质在人身上是好的；由于它们与恶习相对，已经赢得了巨大的声誉。但谁会如此愚昧、如此昏聩，以至于说天主仅仅因人性的卓越而伟大呢？或者说祂因不受恶习玷污而名号伟大超乎万有呢？无论你说什么，无论你在不言的思念中关于天主构想什么，都会流逝并败坏成人的理解，并不能传达其本义，因为它是借我们所使用的、\Emph{只}适合人事的言辞说出的。关于天主的性体，人只能确知一件事：知晓并领悟到，没有任何关于天主的事能以人的语言揭示。\par

20. 那么，这有关形体和性别的事，就是你们——高贵的真理辩护士和虔诚的作者——向你们的神明提出的第一个侮辱。但紧接着的是什么呢？就是你们向我们描述诸神，有的如工匠，有的如医生，有的从事毛纺，有的如水手，有的弹琴吹笛，有的作猎人、牧人，以及，好像再没别的可作了，如同村夫。而那位神，他说，是音乐家，另一位能占卜；因为其他神不能，也不知如何预言将要发生的事，原因是他们缺乏技巧，对将来愚昧无知。一位授以产科之术，另一位则受过医学训练。那么，是否每位神只在自己的领域有能力，而在所求之事属于另一位时，他们便无法提供帮助？这一位口齿伶俐，善于连缀字句；而其他的则愚拙，若要说话，便说不出任何巧妙之言。\par

21. 我还要问，有什么理由、什么不可避免的必然性、什么场合，使得神明们要像卑微的工匠那样了解并熟悉这些手艺呢？因为，难道天上要唱歌奏乐，以便\OriginalTerm{缪斯}{the nine sisters}优雅地协调并和谐音调的休止和节奏吗？难道星辰之山上有森林、树林、小丛林，以便\Person{狄安娜}能在狩猎远征中显得极为强大吗？难道诸神对不久的将来一无所知，并按照命运分配给他们的签运生活度日，以致受灵感的\Person{拉托纳}之子可以解释并宣表明日或下一个时辰给各人带来什么吗？难道他自己也受另一位神灵感，受更大神力的催促和激发，以致可以正确地被称为并被视为附有神灵吗？难道诸神也会染疾，会受伤受害，以致必要时\Place{埃皮达鲁斯}的那位可以前来援助他们吗？他们是否生产、是否分娩，以致\Person{朱诺}可以抚慰，\Person{露西娜}可以缩短可怕的产痛？他们是否从事农耕，或关心战争事务，以致火之主\Person{伏尔甘}可以为他们打造刀剑，或锻造农具？他们是否需要披上长袍，以致\OriginalTerm{特里同少女}{the Tritonian maid}可以以精巧的技艺为他们纺纱织布，并按季节制作合身的袍服，或为三股交织，或为丝绸织物？他们是否进行控告并驳斥，以致\Person{阿特拉斯}的后裔可以因勤勉练习而夺得雄辩的奖赏？\par

22. \Emph{我的对手说}：\Quote{你们错了，受了欺骗；因为诸神本身不是工匠，而是将这些技艺暗示给聪明的人，教导凡人他们该知道的事，使他们的生活方式更文明。} 但凡是将指导给予无知而不情愿的人，并努力使他在某种工作上变得聪明熟练的人，他自己必须先懂得他让别人去实践的那件事。因为没有人能教导一门学问而不通晓他所教导之事的规则，不曾最透彻地掌握其方法。诸神就此成了最初的工匠；无论这是因为他们以知识灌注\Emph{人的心灵}，如你们自己所说，还是因为他们作为不死不生的存在，以其生命之长超越整个地上的族类。那么，这就是问题所在：诸神之间既没有这些技艺的需求，他们的需要和本性也不要求他们有丝毫的巧思或机械技巧，你们凭什么说他们各有精通，一位擅长此艺，另一位擅长彼艺，而且个别人物在特定的部门特别杰出，因精通各门学问而著称呢？\par

但你们也许会说，诸神并非工匠，而是执掌这些技艺，\Emph{且}加以监督；不，一切由我们经营实施的事都置于他们的照顾之下，他们的眷顾看顾其美满顺利的结局。倘若我们在人类事务中所从事、所施行、所尝试的一切，均如我们所愿、所期，那么这话听起来倒真有几分道理与可信。可是，既然事与愿违乃每日所见，行动的成果与意志的意图不相符合，那么说我们有出于迷信幻想而臆造、而非确凿把握的神作为守护者，就是荒谬的。\OriginalTerm{波尔图努斯}{Portunus}赐予水手在海上航行的绝对安全；但为何怒海仍抛出那么多惨遭粉碎的残骸？\OriginalTerm{孔苏斯}{Consus}向我们的心灵提示稳妥有益的路线；为何事态总出乎意料地逆转，结果每与预期相悖？\OriginalTerm{帕勒斯}{Pales}和\OriginalTerm{伊努斯}{Inuus}被立为牲畜的守护者；为何他们以有害的怠惰，不思保护夏日牧场上免受残忍、传染且毁灭性的疾病？娼妓\OriginalTerm{佛洛拉}{Flora}，在淫秽游戏中受敬拜，确保田野繁花盛开；为何嫩芽与幼苗每日被最有害的霜冻摧残殆尽？\OriginalTerm{朱诺}{Juno}掌管分娩，帮助临产的母亲；为何每天有成千的母亲在致命的产痛中丧生？火在\OriginalTerm{伏尔甘}{Vulcan}的看护下，火源受他控制；为何他常常听任神庙和城市街区被火焰吞噬，化为灰烬？占卜者从\OriginalTerm{皮提亚神}{Pythian god}领受其术的知识；为何他经常给予模棱两可、充满疑窦、陷于幽暗蒙昧的回答？\OriginalTerm{埃斯库拉庇乌斯}{Aesculapius}掌管医药的职责与技艺；为何人们不能在种种疾病中恢复健康与身体健全？相反，在医生手下反而更糟。\OriginalTerm{墨丘利}{Mercury}致力于战斗，掌管拳击与角力比赛；为何不使他护佑下的所有人所向无敌？既然被指派同一职司，为何让某些人得胜，却任其他人因可耻的软弱而受嘲笑？\par

我的对手说，没有人向\OriginalTerm{守护神}{tutelar deities}祈求，因此他们收回惯常的恩惠与援助。难道神除非得到熏香和祝圣的祭品，就不能行善吗？难道他们若非见到祭坛被牛血膏抹，就放弃职责吗？而我刚才还以为神的仁慈是出于自由意志，仁爱之礼不求自来呢。难道宇宙的君王，须凭奠酒或祭献才能将生命诸般安慰赐予人类各族吗？天主岂非将太阳育生的温暖、夜晚的时令、风、雨、果实，一律赐予善人与恶人、不义与义人、自由人与奴隶、贫与富吗？因为此乃真实而大能的天主之属性：不待人求，便将仁慈施予困倦软弱、恒常为诸般苦难所困者。因为凭祭品应允祈求，并非向求者施助，而是在兜售其恩惠的财宝。我们人类在如此大事上虚掷光阴，行事愚昧；并且忘却天主是什么、祂圣名的尊威，而将病态轻信所能想象的任何贱劣之事，都归于守护神。\par

\Emph{我的对手说}，\OriginalTerm{恩克西亚}{Unxia}掌管\Emph{门框}的膏抹；\OriginalTerm{辛克西亚}{Cinxia}掌管腰带的解开；至尊贵的\OriginalTerm{维克塔}{Victa}和\OriginalTerm{波图亚}{Potua}司掌饮食。何等稀罕而令人赞叹的对神圣大能的解释啊！若无新妇用油腻的香膏涂抹丈夫的门框；若非丈夫急切近前时解开少女腰带；若人不吃不喝，诸神难道就没有名号了吗？再说，你们不满足于将诸神牵入如此不雅的挂虑中，还赋予他们凶残、冷酷、野蛮的性情，总是因人类的苦难和毁灭而欢欣。\par

我们在此不提盗贼女神\OriginalTerm{拉维尔娜}{Laverna}，以及\OriginalTerm{贝罗纳}{Bellonae}、\OriginalTerm{狄斯科耳狄亚}{Discordiae}、\OriginalTerm{孚里埃}{Furiae}，并且完全缄默不言你们所立那些不祥的神明。我们要提出\OriginalTerm{玛尔斯}{Mars}本人，以及美丽的\OriginalTerm{欲望之母}{mother of the Desires}；你们将战争托付于前者，将爱恋与热切欲望托付于后者。我的对手说，玛尔斯掌管战争；他平息正炽的战火，还是重燃已歇的战事，抑或在和平期间挑起战争？如果他主张战争的疯狂，为何战争每日肆虐？如果他是战争的制造者，那么我们就要说，这神明为了满足自己的倾向，使全世界卷入争斗；在彼此远隔的民族间播下不和与分歧的种子；从各处召集成千上万的人，旋即用死尸堆满原野；使血流成河，将根基稳固的帝国扫荡殆尽，使城邑化为尘土，剥夺自由人的自由，使之沦为奴隶；乐于内乱，乐于兄弟在冲突中流血而亡，最终更乐于子女与父亲间那惨烈、谋杀的搏斗。\par

27. 我们也可以完全相同的方式将这一论证应用于\Person{维纳斯}。因为如果像你们所主张和相信的那样，她使人的心灵充满淫欲的念头，那么随之而来的结论必然是，由这种疯狂所产生的任何耻辱和恶行，都应归咎于\Person{维纳斯}的煽动。那么，难道是在这位女神的强迫下，高贵者也常常将自己的名誉出卖给毫无价值的妓女？难道婚姻的牢固纽带被破坏？近亲因乱伦的欲望而燃烧？母亲对孩子燃起疯狂的激情？父亲把女儿的欲望转向自己？老年人令白发蒙羞，却以少年的热情叹息着要满足污秽的欲望？明智而勇敢的人，在柔弱中丧失了男子气概的力量，漠视坚贞的命令？套索缠绕着他们的脖子？燃烧的柴堆被登上？在不同地方，人们自愿纵身跳下极高极深的悬崖？\par

28. 难道有人，哪怕接受了理性最基本的原则，竟会以如此可耻的道德来玷污或损毁那不变的神性吗？竟然将那些连人类的仁慈也常常在田野的走兽中消除和节制的性情归于诸神吗？我要问，怎能说诸神远离任何激情？说他们温柔、爱好和平、温良？说他们在卓越的完全中达到了完美的顶峰，也达到了最高的智慧？或者，如果我们发现他们正是我们日日所受的一切灾祸的制造者，我们又为什么要祈求他们替我们消除不幸与灾难呢？你们尽可称我们为不虔敬，藐视宗教者，或无神论者，你们永远无法使我们相信存在爱神与战神，无法相信有神散播纷争，并以复仇女神的刺螫扰乱心灵。因为要么他们确实是神，而且不做你们所说的那些事；要么如果他们做了你们所说的事，那么毫无疑问，他们\Emph{根本}不是神。\par

29. 然而，若不是你们自己提出了许多关于诸神的说法，如此自相矛盾且彼此摧毁，迫使我们不敢轻信，我们或许还能接受你们这些充满邪恶谎言的想法。因为当你们各自努力以更深奥的知识在声望上超越彼此时，你们既推翻了你们所信奉的那些神，又用明显不存在的其他神来取代他们；对同样的问题，不同的人给出不同的意见，而你们又写道，那些众人一直公认是单一个体的神，其实是无限的。那么，就让我们从父神\Person{雅努斯}开始履行职责吧，你们中有些人宣称他是世界，有些人宣称他是年份，有些人则宣称他是太阳。但如果我们相信这是真的，那么自然就应该明白，\Person{雅努斯}从未存在过；人们说，他是由\Person{科卢斯}和\Person{赫卡忒}所生，首先在\Place{意大利}为王，建立了\Place{雅尼库卢姆}城，是\Person{福耳斯}的父亲，\Person{武尔图努斯}的女婿，\Person{尤图尔纳}的丈夫；这样你们就抹去了那位在一切祈祷中你们置于首位、并相信他能替你们从诸神那里求得俯听的神的名字。但是，再说，如果\Person{雅努斯}是年，那么他也不能是神。因为谁不知道年是固定的时间段，而由月份的持续和日子的流逝所形成的东西，根本就没有任何神性呢？现在，这同一个\Emph{论证}同样可以应用于\Person{萨图恩}。因为如果在\Person{萨图恩}这一名号下指的是时间，正如希腊思想的解经家们所认为的那样，以至于被认为是\OriginalTerm{克洛诺斯}{Kronos}，其实就是\OriginalTerm{时间}{chronos}，那么就没有\Person{萨图恩}这样的神。因为有谁会如此愚蠢，竟说时间是神，而它不过是在永恒的无尽延续中量出的一段特定空间呢？这样，那位神也将被从不朽者的行列中移除，古人宣布并传给后代说，他由父亲\Person{科卢斯}所生，是大神们的始祖，是葡萄树的种植者，是修剪刀的携带者。\par

30. 但是对于\Person{朱庇特}本人，我们又能说什么呢？智慧者一再断言他就是太阳，驾着有翼的战车，由一群神祗随行；有些人断言他是以太，燃烧着巨大的火焰和不能熄灭的烈火吗？如果这是清楚而确定的，那么按照你们的说法，就根本没有\Person{朱庇特}了；据说他是由父亲\Person{萨图恩}和母亲\Person{俄普斯}所生，为了逃避父亲的愤怒，曾躲藏在\Place{克里特}地区。但是现在，一种类似的思考方式不是也将\Person{朱诺}从神的名单上移除了吗？因为如果她就是空气，正如你们习惯于戏谑而说的那样，把希腊名字的\Emph{音节}颠倒过来念，那么就找不到全能\Person{朱庇特}的妹妹和配偶，就没有\Person{弗路俄尼亚}、没有\Person{波摩娜}、没有\Person{俄西帕吉娜}、没有\Person{费布鲁提斯}、\Person{波普洛尼亚}、\Person{辛克西亚}、\Person{卡普罗提娜}；这样，那个凭借频繁却虚妄的信仰而流传开来的名字的虚构，将变得毫无意义。\par

31. \Person{亚里士多德}，一位智力极强、学问卓越的人，正如\Person{格拉尼乌斯}所说，用看似合理的论证表明\Person{密涅瓦}是月亮，并以博学者的权威加以证明。其他人则说，这同一位女神是深处以太，是最高之处；有些人\Emph{主张}她就是记忆，因此她的名字\Person{密涅瓦}便由此而生，仿佛她是什么记忆女神。但如果相信这一点，那么就没有\Quote{心智}的女儿，没有\Quote{胜利}的女儿，没有从\Person{朱庇特}头颅中诞生的橄榄的发现者，没有精通各种技艺和不同学问分支的\Emph{女神}了。他们说，\Person{尼普顿}之所以获得他的名字和称号，是因为他用水覆盖大地。那么，如果使用这个名字是指遍布的水，就根本没有\Person{尼普顿}这个神；这样，那位\Person{普路托}和\Person{朱庇特}的同父兄弟，手持铁三叉戟的大鱼小鱼之主，深海之王，震撼大地的神，就被撇开，并被从\Emph{我们这里}挪去了。\par

32. \OriginalTerm{墨丘利}{Mercury}也被如此命名，仿佛他是某种中介；因为对话在两人之间传递，并通过他们交换，这个名称所表达的东西就被产生了。如果是这样，\OriginalTerm{墨丘利}{Mercury}并非一位神的名字，而是言语和\Emph{两人之间}交换的话语；因此，著名的\OriginalTerm{基利尼}{Cyllenian}的\OriginalTerm{盘蛇杖}{caduceus}持有者，生于寒冷山顶，制造话语和名称的\Emph{神}，以及主管市场和货物交易与商业往来的神，就被抹去和消灭了。你们有些人说大地是\OriginalTerm{伟大母亲}{Great Mother}，因为它为众生提供食物；另一些人宣称同一\Emph{大地}是\OriginalTerm{刻瑞斯}{Ceres}，因为它产出有用果实的庄稼；而有些人主张它是\OriginalTerm{维斯塔}{Vesta}，因为宇宙中唯独它静止不动，其他部分按其构造则永远运动。现在，如果这是基于可靠的理由提出并坚持的，照你们的解释，三位神就都不存在：刻瑞斯和维斯塔都不应计入诸神之中；最后，众神之母本人，\OriginalTerm{尼基狄乌斯}{Nigidius}认为她曾嫁给\OriginalTerm{萨图恩}{Saturn}，也不能正确地被宣称为女神，如果这些确实都是同一大地的名字，仅由这些称号来指代它的话。\par

33. 我们在此略过\OriginalTerm{伏尔甘}{Vulcan}不表，以免冗长；你们众口一词地声称他是火。\Emph{我们略过}\OriginalTerm{维纳斯}{Venus}，因其得名乃因\Emph{情欲}侵袭万物；又略过\OriginalTerm{普洛塞庇娜}{Proserpina}，因其得名乃因植物逐渐萌发进入光明——由此，你们再次废除了三位神；如果第一位是元素之名，并不表示生命的能力；第二位是众生共有的欲望；而第三位指向种子破土而出，以及生长中的庄稼向上运动。什么！当你们主张\OriginalTerm{巴克斯}{Bacchus}、\OriginalTerm{阿波罗}{Apollo}、\OriginalTerm{太阳}{Sun}是一位神，因使用三个名字而增加数目时，难道在你们的意见中，神的数目未被减少，他们被吹嘘的声望未被推翻吗？因为如果太阳也是巴克斯和阿波罗，那么宇宙中就既无阿波罗也无巴克斯了；于是，\OriginalTerm{塞墨勒}{Semele}之子\Emph{和}\OriginalTerm{皮提亚}{Pythian}的神就被你们抹去\Emph{并}撇弃了——一位是狂欢醉乐之赐予者，另一位是\OriginalTerm{斯敏托斯}{Sminthian}鼠害的毁灭者。\par

34. 你们的一些有学问的人——而且不是\Emph{仅仅}因兴之所至而喋喋不休的人——主张\OriginalTerm{狄安娜}{Diana}、\OriginalTerm{刻瑞斯}{Ceres}、\OriginalTerm{卢娜}{Luna}只是一个三重结合的神；并且不存在三个不同位格，正如有三个不同名字；在所有这些名号中，卢娜被呼求，其他名字则是附加于她名号的一系列绰号。但如果这确实无疑，事实也证明如此，那么刻瑞斯又只是一个空名，狄安娜也一样：于是讨论得出这样的结论：你们引导并建议我们相信，你们认为是地上果实发现者的那位并不存在，而\OriginalTerm{阿波罗}{Apollo}也被夺去了他的妹妹，那位曾有带角的猎人窥见她在一池塘中洗净不洁的肢体，并因他的好奇付出了代价。\par

35. 那些在哲学研究中值得铭记的人，被你们的赞美抬举到最高地位的人，以值得称赞的诚挚，宣布他们的结论说：整个世界，我们都被它的层层包围、覆盖和支撑，是一个具有智慧和理性的动物；然而，如果这是一个真实、可靠且确定的意见，那么你们不久前在其各个部分中无需改名就确立的那些神，也将立刻不再是神。因为正如一个人，在其身体保持完整时，不能被分成许多人；许多人，当他们彼此分离且各不相关时，也不能融合成一个有感觉的个体：所以，如果世界是一个单一的动物，并且由于一个心智的冲动而运动，它就不能分散在多个神之中；如果神祇是它的部分，它们也不能被聚拢并转变为一个活的生物，在其所有部分中具有统一的感觉。月亮、太阳、大地、\OriginalTerm{以太}{Aether}、星辰，都是世界的肢体和部分；但如果它们是部分和肢体，它们本身肯定不是活的生物；因为在任何事物中，部分都不能是整体本身，也不能为自己思考或感觉，因为这不能由它们自己的行动达成，而需要有整个生物的参与；这一点确立并确定后，整个问题归结为：\OriginalTerm{索尔}{Sol}、\OriginalTerm{卢娜}{Luna}、\Emph{以太}、\OriginalTerm{忒勒斯}{Tellus}等，都不是神。因为它们是世界的部分，而非神祇的专名；由此，通过你们对一切神圣事物的扰乱和混淆，世界被树立为宇宙中唯一的神，而其余一切都遭到弃置，且是徒然、无用且毫无真实性地被树立起来。\par

36. 如果我们试图以这么多方式、用这么多论据来颠覆对你们神祇的信仰，没有人会怀疑，你们必会疯狂暴怒，要求用火刑、野兽和刀剑，以及你们惯常用来满足对鲜血极度渴望的其他种种酷刑来对付我们。但当你们自己以聪明和智慧为幌子，几乎摒弃了全体神祇时，你们却毫不犹豫地断言，由于我们，人类遭到了神祇的恶劣对待；尽管事实上，如果它们确实存在于某处，且怒火中烧，那么它们对你们发怒的最好理由，莫过于你们否认它们的存在，并\Emph{声称}它们在宇宙的任何部分都\Emph{找不到}。\par

37. 据\OriginalTerm{姆纳塞阿斯}{Mnaseas}告诉我们，\OriginalTerm{缪斯们}{Muses}是\OriginalTerm{忒勒斯}{Tellus}与\OriginalTerm{科鲁斯}{Coelus}的女儿；另一些人宣称她们是\OriginalTerm{朱庇特}{Jove}与其妻\OriginalTerm{记忆女神}{Memory}或\OriginalTerm{心智女神}{Mens}所生；有人说她们是处女，也有人说她们是已婚妇女。因为我们现在只想简要涉及那些你们因意见分歧而对同一事物做出不同陈述的地方。\OriginalTerm{埃福罗斯}{Ephorus}说她们是三数；我们提到的姆纳塞阿斯说\Emph{她们是四数}；\OriginalTerm{密尔提罗斯}{Myrtilus}提出七数；\OriginalTerm{克拉底斯}{Crates}断言是八数；最后，\OriginalTerm{赫西俄德}{Hesiod}用神祇充斥天地星辰，推出九个名字。\par

如果我们没有弄错，这种不一致标志着那些对真理全然无知的人，并非源于事实的真实状况。因为如果他们的数目清楚已知，所有人的声音就会相同，所有人的一致将会趋向并归结为相同的结论。\par

38. 那么，当你们对众神本身充满错误认识时，又如何能赋予宗教其全部力量呢？或者召唤我们到他们庄严的崇拜中，而你们却没有给我们明确的信息去理解这些神明本身？且不说其他作者，要么第一个人去掉并除灭了六位神圣的\OriginalTerm{缪斯}{Muses}——如果她们确实是九位；要么最后一个人给仅真实存在的三位缪斯添加了六个并不存在的；这样，就无法知道或理解应当添加什么、去掉什么；在举行宗教仪式时，我们便面临危险：要么敬拜并不存在之物，要么忽略了可能存在之物。\Person{皮索}相信\OriginalTerm{诺文西勒斯}{Novensiles}是在\Place{特雷比亚}的萨宾人中设立的九位神。\Person{格拉尼乌斯}与\Person{埃利乌斯}一致，认为他们是缪斯；\Person{瓦罗}教导说他们是九位，因为在做任何事情时，\Emph{那个数目}总是被认为最强大、最伟大；\Person{科尼菲基乌斯}说他们掌管万物的更兴，因为藉由他们的照管，万物在力量上得以更新并持续存在；\Person{曼尼利乌斯}说他们是那九位神，\OriginalTerm{朱庇特}{Jupiter}只把运用他雷电的力量赐给了他们。\Person{钦基乌斯}断言他们是从外地引入的神明，因其新来而得名，因为罗马人习惯有时私下将被征服城市的仪式引入自己家族，有些则公开加以祝圣；并且，为了不致因其数目众多或因无知而遗漏任何神，他们都同样被简短而概括地以一个名字——诺文西勒斯——来呼求。\par

39. 此外，还有一些人断言，那些从人变成神的就是用这个名字来指称的——如\OriginalTerm{赫丘利}{Hercules}、\Person{罗慕路斯}、\Emph{\OriginalTerm{埃斯库拉庇乌斯}{Aeculapius}}、\OriginalTerm{利柏尔}{Liber}、\Emph{\OriginalTerm{埃涅阿斯}{Aeneas}}。显而易见，这些全都是不同的意见；而且在事物的本性上，那些意见不同的人不可能被视为同一真理的教导者。因为如果\Person{皮索}的意见是真实的，\Emph{\Person{埃利乌斯}}和\Person{格拉尼乌斯}说的就是假的；如果他们说的一定是真实的，那么\Person{瓦罗}尽管有各种技巧也错了，他用最琐碎虚无的东西取代了真实存在的事物。如果他们被称为诺文西勒斯是因为数目是九，那么\Person{科尼菲基乌斯}就被证明出了差错，因为他赋予他们本不属于他们的力量和权能，使他们成为更兴神灵的监督。但如果科尼菲基乌斯的见解是正确的，\Person{钦基乌斯}就被发现并非明智，因为他把被征服城市的神灵与\Emph{\OriginalTerm{诺文西勒斯诸神}{dii Novensiles}}的能力联系起来。但如果他们是钦基乌斯所断言\Emph{那样子}，那么\Person{曼尼利乌斯}就被发现说了假话，他把那些掌握别人的雷电者归于此名之下。但如果曼尼利乌斯所持的观点是真实确凿的，那些设想被提升到神圣尊荣、被神化的凡人因位格的新奇而被称为\Emph{这样}的人，就完全错了。但如果诺文西勒斯是那些度过人生之后堪升为星辰的人，那么就根本不存在\Emph{诺文西勒斯诸神}。因为，如同\Quote{奴隶}、\Quote{士兵}、\Quote{主人}并非包含于其下的个人名字，而是官职、等级和职守的名称，同样，当我们说诺文西勒斯是因德行而从人变为神的神的名字时，清楚显明并没有特别指出来一个具体个人，而只是新奇本身以诺文西勒斯这一称号来称谓。\par

40. \Person{尼吉狄乌斯}教导说，\Emph{\OriginalTerm{佩纳特斯诸神}{dii Penates}}是\OriginalTerm{涅普顿}{Neptune}与\OriginalTerm{阿波罗}{Apollo}，他们曾经按约定条件为\Place{伊利乌姆}筑起城墙。他本人在其第十六卷书中，遵循伊特鲁里亚人\Emph{的教导}，又指出有四种佩纳特斯；其中一种属于朱庇特，另一种属于涅普顿，第三种归于下界的幽影，第四种归于必死之人，他作出了某种不可理解的断言。\Person{凯西乌斯}本人也遵循这一教导，认为他们是\OriginalTerm{福尔图娜}{Fortune}、\OriginalTerm{刻瑞斯}{Ceres}、\OriginalTerm{朱庇特之灵}{genius Jovialis}和\OriginalTerm{帕勒斯}{Pales}，但不是通常所接受的女\Emph{神}，而是某位男性的侍从、朱庇特的管家。\Person{瓦罗}认为他们是我们所说的那些在内、并居于天之最深处的神，而且他们的数目和名字都不为人所知。伊特鲁里亚人说这些是\Emph{\OriginalTerm{孔森特斯与孔普利刻斯}{Consentes and Complices}}，如此称呼他们是因为他们一同起落，其中有六位男性，同样数目的女性，其名不为人知，性情冷酷，但他们被视为至高的\OriginalTerm{朱庇特}{Jove}的顾问和君主。也有一些人说，朱庇特、\OriginalTerm{朱诺}{Juno}和\OriginalTerm{弥涅耳瓦}{Minerva}就是\Emph{佩纳特斯诸神}，没有他们，我们便不能生活、不能明智，我们是在理性、情感和思想中由他们内在统御。如你所见，即便在此，也没有任何和谐一致的言说，没有什么是经人人同意而确定的，也没有什么可靠之事可使心灵立足，借着推测几乎接近真理。因为他们的意见如此悬疑，一个猜测被另一个如此推翻，以致要么其中全无真理，要么若有人把它说出来，也在如此众多不同的陈述中无从辨认。\par

41. 如果认为合适，我们也可以简短地谈谈拉瑞斯，大众认为他们是街衢与道路之神，因为希腊人称街道为\Emph{\OriginalTerm{劳拉埃}{laurae}}。\Person{尼吉狄乌斯}在他的著作中不同地方，\Emph{谈到他们}，时而说他们是房屋与住宅的守护者；时而说是\OriginalTerm{库瑞忒斯}{Curetes}，据说他们曾以铙钹的撞击声掩盖了朱庇特的婴啼；时而又说是萨莫色雷斯的五个指头，希腊人\Emph{告诉我们}他们被称作\Emph{\OriginalTerm{伊达山的达克堤利}{Idaei Dactyli}}。\Person{瓦罗}同样犹豫不定，一次说他们是\OriginalTerm{玛内斯}{Manes}，因此拉瑞斯之母名为\OriginalTerm{玛尼亚}{Mania}；又一次，他断定他们是空际之神，称作英雄；再一次，他依随古人的意见，说拉瑞斯是幽灵，好似一种守护灵，是死人的魂魄。\par

42. 逐个分别查考每一种类，并从你们的宗教典籍中显明，你们既不持有也不相信任何一位神明的存在，除非你们对这位神明提出了可疑且不一致的论述，表达出无数种不同的信仰，这乃是一项浩大无尽的工作。然而，为求简洁，避免冗长，刚才所说已经足够；再者，将许多事集合一处也太过繁琐，因为已经以不同方式显明并证实：你们摇摆不定，对于自己所断言的事，从未说过任何确定的话。但你们或许会说，即便我们并不亲自认识\OriginalTerm{拉尔}{Lares}、\OriginalTerm{诺文西勒斯}{Novensiles}、\OriginalTerm{珀那忒斯}{Penates}，可我们众多作者的普遍共识却证明了他们的存在，且这类神祇位列于属天诸神之中。然而，如果某一位神是什么完全未知，又怎能知道是否有这样一位神呢？如果每次求问时，该呼求哪一位尚未确定和指明，求取恩惠又有何用呢？因为每一个寻求从任何神明获得答复的人，都必须知道自己是向谁祈求，呼求谁，为人生事务与际遇向谁求助；尤其你们自己也宣称，并非所有神明拥有一切权能，且每位的震怒与忿恨需以不同仪式平息。\par

43. 因为，\Emph{如果}这位\Emph{神}要求黑色，那位要求白色皮肤；\Emph{如果}向这位献祭时必须蒙首，向那位则必须露首；这位受人咨询婚姻，那位则解除困苦——那么，知晓某个是\OriginalTerm{诺文西勒斯}{Novensills}还是另一个，难道不是有些重要吗？因为对事实的无知和人称的混淆会触怒众神，并必然招致罪咎。再假设我本人，为避某些不便与危险，向这些神祇中的任何一位呼求，说：\Quote{莅临吧，亲近吧，神圣的\OriginalTerm{珀那忒斯}{Penates}，你，阿波罗，还有你，哦，尼普顿，以你们神圣的仁慈，转离所有那些令我烦恼、困扰、折磨的灾祸：}那么，倘若刻瑞斯、帕勒斯、福尔图娜，或\Emph{\OriginalTerm{尤威守护灵}{genius Jovialis}}，而非尼普顿和阿波罗，是\Emph{\OriginalTerm{珀那忒斯众神}{dii Penates}}，我还能有得蒙他们帮助的任何希望吗？或者，如果我呼求\OriginalTerm{库瑞忒斯}{Curetes}而非\OriginalTerm{拉尔}{Lares}——你们某些作者坚持认为库瑞忒斯就是\Emph{\OriginalTerm{萨莫色雷斯的迪吉提}{Digiti Samothracii}}——那么，我既未以他们自己的名字称呼他们，却以不属于他们的名字称呼了其他神祇，我又怎能享有他们的帮助与恩惠呢？因此，我们的利益要求我们正确认识众神，不可对每位的神能、名号犹豫或怀疑；免得当他们被以不属于他们的仪式和名号呼求时，他们立即\Emph{掩耳不听我们的祈祷}，并使我们陷于不可赦免的罪咎之中。\par

44. 所以，如果你们确信，在崇伟的天宫之中居住着、存在着你们所列举的那一大群神祇，你们就当立定心志于一个确定的见解之上，而不可因分持不同且互相矛盾的看法，反倒摧毁了你们试图确立的那些事本身中的信仰。如果存在一位\OriginalTerm{雅努斯}{Janus}，就让雅努斯存在；如果存在一位\OriginalTerm{巴库斯}{Bacchus}，就让巴库斯存在；如果存在一位\OriginalTerm{苏曼努斯}{Summanus}，就让苏曼努斯存在：因为这才是信靠，这才是持守，才是安居于确切知识的把握之中，而不是像瞎眼迷途之人那样说：\OriginalTerm{诺文西勒斯}{Novensiles}就是\OriginalTerm{缪斯}{Muses}，其实他们是\OriginalTerm{特雷比亚诸神}{Trebian gods}，不对，他们的数目是九，更确切说，他们是已倾覆之城的保护者；并将如此重大的事陷入这般险境：以致当你们移除某些神祇，又另立其他取代其位时，完全可以怀疑他们全体是否在任何地方存在。\par

\SourceRecord{Translated by Hamilton Bryce and Hugh Campbell.；Ante-Nicene Fathers；Vol. 6.；Edited by Alexander Roberts, James Donaldson, and A. Cleveland Coxe.；Buffalo, NY: Christian Literature Publishing Co.,；1886.；<http://www.newadvent.org/fathers/06313.htm>.}

% structure: p0001 source=h1 target=section reason=page-title

\section{\WorkTitle{驳异教徒（第四卷）}}

1. 我们问你们，尤其是你们，啊，罗马人，世界的君主与元首，你们是否认为\OriginalTerm{虔敬}{Piety}、\OriginalTerm{和谐}{Concord}、\OriginalTerm{安全}{Safety}、\OriginalTerm{荣誉}{Honour}、\OriginalTerm{美德}{Virtue}、\OriginalTerm{幸福}{Happiness}以及其他这类名号（我们见到你们为它们兴建祭台和辉煌的庙宇）拥有神圣力量，居住在天上？或者，如常见的那样，你们是否仅仅出于形式而将它们列入诸神，因为我们都愿望并祈盼这些福祉降临于我们？因为，如果你们一面认为它们是空洞无实的虚名，一面却又以神圣的荣耀尊奉它们为神，你们必须考虑那究竟是幼稚的嬉戏，还是会将你们的诸神置于被轻视的境地，因为你们把那些空洞伪造的名号等同起来，加进他们的行列。但如果你们已为它们负载了庙宇和卧榻，较为确信地认为它们也是神明，那么我们祈求你们，指教\Emph{我们这些无知者}，凭借何种方式、以何种途径，\OriginalTerm{胜利}{Victory}、\OriginalTerm{和平}{Peace}、\OriginalTerm{公正}{Equity}以及被列于诸神中的其他名号，才能被理解为神明，属于不朽者的团体？\par

2. 对于我们——不过，或许，你们夺走了并剥夺了我们的常识——我们感觉并觉察到，这些无一拥有神圣力量，或具有本身的形体；但\Emph{相反}，它们是：男子气概的优秀、安全者的安全、受敬者的荣誉、征服者的胜利、盟友的和谐、虔敬者的虔敬、守礼者的忆念，以及那生活幸福、不惹嫉妒者的好运。如今容易看出，我们这样说极有道理，当我们观察\Emph{与它们}相对的相反品质时：厄运、不和、遗忘、不义、不虔、精神的卑下，以及身体的不幸虚弱。因为既然这些事是偶然发生，并依赖于人的行为\Emph{和}偶然的情绪，那么它们的反面，即以更可悦的品质命名的，必存在于其他事物中；而这些被臆造的名号，便是由此而生。\par

3. 至于你们向我们举出的其他一群群未知的神明，我们无法确定你们是郑重行事，出于对确定性的信仰；还是\Emph{仅仅}玩弄虚空的虚构，纵情于毫无节制的想像。你们依据\Person{瓦罗}的权威告诉我们，\OriginalTerm{卢珀卡}{Luperca}女神得名，是因为那头凶猛的狼没有伤害弃婴。那么，那女神不是由她自己的能力，\Emph{而是}由事态的进程，才被揭示出来的吗？并且，是否\Emph{只有}在那野兽收住其残忍的利齿之后，她才开始存在并获得其名号？抑或，如果早在\Person{罗慕路斯}及其兄弟出生很久以前，她就已经是女神，那么请告诉我们她的名号和称谓是什么。据你们说，\OriginalTerm{普勒斯塔娜}{Praestana}得名，是因为在投掷标枪时，\Person{奎里努斯}力压群雄；而女神\OriginalTerm{潘达}{Panda}，或称\OriginalTerm{潘提卡}{Pantica}，得名是因为\Person{提图斯·塔提乌斯}获准开辟并修通一条道路，以便攻取\Place{卡彼托山}。那么，在这些事件之前，这些神明从未存在过吗？而如果\Person{罗慕路斯}在投掷标枪中没有赢得第一，如果那萨宾国王未能攻取\Place{塔耳珀伊亚岩}，便不会有潘提卡，不会有普勒斯塔娜吗？如果你们说她们在那些赋予名号的事件之前就已存在——这个问题在前一部分已经讨论过——那么也请告诉我们她们那时被称作什么。\par

4. \OriginalTerm{佩洛尼亚}{Pellonia}是一位有能力逐退敌人的女神。请问，若方便的话，逐退谁的敌人？两军相遇，短兵相接，决出胜负；对一方，这一边是敌人，对另一方，那一边是敌人。那么，既然双方都要战斗，佩洛尼亚会使哪一方溃逃呢？或者她会倾向哪一方，因为她理应将自己名号的力量与援助给予双方？然而，如果她果真如此行事，那就是说，如果她将自己的善意和恩惠给予双方，她便会毁坏自己名号的含义，因其名号本是就击退一方而形成的。但你们也许会说，她仅仅是罗马人的女神，只站在奎里特人一边，总是乐意施恩帮助他们。我们倒真愿这是事实，因为我们喜欢这个名字；但这却是极为可疑的事。什么！难道罗马人拥有专属自己、不帮助其他民族的诸神吗？如果他们对天下万民不行使其神圣权力，他们又如何能成为神呢？而且，我请问，许久以前，当国家荣誉在\Place{考迪昂岔口}被置于轭下时，这位佩洛尼亚女神在哪里？当\Place{特拉西美诺湖}的溪流被鲜血染红时？当\Place{狄奥墨得斯}的平原上堆满罗马人的尸体，在无数不幸的战斗中遭受千百次其他打击时？她是在打鼾沉睡吗？还是如同卑劣者常做的那样，逃到敌人的营地去了？\par

5. \OriginalTerm{左侧的神祇}{sinister deities}只掌管左手边的区域，与右手边的神祇相对立。但这话出于什么理由，或有什么含义，我们自己都不明白；而我们确信，你们也完全无法使之清晰而普遍地被理解。因为首先，世界本身既无左右，也无上下区域，也无前后\Emph{部分}。因为任何圆形、被坚固球体的圆周从各个方向限定的东西，都没有起始，没有终结；在没有终结和起始之处，没有任何部分可以有自己的名称并形成起始。因此，当我们说\Quote{这是右边，那是左边}时，我们指的并不是世界中的任何东西，因为世界各处都大致相同，而是指我们自身的位置和方位，我们被造成这样：我们说某些东西在右手边，另一些在左手边；然而，正是我们称为左的这些事物，以及\Emph{我们称为右}的那些事物，在我们身上并无恒常，并无固定，而是从我们的侧面获得其形式，全凭机遇和瞬间的偶然如何安置我们。若我面朝初升的太阳，北极和北方就在我的左手边；若我将脸转向那里，西方就会在我的左边，因为它会被视为在太阳的背后。然而，若我再将目光转向西方的区域，南方来风与地域就被说成在我的左边。而若我因眼下必需的事务转向这一侧，结果就是，由于位置的进一步改变，东方被说成\Emph{是}左边——由此极易看出，没有什么东西在本性上是或在我们的右边或在我们的左边，而是出于位置、时间，并依照我们相对于周围事物的身体姿势而定。但在这种情况下，何以能有左侧区域的神祇呢？因为显然，相同的区域时而右侧，时而左侧。或者，右边的区域对不朽诸神做了什么，竟理应无人关照它们，而诸神却命定这些区域幸运，且永远有吉祥之兆\Emph{伴随}？\par

6. 如你们所说，\OriginalTerm{拉泰拉努斯}{Lateranus}是灶神与灶的\OriginalTerm{守护神}{genius}，他得此名是因为人们用未烧制的砖建造那种火炉。那么，如果灶是用烧制的黏土或其他任何材料建造的，就没有守护神了吗？而拉泰拉努斯，不管他是谁，会因为他的王国不是由黏土砖砌成，就放弃他作为守护者的职责吗？我还要问，那位神究竟为何接受照料灶的职责？他奔忙于人们的厨房，检查并发现他们的火是用哪种木柴产生热量的；他赐予陶器力量，使它们不致因火焰的猛烈而破碎；他确保未变味的美味佳肴以其自身的鲜美抵达味觉，并扮演品尝者的角色，检验酱汁是否调制得当。难道这不是不得体的，不仅如此——更真实地说——可耻的、不虔敬的，仅仅为此而引入某些虚构的神祇，不以合宜的敬礼尊崇他们，却派他们管理卑贱的事务和可耻的行为吗？\par

7. 还有，\OriginalTerm{战士维纳斯}{Venus Militaris}是否也掌管军营的恶行和年轻男子的放荡？众多神祇中是否也有一个\OriginalTerm{佩尔菲卡}{Perfica}，她使那些卑贱肮脏的欢愉以不间断的快乐达到终点？是否也有一个\OriginalTerm{珀尔通达}{Pertunda}，她掌管婚姻之床？是否也有一个\OriginalTerm{图图努斯}{Tutunus}，你们认为在他的巨大肢体和令人惊骇的\OriginalTerm{法西努姆}{fascinum}上，让你们的妇女被背负是吉祥的，并渴望如此？如果事实本身极不足以让你们对真理有正确的理解，难道你们就连从这些名字本身，也无法明白这些都是最无意义的迷信所臆造，是幻想的假神吗？你们说，\OriginalTerm{普塔}{Puta}掌管树木的修剪，\OriginalTerm{佩塔}{Peta}掌管祈祷；\OriginalTerm{涅美斯特里努斯}{Nemestrinus}是丛林之神；\OriginalTerm{帕特拉娜}{Patellana}是一位神祇，而\OriginalTerm{帕特拉}{Patella}是另一位，前者被派去管理已显露的事情，后者管理有待揭示的事情。\OriginalTerm{诺杜提斯}{Nodutis}被视为一位神，因为他使播种的进入结节；那位掌管谷物脱粒的女神是\OriginalTerm{诺杜特伦西斯}{Noduterensis}；女神\OriginalTerm{乌皮比利亚}{Upibilia}使人免于偏离\Emph{正确}的道路；丧子的父母由\OriginalTerm{奥博娜}{Orbona}照管——那些濒临死亡的人则由\OriginalTerm{纳尼亚}{Naenia}照管。再者，\OriginalTerm{奥西拉戈}{Ossilago}本人被提及，\Emph{她}赐予幼儿骨骼的坚固与结实。\OriginalTerm{麦洛尼亚}{Mellonia}是一位女神，在蜜蜂方面强大有力，照管并守护其蜜的甜美。\par

8. 请，我恳求你们——愿佩塔、普塔、帕特拉施恩于你们——假若那时地球上根本没有蜜蜂，或者我们人类像某些蠕虫一样生来没有骨头，难道就不会有女神麦洛尼亚？或者赐予骨骼坚固的奥西拉戈，会没有她自己的名字吗？我真诚地询问，并急切地探问，你们是否认为，神，或人，或蜜蜂、果实、枝条以及其他，在本性、时间、久远上，哪个更古老？无人会怀疑你们说，诸神比所有事物都早无数年代和世代。但如果真是这样，按事物本性，何以可能，从后来产生的事物那里，他们获得了那些在时间上更早的名字？或者，诸神怎会被赋予照料那些尚未产生、被指定为人所用的事物的职责？又或者，诸神是不是长久没有名字；只是在万物开始出现、在地上存在之后，你们才认为他们应以这些名字和称号称呼？而你们又从何得知该给每个神取什么名字，既然你们完全不知道他们的存在？或者，他们拥有\Emph{任何}固定的能力，既然你们同样不知道他们中谁有任何能力，以及为适合其神圣大能，他应被派往何处？\par

9. 那么，又如何呢？你们说；你们是否宣称这些神明在世间无处存在，是由虚妄的幻想造出来的？不仅我们这么说，真理本身、理性，以及人类共有的常识也这么说。因为有谁会相信存在收益之神，并且他们主宰着收益的获得，而收益往往来自最卑鄙的营生，并且总是以牺牲他人为代价？有谁会相信\OriginalTerm{利本蒂娜}{Libentina}或\OriginalTerm{布尔努斯}{Burnus}掌管着那些智慧要求我们避免的\Emph{那些}情欲，而邪恶污秽之徒却以千百种方式去尝试和践行这些情欲？有谁会相信\OriginalTerm{利门提努斯}{Limentinus}和\OriginalTerm{利玛}{Lima}负责守护门槛，并履行看守者的职责，而我们每天却看到庙宇和私人住宅的\Emph{门槛}被毁坏和推倒，而且通往妓院的臭名昭著的入口也并非没有它们？有谁会相信\OriginalTerm{利弥}{Limi}照管着歪斜之物？有谁会相信\OriginalTerm{撒图尔努斯}{Saturnus}掌管着播下的庄稼？有谁会相信\OriginalTerm{蒙提努斯}{Montinus}是山岳的守护神；\OriginalTerm{穆尔西亚}{Murcia}是懒汉的守护神？最后，有谁会相信\OriginalTerm{金钱女神}{Money}是一位女神，正如你们的著作宣称的那样，\Emph{她}似乎是至高无上的神明，能赐予金戒指、竞技和演出的前排座位、大量的荣誉、显赫的官职，以及懒汉们最喜爱的东西——因财富而带来的无忧无虑的安逸。\par

10. 但是如果你们坚持认为，骨头、\Emph{不同种类的}蜂蜜、门槛，以及所有其他我们匆匆一提或为避免冗长而完全略过的事物，都有它们各自的守护神，那么我们同样可以引进成千上万其他神明，来照料和守护无数的事物。因为神明为什么只掌管蜂蜜，而不掌管葫芦、油菜、香草、水芹、无花果、甜菜、卷心菜呢？为什么只有骨头得到保护，而指甲、头发，以及所有那些位于我们羞于启齿的隐藏部位和器官，由于暴露在众多事故中而更需要神明的眷顾和照料的东西，却得不到保护呢？或者说，如果你们声称这些部位也由各自的守护神来照料，那么有多少事物就会有多少神明；而且也无法说明，如果你们说某些事物由神明主宰和照料，为什么神明的眷顾却不保护所有事物。\par

11. 新宗教的创始者们和掌权者们啊，你们有什么话说？你们是否高声抗议，抱怨这些神明受到我们的羞辱，被我们以亵渎的轻蔑所忽视，即：\OriginalTerm{拉特兰努斯}{Lateranus}，炉灶的守护神；\OriginalTerm{利门提努斯}{Limentinus}，门槛的主宰者；\OriginalTerm{佩尔通达}{Pertunda}、\OriginalTerm{佩尔菲卡}{Perfica}、\OriginalTerm{诺杜特伦西斯}{Noduterensis}；你们是否声称，因为我们对\OriginalTerm{穆图努斯}{Mutunus}和\OriginalTerm{图图努斯}{Tutunus}不低首祈求，事物已陷入毁灭，世界本身也改变了它的法则和结构？但现在请你们看看，恐怕当你们想象出如此怪异的事物，形成如此的观念时，你们已经冒犯了那些确实存在的神明——如果确实有任何神明配得上并持有那最崇高的称号的话；而你们所说的那些灾祸，由于没有其他原因，才日益猖獗，与日俱增。那么，或许你们中的某人会说：你们为何断言这些神明确非真实？当被占卜者呼求时，他们难道不听从召唤，应自己的名号而来，并向求问者给出可信的答复吗？我们可以证明这种说法是虚假的，或者因为在整个这件事上，有最大的怀疑空间，或者因为我们每天看到他们的许多预言要么被证实为假，要么期望被挫败，\Emph{去迎合}相反的结果。\par

12. 但即使如你们所坚持的，它们是真实的，你们难道也要我们相信，比如\OriginalTerm{梅洛尼亚}{Mellonia}，将自己引入内脏之中，或者\OriginalTerm{利门提努斯}{Limentinus}，并且他们设法让人知道你们想要了解的事情？你们可曾见过他们的面容，他们的举止，他们的神色？或者这些甚至在肺脏或肝脏中能被看见？难道不可能，不会发生——尽管你们狡猾地掩饰——此神取代彼神，迷惑、嘲弄、欺蒙，并呈现出被呼求的\Emph{神明}的样子？如果\OriginalTerm{术士}{magi}们，\Emph{他们}与占卜者极为相似，叙述说，在他们的咒语中，伪神常常取代那些被呼求的神明而潜入；而且这些中的一些，还是粗俗物质构成的灵体，他们假装是神明，用谎言和欺骗来误导无知者——那么我们为什么不能同样相信，在这里，也有其他灵体取代了那些并不存在的神明，以便既能巩固你们的迷信信仰，又能因以不属于他们的名号向自己献祭而欢欣呢？\par

13. 或者，如果你们因这说法新奇而拒绝相信，那你们如何能知道是否没有某个存在，取代了所有你们呼求的神明，在世界各处顶替，向你们显明看似许多神明和力量的东西？那是谁呢？有人会问。我们或许，受教于真实可靠的作者，能够说出；但是，恐怕你们不愿相信我们，就让我的对手去问埃及人、波斯人、印度人、迦勒底人、亚美尼亚人，以及所有那些在更隐秘的技艺中见识并熟稔这些事的人吧。那时，你们确实将得知，谁是那独一的天主，或者祂之下的众多神明是谁，他们假装成神，并戏弄人的无知。\par

即便是此刻，我们仍羞于触及那一点，那一点不仅让年轻轻佻的男孩，也让庄重的人，乃至那些\Emph{已然变得}严厉冷峻的人，都忍不住发笑。因为我们全都曾听到我们的师长谆谆教导：在神祇\Emph{名字}的变格中，没有复数形式，因为诸神都是个体，每个名字的所有权不能为众多神祇所共有；你们却因忘却，将幼时所学抛诸脑后，既将相同的名字赋予多位神祇，尽管你们在其他地方对神祇数目尚算节制，却又通过共用名字的方式使其增多了；关于这个话题，确实，一些眼光敏锐、才智卓越的人先前已用拉丁文和希腊文论述过了。那本可减少\Emph{我们的劳作}，若非与此同时我们看到有人对这些著作一无所知；此外，我们既已开始的讨论，也迫使我们提出一些关于这些主题的东西，尽管\Emph{已}为那些\Emph{作家}所着手并讲述过了。\par

14. 那么，你们的神学家和那些论述远古未知之事的作者们说，宇宙间有三位\Person{朱庇特}，一位之父为\Person{以太}；另一位之父为\Person{凯路斯}；第三位之父为\Person{萨图恩}，生于\Place{克里特}岛并葬于该岛。他们\Emph{说到}五个太阳和五位\Person{墨丘利}——据他们所述，其中第一位太阳被称为\Person{朱庇特}之子，被视为\Person{以太}之孙；第二位\Emph{也}是\Person{朱庇特}之子，生他的母亲是\Person{许珀里翁娜}；第三位是\Person{武尔坎}之子，不是\Place{利姆诺斯}的\Emph{武尔坎}，而是\Place{尼罗河}之子；第四位，在英雄时代由\Person{阿坎托}在\Place{罗得}岛生下，\Emph{是}\Person{伊阿吕索斯}之父；\Emph{而}第五位被视为一位\Place{斯基泰}国王和狡黠的\Person{喀耳刻}之子。再者，第一位\Person{墨丘利}，据说曾贪恋\Person{普洛塞庇娜}，是\Person{凯路斯}之子，\Emph{他}是至高无上的。第二位在地下，自诩为\Person{特洛福尼乌斯}。第三位\Emph{乃}其母\Person{迈亚}与第三位\Person{朱庇特}所生；第四位是\Place{尼罗河}的后裔，其名\Place{埃及}人畏惧而不敢提及。第五位是\Person{阿尔戈斯}的斩杀者，一个逃亡者和流放者，也是\Place{埃及}文字的发明者。但他们说，也有五位\Person{密涅瓦}，正如\Emph{有五个}太阳和五位\Person{墨丘利}一样；其中第一位并非童贞女，而是\Person{武尔坎}所生\Person{阿波罗}的母亲；第二位是\Place{尼罗河}的后裔，据称是\Place{埃及}的\Person{萨伊斯}；第三位是\Person{萨图恩}的后代，即那位设计使用兵器者；第四位出自\Person{朱庇特}，\Place{墨塞尼亚}人称她为\Person{科律法西亚}；第五位则是杀死其好色父亲\Person{帕拉斯}的那个。\par

15. 而为了避免逐一细说显得冗长乏味，同一批神学家说还有四位\Person{武尔坎}和三位\Person{狄安娜}，同样数目的\Person{埃斯库拉庇俄斯}和五位\Person{狄俄倪索斯}，六位\Person{赫丘利}和四位\Person{维纳斯}，三组\Person{卡斯托尔}以及同样数目的\Person{缪斯}，三位有翼的\Person{丘比特}，和四位被冠以\Person{阿波罗}之名的神；他们同样地提及诸神的父亲、母亲\Emph{及}出生地，并指出每一位的起源与家族。然而，若这一切是真实确凿的，且被认真当作\Emph{尽人皆知}之事，那么要么它们并不全都是神，因为如我们被告知的那样，同一名下不能有多个存在；要么其中若有一位是真神，他将无法被认识和辨认，因为极为相似的名字造成的混淆使他变得模糊不清。如此，从你们自己的行为中产生的后果——无论你们多么不愿意——便是宗教陷入了困境与混乱，没有任何可以依止的固定目标，只能任由模棱两可的幻象所戏弄。\par

16. 设想我们受你们适宜的影响或强暴的恐惧所驱，前来朝拜\Person{密涅瓦}，例如，以你们视为神圣的仪节与惯常的典礼：如果当我们预备祭品，走近为她指定的燃烧的祭台前献上\Emph{祭品}时，所有的\Person{密涅瓦}都飞聚而来，为那名号的权利而争斗，每一位都要求将预备的祭品归于自己；我们该把哪一只作祭牲的动物放在她们中间？或者，我们应向谁尽当行的圣职？因为我们所说的第一位或许会说：\Quote{\Person{密涅瓦}这名号是我的，神圣的威严也是我的，我生了\Person{阿波罗}和\Person{狄安娜}，以我胎中的果实丰富了天庭，增添了诸神的数目。} 第五位会说：\Quote{不，\Person{密涅瓦}，你是在说话吗？你身为妻子，且屡为母亲，已失去了无玷纯净的圣洁。你岂不见在所有庙宇中，\Person{密涅瓦}的肖像皆是童贞女，一切艺人都避免赋予她们主妇的体态？因此，停息吧，别再擅取不属于你的名号。因为我是\Person{密涅瓦}，由父\Person{帕拉斯}所生，全体诗人都为我作证，称我为\Person{帕拉斯}，这别号源自我的父亲。} 第二位听闻此言将大声说：\Quote{你说什么？你，一个无耻的弑亲者，一个被淫乱欲火污染的人，竟敢号称\Person{密涅瓦}？你以胭脂和娼妓的诡计装饰自己，激起你父亲满溢狂欲的情欲。那么，再往前去吧，为自己另找一个名号；因为这名号属于我，是众河之首的\Place{尼罗河}从流水之中生了我，由水汽凝聚而使我有处女之身。但若你询问此事的可信性，我也将引出埃及人作证，在他们的语言中我被称为\Person{奈特}，正如\Person{柏拉图}的\WorkTitle{蒂迈欧篇}所证实的。} 那么，我们料想结果会如何呢？那位名为\Person{科律法西亚}的，她究竟会停止自称\Person{密涅瓦}吗？——这名字若非为了标明她的母亲，就是因她自\Person{朱庇特}头顶跃出，持盾牌，佩带兵革的恐怖。\Emph{或者我们是否应当认为}，第三位会静静地放弃这名字，而不以如下的话来争辩并抗拒前\Emph{两个}的僭取？\Quote{\Place{赛斯}啊，你这由溪流淤泥与漩涡而生，在泥沼中成形之辈，竟敢如此擅取我名的尊荣？或者，你这妄称自己乃从\Person{朱庇特}头颅中诞生的女神，并说服极愚之人相信你就是理性，你是在篡夺他人的品位吗？他岂能从自己的头中怀孕生子？你所佩带的那兵器可以被铸造成形，难道在他的头壳内竟有铁匠的作坊？\Emph{那里有铁砧}、大锤、熔炉、风箱、煤炭与钳子吗？或者，如你所坚称，若你确是理性，那就停止为自己索取属我的名字；因为你所说的理性，并非某种神明的形态，而是对高深问题的理解。} 那么，如我们所说，在我们尝试献祭时，若有五位\Person{密涅瓦}遇见我们，并为这名号属于谁而争执，每一位都要求把熏香的烟气献给她，或是从金杯中倾倒祭酒；我们该凭哪一位仲裁者，哪一位判官，来处理如此大的争端？或者说，会有哪一位审查者、哪一位裁判者，胆敢对这样的角色作出公正裁决，或宣布她们的理由不基于正义呢？他岂不更愿遁回家中，置身事外，以为与此等事无涉更为安全，免得他因将属于众人之物给了一人而树敌于其余，或因将应属一人的财物屈服给众人而被斥为愚拙？\par

17. 同样的话我们也可用于\Person{墨丘利}们、太阳神们——事实上，一切其他被你增添并繁衍了数目的神祇。但从一件事便足以明白，同样的原则适用于其余；免得我们的冗长使听众厌倦，我们将停止逐一处理，免得我们在指责你过度时，自己也被责以饶舌之过。你们怎么了？你们这些人，借身体酷刑之恐吓，催促我们敬拜诸神，强使我们承行你们神祇的服事。我们本可轻易降服，只要你们能向我们展示某种与如此伟大种族的构想相称之物。向我们展示\Person{墨丘利}，但仅一位；给我们\Person{巴克斯}，但仅一位；一位\Person{维纳斯}，同样地，一位\Person{狄安娜}。因为你们绝不能使我们相信有四位\Person{阿波罗}，或三位\Person{朱庇特}，即使你们召来\Person{朱庇特}本人作证，或以\Person{皮提亚神}为凭借。\par

但反方有人会说：\Quote{我们怎么知道神学家们所写的是确实而众所周知的，还是按他们自己的想法和判断编造的放肆虚构？}这与事情无关；你们论证的合理性也不取决于此——事实是否如\OriginalTerm{神学家}{theologians}的著作所陈述，还是另有不同且明显相异。因为对我们而言，谈论那些昭示于公众眼前之事便已足够；我们\Emph{无须}探究何为真实，而\Emph{只需}反驳并驳倒那显明于众人之前、且已普遍为人们\Emph{的}思想所收纳的事。但如果他们是撒谎者，你们务必自己宣告真理为何，并揭示那无懈可击的奥秘。但当文人学士的协助被搁置一边时，这又如何能做到？因为关于不朽之神，有什么是能说出、而非藉由人们就这些主题所写下的文字才达于人们思想的呢？或者，关于他们的权能与礼仪，你们自己能讲述任何未曾被记载于书籍、并经由作者们的著述而传扬的事吗？如果你们认为这些无关紧要，那就让所有由\OriginalTerm{神学家}{theologians}、\OriginalTerm{司祭}{pontiffs}\Emph{以及}甚至一些致力于哲学研究的人为你们编写的关于众神的书籍统统销毁；不仅如此，我们不妨更假设自世界奠基以来，从来没有人写过任何关于众神的事：我们想要探明，也渴望知道，当你们提及众神时，你们能否嘟囔或喃喃而语，或在思想中构想出那些没有任何书本曾在其心中赋予形象的诸神。然而，当你们显然是藉由书本的提示才得知他们的名字和权能时，却否认这些书籍的可靠性，就是\Emph{不义的}了，因为正是藉着它们的见证和权威，你们才得以确立你们所说的。\par

但是，也许这些事结果会是虚假的，而你们所说的倒是真实的。藉什么证据，藉什么明证，它才能\Emph{被显明}呢？因为双方都是凡人，无论是那些说了一件事的人还是那些说了另一件事的人，双方的讨论都是关于可疑之事；断言自己所以为的是真实的，而冒犯你们感觉的那些则显为放荡与虚妄，这便是傲慢了。按人类的法则，按有死之生物的本性，当你们读到或听到\Quote{那神是由这位\Emph{父亲}和那位母亲所生}时，你们心中难道不觉得所说的事是属于人的，且与我们尘世种族的鄙陋相关吗？或者，当你们这样想的时候，难道你们毫不焦虑，唯恐会以什么方式冒犯众神——不论他们是谁——因为你们相信他们是由于污秽的交合……才得以借着淫荡而达于他们原本不知晓的光明？因为我们，以免有人或许认为我们无知、不了解那圣名所应有的尊威，确实认为众神不应知道诞生；或者，就算他们确实被生，我们秉持并认为：宇宙的主和君宰，以其自身所知的方式，把他们无玷、至洁、未受玷污地差遣而来，不晓性欲之污，且在受生伊始便达致其性体的圆满成全？\par

但你们，相反地，忘记了他们的尊严与伟大是何等崇高，却将诞生与他们相连，并\Emph{将}一种出身\Emph{归}于他们，而这种出身是任何稍有雅致情感的人都视为既可憎又可怖的。你们说，\Person{狄斯丕特}与他的兄弟们乃是由其母亲\Person{俄普斯}和父亲\Person{萨图恩}所生。那么，众神竟有妻子，并且在事先安排婚事之后，他们便受制于婚姻的束缚吗？他们是否借由时效取得、斯佩尔特小麦饼以及假意出卖，而承担起婚床的契约？在既定条件下，他们可曾有情妇、未婚妻和已订婚的新娘？而且，关于他们的婚姻，我们又该说什么呢？你们确实说，有些神庆祝了他们的婚礼，款待了欢愉的人群，女神们也在其间嬉戏；而且，有些神因为没有参与\Emph{歌唱}\OriginalTerm{费斯肯尼诗体}{Fescennine verses}而纷争不已，将一切抛入彻底的混乱，并给下一代人类带来了危险与毁灭。\par

但是，或许这种污秽的污染在其余的事上不太明显。那么，那位天界的统治者、众神与众人之父，他仅凭眉毛的颤动和点头，便摇撼整个天穹并使之战栗——他是从男人和女人那里找到了自己的起源吗？而且，除非两性在肉欲的拥抱中\Emph{自己}陷溺于堕落之乐，那最伟大的\Person{朱庇特}就不会存在；甚至到今日，众神仍无君王，天穹也仍无主宰吗？我们又为何惊讶于你们说\Person{朱庇特}出自妇人的子宫，因为你们的作者不但说他曾有一位乳母，而且还靠着从异母之胸\Emph{汲取}的营养来维持他所获得的性命？你们这些人啊，说什么呢？我要再说一遍，\Emph{那位神}，发出雷霆霹雳、发出闪电、投掷雷电、聚集可怖乌云，竟然吮吸乳房的汁液，像婴孩一样啼哭，四处爬行，而且为了\Emph{被劝服停止}他那极为愚蠢地持续不断的哭泣，竟被哗啷棒的响声弄得安静下来，躺在极柔软的摇篮中入睡，并被不成词的语句抚慰？多么虔敬的宣称啊，关于众神之\Emph{存在}，指出并宣示了他们可畏尊威的可敬庄严！按你们的看法，请问，天界的高超能力就是这样产生的吗？你们的神就是通过这样的出生方式来到光明中的吗？——这些方式，也是驴、猪、狗，是所有这些地上不洁的蓄类受孕与产出的方式？\par

22. 你们不仅把肉欲的结合归于可敬的\OriginalTerm{萨图恩}{Saturn}，还宣称世界之王本人所生的子女比他自己的出生和受生更为可耻。据你们说，他的母亲是\OriginalTerm{许珀里翁那}{Hyperiona}，\OriginalTerm{朱庇特}{Jupiter}——挥舞雷霆者——生下了金色炽热的太阳；他与\OriginalTerm{拉托娜}{Latona}生了\OriginalTerm{得洛斯}{Delian}弓箭手和\OriginalTerm{狄安娜}{Diana}，唤醒林间；他与\OriginalTerm{勒达}{Leda}生了希腊语称为\OriginalTerm{狄奥斯库里}{Dioscori}者；他与\OriginalTerm{阿尔克墨涅}{Aclmena}生了\OriginalTerm{底比斯}{Theban}的\OriginalTerm{赫拉克勒斯}{Hercules}，以棍棒和兽皮护身；他与\OriginalTerm{塞墨勒}{Semele}生了\OriginalTerm{利贝尔}{Liber}，又名\OriginalTerm{布罗弥俄斯}{Bromius}，从父亲大腿第二次出生；他又与\OriginalTerm{迈亚}{Main}生了\OriginalTerm{墨丘利}{Mercury}，善于辞令，携无害之蛇。还能对你们的朱庇特施加更大的侮辱吗？还有什么能更彻底地败坏、毁灭这位众神之首的声誉，超过你们相信他曾屡次为邪恶的享乐所征服，因欲望之火炽热而追逐女性？萨图恩的后代又何必介入奇异的婚姻？\OriginalTerm{朱诺}{Juno}难道不足以满足他吗？难道他不能将自己的欲望之力约束在这位女神之后身上，纵然她拥有如此卓越的佳丽、\Emph{这般}美貌、威仪的容颜以及雪白如玉的手臂？还是说，他并不满足于一妻，而是耽于妾侍、情妇和娼妓，这位淫佚之神如放荡少年一般四处放纵无度，并在年迈之时，与无数人交合之后，又重燃起对\Emph{已然}乏味的享乐的渴求？你们这些亵渎者，你们说些什么？你们对自己所爱的神又构想出何等卑劣的念头？难道你们没有看到、没有察觉你们用何等耻辱来玷污他吗？你们使他成为何等恶行的始作俑者？你们将何等罪污、何等骂名堆积于他？\par

23. 凡人虽易陷于情欲，且因性格软弱而倾向于\Emph{屈服于}感官享乐的诱惑，然而法律仍惩罚通奸，并对那些强行侵犯他人婚床、占据他人权利者处以死刑。然而，\Emph{你们却告诉我们，}这位最伟大的君王竟不知诱奸者和淫夫的身份是何等卑劣、何等可耻；而且，据说那位审察我们功过的神，竟然因为其堕落的内心之推理，看不出他\Emph{自己}应决定采取什么恰当的行动。可是，如果你们至少让他与同等地位者结合，\Emph{如果}你们使他成为不朽女神的情夫，这种不当行为或许还能容忍。但我问你们，人类的身体中有什么美丽，有什么优雅，能够打动、能够吸引朱庇特的目光呢？皮肤、内脏、痰液，以及肠衣覆盖下的一切污秽之物，不仅\OriginalTerm{林叩斯}{Lynceus}以其锐利的目光会为之战栗，任何其他人即使仅仅想一想也会\Emph{不由得}转身避开。\par

24. 如果你们睁开心灵的眼睛，不存任何私心地观看真实真相，就会发觉，如你们所说，人类长期以来遭受的一切灾祸，其原因都源自你们昔日对众神所持的那些信仰；而尽管真理摆在眼前，你们仍拒绝改正。请问，关于他们，我们可曾设想过任何不当之事，或者写出过可耻的文字，以致人们用袭击人类的苦难和丧失生活福祉来煽动对我们的偏见？我们可曾说过，某些神是从蛋中产生的，就像鹳鸟和鸽子那样？\Emph{我们可曾说过}，光彩照人的\OriginalTerm{塞西拉的维纳斯}{Cytherean Venus}是从海水泡沫和\OriginalTerm{科俄路斯}{Coelus}被割下的生殖器中成形的？萨图恩因弑亲而被戴上锁链，只在自己的节期才得以卸下重负？朱庇特因\OriginalTerm{库里特斯}{Curetes}的服侍而免于一死？他将其父赶下权位，以暴力和欺诈占据了本不属于他的王权？我们可曾说过，他那被驱逐的年迈父亲隐匿在\OriginalTerm{意大利}{Itali}的地域，并把名字赠予\OriginalTerm{拉丁姆}{Latium}，因为他在\Emph{那里}受到了儿子之外的庇护？我们可曾说过，朱庇特本人乱伦地娶了自己的姊妹？或者，在受邀赴宴时，不曾享用猪肉，却在无知中吃了\OriginalTerm{吕卡翁}{Lycaon}的儿子？\OriginalTerm{伏尔甘}{Vulcan}跛着脚，在\OriginalTerm{利姆诺斯岛}{Lemnos}上以铁匠为业？贪得无厌的\OriginalTerm{埃斯库拉庇乌斯}{Aeculapius}被雷电击穿，正如\OriginalTerm{维奥蒂亚}{Boeotian}的\OriginalTerm{品达}{Pindar}所吟唱的？\OriginalTerm{阿波罗}{Apollo}富甲一方，却以模棱两可的神谕欺骗了那些用财宝和礼物使他致富的国王？我们可曾宣称墨丘利是一个窃贼？\OriginalTerm{拉维娜}{Laverna}同样如此，并与他一同主宰着隐秘的欺诈？称\OriginalTerm{缪斯}{Muses}是\OriginalTerm{马卡路斯}{Macarus}之女\OriginalTerm{墨伽尔孔}{Megalcon}的侍女的作家\OriginalTerm{米尔提洛斯}{Myrtilus}，是我辈中人吗？\par

25. 是我们说过\Person{维纳斯}是个妓女，被一位名叫\Person{喀倪剌斯}的塞浦路斯国王神化吗？是谁报道说\OriginalTerm{帕拉狄翁}{palladium}是由\Person{珀罗普斯}的遗骸制成的？岂不就是你们吗？谁说\Person{玛尔斯}是斯巴达人？岂不是你们的作家\Person{厄庇卡摩斯}吗？谁说他出生在\Place{色雷斯}境内？岂不是雅典人\Person{索福克勒斯}，且获全体观众赞同吗？谁\Emph{说他出生}在\Place{阿卡狄亚}？岂不是你们吗？谁说他曾被囚禁了十三个月？岂不是\Place{墨勒斯河}之子吗？谁\Emph{说}卡里亚人向他献祭狗，而斯基泰人献祭驴？岂不是尤其是\Person{阿波罗多洛斯}及其余众人吗？谁说他因玷污他人婚床而被网罗缠住捉住？岂不是你们的作品、你们的悲剧吗？我们可曾写过众神为报酬而忍受奴役，如\Person{赫拉克勒斯}在\Place{萨迪斯}为淫欲和放荡；如\Place{提洛}的\Person{阿波罗}\Emph{曾侍奉}\Person{阿德墨托斯}，如\Person{朱庇特}的兄弟\Emph{曾侍奉}特洛伊人\Person{拉俄墨冬}，而\Person{皮提亚}也曾\Emph{一同侍奉}，却是与其叔父；如\Person{密涅瓦}为秘密情人照明、修剪灯芯？那将\Person{玛尔斯}和\Person{维纳斯}描绘成被凡人之手所伤的，岂不就是你们的一个诗人吗？那叙述父\Person{狄斯}和天后\Person{朱诺}被\Person{赫拉克勒斯}所伤的，岂不就是你们的\Person{帕尼亚西斯}吗？你们的\Person{波勒谟}的作品岂不说\Person{帕拉斯}被杀死，浑身是血，被\Person{奥尔倪托斯}压倒吗？\Person{索西比俄斯}岂不宣称\Person{赫拉克勒斯}自己因\Person{希波科翁}之子们所造成的创伤和痛苦而受苦吗？是我们促使人说\Person{朱庇特}在\Place{克里特}岛被葬入坟墓吗？是我们说那在摇篮中便结合在一起的兄弟俩被葬在\Place{斯巴达}和\Place{拉刻代蒙}的领土内吗？那在其作品中被称为\Place{图里伊}人\Person{帕特洛克勒斯}，叙述\Person{萨图恩}的坟墓和遗骸在\Place{西西里}被发现的那位作家，是我们的一员吗？\Place{喀罗尼亚}的\Person{普鲁塔克}被算作我们中的一员吗？他曾说\Person{赫拉克勒斯}因癫痫丧失力量后，在\Place{俄塔山}山顶被焚化为灰烬？\par

26. 但我又该怎样述说那些欲望呢？据你们书中所写、你们作家所载，圣洁的不朽者竟贪恋女色。是我们断言海中之王在疯狂情欲的炽热中，夺去了\Person{安菲特里忒}、\Person{希波托厄}、\Person{阿密摩涅}、\Person{墨那利珀}、\Person{阿洛珀}的童贞吗？那无瑕的\Person{阿波罗}，\Person{拉托娜}之子，最是贞洁纯净，竟被不受理性支配的胸中情欲所俘，贪求\Person{阿尔西诺厄}、\Person{埃图萨}、\Person{许普西皮勒}、\Person{玛耳佩萨}、\Person{宙克西珀}、\Person{普罗托厄}、\Person{达佛涅}、\Person{斯忒洛珀}吗？我们的诗中可曾表明，那早已满头白发、因年迈而心冷的年迈的\Person{萨图恩}，在通奸时被妻子捉住，便化身为下等动物，\Emph{高声}嘶鸣，以野兽之形逃遁吗？你们岂不指控\Person{朱庇特}本人曾变幻无数形象，以卑劣的欺骗来掩饰其淫欲的炽热吗？我们可曾写过，他通过诈骗得偿所愿，一时变为黄金，一时变为嬉戏的羊足怪；变成蛇、鸟、牛；甚至越过一切羞耻的界限，变成一只小蚂蚁，好让\Person{克利托尔}的女儿在\Place{忒萨利亚}成为\Person{密耳弥冬}的母亲？是谁描绘他彻夜不眠地看守\Person{阿尔克墨涅}凡九夜？岂不就是你们吗？——他竟擅离天职，慵懒地放纵情欲？岂不就是你们吗？的确，你们归给他的恩惠还不算小；因为据你们看，神\Person{赫拉克勒斯}生来便在诸如此类的事上超越并胜过其父的能力。他费了九夜才勉强生得一子；但\Person{赫拉克勒斯}，一位圣洁的神，竟在一夜之间使\Person{特斯提俄斯}的五十个女儿同时放下处女的名号，肩负起母亲的重担。更有甚者，你们不满足于将恋女之欲归于诸神，还说他们也对男色起淫念吗？有人恋慕\Person{许拉斯}；另一人迷恋\Person{雅辛托斯}；那一位为\Person{珀罗普斯}燃起欲火；这一位为\Person{克律西波斯}更为炽烈地叹息；\OriginalTerm{卡塔弥图斯}{Catamitus}被掳去作嬖童和司酒者；而\Person{法比乌斯}，为得称\Person{朱庇特}的爱宠，竟在柔软部位受烙印，于臀后留有标记。\par

27. 但在你们中间，难道只有男性贪欲，而女性却保守了贞洁吗？你们的书中岂不证实，\Person{提托诺斯}为\Person{奥罗拉}所爱；\Person{卢娜}渴慕\Person{恩底弥翁}；\Person{涅瑞伊得斯}贪恋\Person{埃阿科斯}；\Person{忒提斯}贪恋\Person{阿喀琉斯}之父；\Person{普洛塞庇娜}贪恋\Person{阿多尼斯}；她的母亲\Person{刻瑞斯}贪恋某个乡野的\Person{伊阿西翁}，又相继贪恋\Person{伏尔甘}、\Person{法厄同}、\Person{玛尔斯}；\Person{维纳斯}本人，\Person{埃涅阿斯}之母，罗马权势的奠基者，竟嫁与\Person{安喀塞斯}吗？因此，当你们无一例外地指控，不是指名道姓地指控某一个，而是指控你们所信奉的全体神灵，犯下如此异常可耻和下贱的行为时，你们竟毫不害臊地敢说，要么是我们不敬虔，要么是你们敬虔么？尽管他们从你们那里，因你们堆积在其名誉上的所有这些可耻行为，而获致比涉及他们神威、尊荣与崇拜所要求的服侍与义务远为重大的冒犯。因为，要么你们所提出的这些关于他们各人的事全是虚假的，损及他们的信誉和声望；果真如此，那么众神应当将人类尽行毁灭，这倒是相当合宜的；要么，如果这些事真实无疑，且被确认毫无可疑之理，那么，无论你们多么不情愿，其结论便是：我们相信他们并非来自天上，而是生于地上。\par

28. 因为哪里有婚礼、婚姻、生育、奶妈、技艺和软弱；哪里有自由和奴隶制；哪里有伤害、杀戮和\Emph{流血}；哪里有贪欲、情欲和感官的享乐；哪里有从可憎的情感所生发的一切内心的激情——那里必然没有神圣性；那些属于会逝灭的族类和大地脆弱之物的东西，也不能贴近一个优越的本性。因为谁若仅仅认识到并理解那种能力的本性，还能够相信：一位神曾有生殖器官，并且通过极为卑劣的手术被剥夺了它；或者他曾一度阉割自己所生的儿子，并因此受罚而被监禁；或者他在某种意义上对他的父亲发动内战，并剥夺其统治权；或者他在被打败后，因恐惧一个更年轻者而转身逃跑，如同一个逃亡者和流放者，躲藏在遥远的荒野中？我说，谁能相信：那位神靠在人的餐桌旁，因贪欲而烦恼，以模棱两可的回答欺骗祈求者，精于盗贼的诡计，犯下奸淫，充当奴隶，受伤，并且恋爱，在各种形式的色欲中屈服于不洁欲望的诱惑？但是你们却宣称所有这一切在你们的诸神身上既曾有过，也依然存在；并且你们在你们幻想的恣肆中，不放过任何形式的邪恶、恶行、错误，不把它们提出来作为对诸神的非难。因此，你们必须要么找出另外的神，这些\Emph{非难}加不到他们身上，因为只有属人类和尘世的族类才适合这些非难；要么，如果只有你们所列述名字和特征的那些神，那么你们的信仰就废除了他们：因为你们所说的所有事情都适用于人。\par

29. 而且在此，我们确实可以指出，你们向我们呈现并称为诸神的那些人，不过是人，通过引用以下作者即可：\Person{阿克剌伽的欧赫墨罗斯}，其著作由\Person{恩尼乌斯}翻译成拉丁文，以使所有人都能透彻了解它们；或是\Person{塞浦路斯的尼加诺尔}；或是\Person{佩拉的莱翁}；或是\Person{昔兰尼的狄奥多罗斯}；或是\Person{希波}和\Person{米洛斯的狄亚戈拉斯}；或是成千上百的其他作者，他们以高尚的坦率，精细地、孜孜不倦地、仔细地将隐秘之事揭示出来。我再说一遍，我们可以随意宣告\Person{朱庇特}的事迹，\Person{密涅瓦}和童贞\Person{狄安娜}的战争；\Person{利柏耳}以何种计谋努力使自己成为印度帝国的主宰；\Person{维纳斯}的处境、职分和收益是什么；\Person{大神母}被许配给了谁；\Person{阿提斯}的俊美在她里面激起了何种希望、何种欢乐；埃及的\Person{塞拉皮斯}和\Person{伊西斯}来自何处，或是出于什么理由甚至他们的名字本身也被编造出来。\par

30. 然而在我们当前所进行的讨论中，我们不承担这种麻烦或辛劳，去指明和宣布所有这些人是谁。\Emph{但}我们所定意的是：你们称我们为不虔敬和无宗教的，\Emph{而}另一方面，你们自称是虔敬且事奉诸神的，我们就要证明并显明，他们所受的待遇，没有比你们所加的更不恭敬的了。如果借由这些侮辱本身就证明了这一点，那么，结果必然能被理解的是，正是你们激起了诸神猛烈而可怕的愤怒，因为你们或是听从、或是相信、或是自己捏造了关于他们的如此贬损的故事。因为不是那正在苦思宗教仪式，宰杀无瑕的牺牲，堆积成堆的香投入火中燃烧的人，必须被认为是在敬拜神明，或是单独履行宗教的义务。真正的敬拜在于内心，和对诸神相配的信念；倘若你相信关于他们的事不仅远离而且不似他们的本性，甚至在某种程度上玷污和羞辱了他们的尊严和美德，那么带来血和凝血便毫无用处。\par

31. 因此，我们想要质问你们，并邀请你们回答一个简短的问题：你们是否认为，向神明献上他们所不愿和不渴望的祭品，抑或以污秽的信念，持有关于他们的极其贬损的意见，竟至于能激起任何人的精神去疯狂地渴望报复，哪一个更是更大的冒犯？如果权衡这些事情的相对重要性，你们将找不到一个带有偏见的法官，竟会不相信，通过明显的侮辱来诽谤任何人的声誉，比起以沉默的忽视来对待它，是更大的罪行。因为后者，或许可以出于对理性的尊重而被认为并相信是可行的；\Emph{但}前者却显露出不虔敬的精神，以及一种虚构中无可救药的盲目。如果在你们的仪式和典礼中，因忽略祭品和赎罪的供奉而被追讨，罪责便被说成已经背负；如果由于一时的疏忽，任何人在说话或斟酒时犯了错误；或者，如果在庄严的竞技和神圣的赛跑中，舞者突然停顿，或者乐师突然缄默——你们所有人便立即喊叫起来，认为某些事情的发生，违反了仪式的神圣性；或者，如果那被称为\OriginalTerm{双亲健在的童仆}{patrimus}的男孩，因无知而放开了皮带，或者不能把它持贴地面：\Emph{然而}，你们还敢否认，诸神正在被你们以如此严重的罪过所触犯，而你们自己却承认，在更小的事情上，他们时常发怒，甚至导致国家的毁灭吗？\par

32. 但他们说，这一切都是诗人的虚构，以及为取乐而编排的游戏。的确，难以相信，那些绝非轻率的人，他们试图探究最远古时代的特征，要么没有将存留于人们心中和普遍谈论之中的寓言插入他们的诗作；要么他们会为自己攫取如此大的自由，以至于愚蠢地臆造近乎彻底的疯狂之事，而这可能使他们有理由畏惧众神，并使他们陷入人际间的危险。但让我们承认，如你们所说，诗人是如此可耻故事的编造者和作者；然而，即便如此，你们也不能免除羞辱神明的罪责，这些神明要么疏于惩罚这类冒犯，要么没有通过颁布法律，并以严厉的惩罚，来反对如此轻率的行为，并规定此后任何人不可说有损神明荣耀或不配于神明荣耀的话。因为无论谁容许作恶者犯罪，都是助长其大胆；而且以虚假的指控诋毁和标记任何人，比揭发和指责他们真实的罪行更带侮辱性。因为被称为你所是的，以及你感觉自己是的那种人，较少冒犯性，因为\Emph{你的怨愤}受到私下省察你的生活时提供的证据所抑制；但那种伤害非常剧烈，其诋毁无辜者，并败坏一个人的名誉和声望。\par

33. 据记载，你们的神在天上的卧榻上用餐，在黄金的房间里饮酒，最后被竖琴和歌声所抚慰。你们为他们配备不易疲倦的耳朵；你们不认为将支撑地上肉体并因萎靡不振的柔弱精神所追求的快乐分配给神明是不合适的。他们中的一些以恋人和破坏贞洁者的形象出现，不仅与女人，也与男人犯下可耻可鄙之事。你们不在乎关于如此重要事情的说法，也不至少以对惩罚的任何恐惧来约束你们放荡文学的恣肆；其他的，由于疯狂和狂乱，自戕其身，屠杀自己的亲属，使身上溅满鲜血，仿佛那是敌人的血。你们对这些崇高表达的不虔敬啧啧称奇；那些本应遭受一切惩罚的事，你们却用赞美加以颂扬，激励他们，以激起更大的猛烈。他们悲痛丧亲之伤，以不合宜的哀号抱怨残酷的命运；你们惊异于他们雄辩的力量，仔细研究\Emph{并}牢记那些本应完全从人类社会摒弃的内容，并热切希望它们不因任何遗忘而湮灭。他们被刻画为受伤、受虐待，互相进行激烈而狂怒的争斗；你们欣赏这种描写，并且为了替作者如此大胆的行为辩护，假装这些是寓言，包含着自然科学原理。\par

34. 但我为何抱怨你们轻看对其他神祇所加的侮辱呢？那位\OriginalTerm{朱庇特}{Jupiter}，你们本不该在全身战栗之前提到他的名字，却被描述为在屈服于对他妻子的情欲时坦白自己的过错，并且在无耻中变得刚硬，仿佛疯狂而愚昧地透露他所偏爱胜于自己配偶的情妇，他所偏爱胜于自己妻子的妾室；你们说，那些说出如此奇异事情的人是诗人中的首领与君王，被赋予了如神般的才华，是最圣洁的人；你们在自己提出的宗教事务上全然忘却了本分，以至于在你们看来，言语比不朽者的被亵渎的威严更为重要。那么，只要你们对众神有任何畏惧，或带着确信与毫不犹豫的肯定相信他们存在过，难道你们不应该通过法案、民众投票、对元老院法令的畏惧，阻止、防备\Emph{并}禁止任何人随意谈论众神而不以虔诚的方式吗？他们甚至没有从你们手中获得这种荣誉，即你们应该用你们用来保护自己免受侮辱的同一法律来抵挡对他们的侮辱。在你们中间，那些对你们的君王低声说过任何坏话的人被指控犯了叛逆罪。贬低一位行政官，或对元老使用侮辱言语，你们已通过法令将其定为\Emph{一项罪行}，并附以最严厉的惩罚。写一首讽刺诗，对别人的名声和品格造成污蔑，你们根据十人委员会的裁定，决定这不该不受惩罚；并且为了不让人以过于放纵的辱骂冒犯你们的耳朵，你们为严重的冒犯确立了程式。唯有在你们这里，众神是不受尊敬的、可鄙的、卑微的；你们允许任何人随意说出他想说的任何话，指控他们犯了其情欲所虚构和设计的卑劣行为。\Emph{然而}你们却不脸红地向我们提出对如此声名狼藉的神祇缺乏尊重的指控，尽管不相信众神的存在远比认为他们是这样的，并有这样的名声要好得多。\par

然而，你们认为只有诗人可以编造关于神祇的不雅故事，并可耻地将它们化为戏谑吗？你们的哑剧演员、演员、那群模仿者与奸夫又怎样呢？他们难道不是滥用你们的神祇来为自己牟利吗？而\Emph{其他人}难道不是在神祇所受的羞辱与损害中寻求诱人的快感吗？在公共赛会上，所有祭司与行政长官的团体也各就其位：最高大祭司，以及\OriginalTerm{库里亚大祭司}{sacerdotes curionum}；头戴月桂花环的\OriginalTerm{十五人司祭团}{Quindecemviri}，佩戴头巾的\OriginalTerm{朱庇特祭司}{flamines Diales}；揭示神明心意和意志的\OriginalTerm{占卜官}{augures}；同样就座的还有守护永燃之火的贞女；全体人民和元老院议员就座；曾任执政官、地位仅次于神祇并最应受敬重的元老们也在此列；而说来可耻，\Person{维纳斯}，\Person{玛尔斯}种族的母亲，帝国人民的祖先，被用手势表演恋爱，并被无耻地模仿成\OriginalTerm{酒神狂女}{Bacchanal}般疯狂，带有卑贱妓女的所有情感。同样，佩戴神圣飘带的\OriginalTerm{伟大母亲}{Magna Mater}，被用舞蹈表演；那\OriginalTerm{佩西努斯的丁迪梅涅}{Pessinuntia Dindymene}，不顾其年龄的荣誉，被表演为在拥抱一个牧人时，带着可耻的欲望，使用激情的手势；此外，在\Person{索福克勒斯}的\WorkTitle{特拉基斯少女}中，\Person{朱庇特}之子\Person{赫拉克勒斯}，陷入致命衣袍的网罗，被展现为发出可怜的哭喊，被剧烈痛苦所征服，最终衰弱并耗尽，当他的肠子软化并溶解时。但在\Emph{这些}故事里，甚至天界至高统治者本人也被搬出，毫无对他名号与威严的敬畏，扮演通奸者的角色，为了诱惑而改变面容，好使他能用狡计夺去他人妻子的贞洁，假扮成她们的丈夫，化身为另一人的模样。\par

但是，这罪行还不够：最神圣神祇的身份也被混入闹剧和下流剧之中。为使闲散观众被激发得哄笑欢腾，神祇在滑稽的俏皮话中被嘲弄，观众呼喊起立，整个剧场回荡着拍手鼓掌。嘲弄神祇的放荡者们被赐予礼物赏赐、安逸、免于公共负担、豁免和减免，连同\OriginalTerm{凯旋花环}{coronae triumphales}——这是一项任何道歉都无法弥补的罪行。此后，你们竟敢惊讶于这些人世间的灾祸从何而来，它们毫无间断地淹没并压倒人类，同时你们每天重复并记诵所有这些事情，其中混杂着对神祇的诽谤和中伤言论；当你们希望你们闲散的心智被无用的梦想占据时，要求给你们日子，并毫无间断地呈献表演？但如果你们为你们的宗教信仰感到任何真正的愤慨，你们早就该焚烧这些著作，摧毁你们的那些书籍，并推翻这些剧院，在那里每天以可耻的故事公开对你们的神祇进行邪恶报导。因为，究竟为什么，我们的著作理应被付之一炬？我们的聚会理应被残酷地解散，在聚会中我们向\OriginalTerm{至高天主}{Deus Summus}祈祷，为所有掌权者、士兵、国王、朋友、敌人，为那些仍在生者，以及那些脱离肉躯束缚的人，祈求和平与宽恕；在聚会中所说的一切，都是为了使\Emph{人}成为人道、温和、谦逊、有德、贞洁、慷慨处理财物，并与我们\OriginalTerm{兄弟情谊}{fraternitas}中所有拥抱的人不可分离地联合？\par

但事实是，由于你们在战争和军事力量上极其强大，你们认为自己在真理知识方面也卓越，并且你们在诸神面前显得虔诚，然而你们却是最先用污秽的想象玷污他们威能的人。在此，如果你们的凶猛允许，疯狂容许，我们请你们回答这个问题：你们是否认为愤怒在\OriginalTerm{神性}{natura divina}中占有一席之地，还是\OriginalTerm{神福}{beatitudo divina}远离这类\OriginalTerm{情感}{passiones}？因为如果他们像你们的想象所暗示的那样，受如此狂怒的情感支配，并被暴怒的情绪激发——因为你们说他们常常以咆哮震动大地，给人类带来悲惨的苦难，以瘟疫般的传染败坏时代的风气，都是因为他们的赛会庆祝得不够仔细，因为他们的祭司未被悦纳，因为一些小空间被亵渎，因为他们的礼仪未按规执行——那么必然得出，他们对所提及的那些意见感到不小的愤怒。但是，如果必然得出的推论是，所有这些长期压垮人类的苦难都源于这些虚构，如果神祇的愤怒是由这些原因激起的，那么你们就是这些可怕灾祸的肇因，因为你们从不停止刺激神祇的情绪，激起他们强烈的报复渴望。但是，另一方面，如果神祇不受这类情感影响，也根本不知道何为愤怒，那么确实没有理由说，这些不知道愤怒为何物的神祇对我们感到愤怒，*并且他们不受其存在影响，也不受\Emph{它所引起的}混乱影响。因为在事物本性中，一不能变为二；那本性上非复合的\OriginalTerm{统一体}{unitas}也不能分裂并分离为不同的事物。\par

\SourceRecord{Translated by Hamilton Bryce and Hugh Campbell.；Ante-Nicene Fathers；Vol. 6.；Edited by Alexander Roberts, James Donaldson, and A. Cleveland Coxe.；Buffalo, NY: Christian Literature Publishing Co.,；1886.；<http://www.newadvent.org/fathers/06314.htm>.}

% structure: p0001 source=h1 target=section reason=page-title

\section{驳异教徒（卷五）}

1. 承认诗人们仅仅出于戏谑而编造了所有这些对不死的神祇造成羞辱的事，那么，\Emph{至于}那些载于严肃、认真、谨慎的史书中，并由你们在隐秘的奥迹中流传下来的事，又如何呢？它们也是诗人放荡的幻想编造出来的吗？如果它们对你们而言显得如此荒谬，你们就不会在日常使用中保留其中某些故事，不会年复一年地举行庄严的节庆来庆祝，也不会在你们的圣礼中把它们当作真实事件的影子来维护。我将严格节制自己，只举出这众多故事中的一个；在这个故事里，\Person{朱庇特}本人被搬上台，显得愚蠢而轻率，被词语的歧义所欺骗。在\Person{安提阿斯}的第二卷中——以免可能有人以为我们在恶意捏造罪名——写着这样一个故事：——\par

著名的国王\Person{努玛}不知道如何禳解雷所预示的灾祸，又渴望得知其法，便依\Person{厄革里亚}的建议，在一口泉边藏匿了十二名带着锁链的贞洁青年；这样，当\Person{法乌努斯}和\Person{玛尔提乌斯·皮库斯}来到此处饮水时——因为他们惯常到这里取水——他们便可冲上前去，抓住并捆绑他们。然而，为了使这事做得更为迅捷，国王以酒和蜜酒斟满许多杯盏，并将它们放置在通往泉水的通道上，使其显而易见——这是为来者设下的狡诈圈套。他们一如往常，因口渴难耐，来到熟悉的出没之处。但一看到那些散发着甜香的杯子，他们便喜新厌旧；急切地扑向它们；为那饮品的甘美所迷醉，喝得过多；酩酊大醉，沉沉入睡。于是，那十二个\Emph{青年}扑向酣睡者，\Emph{并}用锁链捆绑那些醉倒在酒中的人；他们惊醒后，立即教给国王用什么方法和祭品能把\Person{朱庇特}唤到地上。国王凭借这知识在\Place{阿文蒂诺}山上举行了神圣的仪式，将\Person{朱庇特}引到地上，并向他请教应有的禳解方式。\Person{朱庇特}犹豫良久，说道：\Quote{你应当用一个头来禳解雷所预示的灾祸。}国王答道：\Quote{用一个洋葱。}\Person{朱庇特}又说：\Quote{用一个男人的头。}国王回说：\Quote{但要用头发。}神明反过来：\Quote{用生命。用一条鱼，}\Person{庞皮里乌斯}接话。于是，\Person{朱庇特}被所用的含糊言辞所困，说出了这些话：\Quote{你骗过了我，\Person{努玛}；因为我本已决定，要用\Emph{人头的祭献}来禳除雷所示的灾祸，而不是用头发\Emph{和}洋葱。然而，既然你的诡计胜过了我，就按照你所希望的方式吧；今后你永远要用你所换得的这些东西来进行雷霆之兆的禳解。}\par

2. 心灵应当首先思考什么，最后思考什么，或应默然略过什么，这不容易说清，任何思考也弄不明白；因为所有这一切都是如此设计和编造以引人发笑，以致你们应当努力让人相信它们是虚假的——即便它们是真实的——而不是让其作为真实之事流行开来，从而仿佛提示出某种非凡之事，并给神性本身带来蔑视。那么，你们说，你们啊——？我们岂能相信，\Person{法乌努斯}和\Person{玛尔提乌斯·皮库斯}——若他们是神祇之数，具有永恒不朽的本体——曾一度口干舌燥，去寻涌流的泉水，好用水冷却他们灼热的血管？我们岂能相信，他们为酒所困，为蜜酒的甘甜所惑，那么久地沉溺于那些诱人的杯盏，以至于甚至有醉倒之虞？我们岂能相信，他们沉睡不起，坠入深沉睡眠的遗忘中，竟给了地上的造物捆绑他们的机会？那么，那些束缚和锁链是搭在什么部位上的呢？他们有任何固体实质吗？或者他们的双手是由坚硬的骨头做成，以致能用绳索捆绑并就紧紧地系住的结扣将他们牢牢捉住？因为我并不询问，也不探究，他们在醉后的胡言乱语中摇晃不定时，是否能说出什么；或者，当\Person{朱庇特}不愿意，或者不如说是不知情时，是否有人能够使人知晓将他唤到地上的方法。我只想听到这一点：如果\Person{法乌努斯}和\Person{皮库斯}具有神圣的起源和权能，他们为什么不自己向询问的\Person{努玛}宣告他甘冒更大风险想从\Person{朱庇特}本人那里学到的东西呢？难道只有\Person{朱庇特}才知道——因为雷霆是由他投下的——如何藉某种知识的训练来逃避即将来临的危险？或者，尽管他亲自投掷这些火矢，平息他的愤怒和忿恨的方法却要由别的神知道？因为，认为他亲自指定了可以通过哪些方法来避免他决定借雷霆投下、降于人类的灾祸，这真的荒谬至极。因为这就等于说，借着这般的仪式，你将转移我的愤怒；而且，倘若我曾在任何时候用闪电预示过灾祸即将来临，你就做这做那，以致我所定意要做的事全然落空，并藉着这些仪式的力量而悄然消逝。\par

3. 但是，我们姑且承认，如他们所说，\Person{朱庇特}自己曾为自己指定了方法和途径，借此可以适当地反对他自己宣布的意图：我们是否也要相信，一位如此威严的神明竟被拖到地上，与一个侏儒站在一座小丘上，陷入一场口角之争？请问，是什么魔法迫使\Person{朱庇特}离开那对宇宙至关重要的治理工作，而应凡人之命显现？是祭食、馨香、血、燃烧的月桂枝的气味，和喃喃的咒语吗？难道这些东西都比\Person{朱庇特}更强大，以至于能强迫他不情愿地履行所命之事，或是他自己心甘情愿地委身于这些狡诈的把戏？什么！我们还能相信以下的事吗：\Person{萨图恩}的儿子竟如此缺乏远见，以至于他提出的条件本身就因含糊不清而使他自陷其中，或者他竟不知道将要发生什么，不知道一个凡人的诡计和狡猾将如何胜过他？他说，当雷电劈落时，你要用一个头来赎罪。这话还不完整，意思还没有充分表明和界定；因为必然需要知道，\Person{狄埃斯皮特}规定这赎罪要用一只阉公羊头、一头母猪头、一头公牛头，还是任何别的畜生的头？可是，他既没有对此作出明确的规定，他的决定既仍悬而未决，也未确定，那么\Person{努马}怎会知道\Person{朱庇特}将说的是一个人头，从而预知\Emph{并且}阻止\Emph{他}，并把他那含糊不明的话变成\Quote{一颗洋葱头？}\par

4. 但是，你或许会说，国王是一位占卜者。他能比\Person{朱庇特}本人更会占卜吗？对于一个凡人而言，预知\Person{朱庇特}——他欺瞒了\Person{朱庇特}——将要说的话，难道这位神明竟不知道一个人正准备用什么方法欺瞒他吗？那么，这岂不是清楚又明显，这些都是幼稚而虚构的捏造，借此将敏锐的机巧归给\Person{努马}，却将最大的缺乏远见归给\Person{朱庇特}吗？因为还有什么比这更缺乏远见呢：承认自己被一个人智力的狡黠所诱骗，又在受骗恼怒时，向那战胜了你的人让步，并放弃你所提出的手段？因为，如果之所以被雷击的东西须用人头作赎罪祭，是有理可据、有某种自然合宜性的，那么我看不出国王为什么提出一颗洋葱头；但如果也可以用一颗洋葱来完成，那么对人血就有了一种贪婪的嗜血欲望。而且，这两部分自相矛盾：因此，一方面，\Person{努马}被表明不想知道他所确实想知道的事；而另一方面，\Person{朱庇特}被表明曾是残忍无情的，因为他说他希望用人的头来赎罪，而这件赎罪祭本可由\Person{努马}用一颗洋葱头来完成的。\par

5. \Person{蒂莫修斯}绝非凡俗的神话学者，其他同样学识渊博的人亦然，他们详细叙述了\OriginalTerm{大神母}{Great Mother}的诞生及其仪式的起源，据他本人撰写和所述——这一切都取自古代博学的典籍，以及他对最隐秘奥秘的\Emph{通晓}：——他说，在\Place{弗里吉亚}境内，有一块在各方面都荒僻得闻所未闻的岩石，其名为\Place{阿格杜斯}，是当地原住民如此称呼的。按\Person{忒弥斯}凭其神谕所吩咐的，\Person{丢卡利翁}和\Person{皮拉}从这块岩石上取下石头抛在地上，那时大地已空无一人；从这些石头中，这位被称为大神母的，也连同其他石头一起被塑造出来，并由神明赋予了生气。当她被置于岩石顶端安歇入睡时，\Person{朱庇特}怀着最淫荡的欲望攻击了她。但是，经过长时间的挣扎，他仍然无法达到自己的目的，于是一筹莫展，便将情欲发泄在石头上。岩石承受了这一切，并伴随着阵阵呻吟，在第十个月生下了\Person{阿克得斯提斯}，他的名字得自其母岩。在他身上，有不可抗拒的力量，有无法控制的凶暴性情，有狂暴的欲望，且\Emph{源自}两性。他大肆劫掠，肆意毁坏；在他凶残性情所到之处，散布毁灭；他既不顾神，也不理人，也不以为有任何比自己更强大的东西；他藐视大地、苍天和星辰。\par

在众神的会议中，此事常被考虑，如何才有可能削弱或抑制他的胆大妄为，\Person{利柏耳}（其余神祇退缩不前）担起了这项任务。他用最烈的酒在\Person{阿刻得斯提斯}常去的一处泉水里下了药；那地方本是他在运动和狩猎后平息\Emph{他内心}所激起的燥热与焦渴之处。\Person{阿刻得斯提斯}一感到需要，便跑到那里去喝；他过于贪婪地将那饮料吞入干渴的血管。由于全然不习惯这酒力，他很快便被制服，沉沉入睡。\Person{利柏耳}就守在\Emph{他所设的}圈套旁；他把一条用毛发精心编成的套索一端套在他的脚上，另一端则抓住他的私处。酒劲过后，\Person{阿刻得斯提斯}狂怒地跳起，他的脚拖动绳套，靠自己的力量剥夺了自己的性别；\Emph{这些}部位撕裂时，血流如注；两者都被大地吞没并吸收；从它们那里突然长出一棵挂满果实的石榴树，\Person{珊伽里乌斯}国王（或河神）之女\Person{娜娜}见了它的美丽，满怀赞叹，摘\Emph{了一些果实}放在怀中。她因此怀孕；她的父亲把她关了起来，以为她行为不检，并设法让她饿死；她却被\OriginalTerm{诸神之母}{mother of the gods}用苹果和其他食物养活，\Emph{并}生下一个孩子，而\Person{珊伽里乌斯}却下令将他丢弃。一个名叫\Person{福耳巴斯}的人发现了这孩子，带回家中，用山羊奶喂养他；又因为\Place{吕底亚}地方把俊美的男子如此称呼，或是因为弗里吉亚人依其方言称山羊为 \OriginalTerm{阿塔吉}{attagi}，这男孩便因此得名\Person{阿提斯}。\OriginalTerm{诸神之母}{mother of the gods}因他容貌出奇地俊美而极其钟爱他；他的同伴\Person{阿刻得斯提斯}，随着他长大，爱抚他，并以如今唯一可能的方式，用淫邪的顺从与他相连，领他穿过林木茂密的林间空地，送给他许多猎获的野兽，小\Person{阿提斯}起初还吹嘘说那是自己辛苦劳力所得。后来，在酒力影响下，他承认自己既受\Person{阿刻得斯提斯}宠爱，又被他用林中得来的礼物尊崇；因此，凡是因\Emph{饮酒}而受玷污的人，都不得进入他的圣所，因为酒揭露了他的秘密。\par

随后，\Place{佩西努斯}的国王\Person{弥达斯}，意欲使少年脱离这如此可耻的亲密关系，决定将自己的女儿许配给他，并下令关闭\Emph{城门}，以免任何不祥之人扰乱他们的新婚之乐。然而，\OriginalTerm{诸神之母}{mother of the gods}知晓这年轻人的命运，知道他若要保持在人间的安全，\Emph{唯有}不结婚姻，为使灾祸不致发生，她进入了紧闭的城，用头抬起城墙，她的头因此开始被塔楼环绕。\Person{阿刻得斯提斯}因男孩被夺走、且被带去娶妻而勃然大怒，使所有宾客陷入疯狂错乱：弗里吉亚人因神祇显现而惊恐尖叫；奸夫\Person{加卢斯}的女儿割下自己的乳房；\Person{阿提斯}夺过那催人发狂之人所持的笛子；他自己也此刻充满暴怒，疯狂咆哮\Emph{且}四处冲撞，最后将自己摔倒在地，在一棵松树下自残，并说：\Quote{收下这些吧，\Person{阿刻得斯提斯}，你为它们挑起了如此巨大而可怖的骚动。}生命随着涌流的鲜血逝去；但\OriginalTerm{诸神的大母}{Great Mother of the gods}将被割下的部分收集起来，先用死者的衣服覆盖包裹，再撒上泥土。从涌出的鲜血中长出一种花——紫罗兰，并用它装点了那棵松树。从此便开始了这种习俗，你们至今还用花朵遮盖并缠绕那神圣的松树。那曾作新娘的贞女，据\OriginalTerm{大祭司}{pontifex} \Person{瓦勒里乌斯}记载，名叫\Person{伊阿}，用柔软的羊毛遮盖了\Emph{那少年}毫无生气的胸膛，与\Person{阿刻得斯提斯}一同流泪，然后自杀身亡。她死后，她的血化为紫色的紫罗兰。\OriginalTerm{诸神之母}{mother of the gods}也流泪，泪水长出一棵杏树，象征死亡的苦涩。然后她将那棵\Person{阿提斯}在其下自阉的松树带回自己的洞穴；\Person{阿刻得斯提斯}加入她的悲号，她则绕着已静止的树干\Emph{踱步}，捶打着自己的胸膛。\Person{阿刻得斯提斯}恳求\Person{朱庇特}让\Person{阿提斯}复活：他不允许。然而，命运所容许的，他欣然应允：他的身体不会朽坏，他的头发将永远生长，他的最小一根手指将存活并保持永远活动；\Emph{据称}，\Person{阿刻得斯提斯}满足于这些恩惠，便在\Place{佩西努斯}将他的遗体祝圣，并以每年的礼仪和司祭服事加以尊崇。\par

8. 如果某人，鄙视神祇，且带着一种野蛮渎神的烈怒，决意亵渎你们的神祇，他敢说出比这个故事所叙述的更为严厉的话攻击他们吗？你们却将之编成形式，仿佛\Emph{它就是}一个奇妙的故事，并不断加以尊崇，以免时间的冲刷和年代的久远使之被遗忘。因为在这个故事中，有什么关于神祇的断言或记载，若是用在一个人，培养出坏习惯和相当粗野的训练的人身上，不会使你们因冤枉和侮辱他而招致谴责，并引起憎恨和厌恶，伴随着无法平息的怨恨呢？你们说，众神之母是由\Person{丢卡利翁}和\Person{皮拉}所掷的石头产生的。啊，神学家们，你们说什么？啊，天上权能的祭司们，你们说什么？难道众神之母的诞生完全是因为洪水吗？若不是暴雨冲走了全人类，她的出生就没有原因或开端吗？那么，她因人类才感到自己存在，她多亏\Person{皮拉}的恩惠才被当作一个真实的存在；但如果这是真的，那么这也必然不假：她是人，不是神。因为如果人最初是由投掷石头产生的，那么必须相信她也是我们中的一员，因为她是由同样的原因产生的。因为不可能，自然不容许这样，从同一种石头，和同一种投掷\Emph{它们}的方式，有些被塑造成列于不朽之列，另一些则具有人的境况。那位著名的罗马人\Person{瓦罗}，以其学识广博和孜孜不倦地研究古代而闻名，在他所遗留的四卷关于罗马人民族裔的著作的第一卷中，通过精确计算表明，从我们之前提到的洪水时代，直到\Person{希尔提乌斯}和\Person{潘萨}的执政官任期，几乎不到两千年；如果他是可信的，那么也必须说，大母神的整个生命被这个数字的界限所限定。于是，事情就发展到这个地步：那位被称为所有神祇之母的，不是他们的母亲，而是他们的女儿；不，她倒不如说\Emph{仅仅是个孩子}，一个小女孩，因为我们承认，在无尽的时间序列中，诸神既无开始也无终结。\par

9. 然而，我们为何还要说你们以尘世的污秽玷污了众神的大母神，而你们竟连片刻都未曾停止说\Person{朱庇特}本人的坏话呢？当时众神之母正在\Place{阿格杜斯}的最高峰上熟睡，你们说，她的儿子在她睡着时试图悄悄地夺取她的贞洁。在夺去无数处女和主妇的贞洁之后，\Person{朱庇特}难道指望在他母亲身上满足他可恶的欲望吗？难道他不能因为自然本身不仅在人类，而且在某些\Emph{其他动物}中，以及普遍的感受所激起的恐惧，而从他猛烈的欲望中回头吗？身为庙宇之首的他，难道不顾虔敬与荣誉吗？在他心智狂乱激动之时，难道他不能反思或察觉他的欲望是多么邪恶吗？但是，事实却是，他忘记了他的威严与尊贵，爬向前去偷那卑贱的快乐，因恐惧而颤抖战栗，屏住呼吸，踮着脚尖惊恐行走，在盼望与恐惧之间，触摸她的隐秘部位，试探他母亲睡得有多沉，以及她将忍受什么。啊，可耻的描绘！啊，\Person{朱庇特}不光彩的困境，预备尝试肮脏的争斗！于是，当他因疏忽和匆忙而未能用突袭偷取之时，世界的统治者是否就诉诸暴力？当他不能以狡诈手段攫取快乐时，他是否以暴力攻击他的母亲，并毫不遮掩地开始破坏他本应尊崇的贞洁？然后，在她长时间反抗后，他是否败退，被征服，被制服，并离去？那个虔敬都无法阻止其对自己母亲产生可憎淫欲的人，他那耗尽的淫欲是否就离开了他？\par

10. 但你们或许会说，人类是厌恶并诅咒这类结合的；在众神之间不存在乱伦。那么，\Emph{那么}，他的母亲在她儿子向她施暴时，以最强烈的态度反抗？她为什么要躲避他的拥抱，仿佛在逃避不洁的亲近？因为如果那事没有不对，她就应该毫不勉强地满足他，正如他热切地希望满足自己的淫欲一样。而在这里，确实，那些非常节俭的人们，甚至在可耻的行为上也斤斤计较，为了不让那神圣的种子看似白白泼洒——那岩石，据说，吸收了\Person{朱庇特}肮脏的失禁。接下来发生了什么，我问你们？说吧。就在岩石的中心，在那燧石般的坚硬中，一个孩子形成并被赋予了生命，成为伟大\Person{朱庇特}的后代。要反对这样不自然且如此奇异的观念并不容易。因为正如你们所说，人类是从石头中产生并繁衍出来的，那么我们必定要相信，那些石头既有生殖器官，又吸收了洒在它们身上的种子，并在满月时怀孕，最后如妇人分娩般痛苦地生产。这促使我们的好奇心去追问：既然你们说分娩发生在十个月后，那么在那段时间里，他被封闭在岩石的哪个子宫里？靠什么食物，什么汁液供养？或者他从坚硬的石头中能汲取什么来支撑自己，就像未出生的婴儿通常从母亲那里\Emph{汲取}一样！他尚未见到光，\Emph{我的报信者}说道；但他已经在吼叫，模仿他父亲的雷声，重现了那种\Emph{声音}。当他能见到天空和白昼的光时，便攻击所遇见的一切，肆意破坏，并确信自己能够将众神从天上推下来。哦，谨慎而又有先见之明的诸神之母啊，她为了不承受如此傲慢儿子的恶意，或者为了不让他还在未出生时的吼叫打扰她的睡眠或打破她的安歇，便退避开来，将那极其有害的种子远远送走，给了粗糙的岩石。\par

11. 在诸神的会议中，对于如何制服那顽固而凶暴的力量产生了疑惑；当别无他法时，他们采取了一种手段，就是让他喝下大量酒，并割去他的肢体，使其丧失男性器官。\Emph{仿佛}，确实，那些失去这些部分的人会变得不那么傲慢，仿佛我们不是每天见到那些自阉的人变得更加放荡，轻看所有贞洁和端庄的约束，一头栽进污秽的卑贱中，将他们的可耻行为公之于众。然而，我倒是想看看——倘若我蒙允出生在那时——父\Person{利柏尔}，他战胜了\Person{阿克戴斯提斯}的凶猛，从诸天的高峰，在诸神极其庄严的集会之后滑落下来，割下马尾，编织柔软的缰绳，用大量烈酒迷醉那些无害而纯净的水，并在饮酒导致的沉醉之后，仔细地伸出双手，摆弄那沉睡者的肢体，巧妙地将他的关怀导向那些将要毁灭的部分，好让套在那\Emph{些部分}上的套索能完全环绕它们全部。\par

12. 有谁会这样谈论诸神，哪怕对他们评价极低？或者，如果他们被这类事务、考虑、忧虑所占据，任何智慧之人会相信他们是神，或甚至将他们算作人吗？请说，那个\Person{阿克戴斯提斯}，割下他可耻的肢体是为了给不朽的众神带来安全感，\Emph{他是个}地上的造物，还是位神，拥有不朽？因为如果他被认为\Emph{属于}我们的命运，处在人的境况中，那他为何使众神如此恐惧？但如果他是神，那他怎会被欺骗，或者从神的身体上怎能割下什么？但我们在这一点上不提出异议：他可以具有神性的出身，也可以是我们中的一员，如果你们认为这样说更准确。另外，从那流淌的血和割下的部分中，是不是也长出了一棵石榴树？还是说，当那个器官被藏在大地的怀里时，它用根抓住土地，长成一棵大树，开出缀满花朵的枝条，并在瞬间结出完全成熟的甜美果实？而且，因为它们源于红血，所以它们的颜色是明亮的紫色，并带有一抹黄色？再进一步说，它们多汁，带有葡萄酒的味道，因为它们源于一个充满葡萄酒之人的血，这样你的故事就首尾一贯了。哦，\Place{阿布德拉}，\Place{阿布德拉}，如果这样一个故事是你们编出来的，那会给人们提供多少嘲笑你们的把柄啊！所有父辈都将它传述，傲慢的城邦将它研读；而你们却被视为愚笨，全然迟钝又愚蠢。\par

13. 据说，\Person{娜娜}通过一只苹果从她的腹中怀孕，生下了一个儿子。这种说法是自相一致的；因为既然岩石和坚硬的石头都能生育，那么苹果也必定有其生育的时节。\Person{贝雷辛提亚女神}用坚果和无花果喂养被囚禁的少女，这既恰当又正确；因为一个因苹果而成为母亲的女子，理应以苹果为食。她的孩子出生后，\Person{桑伽里乌斯}下令将他远远抛弃：他早已认为这孩子是神性受孕的，却又不会称其为自己孩子的后代。婴儿依靠公山羊的奶被抚养长大。哦，这故事总是与男性为敌，极其憎恨男性；在其中，不仅人放弃了男子的能力，就连雄性的野兽也变成了母亲！他以美貌著称，并以非凡的俊秀而出众。最令人惊奇的是，山羊的恶臭竟没有使人躲避和逃离他。\Person{大母}爱他——如果像祖母爱孙子那样，那倒无可非议；但如果像戏剧中所描述的那样，她的爱便是不名誉和可耻的。\Person{阿克得斯提斯}也最爱他，用猎人的礼物使他富足。\Emph{你们说}，从一个被阉割的人那里，他的纯洁不会受到任何危险；但\Person{弥达斯}所惧怕的是什么，却不容易猜到？\Person{那位母亲}进城时竟然背负着城墙。我们在这里确实惊叹于神祇的大能与力量；但我们也再次指责她的粗心大意，因为当她记起命运的谕令时，却漫不经心地让城池暴露于敌人面前。\Person{阿克得斯提斯}使那些庆祝婚誓的人陷入狂怒和疯狂。如果\Person{弥达斯}王得罪了\Emph{他}——那个正把这青年与一位妻子结合的人，那么\Person{加卢斯}和他的妾的女儿又犯了什么罪，以至于他要自残男根，她要自割乳房？他说：\Quote{拿住并收好这些，}又说：\Quote{就是因为你们激起了如此骚乱，用恐惧压倒了\Emph{我们}的心神。} 要不是那个少年为了平息他的怒火，把割下的部分扔给了他，我们谁也不会知道疯狂的\Person{阿克得斯提斯}究竟想要他情妇身体上的什么。\par

14. 啊，你们这些种族和民族，竟沉溺于这样的信仰，你们还有什么可说的？当这些事情被提出来时，你们难道不感到羞耻和困惑，竟然说出如此不体面的话吗？我们本想从你们那里听到或学到一些与神祇相称的事情；但你们相反地，却向我们提出割去乳房、砍掉男子器官、狂怒、流血、疯狂、少女自尽，以及从死者的血中生出的花朵与树木。再说，那时\Person{众神之母}，难道在悲痛中，小心翼翼地亲手将散落的生殖器连同流出的血收集起来了吗？她是否用自己神圣的、属神的手，触摸并捧起那从事可耻和不体面工作的器具？她是否也将它们埋入土中，使之不见天日；又生怕它们裸露时在大地的怀抱中散失，是否真的在将它们用他的衣服包裹掩盖之前，先用芳香的膏油清洗并涂抹它们？因为若不是添加了那些膏油，改变那器官的腐臭气味，紫罗兰的芬芳又从何而来？请问，当你们读着这样的故事时，难道不觉得自己仿佛是在听织机旁的少女消磨冗长的劳作时光，或者是在听老妇人给轻信的孩子找些消遣，并且以真理的外衣编造出各种虚构吗？\Person{阿克得斯提斯}向\Person{朱庇特}恳求，让他的情侣复生：\Person{朱庇特}不肯答应，因为他受制于比\Emph{他自己}更强大的命运；而为了不在各方面都显得铁石心肠，他赐下了一个恩惠——身体不会因任何腐朽而分解；头发将永远生长；身体上仅最小的那根手指能存活，并且单独保持永恒的运动。谁会认可这一点，或者毫不犹豫地支持它——头发在死尸上生长——那部分已消亡，而\Emph{他身体其余}有死的部分，却不受腐朽律的约束，竟依然存留至今？\par

15. 我们早就可以敦促你们思考这一点，若非要求证明这类事情如同讲述它们一样愚蠢。但这故事是虚假的，全然不实。事实上，因为你们坚持认为诸神是因何人被逐出大地，对我们而言，这故事是否前后一致并且有确实根据，还是完全出于捏造和虚构，都无关紧要。因为对我们来说——我们今日已表明要阐明这一点——这就足够了：你们所提出的那些神灵，如果他们存在于世上某处，并且燃烧着愤怒的火焰，那么他们因凶猛仇恨而发的怒火，主要原因不在我们而在你们；而那个\Emph{故事}，已被你们列为事件并记录下来，被你们每天乐意地诵读，又按序传承以启迪后世。现在，如果这个\Emph{故事}确实是真实的，我们看不出其中有任何理由可以断言天上的神明对我们发怒，因为我们既没有说出那么多令他们蒙羞的事，也从未将它们记录下来，更未透过神圣仪式的举行将其公之于众；但倘若，正如你们所认为的，这故事并不真实，由虚妄的谎言构成，那么无人能够怀疑你们才是触犯神明的原因，因为你们要么容许某些人写下这样的故事，要么任由这些写下的\Emph{东西}在世代记忆中留存。\par

16. 然而，当你们全年所举行的那些仪式本身就证明你们相信\Emph{这些事}是真的，并且认为它们完全可信时，你们又怎能断言这故事是虚假的呢？因为在固定的日子里你们总是带进\OriginalTerm{众神之母}{mother of the gods}的圣所的那棵松树，是什么意思呢？难道这不是模仿那棵树吗？在那棵树下，那狂暴而不幸的青年对自己下手，而众神的母亲将\Emph{这棵树}祝圣，以减轻她的悲伤。你们用来绑扎并缠绕树干的那一团团羊毛，是什么意思呢？难道不是让人回想起\OriginalTerm{拉}{la}用羊毛覆盖那垂死的\Emph{青年}，并且以为能为他那因寒冷而\Emph{迅速}僵硬的肢体带来一些温暖吗？那缠绕并装饰着紫罗兰花环的树枝，是什么意思呢？难道不是表明，众神之母如何用早春的花朵装饰了那棵指示并见证这悲惨事故的松树吗？那些\Emph{\OriginalTerm{伽利祭司}{Galli}}披散着头发，用手掌拍打胸膛，是什么意思呢？难道不是使人想起那戴塔冠的众神之母，与哭泣的\OriginalTerm{阿克得斯提斯}{Acdestis}一起，高声哀号，追随着那少年吗？你们称为\Emph{\OriginalTerm{禁食}{castus}}的戒食面包，是什么意思呢？难道不是在模仿那位女神因深切悲伤而禁食\OriginalTerm{刻瑞斯}{Ceres}的果实的时候吗？\par

17. 或者，如果我们所说并非如此，那么请你们自己——我们在你们这位神明的神圣仪式中看到的那些阴柔娇气的\Emph{男人们}——说说，他们在那里有什么事、\Emph{什么}挂虑、\Emph{什么}关切；为什么他们像哀悼者一样弄伤自己的臂膀和胸膛，表现得仿佛身处悲惨境遇中呢？那些花环\Emph{是什么意思}，那些紫罗兰\Emph{是什么意思}，那些绷带、那些柔软羊毛的覆盖物\Emph{是什么意思}？最后，为什么那棵不久之前还在灌木丛中来回摇摆的松树，一根完全惰性的木头，被立在众神之母的神庙中，就像某种吉祥而非常可敬的神明一样呢？因为，要么这就是我们在你们的著作和论文中找到的原因，\Emph{那样的话}，显然你们不是在举行神圣的仪式，而是在表演悲惨的事件；要么，如果另有某种奥秘的隐晦瞒过了我们的原因，那也必定牵连在某种可耻行为的丑名之中。因为谁会相信，无用的\Emph{\OriginalTerm{伽利祭司}{Galli}}着手，而阴柔的淫荡之徒完成的事，会有什么尊荣呢？\par

18. 主题的重大，以及对那些受辩护者的责任，也要求我们同样地去搜寻其他可耻的形式，无论是古代历史记载的那些，还是包含在名为\Emph{\OriginalTerm{入会礼}{initia}}的神圣奥迹中的那些——这些奥迹并不公开泄露于所有人，只对少数人的缄默透露；但是你们不可胜数的神圣仪式，以及它们的全部令人憎恶之处，不允许我们逐一详细数算：而且，说实话，我们故意有意地避开了一些，免得在力图揭露一切时，我们自己反倒在论述中因沾染污秽而被玷污。因此，让我们略过\OriginalTerm{福娜·法图阿}{Fauna Fatua}，她被称为\OriginalTerm{善德女神}{Bona Dea}，据\Person{塞克斯图斯·克洛狄乌斯}在他关于诸神的希腊文第六卷书中宣称，她因为在她丈夫不知道的情况下喝了一整罐酒，而被桃金娘枝鞭打致死；这便是一个证据，即当妇女们向她献上神圣的敬礼时，会放置一罐酒\Emph{在那里，却}遮掩起来不让人看见，而且不允许带入桃金娘的枝条，正如\Person{布塔斯}在他的\WorkTitle{原因论}中所提到的。但是让我们以同样的忽略略过\Emph{\OriginalTerm{掩埋之神}{dii conserentes}}吧，根据\Person{弗拉库斯}及其他人的叙述，他们把自己掩埋起来，在一锅\OriginalTerm{内脏}{exta}下面的灰烬中变成了\Emph{\OriginalTerm{类似于人阴茎的形状}{in humani penis similitudinem}}。当精通\Place{伊特鲁里亚}术数的\Person{塔娜奎尔}搅动了这些时，那些神明挺立起来，变得僵硬。她于是命令一名来自\Place{科尔尼库鲁姆}的女俘去学习和理解这是什么意思：最具智慧的女子\Person{奥克里西亚}，\Emph{\OriginalTerm{诸神插入生殖器，展开确定的动作}{divos inseruisse genitali, explicuisse motus certos}}。然后，那圣洁而炽热的神明倾泻出\Person{卢基里乌斯}的力量，\Emph{这样}，罗马国王\Person{塞尔维乌斯}便出生了。\par

19. 我们也将略过那狂野的\OriginalTerm{酒神节}{Bacchanalia}，希腊语中称之为\OriginalTerm{生肉宴}{Omophagia}，在其中，你们带着假装的疯狂和丧失理智，用蛇缠绕自己；并且，为了显示自己充满神明的神性与威严，用血口撕裂嘶叫的山羊的肉。我们也略过那\OriginalTerm{塞浦路斯的维纳斯}{Cyprian Venus}的隐秘奥迹，据说其创立者是\Person{辛尼拉斯}王；在其中，受洗者如同对妓女一样缴纳规定的费用，并带走\OriginalTerm{阴茎像}{phalli}，作为吉祥神明的标记。也让\OriginalTerm{科里班忒斯祭司}{Corybantes}的仪式被忘却吧，其中揭示了那神圣的奥迹：一个兄弟被他的兄弟们杀害，欧芹从被杀者的血中生出，这种蔬菜被禁止放在餐桌上，以免死者的\OriginalTerm{亡灵}{manes}受到不可平息的冒犯。但我们拒绝公开宣扬那些其他的\OriginalTerm{酒神节}{Bacchanalia}，其中向受洗者揭示并教导一个不可言说的秘密：\OriginalTerm{利伯尔}{Liber}如何因童稚的游戏而被\OriginalTerm{泰坦}{Titans}们撕碎；又如何被他们切成碎块，扔进锅里烹煮；\OriginalTerm{朱庇特}{Jupiter}如何被诱人的香气吸引，不请自来地冲向餐席，发现所发生的事，便用他可怕的雷电压倒了那些狂欢者，并将他们抛入\OriginalTerm{塔尔塔罗斯}{Tartarus}的最底层。作为其证据和证明，那色雷斯的\Emph{吟游诗人}在他的诗篇中传下了骰子、镜子、陀螺、环圈、光滑的球，以及取自童贞女\OriginalTerm{赫斯珀里得斯}{Hesperides}的金苹果。\par

20. 我们原打算也不提那些弗里吉亚受传授的秘仪，以及所有那个种族，若不是他们\Emph{已引进的}\Person{Jupiter}的名字不容我们匆匆掠过对他所加的伤害和侮辱；并非我们喜欢讨论如此污秽的秘仪，而是为了一再向你们显明，你们把何等错谬加在你们自认是守护者、捍卫者、敬拜者的那些神身上。他们说，从前\Person{Diespiter}对母亲\Person{刻瑞斯}燃起邪情和可耻的欲望——因那地方的土人说，她\Emph{是}\Person{Jupiter}的母亲——却不敢公然用强索取他无耻贪恋之事，便想出一条诡计，好夺取那毫无防备的母亲贞操。他变成公牛，而非原本的神；在野兽潜伏的表象下遮掩意图和胆大妄为，他疯狂地以突发的暴力冲向那毫无所思、毫无察觉的她，遂了乱伦的欲望；此事因其淫欲而暴露后，他便逃走，已被认出和发现。他的母亲燃烧、口吐白沫、喘息、沸腾着暴怒和愤慨；既不能抑制内心风暴和怒涛，此后因她那永不止息的激情而得名为\OriginalTerm{布里摩}{Brimo}：且她别无他愿，唯求尽其所能惩罚儿子的胆大妄为。\par

21. \Person{Jupiter}十分烦恼，恐惧万分，找不到办法平息受辱母亲的愤怒。他倾吐祈祷，哀求恳请；她的耳朵因悲痛而紧闭。众神全体都被派去为他求情；无人有足够分量能得倾听。最后，儿子寻求补偿之法，想出此策：\OriginalTerm{他挑选了一只高贵而睾丸硕大的公羊，亲手将其阉割，并从毛茸茸的外皮中取出}{Arietem nobilem bene grandibus cum testiculis deligit, exsecat hos ipse et lanato exuit ex folliculi tegmine}。他悲哀地低垂着头走近母亲，仿佛自己已定自己的罪，把这些投扔到她怀中。当她看到他的信物，便稍觉缓和，容许自己被唤回关照她所怀的后代。十个月后她生下一个女儿，形容美丽，后世或称她为\Person{利柏拉}，或称\Person{普洛塞庇娜}；当\Person{Jupiter Verveceus}看到她健壮、丰满而娇艳，竟然忘了他刚才陷入的种种邪恶、罪行和何其大的鲁莽，又重回故态；且因为觉得父亲公然与女儿结为夫妻似乎太过邪恶，他便化作可怕的龙形：将那惊恐的少女盘绕在巨大蛇身中，以凶猛的样貌嬉戏爱抚\Emph{她}，用最温柔的拥抱。她也因而充满大能\Person{Jupiter}的种子，却不像她母亲\Emph{那样}，因为她生了一个像自己的女儿；而从这少女生的却是一件类似牛的东西，证明她曾被\Person{Jupiter}诱奸。若有人问是谁讲述此事，那我们将引用古人咏唱的一位塔兰托诗人的著名六音步诗句：\Quote{公牛生龙，龙生公牛。}最后，圣礼本身，乃至名为\OriginalTerm{塞巴迪亚}{Sebadia}的传授仪式，也可证实其真；因为在仪式中，一条金蛇被放入入会者的怀中，又从下身取出。\par

22. 我在此也不必用许多话逐节讨论，并指出每一细节中有多少下贱不当之处。因为有哪一个凡人，只要对人之本分稍具意识，会不自己清楚看出这一切事的性质是\Emph{何等邪恶}、卑劣，并看出他们的秘仪中的种种礼节，以及其仪式可耻的起源，给神明带来了怎样的耻辱？据说，\Person{Jupiter}贪恋\Person{刻瑞斯}。我要问，\Person{Jupiter}如何得罪了你们，竟至于你们将一切羞辱、恶名昭彰的奸淫，像堆在什么卑贱无用之人身上一般，全都堆在他的头上？\Person{勒达}对她的婚姻誓言不忠；\Person{Jupiter}被指为罪过之因。\Person{达娜厄}守不住童贞；偷窃据说是\Person{Jupiter}所为。\Person{欧罗巴}急于成为妇人；他再次被宣布为她贞操的攻击者。\Person{阿尔克墨涅}、\Person{厄勒克特拉}、\Person{拉托娜}、\Person{拉奥达弥亚}，以及成千其他贞女、成千主妇，连同童子\Person{卡塔米图斯}，都一并被夺去了荣誉和贞节。到处是同一故事——\Person{Jupiter}。也再没有什么卑劣之事，你们不把他的名字与放纵情欲相连；以致这可怜的家伙似乎降生只为作为罪恶的沃土、毁谤的契机、某种敞开之地，好从戏台的污秽中汇聚一切肮脏。然而，你们若说他与陌生妇人交合，那诚然是不敬，但诽谤他的罪过尚可容忍。难道他也对自己母亲怀着疯狂欲望，乃至对亲女亦然；难道父母之身的神圣、对父母之身的敬畏、甚至对自己所出的子女的\Emph{毫无}退缩，都不能阻止他构想如此可憎的计划？\par

23. 因此，我愿看见那位始终掌管世界与人类的诸神之父朱庇特，头戴牛角，摇晃着毛茸茸的耳朵，双足缩成蹄形，咀嚼着青草，身后拖着尾巴，臀部和脚踝涂满了柔软的粪便，遍身沾污着排出的秽物。我愿——我必须一再重复——看见那位运转星辰、吓倒并倾覆那些惊恐万状的民族的神，他追赶着阉羊群，\Emph{\OriginalTerm{审视公羊的睾丸}{inspicientem testiculos aretinos}}，用那只严厉而神圣的手——他惯常用它抛出耀眼的闪电、在暴怒中投掷雷霆——夺走这些东西。然后，我更愿看见他用灼热的刀剖开它们的最深处；在所有目击者被遣走之后，撕下\Emph{\OriginalTerm{包裹胎儿的膜}{circumjectas prolibus}}，把它们带给他的母亲，他依然怒火中烧，将这些膜当作一种祈求怜悯的束发带，他面容沮丧，苍白，受伤，假装在痛苦中煎熬；而且，为使人相信，他沾染着雨的鲜血，并用羊毛和麻布包扎他那假装的伤口。难道这些事竟能在世间被人听闻和阅读，而讨论这些事的人竟希望自己被视为虔诚、圣洁、宗教的捍卫者？难道有比这更大的亵渎，或者能找出如此深染不虔思想的心灵，竟相信这类故事，或接受它们，或在最秘密的宗教仪式中传承它们？如果你们所说的那位朱庇特，无论他是谁，当真存在，或曾受到任何不公的感触，难道不理所当然，他应该被激怒，将大地从我们脚下移走，熄灭日月的光芒；不但如此，还应将万物混成一团，如同古时一样？\par

24. 但是，\Emph{我的对手会说}，这些并非我们国家的仪式。请问，谁这样说，或谁重复这话？他是罗马人、高卢人、西班牙人、阿非利加人、日耳曼人，还是西西里人？如果这些故事不是你们的，而编造它们的人却站在你们一边，那对你们的事业又有何益处？或者，你们是否赞同它们，这又有什么关系，既然你们自己所讲的，被证明要么同样污秽，要么更加卑劣？你们是否希望我们细察那些希腊人称为\OriginalTerm{塞斯摩福里亚}{Thesmophoria}的奥秘和仪式？在这些仪式中，雅典人\Emph{为女神}举行了圣夜的守斋和庄严的守夜。我说，你们是否愿意我们查考它们的起源、原因，以便证明，连在艺术和文明追求方面如此卓越的雅典，也同样讲出对其他神祇的侮辱之辞，而且在那里，假托宗教名义公开传述的故事，与你们其他人\Emph{丢给我们的}一样可耻？据说，有一次，普罗塞庇娜——尚未成年，仍是一位少女——在西西里的草地上采摘紫色的花朵，她热切采摘，四处走动，忽然，冥王从一处不知多深的裂口中跃出，抓住并掳走了这少女，随后重又隐没在地腹之中。那时，刻瑞斯不知发生了何事，全然不知女儿究竟去了何方，便决意寻遍全世界，找回失落的女儿。她抓起两把在埃特纳火山之焰上点燃的火炬，藉着它们的光照，在大地的每一角落进行她的搜寻。\par

25. 在她颠沛流离的寻找途中，她来到了厄琉西斯以及其他地区的边界——那是阿提卡的一个乡区。当时，这些地方居住着原住民，名为包波、特里普托勒摩斯、欧布勒俄斯、欧摩尔波斯、狄索勒斯：特里普托勒摩斯是驾牛耕田者；狄索勒斯是牧羊人；欧布勒俄斯养猪；欧摩尔波斯牧羊，欧摩尔波斯族即是由他而来，\Emph{也由他}衍生出那在雅典人中闻名的名字，以及后来兴起的\OriginalTerm{神使}{caduceatores}、\OriginalTerm{圣显司祭}{hierophants}和传令官。那么，我们说过的那个住在厄琉西斯乡区的包波，殷勤地接待了被种种不幸折磨得疲惫不堪的刻瑞斯，用讨好的关心围绕着她，恳求她不要忽略身体的歇息，又端上用小麦面浆稠和的酒——希腊人叫作\OriginalTerm{曲刻翁}{cyceon}——为她解渴。悲伤的女神掉头避开那善意递来的慰藉，拒绝了\Emph{它们}；她的不幸不容她想起身体常需之物。包波则按这类灾难中的常情，恳求并劝勉她，不要轻视她的人情；刻瑞斯全然不为所动，固执地保持着不可征服的严苛。但这样反复多次之后，她的坚执心意仍无法被任何殷勤消磨，包波便改变计划，决定用奇特的戏谑来逗乐她无法。以庄重赢得的人。那女子们藉以生育子女、获得母亲之名的身体部分，她从此不再长期疏忽：她使它外表更显洁净，变得如同孩童般光滑，尚未因毛发而坚硬粗糙。她就这样回到悲伤的女神面前；正当她尝试那些常用于打破并缓和悲痛之力的普通办法时，她敞开自己，裸露下体，展示出所有为羞耻所遮掩的部位；于是女神的双眼凝视其上，对这奇特的安慰方式感到欣喜。大笑之后，她神色更为开朗，端起并且饮尽了\Emph{先前}被拒的饮料，而那无耻的亵露行为，竟强行成就了包波端庄举止长久未能赢得的事。\par

26. 若有人碰巧以为我们在讲邪恶的诽谤，就让他拿起那些被你们称为具有神圣古老性的色雷斯占卜者之书；他将发现，我们既非狡猾地编造任何事，也非设法将神祇的神圣性沦为笑柄，而确实如此：因为我们将引出卡利俄佩之子用希腊语述说的诗句，这诗句他曾在他的歌中向全人类传递，历经万代：——\par

她一边说着这些话，一边从最下边的\Emph{衣边}撩起衣裳，\newline{}\Emph{\OriginalTerm{那形似私处之物}{formatas inguinibus res}}暴露于眼前，\newline{}\Person{鲍波}用空手握住了它，因为\newline{}它形如婴孩，她轻轻击打，温柔抚摸。\newline{}于是那女神，将威严光明的双眸凝视，\newline{}心软下来，暂且将心中的悲伤搁置；\newline{}随后她手拿杯子，大笑着，\newline{}欣然将整杯\OriginalTerm{西塞昂饮料}{cyceon}一饮而尽。\par

啊，\Person{厄瑞克透斯}的智慧子孙们，你们怎么说？\Person{弥涅耳瓦}的公民们，你们又怎么说？人们急切地想知道，你们将用什么样的言辞来捍卫那危险至极的主张，或者你们有什么样的手段，能为那些受到致命伤害的人物和事业提供安全。这并非毫无根据的猜疑，你们也不是遭到诬陷的指控：你们\OriginalTerm{厄琉息斯秘仪}{Eleusinia}的丑行，既被它们卑劣的起源所证实，也被古代文献的记录所证实，最后，更被你们在接受圣物受询问时所用的标记所证实——\Quote{我禁食了，饮了那饮料；我从那\Emph{神秘的圣物箱}中取出，放入柳条筐里；我又取回，并转移到了小箱子里。}\par

难道你们的众神会被武力掳走，又像他们神圣隐秘的秘仪所讲述的那样，用暴力强夺吗？他们是否以偷摸和欺骗的方式缔结婚姻？是否从抗拒不从的童女那里夺去了贞操？他们对即将临头的伤害一无所知，对那些被强行掳走者的遭遇毫无了解吗？他们走失时，如同凡人一般被寻找吗？他们在阳光最灿烂之时，也要举着灯烛火把，遍寻大地广袤的角落吗？他们感到痛苦吗？他们心绪烦乱吗？他们穿上哀悼者那肮脏的衣服，显出悲惨的标记吗？而为了能让他们将心思转向食物和进食，所用的不是理智，不是恰当的时机，不是某些郑重的话语或恳切的礼遇，而是将那羞耻不雅的部位展示出来吗？那些部位，人人皆感羞耻，自然的廉耻法则命我们遮掩，未经许可不宜在纯洁的耳朵前道出，还得说一句\Quote{请恕冒昧？}我问你们，在那样的景象中，在\Person{鲍波}的私处里，究竟有什么，能让一位同性、且拥有相似部位的女神感到惊异和发笑？究竟有什么，一旦呈现于神明眼前，就能让她忘却苦难，并陡然欢欣，心绪转佳？哦，若非出于对读者的尊重，顾及文学的尊严，我们本可尽情讥嘲揶揄！\par

我承认，我久久踌躇，左顾右盼，推搪闪躲，倍增\Person{泰莱内}式的困惑；因为我羞于提及那些\OriginalTerm{阿利蒙提亚}{Alimontian}秘仪，在希腊，人们为尊崇\Person{巴克科斯}父神而树立\Emph{\OriginalTerm{阳具像}{phalli}}，整个地区遍布男子\Emph{\OriginalTerm{生殖器像}{fascina}}。这事的含义或许幽晦不明，有人会问为何要这样做。凡对此无知者，理应学习，并于惊叹如此要事之余，以一颗纯洁之心，恒常虔敬持守。话说，生于\Place{倪萨山}、身为\Person{塞墨勒}之子的\Person{利柏尔}尚在人间时，据传，他愿结识冥界幽影，探询\Place{塔耳塔洛斯}之事；但这一愿望因某些困难受阻，他不识路途，不知由何处前往且行进。此时冒出一个名叫\Person{普洛绪姆诺斯}的人，他是这位神明下贱的爱慕者，且是一个极易陷入邪淫欲念的\Emph{家伙}，他许诺指出\Person{狄斯}的门径，以及通往\Place{阿刻戎河}的入口，条件是该神须满足他的欲望，容许他\Emph{\OriginalTerm{从自己身上获取夫妻之乐}{uxorias voluptates ex se carpi}}。神明毫不犹豫地立誓，将自己交于他掌控处置，但只等他从地府返回之后，心愿已遂且欲念得偿。\Person{普洛绪姆诺斯}客气地告知道路，将他安置在地府的门槛之处。就在\Person{利柏尔}悉心审视\Place{斯提克斯河}、\Person{刻耳柏洛斯}、\Person{复仇女神}及其他种种之际，那位指路人却已从生者之列消逝，按照凡人的方式入土安葬。\Person{欧伊俄斯}由地府上来，方知其向导已死。但为履行承诺，解脱誓言之缠累，他来到埋葬之地，并且——\OriginalTerm{从无花果树上砍下一根最粗壮的树枝，削、刨、磨光，做成男子阳具的形状，插在坟丘之上，然后从后面褪去衣裳，走上前来，蹲下身子，跨坐其上，继而任那饥渴的淫欲放纵开来，臀部前后扭动，一心从木棍上获取他早前真实许诺过的那番感受。}{ficorum ex arbore ramum validissimum praesecans dolat, runcinat, levigat et humani speciem fabricatur in penis, figit super aggerem tumuli, et postica ex parte nudatus accedit, subsidit, insidit. Lascivia deinde surientis assumptâ, huc atque illuc clunes torquet et meditatur ab ligno pati quod jamdudum in veritate promiserat.}\par

29. 现在，为避免有人以为我们杜撰了如此不虔敬的事，我们并不要求他相信\Person{赫拉克利特}作证，也不要求他接受\Person{赫拉克利特}记述对此等密仪的感受。让他去问全希腊，那遍及乡野城镇，由古俗竖立敬拜的\OriginalTerm{\Emph{男根像}}{phalli}是何含义：他必发现其原因正如我们所说；或者，他们若羞于诚实地吐露实情，那么遮掩、隐藏仪式的起因与根源又有何用，既然崇拜行为本身已受指控？各民族啊，你们怎么说？你们这些忙于神殿事奉、\Emph{献身其中}的列国啊，怎么说？难道就是为这些仪式，你们用火刑、流放、屠杀以及各种惩罚，用残酷折磨的恐惧，驱赶我们吗？这些就是你们带给我们的神明，强行塞给我们的神明，你们自己不愿像他们那样，也不愿任何血亲挚友像他们那样吗？你们能向尚着童装的无须少子宣告，\Person{利柏耳}与其情人所订的密约吗？你们能敦促儿媳，甚或自己的妻子，去\Emph{展示}\Person{鲍波}的羞态，去\Emph{享受}\Person{刻瑞斯}的贞洁之乐吗？你们愿你们的青年知晓、听闻、\Emph{并}学到，\Person{朱庇特}甚至在不止一位贵妇面前显出何种模样吗？你们愿你们的成年闺女和精力尚旺的父亲学习，这同一位神如何与自己的女儿嬉戏吗？你们愿情欲已炽的同手足兄弟、同父同母的亲姊妹，听闻他又一次不拒其姊妹的拥抱与床榻吗？这样，我们还不该远远逃离这些神明吗？还不该全然塞住耳朵，免得这至洁宗教的污秽渗入心灵吗？因为，有哪个受过如此纯洁道德教养的人，神明的榜样不会激起他类似的疯狂呢？或者，谁能约束对亲属、对应当敬畏之人的欲望，当他看到在至高的神明当中，因情欲的混乱，一切都不被视为神圣？因为既然那最初而完美的本性，都无法将激情限制在正当界限内，那么人，既受本性的脆弱所驱驰，又得神圣神明的教导所襄助，为何不该不加分别地放纵欲望呢？\par

30. 我承认，每当心中反思这些怪异故事，我长久以来便惯于诧异，你们竟敢把那些要么完全否认神明存在、要么怀疑其存在、要么断言他们是人、因某种功德与大恩而被列于神位的人，称为无神论者、不虔敬者和亵圣者；因为，若公正地审察，最该被如此称呼的，无过于你们自己，你们佯装向神明示敬，却在此过程中堆叠更多的侮辱与指控，甚至比公开承认予以辱骂更甚。那怀疑神明存在、甚或全然否认的人，虽似凭其臆测之大胆而持怪异见解，却未加侮辱而拒绝相信幽晦之事；而那断言他们是凡人的人，虽将他们从高居天堂者的尊位拉下，却因推想他们因功业、因令人仰慕的德性而被擢升于神列，反将其他荣耀加于其身。\par

31. 然而你们，自称是神明不朽的捍卫者与宣扬者，你们曾否放过、曾否略过他们中的任何一位，而不以辱骂攻击他？抑或，有哪种在众人眼中如此可憎的侮辱，竟使你们不敢用之其身，哪怕因其名号的尊严而有所顾忌？谁宣称神明爱恋脆弱必朽之身？\Emph{难道不是你们吗？}谁宣称他们在他人床榻上施行了那些极诱人的窃玉偷香？\Emph{难道不是你们吗？}谁宣称子女与母亲交合，反之，父亲与贞女女儿交合？\Emph{难道不是你们吗？}谁宣称面容俊美的童子、乃至容貌极佳的成年男子，被不正当觊觎？\Emph{难道不是你们吗？}谁宣称他们被伤残、被玷污、精于伪装、作盗贼、囚于锁链桎梏，最终遭雷霆轰击、受伤而亡，乃至在世间有坟墓？\Emph{难道不是你们吗？}那么，你们既已提出如此众多而严重的指控，伤害了神明，岂敢断言神明是因我们而不悦，而久已显明，你们才是这类恼怒的罪魁，神明震怒的起因？\par

\Emph{我的对手}说，但是你们错了，你们被误导了，甚至在批评\Emph{这些}事情时，显示出你们相当无知、没有学识、粗俗不堪。因为所有那些在你们看来不光彩、有损于诸神声誉的故事，其中都包含着神圣的奥秘，精妙而深奥的理论，绝非任何人仅凭理解力就能轻易知晓的。因为故事表面所写和所说的，并不是其真正的意图和含义；相反，这一切都应从寓意层面来理解，并借助私下传授的秘密解释。因此，说\Person{朱庇特}与母亲同寝的人，并不是指\Person{维纳斯}那种乱伦或可耻的拥抱，而是以\Person{朱庇特}指代雨，以\Person{刻瑞斯}指代大地。再者，说他与女儿淫乱的人，也并非谈论秽恶的享乐，而是以\Person{朱庇特}指代雨水，以他的女儿指代所播种的庄稼。同样，说\Person{普洛塞庇娜}被父神\Person{狄斯}掳走的人，也并非如你所猜想的那样，是说那少女被掳以\Emph{满足}最卑下的欲望；而是因为我们用土块覆盖种子，他借此表示女神沉入地下，与\Person{俄耳库斯}结合以结出果实。同样，在其他的故事中，所说的确是一回事，而所理解的则是另一回事；在通俗直白的表达之下，隐藏着秘密的教义和幽深的奥秘。\par

这些显然都是诡辩和遁词，他们惯于用这种手法在陪审团面前为站不住脚的案子辩护；不，更确切地说，这些都是借口，如同在诡辩推理中所用的那类借口，不是追寻真理，而总是追寻真理的影像、表象和影子。因为既然接受那准确的叙述为真实是羞耻和不妥的，你们便求助于这种权宜之计，用一事物替代另一事物，并使得原本可耻的事物，在解释之下，硬被装点出道貌岸然的样子。然而，这些空洞的故事背后是否有别的含义和意义，与我们何干？因为我们断言诸神被你们邪恶而不敬地对待，我们只需要接受所写下的和所说的，无需关顾那被隐藏起来的，因为对神明的侮辱不在于隐藏于其含义中的观念，而在于那些字面所表达的意义。不过，为了不显得我们不愿审查你们的说法，我们首先问你们，只要你们能忍受我们，你们是从谁那里得知，或是谁让你们明白，这些故事是寓意性地写成的，或者应当以这种方式来理解？是写作者们召唤你们去与他们商量吗？还是当他们在不顾真相而用此物代彼物时，你们当时就躲在他们怀里呢？如果说，他们出于宗教敬畏和恐惧，无论为何原因，选择将那些奥秘包裹在晦暗之中，那么你们竟然想要理解他们不愿让人理解的，自己去知晓并让所有人都知道他们徒劳地试图用不暗示真相的言词来隐藏的东西，这显示出何等的鲁莽！\par

但是，就算我们同意你，在所有这些故事中，用鹿代替了\Person{伊菲革涅亚}，然而，当你们解释或展开这些寓意时，你们如何确定你们给出的解释，或者你们持有的观念，正是写作者们自己在默然的思想中所怀有的、却用不适合于其本意而适合他物的言辞来表述的呢？你们说，雨水落入大地的怀抱被说成是\Person{朱庇特}与\Person{刻瑞斯}的结合；另一个人或许会以更精巧的方式设计，并以某种可能性推测出别的什么；第三个人、第四个人也会如此；并且随着思考者心智特征的不同，每件事都可能以无限多的方式被解释。因为，既然所有所谓的寓意，都取自特意弄得晦涩的叙述，且没有任何确定的界限，能使所谓故事的含义牢固固定、不可更改，那么每个人都可以随心所欲地把意义附加进去，并断言他所思想与推测出的意义就是被采纳的。既然情况如此，你们又如何能从可疑的东西中获得确定性，并为你们看到可用无数不同方式解释的一个说法强加上唯一的含义呢？\par

35. 最后，若你以为合理，回到我们的探究，我们向你请问：你认为关于诸神的所有故事，也就是无一例外地，全都贯穿着双重意义与含义，并且在某种程度上允许多种解释；还是说其中有些部分毫不含糊，相反，另一些则具有多重意义，并被包裹在环绕其上的寓意之幕中？因为如果叙述的整体结构从头到尾都被寓意之幕所笼罩，那么向我们解释，告诉我们，对于每个故事所说的每一事物，我们应当用什么来替代，又应当将它们归于哪些其他的事物与含义。因为，举例来说，正如你们希望用\Person{朱庇特}来代替雨，用\Person{刻瑞斯}代替大地，用\Person{利贝拉}和父\Person{狄斯}来代表种子陷入和播入地中，同样，你们也应当说明，对于公牛我们应当理解为什么，对于\Person{刻瑞斯}的愤怒和暴怒应理解为什么；\OriginalTerm{布里摩}{Brimo}一词意味着什么；\Person{朱庇特}焦急的祈祷，他差遣向他说情却不被听从的神们，意味着什么；阉割的公羊意味着什么；阉割的公羊的部位意味着什么；用这些所做的赔罪意味着什么；对他女儿进一步的行为，其淫欲更加无耻，意味着什么；同样，在另一个故事中，\Place{赫纳}的树林和花朵是什么；从\Place{埃特纳}取来的火以及用它点燃的火把是什么；带着这些火把巡游世界是什么；\Place{阿提卡}地方、\Place{厄琉西斯}地区、\Person{包波}的小屋以及她乡间的款待是什么；饮下\OriginalTerm{库刻翁}{cyceon}意味着什么，拒绝饮它意味着什么，剃掉和暴露私处意味着什么，那景象的羞耻诱惑，以及由此产生的对她丧亲之痛的遗忘意味着什么。现在，如果你们能指出所有这些东西应当用什么来代替，以一个替换另一个，我们就承认你们的断言；但如果你们既不能为每一情况提出另一种假定，也不能求助于整体的上下文，那么，为什么你们要利用看似冠冕堂皇的寓意，将那些已被清楚言说、向所有人理解敞开的事物弄得晦暗不明呢？\par

36. 但你们或许会说，这些寓意并非存在于整个故事的全体之中，而是有些部分写得让所有人都能明白，另一些部分则有双重含义，并被包裹在模糊之中。这是一种精巧的微妙，即使最愚钝的人也能看穿。因为你们很难将所有已说之事移换、逆转并转为他意，于是你们挑选出适合自己目的的一些东西，并借助这些努力维持那些虚假伪造的说法，将其散布于隐藏在其下的真理之上。然而，纵使我们向你们承认事情正如你们所言，你们又如何知道，或从哪里得知，故事的哪一部分是毫无双重含义地写下的，哪一部分又是被不协调和不相干的含义所覆盖的呢？因为可能你们所相信的如此之事并非如此，而你们所相信的不同之事却是以不同甚至相反的表达方式产生的。因为在一致的整体中，当一部分被说成是寓意地写成，另一部分则以明白可信的语言写成，而事物本身中没有任何标记来指明模棱两可之言与简洁之言的区别时，那简洁之言很可能被认为具有双重含义，正如那模棱两可写成的东西被相信是包裹在隐晦之中一样。但事实上，我们承认我们完全不理解，这件事是被谁做成的，或者怎么可能被相信能够做到。\par

37. 那么，让我们来检视一下按这种方式所说的话。在\Place{赫纳}的树林中，我的对手说，少女\Person{普洛塞庇娜}曾经在采花：这至今未受污染，且以直白的方式讲述，因为所有人都毫无疑义地知道树林和花朵是什么，\Person{普洛塞庇娜}是什么，以及少女是什么。\Person{苏曼努斯}从地中跃出，驾着四马战车行驶：这也同样简单，因为四马队、战车和\Person{苏曼努斯}无需解释者。他忽然抓走了\Person{普洛塞庇娜}，并带她一同进入地下：我的对手说，种子的掩埋就是由对\Person{普洛塞庇娜}的抢夺来意指的。请问，发生了什么，使得故事突然转向了别的事物？\Person{普洛塞庇娜}竟被称为种子？她长期以来被认为是一个采花的少女，在她被强行带走和掠走之后，竟开始意指那播下的种子？我的对手说，\Person{朱庇特}将自己变成公牛，渴望与他的母亲\Person{刻瑞斯}交合：正如先前解释的，在这些名字下，谈论的是大地和落雨。我看到寓意之法则借着晦暗和模棱的词语表达出来。\Person{刻瑞斯}被激怒且生气，并接受了公羊的（睾丸和）赔偿，作为复仇所要求的惩罚：这我再次看到是以通常的语言表达的，因为愤怒和（睾丸以及）赔偿都是以它们通常的情形谈论的。那么，这里发生了什么——从被用来指雨的\Person{朱庇特}和被用来指大地的\Person{刻瑞斯}，故事竟转到了真正的\Person{朱庇特}，以及对事件的最直白的叙述上？\par

38. 那么，要么它们全都是用\OriginalTerm{寓意}{allegory}写成并提出的，并且整个寓意应该向我们指出；要么没有一样是这样写的，因为所谓的\Emph{寓意}似乎不像是叙述的一部分。你说，这些全都是用寓意写的。这似乎绝不肯定。你问是什么理由、什么原因？我回答说，因为所有已经发生并被明确记载在任何书中的事，都不能被转变成寓意，因为已经成就的事不能被撤销，一个事件的性质也不能变成完全不同的另一种性质。难道特洛伊战争能被变成\Person{苏格拉底}的定罪吗？或者\Place{坎尼}战役能变成\Person{苏拉}的残酷剥夺公权吗？的确，如\Person{图利乌斯}开玩笑所说，剥夺公权可以被说成是一场战斗，并被称为坎尼的那场战斗；但是已经发生的事，不能同时既是一场战斗又是一次剥夺公权；因为，如我已说过的，已经发生的事只能是已发生的事，不能是其他任何事物；而且那已被牢固确定在其自身本性和特有状态中的事物，也不能转变为与它异质的实体。\par

39. 那么，我们从何证明所有这些叙述都是事件的记录呢？从那些庄严的仪式和入教秘仪来看，这是清楚的，无论是在固定的时间和指定的日子举行的那些，还是那些被异教徒秘密传授而从不中断其习俗遵行的那些。因为不能相信这些仪式毫无起源，没有道理和意义地被实行，且与它们最初的开始没有任何原因相连。那棵通常被带入伟大母亲圣所的松树，难道不是模仿那棵在其下\Person{阿提斯}自残并自阉的树，而且他们说女神为缓解悲伤而祝圣了那棵树吗？那种每年在希腊受崇拜并在仪式中庆祝的\OriginalTerm{\Emph{阳具}}{phalli and fascina}竖立，难道不使人想起\Person{利柏尔}还债的那件事吗？那些\OriginalTerm{厄琉息斯秘仪}{Eleusinian mysteries}和秘密仪式又叙述了什么呢？难道不是叙述了\Person{刻瑞斯}在寻找女儿时疲惫不堪，来到\Place{阿提卡}边界时，带来了小麦，用鹿皮装饰了内布里德家族，并对\Person{包玻}腹股沟处最奇妙的景象发笑的那次漫游吗？或者如果有其他原因，那对我们也一样，只要它们都是由某种原因产生的。因为不能相信这些事是在没有任何原因在前的情况下发起的，否则必须认为阿提卡的居民是疯狂的，才会接受一个毫无理由被制造出来的宗教仪式。但如果这是清楚而确定的，即，如果这些秘仪的原因和起源可追溯到过去的事件，那么无论如何它们也不能被变成寓意的形象；因为已经成就的事，即已经发生的事，在事物的本性中是不能被撤销的。\par

40. 然而，即使我们承认你所说的这种情况，即，即使这些叙述在言辞上表达一件事，\Emph{却}意指另一件事，像是疯癫的先知那样，难道在这种情况下，你难道没有观察到，你难道没有看到，这种所谓的做法对众神是多么不敬、多么侮辱吗？或者说，还能设计出什么更大的错误，胜于将大地和雨水，或其他任何东西——因为解释上做什么改变无关紧要——称作\Person{朱庇特}与\Person{刻瑞斯}的交合？并且借对众神的控告来指称雨水自天降下和大地的湿润吗？还有什么比这更不虔敬的事能想到或相信吗：说\Person{普洛塞庇娜}的劫夺说的是种子埋入土中，或其他任何东西——因为同样无关紧要——并且说它谈论农业的操劳，却使父亲\Person{狄斯}蒙羞？难道不是比这好上一千倍，成为哑巴和不能说话的，并失去那种话流和嘈杂而不雅的饶舌，强似用众神的名字称呼最卑贱的事物；甚至，更甚者，用众神的卑贱行为来指称平常的事物吗？\par

41. 从前，在寓意地说话时，人们习惯于将那些公开谈论起来卑鄙而可怕的事，隐藏在完全正派的观念下，并披上端庄可敬的外衣；但现在，由于你的缘故，可尊敬的事被说得卑鄙，十分纯洁的事被用肮脏的语言讲述，以致从前罪恶出于羞耻而隐藏的事，现在却被卑贱而下流地谈论，合适的说话方式已被改变。在谈论\Person{玛尔斯}和\Person{维纳斯}被\Person{武尔坎}的计谋捉奸在床时，我的对手说，我们谈论的是情欲和愤怒，它们被理性的力量和意图所约束。那么，什么阻碍了你，什么阻止了你用每件事所特有的语词和术语来表达呢？不仅如此，当你已决定通过论文和著作来宣布某些事时，又有什么必要去决定那不该是你所指的意义，并在一个叙述中同时采取相反的立场——既有希望教导者的热心，又有不愿公开者的吝啬呢？谈论众神的不贞就没有风险吗？我的对手说，提到情欲和愤怒，可能用污秽的传染玷污舌头和嘴。但是，当然，如果这事做成了，且寓意的晦涩面纱被揭开，事情就会容易理解，而且同样众神的尊严也会毫发无损地得以维持。但现在，当恶习的约束被说成是由捆绑\Person{玛尔斯}和\Person{维纳斯}来指称时，两件完全不协调的事同时发生了；以致于，一方面，对某种卑鄙之事的描述暗示了一种可敬的意义，另一方面，卑鄙先于任何对宗教的尊重而占据了心灵。\par

42. 但你们或许会说——因为这是你们唯一会想到用来辩驳的了——天上的众神不愿让凡人知道他们的奥秘，因此这些叙述才被写成寓意隐晦的。你们从何得知，天上的众神不愿让它们的奥秘公之于众？你们从何熟悉了这些？或者，你们为何急于通过寓意解释来揭示这些？最后，归根结底，众神的意思究竟是什么：他们不愿人谈论可敬之事，却容许人谈论他们身上无礼的、甚至最卑劣的事？当我提起\Person{阿提斯}，\Emph{我的对手}说，我们指的是并且是在谈论太阳；可是如果\Person{阿提斯}就是太阳，正如你们所估算和声称的那样，那么你们的书中所记载和声称的那个出生在\Place{弗里吉亚}、经历了某些事、也做了某些事的\Person{阿提斯}，又会是谁呢？所有的剧场都在舞台表演中认识他，每年我们都在其他\Emph{宗教}仪式中，看到人们指名道姓地向他献上神圣的敬礼。这个名字究竟是从太阳转移到人身上，还是从人身上转移到太阳那里？因为如果这个名字最初源于太阳，那么，请问，金灿灿的太阳对你们做了什么，你们竟让他跟一个阉人共用一个名字？但是，如果它\Emph{源于}一只山羊，并且是\Place{弗里吉亚}的，那么，\Person{法厄同}的父亲，这光明与辉煌的本源，又犯了什么罪过，竟让自己看来配得上以被阉割的男人命名，并在被人以一个阉割之身的名称称呼时，变得更可敬了呢？\par

43. 这其中的含义，已经众人皆知了。因为你们为这样的作者和历史感到羞耻，却又不明白，那些曾以污秽的言语被记载下来的事，是无法消除的，于是你们竭力将卑鄙的事说成可敬的，并用尽种种诡辩，为了假借义的解释，歪曲败坏了词语的真实含义；并且，像常发生在病人身上的情形那样，他们的知觉和理解力被混乱的疾病势力逐散，你们也投出混乱而不确定的\Emph{猜测}，在空洞的虚构中胡言乱语。\par

就算\Emph{承认}，灌溉大地是由\Person{朱庇特}与\Person{刻瑞斯}的结合所意指的，种子的埋藏是由父\Person{狄斯}掳走\Emph{\Person{普洛塞庇娜}}所意指的，酒洒遍大地是由\Person{利柏耳}被\Emph{\Person{提坦}}们撕裂的肢体所意指的，那对贪欲与鲁莽的约束也被说成是捆绑奸淫的\Person{维纳斯}和\Person{玛尔斯}。\par

44. 但如果你们得出结论，认为这些神话是以寓意方式写成的，那么剩下的那些——我们看到的无法被强行转化为这类\Emph{意义}转变的神话——又该如何处理呢？当\Person{塞墨勒}之子那淫欲的燥热迫使他伏在坟冢上扭摆时，我们该拿什么来替换那些扭摆呢？对那些被掳去并被派去掌管淫欲之事的\Person{伽倪墨得斯}们，我们又该拿什么来替换？对那化身为蚁的变形，\Person{朱庇特}，\Emph{众神中}最伟大的一位，竟将庞硕身躯的轮廓缩小为蚁，又该拿什么来替换？天鹅与\Person{萨堤尔}呢？那同一位充满诱惑力的\Emph{神}以诡诈的欺骗穿上的黄金雨呢，他以形体的变换为乐？为使我们的批评看似不单单针对\Person{朱庇特}：其他诸神的情爱之事，又能有什么寓意呢？他们那些受雇为仆、为奴的情形，又能有什么寓意呢？他们的捆绑、丧亲、哀泣，能有什么寓意？他们的痛苦、伤痕、坟墓，又能有什么寓意呢？如今，你们在这方面，或可因以这种方式书写众神而被判有罪，但你们又以神明之名来称呼卑鄙之事，接着又借\Emph{赋予神明}丑事之名而亵渎了神明，这罪过更是大大加增了。然而，如果你们深信不疑地相信，他们就在近旁，或者在任何地方，那么恐惧便会阻止你们提到他们，你们的信仰和不变的意念，就应当恰好如同他们在聆听你们、听见你们的话一般。因为，在那些献身于宗教敬礼的人中间，不仅神明本身，就连神明的名字也应受到敬畏，这些名字中应具有与其名下所思想的存在本身同样多的威严。\par

45. 公正地判断吧，你们在这方面是应受指责的：在你们的日常谈话中，当你们意指战斗时，你们却称其为\Person{玛尔斯}；当你们意指海洋时，你们却称其为\Person{涅普顿}；当你们意指面包时，你们却称其为\Person{刻瑞斯}；当你们意指纺织时，你们却称其为\Person{弥涅耳瓦}；当你们意指污秽的贪欲时，你们却称其为\Person{维纳斯}。有什么理由，当事物都可以归入它们自己的名下时，却偏要以神明的名字来称呼它们？并且，这样的侮辱竟加予神明，即使是我们凡人，若有人将我们的名字用于琐碎之物，我们也是无法容忍的。但是，\Emph{你们说}，言语若是被这样的词汇玷污，就变得可鄙了。啊，多么值得称赞的庄重！你们羞于说出面包和酒，却毫不畏惧地用\Person{维纳斯}来指称肉欲的交合！\par

\SourceRecord{Translated by Hamilton Bryce and Hugh Campbell.；Ante-Nicene Fathers；Vol. 6.；Edited by Alexander Roberts, James Donaldson, and A. Cleveland Coxe.；Buffalo, NY: Christian Literature Publishing Co.,；1886.；<http://www.newadvent.org/fathers/06315.htm>.}

% structure: p0001 source=h1 target=section reason=page-title

\section{驳异教徒（卷六）}

1. 在简略指出你们对所信奉神祇形成的观点是\Emph{何等不虔敬且可耻}之后，我们现在必须谈论他们的庙宇、他们的偶像以及祭献，还有与之紧密相连的其他事物。因为你们惯于就此严厉指责我们为不虔敬：因为我们不为敬拜仪式建造庙宇，不竖立任何神祇的雕像与偶像，不修建祭台，不献上在\Emph{祭献}中被宰杀生物的血、\Emph{乳香}和祭餐，最后也不从神圣的杯中倾流奠酒；而这些建造和行为我们确实未做，并非因为我们心怀不虔敬和邪恶的性情，或对神祇产生了任何疯狂绝望的轻蔑之意，而是因为我们思量与相信，他们——倘若他们真是神，且堪当此崇高职分——或轻视此种敬礼，若他们流露轻蔑，或愤怒忍受\Emph{它们}，若他们被暴怒所激动。\par

2. 为了让你们得知我们对于这类存在的观念和意见——我们思量，他们——倘若他们真是神（这同一件事我们可以反复说，直到你们听厌为止）——应具备一切完善的德性，应是睿智、正直、可敬的——只要我们将人类荣誉加诸他们并非罪行——内在充满卓绝之善，且不应依赖外在支撑，因为他们那不间断的幸福之圆满得以完全；\Emph{应}远离一切纷扰不休的激情；不应怒火中烧，不应因任何欲望而激动；不应降祸于任何人，不应以人类之不幸为残忍之乐；不应以异兆恐吓，不应显示怪异以制造恐惧；不应使人为他们应偿的誓愿负责并受罚，也不应以凶兆相威胁而要求赎罪祭献；不应败坏空气而带来瘟疫疾病，不应以干旱烧毁果实；不应参与战争屠杀与毁灭城池；不应盼一方受损而促另一方成功；而应如伟大心灵之所当行，以公正天平权衡一切，平等对待所有人。因为，以别的方式行事，乃属有死的人间族群与人类软弱；智者的格言与断言明确声称，凡受激情触动者，过着受苦的生活，且\Emph{因忧伤而软弱}，且凡被交付于烦乱情感者，不可能不被必死的法则所束缚。既然如此，我们怎能被误认为轻视神祇，其实我们说的是，他们若不是正义且配得上伟大心灵所激发的钦佩，便不是神，也无法与天界大能相连？\par

3. \Emph{有人告诉说}，我们不给他们建造庙宇，也不敬拜他们的偶像；我们不宰杀牺牲献祭，我们不献乳香与奠酒。而我们能赋予他们什么更大的敬礼与尊荣，胜过将他们置于与宇宙的元首和主相同的地位？那些神祇与我们一同亏欠祂，才得以意识到自己存在，且拥有生命。难道我们以神龛或建造殿宇来敬拜祂吗？我们竟也向祂宰杀牺牲吗？我们献给祂其他物品，取用并浇奠这些，所表明的不是对理性的谨慎关切，而是对仅仅由习俗所保持的惯例之听从？因为，以你们的需求来衡量更伟大的权能，把对你们有用的东西献给那赐予万有的神祇，并且把这当成一种敬礼而非侮辱，实在是彻底的愚昧。因此我们要问，你们说建造神庙是为神祇提供什么服务，或满足什么需求，并认为应当继续建造呢？难道他们会感受冬日的寒冷，或被夏日炎阳灼烧？暴雨会淋在他们身上，旋风会摇撼他们吗？他们是否有遭受敌人袭击或野兽猛攻的危险，以致适宜且妥当将他们关在安全处所，或用石墙保护他们？这些神庙究竟是什么呢？若你们询问人类的软弱——则它们是宏大宽敞之物；若你们考虑神祇的能力——则它们仿佛是渺小的洞穴，甚至更真实地说，是以鄙陋判断建造和设计的极其狭窄的地窖。现在，若你们要人告诉谁是它们最初的创立者和建造者，有人会对你们提起\Person{福罗纽斯}或埃及的\Person{墨洛普斯}，或者如\Person{瓦罗}在他的\WorkTitle{论奇事}中所记述的，\Person{朱庇特}的后裔\Person{厄阿科斯}。然而，即使这些神庙由大堆大理石建造，或以镶金的天花板璀璨生辉，即使宝石在此闪耀，如同星辰疏密有致地闪烁着，这一切都不过是尘土和更低劣物质的糟粕所构成。因为，即便你们更加看重这些，也不可相信神祇会喜爱它们，或相信他们不拒绝、不鄙夷将自己封闭于这些界限内。\Emph{对手说}，这是\Person{马尔斯}的庙，这是\Person{朱诺}与\Person{维纳斯}的庙，这是\Person{赫拉克勒斯}的庙，是\Person{阿波罗}的，是\Person{狄斯}的。这不是等于说，这是\Person{马尔斯}的家，这是\Person{朱诺}与\Person{维纳斯}的家，\Person{阿波罗}住在这里，\Person{赫拉克勒斯}居于此间，\Person{苏曼努斯}住在那处吗？那么，将神祇囚禁于住所，给他们小屋，建造上锁的处所与密室，并认为这些是人、猫、蚂蚁、蜥蜴，以及颤抖胆怯的小鼠所需之物，这难道不是最大的侮辱吗？\par

\Emph{我的论敌}说，我们之所以为诸神修建庙宇，并非因为想为他们挡开倾盆暴雨、狂风、阵雨或阳光；而是为了能亲眼、近距离看见他们，走近并与他们交谈，并能在他们某种程度亲自临在时，表达我们的崇敬之情。因为，若是人们在天穹下、\OriginalTerm{以太}{ether}的穹顶下祈求他们，\Emph{我想}，他们便什么也听不见；而除非贴近他们祈祷，他们就会像什么也没说一样，充耳不闻，无动于衷。然而，我们认为，任何一位神——只要他配得上这个称号——都应听见来自世界各地每个人所说的话，如同亲临一般；不但如此，还应无需告知便能预知每个人在秘密沉默的意念中所思所想。就如星辰、日、月，高悬于大地之上，恒久且无处不在地被所有仰望它们的人看见，无一例外；同样，神的耳朵理应不对任何语言封闭，永远触手可及，即便声音从相距遥远的各地同时汇集于他们。因为这才是神特别应有的本性——以他们的能力充满万物，不是部分在某处，而是全部无处不在，不会去与埃塞俄比亚人共进晚餐，并在十二天后返回自己的住所。\par

现在，若非如此，一切获得帮助的希望便化为乌有，你每次以合宜的礼仪举行神圣仪式时，是否被神垂听，便值得怀疑了。为明确起见，我们假设，在\Place{加那利群岛}有某位神明的庙宇，同一神明的另一座庙宇在极远的\Place{图勒}，又在赛里斯人中，在棕肤的加拉曼特人中，以及任何因海洋、山脉、森林和世界四方而相互隔绝不为人知的民族中。倘若他们同时以祭品向神明祈求各自所需，若神不能听到从各处向他发出的呼求，且有任何距离使求救者的言辞无法抵达，那么，我请问，所有人获得恩惠的希望又在哪里呢？因为，他要么根本不在任何地方，如果他可能有时不在某处，要么他只会在一处，因为他无法普遍而无差别地加以关注。于是，结果就是：要么神根本不帮助任何人，如果他忙于某些事而未能迅速垂听他们的呼求；要么只有一人得蒙俯允，而其余的人则一无所获。\par

对此，你还能说什么呢？著作家们的作品证实，许多用金色穹顶和高耸屋顶建起的庙宇，遮盖的却是骸骨与骨灰，实际上是死者的坟墓。难道这不是明明白白的吗：要么你们把必死的人当作不朽的神崇拜，要么对神祇施以不可弥补的侮辱，因为他们的圣所与庙宇竟建在死者的坟墓之上？\Person{安条克}在\WorkTitle{历史}第九卷中记载，\Person{刻克罗普斯}葬于\Place{雅典}的\Person{密涅瓦}神庙；又说，在\Place{拉里萨}卫城的同一女神庙中，据称埋葬着\Person{阿克里西俄斯}，而在\Person{波利阿斯}圣所内则葬着\Person{厄里克托尼俄斯}；\Person{代拉斯}与\Person{伊姆马纳库斯}兄弟则葬于靠近城市的\Place{厄琉西斯}围场。关于\Person{科勒俄斯}未婚的女儿们，你有何话说？她们不是据说葬在\Place{厄琉西斯}的\Person{刻瑞斯}神庙中吗？而在\Place{德洛斯岛}\Person{阿波罗}神庙内所建的\Person{狄安娜}的神龛中，不是埋葬着\Person{许珀洛刻}与\Person{拉俄狄刻}吗？据说她们是从极北族之地被带到那里的。在\Place{米利都}的\Place{狄底马}，\Person{利安德里乌斯}说，\Person{克勒俄科斯}以最后的葬礼荣典安葬。\Place{明杜斯}人\Person{芝诺}公开述说，\Person{琉科弗律涅}的墓碑在\Place{马格尼西亚}的\Person{狄安娜}圣所内。在\Place{忒尔墨索斯}城可见的\Person{阿波罗}祭坛下，难道不是众口一词地记载着预言家\Person{忒尔墨索斯}葬于此吗？\Person{阿革萨尔库斯}之子\Person{托勒密}，在他所著的\WorkTitle{菲洛帕托尔史}第一卷中，根据文献证实，\Place{帕福斯}王\Person{喀倪拉斯}与其全部家人，甚至所有后代，均同葬于\Person{维纳斯}神庙中。若要描述全世界它们所有所在的圣所，那将是一项无穷无尽的工作；也无需过于焦虑，尽管埃及人曾规定，凡泄露\Person{阿匹斯}藏匿之处的人将受惩罚，至于\Person{瓦罗}的\WorkTitle{多人葬所}，它们被哪些庙宇遮盖，上面又堆积了多少沉重的土石，便也无须深究了。\par

7. 但我为何\Emph{谈论}这些琐事呢？有谁不知在帝国人民的\Place{卡比托利欧}山丘上，有着\Person{托卢斯·武尔肯塔努斯}的坟墓呢？我说，有谁不知道从这地基下面滚出一个人头，埋藏不久，要么是单独无躯——有些人这样传说——要么是连同全部肢体？现在，如果你们要求借著作者的证词澄清此事，\Person{萨莫尼库斯}、\Person{格拉尼乌斯}、\Person{瓦勒里安努斯}和\Person{法比乌斯}会告诉你们\Person{奥卢斯}是谁的儿子，属于何种族和民族，他如何被他兄弟的奴隶夺去生命和光明，他对同胞犯下何等罪行，以致被拒绝葬于故土。你们还会知道——尽管他们假装不愿公开——他的头被割下后如何处置，或被封存在何处，以及整件事如何被小心隐藏，以便神明所证实的预兆可以不受干扰、永恒不变、确定无疑。然而，本应被压制、隐藏并随时间遗忘的这个\Emph{故事}，却因组合名称而公开，并以无法摆脱的见证，使其连同原因一起留在\Emph{人们的记忆中}，只要它本身持续存在；而那个\Emph{最伟大}的、崇拜一切神祇的国家，在给神殿命名时，竟不以\Person{朱庇特}之名，却从\Person{奥卢斯}的头颅命名为\Place{卡比托利欧}，而不觉羞愧。\par

8. 因此，我们——如我所想——已充分表明，为不死的神明建造的神殿，要么是枉然的，要么是基于\Emph{所持的}侮辱性看法而建的，有损其荣耀并贬低人们相信\Emph{在它们手中的}权能。接下来，我们要谈一谈雕像和肖像，你们以高超的技艺塑造，并以宗教般的虔诚加以照料——在这件事上，若有任何可信之处，我们无论如何深思熟虑，都无法断定你们是认真、严肃地这样做，还是以幼稚的幻想嘲弄这些事物来自娱。因为如果你们确信自己所设想的那些神祇存在，并且居住在天界的最高处，那么，当你们已有真实的存在可以直接向它们殷切祈祷，并\Emph{从它们那里}在困境中求助时，还有什么理由、什么原因，让你们去塑造这些肖像呢？反之，如果你们并不相信，或者说得克制些，心存怀疑，那么，请问又有什么理由，去塑造和竖立这些有疑问的\Emph{存在物}的肖像，以徒劳的模仿去塑造你们不相信存在的东西呢？你们或许会说，在这些神明的肖像之下，它们的存在仿佛就显示给你们；因为你们未被准许看见神祇，所以以这种方式崇拜它们，并将应尽的义务献\Emph{给它们}。说这话、断言这事的人，并不相信神祇存在；凡必须看见他所持守的，便证明他不信赖自己的宗教，以免那因\Emph{是}隐晦而不可见的事物，或许竟是虚空。\par

9. 你们说，我们藉着肖像崇拜诸神。那么，若没有这些，诸神就不知道它们被崇拜了吗？它们会认为你们没有给予任何荣耀吗？如此说来，它们是通过迂回途径，经由所谓的第三方指派，来接收并接受你们的服侍；而在那些应受服侍的存在体验到之前，你们先向肖像献祭，并如同经过别人的许可，才将些许残余传予它们。还有什么更大的不公、羞辱、苦难，能比这更甚呢：承认一位神明，\Emph{却}向别的东西祈祷——希望从神祇得到帮助，却向没有感觉的肖像恳求？请问，这不正像俗谚中所说的吗：\Quote{击打漂洗工，却砍倒了铁匠；}\Quote{寻求人的忠告时，却向驴和猪请教该做什么？}\par

10. 最后，你们怎么知道你们所塑造并用来取代不朽神明的所有这些像，是否再现神明并与之相似？因为在天上，可能某位神原有胡须，却被你们表现为光洁的面颊；\Emph{另一位}可能年事已高，却被你们赋予青年的外貌；此处的那位是金发\Emph{碧眼}的，而真实的却可能是灰眼；那位鼻孔宽大，却被你们塑造为高鼻梁。因为那种并非源自原初面容特征、不能从明显特征中清晰明确辨认出的像，是不应当称之为像的。因为尽管我们人类肉眼所见的太阳是完美的圆形——这一点不容置疑——你们却将它赋予了人的面貌和可朽坏的身体的特征。月亮总是在运行，每月恢复时呈现三十种面貌：而在你们这些引领者和设计者手中，它被\Emph{表现为}一个女人，仅有一个面容，却每日变化，历经千种不同状态。我们理解，所有的风\Emph{仅}是空气的流动，被以世俗的方式驱动和推动：而在你们手中，它们采取人的形状，用胸中鼓出的气息吹满扭曲的号角。在你们神明的\Emph{像}中，我们看到一张极其严厉的狮子面孔，涂满纯朱砂，它被命名为\Person{\OriginalTerm{弗鲁吉弗}{Frugifer}}。如果所有这些像都是天上神明的肖像，那么天上也必然住着一位神，就像为表现其形态外貌而制作的像一样；当然，正如此处你们那个形象，那里的神明本身也只是一副面具和面孔，没有身体的其他部分，怒张着可怕的大口咆哮着，狰狞可怖，红如鲜血，用牙齿紧紧咬住一个苹果，有时，像疲惫的狗一样，从张开的口中伸出舌头。但如果其实并非如此——正如我们都认为的那样——那么，请问，如此胆大妄为，你们随心所欲地为自己塑造任何形式，并称之为一个你们根本无法证明其存在之神的像，意义何在呢？\par

11. 你们嘲笑古人，因为据说\Place{波斯}人崇拜河流，正如那些将此传于后世的著作中所记载的；\Place{阿拉伯}人崇拜一块未经雕琢的石头；\Place{斯基泰}民族崇拜一把军刀；\Place{特斯皮亚}人崇拜一根树枝以代替\Person{\OriginalTerm{辛克西娅}{Cinxia}}；\Place{伊卡里亚}人崇拜一段未加工的木料以代替\Person{\OriginalTerm{狄安娜}{Diana}}；\Place{培西努}的居民崇拜一块燧石以代替众神之母；\Place{罗马}人崇拜一支长矛以代替\Person{\OriginalTerm{玛尔斯}{Mars}}，正如\Person{\OriginalTerm{瓦罗}{Varro}}的缪斯所指出的；以及，据\Person{\OriginalTerm{埃特利乌斯}{Aethlius}}所述，\Place{萨摩斯}人在掌握雕像艺术之前崇拜一块木板以代替\Person{\OriginalTerm{朱诺}{Juno}}：而你们却不嘲笑自己，因为你们不是向不朽的神明，而是向小人像和人的形象祈祷——不仅如此，你们甚至认为这些小人像就是神，除此之外，你们不相信任何事物具有神圣的能力。你们有什么可说的呢，啊，你们——！那么，天上的神有耳朵、太阳穴、后脑勺、脊梁、腰、肋、腿窝、臀、膝关节、脚踝，以及我们被造时所具有的其他肢体吗？这些已在本书第一部分更详细地提及，并以更丰富的语言引证过。但愿有可能窥见你们内心的感受和最深的隐秘，你们在其中反复思量，陷入最隐晦的考量：我们会发现你们自己也和我们有同样的感受，对神明的形象也没有别的看法。然而，面对固执的偏见，我们能做什么？面对那些用刀剑威胁我们、并不断想出新的刑罚来对付我们的人，我们能做什么？你们在狂怒中坚持一个糟糕的主张，尽管你们对此心知肚明；你们为自己曾无理做过的事辩护，以免显得自己曾无知过；你们认为不被征服，胜于屈服并顺从已被确认的真理。\par

12. 由于这些缘由，在你们的默认下，又出现了这样的情况：艺术家们任性的幻想在\Emph{描绘}神祇的身体并赋予其形态时得到了充分施展，即便是最严肃的人也可能会发笑。因此，\Person{哈蒙}至今仍被塑造和描绘为长着公羊的角；\Person{萨图恩}手持弯镰，如同一个田地的看守人，\Emph{并且}是修剪过于茂盛枝条的人；\Person{迈亚之子}头戴一顶阔边旅行帽，好像准备赶路，躲避阳光与尘沙；\Person{利柏尔}肢体柔弱，而\Person{维纳斯}有着女人般完全放纵自如、线条流畅的身体，赤裸而不挂一丝，仿佛你公然宣称，她将自己那被玷污的身体之美公然展示，并卖给所有来者；\Person{伏尔甘}戴着便帽，持着锤子，右手空着，衣裳束起，如同工人准备做工；\Person{德洛斯神}手握琴拨与里拉，像一名基萨拉琴师和即将歌唱的演员般作着手势；\Person{海神}手持三叉戟，仿佛要在角斗比赛中厮杀一般：任何一位神祇的形象，都无法找到没有其制作者慷慨赋予\Emph{于其上}的某些特征的。且看，如果某个机智狡黠的国王，将\Person{太阳神}从他在大门前的\Emph{位置}移开，把他挪到\Person{墨丘利}的位置上，然后再将\Person{墨丘利}搬走，让他迁入\Person{太阳神}的神殿——因为两者都被你们制作成无胡须的，面容光滑——并给这一个赋予\Emph{光}芒，给\Person{太阳神}的头上戴上一顶小帽，那么你们又将如何能够区分他们，究竟这位是\Person{太阳神}，还是那位是\Person{墨丘利}，因为通常是指出神祇给你们看的，是服饰，而非面孔的独特特征？再者，如果以同样的方式挪移神像之后，他从赤裸的\Person{朱庇特}身上取下\Person{伊里斯}的角，把它们安在\Person{马尔斯}的额上，又剥去\Person{马尔斯}的武装，反过来再把那些武装加给\Person{哈蒙}，那么他们之间又能有何区别呢？因为曾经是\Person{朱庇特}的，也可能被认作\Person{马尔斯}，而曾经是\Person{马约尔斯}的，也可能呈现\Person{朱庇特·哈蒙}的面貌。对塑造这些神像并祝圣其名字的放肆竟到了如此地步，仿佛这些名字\Emph{是}其独有的；因为，若除去他们的服饰，辨认每位神祇的\Emph{方法}也就被消除了，神祇可能被认作神祇，一位可能看似另一位，甚至，两者都可能被认作两者！\par

13. 然而，我为何要嘲笑赋予众神的镰刀和三叉戟呢？为何嘲笑那些角、锤子和帽子呢？因为我知道某些神像具有某些男子的形体，以及声名狼藉的妓女的面容。有谁不知道，雅典人按\Person{阿尔西比亚德}的相貌塑造了赫尔墨斯柱像呢？有谁不知道——若他重读过\Person{波西迪普斯}的话——\Person{普拉克西特列斯}竭尽技艺，模仿妓女\Person{格拉蒂娜}的面容塑造了\Person{克尼多斯的维纳斯}的脸，而这位不幸的男人竟绝望地爱着她？然而，难道这就是唯一一个以妓女的面容赋予了美的\Person{维纳斯}吗？据说，\Person{弗里内}——那位著名的\Place{塞斯比亚}人——正如那些撰写\Emph{关于}\WorkTitle{塞斯比亚事物志}的人所记载的，当她处于容貌、娇美与青春活力的巅峰之时，无论希腊各城邦，还是此地，所有的\Person{维纳斯}神像，凡是\Emph{受}人珍视的，都以她为模型，人们对这种雕像的渴慕和热望也涌流而至。因此，所有当时在世的艺术家，那些其忠实表现赋予了他们描绘肖像最高能力的人，都竞相以全部的细心和热忱，将一个妓女的轮廓移植到\Person{库塞拉女神}的神像上。艺术家们美妙的\Emph{构思}充满了火热；他们彼此争先，互相竞赛，不是为了\Person{维纳斯}变得更庄严，而是为了\Person{弗里内}能够代表\Person{维纳斯}。事情就这样发展到，将神圣的尊荣献给了妓女，而非不朽的神祇，而一种不幸的崇拜体系也因制造雕像而误入歧途。那位著名而最杰出的雕塑家\Person{菲迪亚斯}，当他以极大的劳动和努力塑造了\Person{奥林匹亚的朱庇特}的形象之后，在神像的手指上刻下了\Quote{\Emph{潘塔克斯}\Emph{是}\Emph{美少年}}——\Emph{这}，正是他所爱的一个男孩的名字，且是带着淫欲的爱——他竟没有任何恐惧或宗教敬畏，不敢以妓女之名来称呼神祇；不，他反而将朱庇特的神性和神像祝圣给一个浪荡子。至于塑造这些小型神像，将其当作神祇来崇拜，把种种神性堆砌其上，这种行为竟到了如此放肆和幼稚的地步，我们看到艺术家们自己在塑造它们时以此为乐，并将其树立为自己情欲的纪念碑！因为，请问，有什么\Emph{理由}\Person{菲迪亚斯}会犹豫不去自娱自乐、放肆妄为呢？当他明知，仅仅片刻之前，他所\Emph{制作的}朱庇特本身只是黄金、石块和象牙，无定形、分离、混乱的，是他自己将所有这些组合并固定起来，它们的形状是他模仿所雕刻的肢体而赋予的；而且，比一切都重要的是，\Emph{朱庇特}之被造出并受人崇拜，这完全是他自己的慷慨赠予呢？\par

14. 我们在此，仿佛世上万民都亲临其境，要说一番话，灌入他们所有人的耳中，这话当为众人所共闻：人啊，请问，这是何故？你们竟自愿盲目地欺骗自己？现在驱散黑暗，回归心灵之光，更仔细地看清那正在发生的事，只要你们还保有理智，尚未超出赋予你们的理性和明达所及的范围。那些令你们充满恐惧，在各庙宇中俯伏敬拜的像，不过是骨骸、石头、铜、银、金、黏土、从树上取得的木头，或是掺了石膏的胶。它们或许是从娼妓的装饰品或妇人的饰物、骆驼骨或印度野兽的牙齿、烹锅和小罐、烛台和灯盏或其他更不洁的器皿中堆聚而来，\Emph{并且}被熔化后，铸成这些形状，呈现出你们所见的模样，这些像在陶匠的窑中烧制，经铁砧和锤子打造，被银匠刮削，又用\Emph{普通}锉刀锉平，用锯子、螺丝钻、斧子劈开削平，经钻头钻挖掏空，\Emph{并且}用刨子刨光。这难道不是错谬吗？说得准确些，你们相信自己费心制成的就是\Emph{那}位神明，向那出自你手的作品跪伏颤栗、哀恳祈求，你们明知且确知那是你双手劳苦的成果——却仍俯伏于地，恳求援助，在逆境困苦时请求它施以仁慈和神圣的恩惠来救助\Emph{你}，这难道不是愚妄吗？\par

15. 看，若有人将未经任何\Emph{艺术}加工的块状铜、未经锻造的银块、未塑形的金块，以及木头、石头、骨骸，连同所有其他通常用以制造神像和神明肖像的材料，放在你们面前——更有甚者，若有人将毁损之神的面容、被熔化且破坏的像放在你们面前，并命令你们向这些碎片残块宰杀牺牲，将神圣的敬礼献给无形式的团块——我们要请你们告诉我们，你们是否愿意这样做，还是拒绝服从。或许你们会说，为什么呢？因为没有人会愚蠢盲目到将银、铜、金、石膏、象牙、陶土归入神灵之列，并声称这些东西本身具有并拥有神能。既然如此，这些物体在保持原状、未经雕琢时缺乏神性和上天尊位，\Emph{但}一旦它们获得人的形状、耳朵、鼻子、面颊、嘴唇、眼睛和眉毛，就立刻变成神明，并被列入天上的居民之中，这又有什么道理呢？赋予外形是否给这些物体增添了任何新质，以致你们不得不相信，有某种神圣而庄严的东西已与它们结合？它是否将铜变成金，或迫使无价值的陶器变成银？它是否能使片刻之前尚无感觉的东西活起来、呼吸起来？如果它们先前有任何自然属性，当被堆砌成雕像的形体时，它们仍保留着这一切。认为事物的本质因\Emph{被迫}采用的外形类别而改变，那原本在其原始躯体中呆滞、无理性、对感觉无动于衷的东西，竟从赋予它的外观获得神性——这是何等愚蠢（我拒绝称之为盲目）啊！\par

而你们，作为有理性的受造物、被赋予智慧与审慎之恩赐的人，却如此忘却并疏忽了偶像的实质与起源，在烧制的陶器前俯伏，朝拜铜板，向象牙祈求健康、官职、王权、能力、胜利、获取、收益、丰美的收成以及最丰富的葡萄收成；而尽管你们明显\Emph{且}清晰地在向无感觉的事物说话，你们却自以为被垂听，并甘心情愿地自取其辱，徒然且轻信地自欺。啊，但愿你们能进入某尊雕像！不如说，但愿你们能将那些\Place{奥林匹亚}和\Place{卡比托利欧}的\Person{朱庇特}像拆散、打碎，亲眼观看那些单独并独自构成其整个身躯的所有部分！你们立刻会看到，你们那些\Emph{其}外表的平滑因诱人光泽而赋予庄重外貌的神祇，\Emph{只是}柔韧薄片的框架，是无定形微粒的拼合体；它们全凭鸠尾榫、夹具和铁箍才免于坍塌和毁灭的恐惧；而铅被灌入所有空洞的中央以及接合之处，造成了一种有利于保存它们的延迟。我告诉你们，你们会立刻看到它们\Emph{仅有面孔}而无头的其余部分，有残缺的手而无手臂，腹部和侧身是半边的，脚是残缺的，而最可笑的是，\Emph{它们}的身体构造是被不均衡地拼合起来的，一部分是木头，另一部分却是石头。如今，确实，即使这些因被巧妙掩藏而无法得见，至少那些公开暴露于众人眼前的，本该已教导并训诫你们，你们是在徒劳，是在白向死物献上侍奉。因为，在这事上，你们难道没有看见这些似乎有气息的偶像，就是你们在祈祷时触摸并摩挲其脚膝的，有时因雨水不断滴落而倒塌为废墟，有时因腐烂朽坏而失去各部分的牢固结合——它们如何因祭献的蒸汽和烟雾而变黑、褪色——如何因持续疏忽而失去其姿态与外观，并被锈迹侵蚀？在这事上，我说，你们难道没有看见蝾螈、鼩鼱、老鼠和蟑螂，这些躲避光明的生物，在这些雕像的空洞部位下筑巢而居吗？它们小心翼翼地将各种污秽，以及其它适合其所需的东西——坚硬且被啃了一半的面包、因预想匮乏而拖来的骨头、破布、羽绒和碎纸片——收集其中，以使它们的巢穴柔软，并保持幼崽温暖吗？你们难道没有看见，有时在偶像面孔上，有蜘蛛织成的蛛网和险恶的罗网，为要缠住那些在飞行中嗡嗡作响且不知戒备的飞虫吗？最后，你们难道没有看见满身污秽的燕子，在庙宇的穹顶下飞舞，翻飞嬉戏，时而将污秽拉到神祇的面孔、嘴、胡须、眼睛、鼻子以及所有被排泄物落上的其它部位吗？那么，即使为时已晚，也当脸红，从无言的生灵那里接受真实的道理和观察，让它们教导你们：在偶像中毫无神性，因为它们（这些生物）依照其生命的律法，被无误的本能引导，在这些偶像内拋置不洁之物而毫无恐惧或顾忌。\par

但是你们错了，\Emph{我的对手说}，你们误会了，因为我们并不认为铜、金银或其它制作雕像的材料在它们自身中即为神明和神圣的神祇；而是在它们之中，我们敬拜并尊崇那些因其\OriginalTerm{祝圣}{dedication}被引入、并居住于工匠所造雕像内的神明。这个推理并不邪恶或可鄙，借着它，任何人——无论愚钝者还是最聪明的人——都可以相信：神祇们离弃了自己本然的居所——即天堂——并不退缩回拒进入尘世的居所；更有甚者，在\OriginalTerm{祝圣礼仪}{rite of dedication}的推动下，他们与偶像结合在一起。那么，你们的神祇是居住在石膏和陶土形象中吗？不，不如说，神祇竟是陶土和石膏形象的心智、精神和灵魂吗？而且，为了使最卑微的事物能变得更重要，他们容许自己被关锁、隐藏、禁闭在阴暗的居所里吗？在此，首先，我们渴望并要求你告诉我们这一点：他们这样做是违心的——即，他们进入作为居所的偶像，是被祝圣礼仪拖拽而入——还是他们自愿乐意的？你们难道不是以任何必然性的考虑来召唤他们吗？他们这样做是不情愿的吗？那么，怎么有可能在不损害其尊严的情况下，他们被强逼而\Emph{服从}任何必然性呢？是甘愿同意吗？神祇们在陶土形象中寻求什么，以至于他们宁愿要这些牢笼而不要星空中的住所——以至于，他们几乎被固定其上，就使陶土以及制作偶像的其它物质变得高贵了？\par

18. 那么怎样呢？神祇们是否总是停留在这样的形体中，即使有最重大的事务召唤，也不离往任何地方？还是他们可以自由通行，想往哪里去就往哪里去，并离开自己的庙座与神像？倘若他们不得不留在其中，还有什么比他们更可怜，比被钩子和铅索牢牢缚在基座上更不幸的呢？但\Emph{如果}我们承认他们喜爱\Emph{这些神像}胜过上天和星辰的居所，他们便失去了神性。反之，倘若他们愿意时，便能飞出，完全可以离开神像使之空虚，那么这些神像便将在某个时候不再是神，何时应当奉献祭品、何时应当停止便成为疑问。我们常见工匠们将这些神像时而造得很小，缩至手掌大小；时而又造得极高，建成惊人的尺寸。由此，我们便应理解，神祇们在小雕像中收缩自身，被压缩得如同一个陌生的形体；或者，\Emph{他们}又会伸展至极长，在体积巨大的神像中扩至无限。那么，若果真如此，在坐姿雕像中，神祇也应被说成是坐着的；在站立的雕像中，则是站着的；在那些向前伸展作奔跑状的雕像中，是在奔跑；在那些\Emph{表现为}投掷标枪的雕像中，是在投掷标枪；并使自身适应和塑造得与雕像的面容一致，与\Emph{工匠}所造身体的其他特征相似。\par

19. 神祇居住在神像内——是每位完整地住在一尊像内，还是被分为各部分、各肢体？因为一位神祇既不可能同时存在于多尊神像内，也不可能因切割而被分割成部分。让我们假设普世有一万尊\Person{伍尔坎}的神像：如我所说，一位\Emph{神}怎么可能同时存在于这一万尊像内？我不\Emph{这样}认为。\Emph{你问}为什么？因为本性单一而独一的事物，在其单纯之完整得以保持时，不可能变为多。而且，如果神祇具有人的形体，正如你们的信仰所宣称，他们更不可能\Emph{变成}多个；因为一只手若与头分离，或一只脚若与身体分开，便无法彰显整体的完美；要么必须承认部分可以与整体相同，而整体若非由部分聚集而成便不能存在。再者，如果同一位\Emph{神}被说成是在\Emph{所有神像}中，那么若假定一位神能同时存留于\Emph{它们全部}，真理便丧失了一切合理与健全；或者，必须说每一位神祇都从自身分裂出自己，以致他同时是自身又是另一位，并非由任何区别分开，而是自身与另一位相同。然而，本性排斥、摒弃并嘲笑这种说法，因此必须要么承认并宣称有无数个\Person{伍尔坎}，若我们断定他存在于所有神像中；要么他不在任何一尊像内，因为他的本性不容许在多个形体中被分割。\par

20. 然而，你们啊——如果你们清楚明白地知道神祇活着，天上的居民居住在神像的内部，为什么你们还用最坚固的锁钥、巨大的门闩、铁棍、栓子以及其他此类东西把神像牢牢关锁、保护起来，又用成千上万的男女看守防备，惟恐有什么盗贼或夜窃者溜入？你们为什么在神殿中喂养狗？为什么给鹅提供食物和滋养？其实，如果你们确信神祇就在那里，且从不离开自己的雕像和神像往任何地方去，那就让他们自己看管自己吧，让他们的圣所永远不上锁、敞开着吧；若有人以冒失的欺诈悄悄偷走任何东西，就让他们显示神性的大能，在偷窃和\Emph{邪恶}行径发生的当时，将亵渎的盗贼施以应得的惩罚。因为，让至高神祇的护卫交付给狗来看管，当你们寻求某种恐吓盗贼以使其远离的方法时，不去祈求\Emph{神祇}自身，却依赖于鹅的咯咯叫声，这不但不成体统，也损毁了他们的能力和威严。\par

21. 他们说，\Place{基齐库姆}的\Person{安条克}从神殿取走一尊十肘\Emph{高}的黄金\Person{朱庇特}神像，代之以一尊镀薄金片的铜像。如果神祇确实临在并居住于自己的神像内，朱庇特是被何种事务、何种牵挂所缠扰，以致无法惩罚对自己的伤害，无法报复自己被替换成较贱金属之事？当著名的\Person{狄奥尼修斯}——不过是\Emph{那}较年轻的——夺去朱庇特的金袍，代之以羊毛袍子，并\Emph{且}以戏谑嘲弄\Emph{他}时，他说，\Emph{他所取走的那件}在冬日的霜冻中寒冷，而这件暖和；又说那件在夏天闷热，而这件在炎热天气中又透气——身为世界之王的朱庇特，为何不以某种可怕的行为显示其临在，并以痛苦的折磨让那滑稽的小丑恢复清醒？何必需要我提及\Person{埃斯库拉庇乌斯}的尊严被他嘲弄？因为当狄奥尼修斯夺去他那极为丰盈的胡须——\Emph{那胡须}既重且具哲学家的浓密——时，他说，一位出自\Person{阿波罗}的儿子，父亲却是光滑无须、极像一名纯粹男孩，竟被塑造出如此一部胡须，以致令人无从分辨他们中谁是父亲、谁是儿子，甚或是同一种族和家族，这是不合适的。然而，当这一切发生时，这盗贼以不虔的嘲弄发言，若这位神祇隐藏于奉献给他的名与威仪的神像内，为何他不用公正且应得的惩罚来报复面貌被剃须、面容被丑化之辱，并藉此表明他自身确实临在，且从不间断地看顾着自己的殿宇与神像？\par

22. 但你们也许会说，众神并不为这些损失烦恼，也不认为有足够理由出来惩罚那些犯下不敬亵渎之罪的人。那么，如果情况如此，他们也不希望拥有这些任由人连根拔起和扯走而不受惩罚的偶像；不，恰恰相反，他们明白地告诉我们，他们鄙视这些\Emph{雕像}，因为他们甚至不愿通过任何报复来表明自己受蔑视。\Person{菲罗斯特凡诺斯}在其\WorkTitle{塞浦路斯志}中记述，\Place{塞浦路斯}王\Person{皮格马利翁}热爱一座\Person{维纳斯}像犹如爱女人，这像自古以来被塞浦路斯人视为神圣可敬；他心智、灵魂、理智之光以及判断力都昏昧了；他在疯狂中，仿佛对待妻子一般，把神像抬上他的床榻，与之拥抱、\Emph{面对面}结合，并做其他虚妄之事，他被愚蠢的淫荡想象\Emph{所驱使}。同样，\Person{波西迪浦斯}在他提及\Emph{曾}写下的关于\Place{克尼多斯}及其事务的书中记述，一位出身高贵的青年——但他隐去其名——因爱上那座使克尼多斯闻名的\Person{维纳斯}像，亦在欢娱之榻上与这同一神像行淫荡之事，享受随之而来的快感。再以同样方式追问：如果上方众神的权能潜伏于铜与其他构成偶像的物质中，那么这尊维纳斯和那尊维纳斯究竟在哪里，能从那两个青年身上远远驱走淫荡的放肆，并以可怖的痛苦惩罚他们不敬的触摸？或者，既然女神们性情温和，较为平静，那么她们要做些什么，才能平息那些可怜的人狂怒的喜乐，使他们疯狂的头脑恢复理智呢？\par

23. 但也许，如你们所说，女神们对这些淫荡好色的侮辱感到极大愉悦，并不认为需要采取报复行动，因为这使她们心神舒畅，她们知道这是由她们自己启发给人类欲望的。然而，如果这些女神，即维纳斯们，天生性情相当平和，认为应当对昏了头的青年们的苦难予以恩惠；当贪婪的火焰曾如此频繁地吞噬\Place{卡庇托林山}，并连同其妻女一同摧毁了卡庇托林山的\Person{朱庇特}时，\Person{雷霆之神}当时在哪里，来避免那场灾难性的大火，救他的财产、他自己以及他的全家脱离毁灭呢？当猛烈的火灾在\Place{阿尔戈斯}摧毁了她著名的庙宇及其女祭司\Person{克律西斯}时，庄严的\Person{朱诺}又在哪里？当出于类似的灾难，埃及的\Person{塞拉皮斯}，连同所有奥迹，以及\Person{伊西斯}，他的庙宇被焚为灰烬时，他在哪里？当\Person{自由者利贝尔}在\Place{雅典}庙宇坍塌时，他在哪里？当\Person{狄安娜}在\Place{以弗所}庙宇坍塌时，她在哪里？当\Place{多多纳}的\Person{朱庇特}在多多纳庙宇坍塌时，他在哪里？最后，当先知\Person{阿波罗}遭到海盗和海贼抢劫并纵火，以至于从无数世代积聚起来的那么多磅黄金中，他竟连一微量也拿不出，展示给在他岩洞下筑巢的燕子，正如\Person{瓦罗}在\WorkTitle{梅尼普斯式讽刺诗}中所说的那样，他又在哪里？要写下在整个世界有哪些庙宇因地震和暴风雨而毁坏、有哪些被敌人及君主和暴君纵火焚烧、有哪些被庙宇监督者及祭司们自己剥得精光，即便他们转移了嫌疑——最后，有哪些被窃贼和\Place{卡纳凯尼人}以未知的方式撬开而洗劫一空，那将是一项无止境的工作；确实，如果众神真如所言在保护他们，或对他们自己的庙宇有任何顾念，那么这些庙宇本应安然无恙，不受任何灾祸侵凌。但如今，因为这些庙宇是空的，没有任何内居者保护，命运便有权柄控制它们，它们就如一切无生命之物一样，面临着一切意外。\par

24. 这里，偶像的辩护者们也惯常说这话：古人十分清楚偶像并无神性，其中也没有感觉，但他们为了那些难以管束和无知的群众——这些群众构成了各族各国的大多数——才明智地立了偶像，为的是给群众呈现一种神祇的似是非是的表象，好让他们因恐惧而摆脱粗暴的性情，并设想自己是行在众神面前，从而丢弃不敬的行为，改变作风，学习像人一样行事；并且，为他们寻求金银的庄严形像，不出于其他理由，只因相信这些光辉中有某种力量，不仅眩惑眼目，甚且能以威严闪烁的光芒令心灵本身颤栗。这话也许还能被认为稍有道理，倘若在众神的庙宇建立、他们的偶像树立之后，世界上便没有了恶人，根本没有了任何恶行；倘若正义、和平、诚信占据了人心，世上没有人被称为有罪或无罪，因为所有人都对恶行无知。但是现在，恰恰相反，一切都充满了恶人，无辜之名几乎湮灭；每一刻，每一秒，前所未闻的恶行都因作恶者的邪恶而成千上万地涌现；我们又如何能说，设立偶像的目的是要恐吓群众，而同时，除了数不清的罪行和邪恶外，我们竟看到连庙宇本身都遭到暴君、君主、盗贼和夜贼的袭击，以及看到那些经古人塑造并祝圣以制造恐惧的众神本身，竟被带进盗贼的洞穴，即便有着黄金的可怖光辉？\par

25. 如果你不带偏见地审视真理，他们所说的这些形象之中，究竟有什么庄严伟大之处，竟让古人会有理由希望并认为，只要观看它们，人的恶习就能被制伏，品行与邪恶的行径能受到约束呢？例如，分配给了\Person{撒图恩}的镰刀，难道是要激发凡人的恐惧，使他们愿意平安度日，弃绝恶毒的倾向吗？\Person{雅努斯}，双面，或那使他卓然出众的尖刺钥匙；\Person{朱庇特}，披斗篷、蓄胡须，右手握着一根形如雷电的木头；\Person{朱诺}的腰带，或是藏匿在战士头盔下的少女；诸神之母，手持手鼓；\Person{缪斯}，带着排箫与琴；\Person{墨丘利}，有翼的\Person{阿尔戈斯}杀戮者；\Person{埃斯库拉庇乌斯}，拿着权杖；\Person{刻瑞斯}，胸脯硕大，或是在\Person{利柏耳}右手中晃动的酒杯；\Person{穆耳西柏耳}，身穿工匠的服装；或是\Person{福耳图娜}，手持满装苹果、无花果或秋日果实的丰角；\Person{狄安娜}，大腿半露，或\Person{维纳斯}裸体，挑动淫欲；\Person{阿努比斯}，狗头；或是\Person{普里阿普斯}，还不如自己的生殖器重要：\Emph{难道人们指望这些能使人惧怕吗？}\par

26. 哦，那可怖的恐惧之形、骇人的妖怪，人类曾要因它们而永远麻木，在极度的惊愕中不敢妄动，克制自己不去做任何邪恶可耻之事——小小的镰刀、钥匙、帽子、木块、有翼的凉鞋、棍杖、小手鼓、排箫、琴、突出而硕大的乳房、小酒杯、钳子、装满果实的角、裸体的女人，以及公然暴露的巨大\Emph{\OriginalTerm{阳具}{veretra}}！难道跳舞\Emph{并且}唱歌，不比称之为庄重并装作严肃、讲述如此乏味愚蠢的事——即古人造像以遏制恶行，并\Emph{激发}邪恶与不虔之人的恐惧——更好吗？难道那个时代的人，在理解力上如此缺乏理性和良好的判断力，以至他们就像小男孩一样，被面具那超自然的凶恶、鬼脸和妖怪吓阻，不敢作恶吗？而这一切何以彻底改变了，以致尽管你们各邦中有这么多庙宇，充满了所有神祇的像，但\Emph{即便}有这么多法律和如此可怕的惩罚，仍无法遏止众多的罪犯，他们的胆大妄为无法以任何方式制伏，邪恶的行为一再重复，并且当法律和\Emph{严酷}的判决力图减少残酷行为的数目，并通过惩罚\Emph{所施加}的约束来平息它们时，罪行却反而倍增？然而，如果像果真在人心中引起任何恐惧，法律的制定就会停止，也不会有那么多种酷刑用以对付胆敢犯罪的罪人：但现在，既然已证明并确定，那据称从像中流出的所谓恐惧实际上是虚妄的，人们便转而诉诸法律的条例，借此可在人心中也固定一种\Emph{由惩罚而来的}极确切的畏惧，并确定定罪；这些像本身也正因如此，才得以仍安然矗立，并因人们对其表现出的一些尊敬而得到保护。\par

\SourceRecord{Translated by Hamilton Bryce and Hugh Campbell.；Ante-Nicene Fathers；Vol. 6.；Edited by Alexander Roberts, James Donaldson, and A. Cleveland Coxe.；Buffalo, NY: Christian Literature Publishing Co.,；1886.；<http://www.newadvent.org/fathers/06316.htm>.}

% structure: p0001 source=h1 target=section reason=page-title

\section{\WorkTitle{反异教徒}（卷七）}

1. 由于已经充分表明——只要有机会——制造偶像是多么徒劳，我们论证的进程要求我们接下来尽可能简要地，毫不绕弯地论及祭献、牺牲的宰杀与焚烧、纯酒、焚香，以及在这种场合所提供的一切其他事物。为此，你们惯于激起对我们最猛烈的仇视，称我们为无神论者，并以死刑惩罚我们，甚至用野兽残暴地撕碎我们，理由是我们在极小的程度上敬重神祇；当然，我们承认这一点，不是出于对神圣的蔑视或鄙视，而是因为我们认为，这类大能者不需要任何这类事物，也不为对这些事物的欲望所支配。\par

那么，有人会问，你们是不是认为完全不应该举行任何祭献呢？我们不用自己的想法，而用你们的\Person{瓦罗}的意见来回答你们——不举行任何祭献。为什么这样呢？因为，如他所说，\Quote{真神既不希望也不要求这些；而那些用铜、陶、石膏或大理石造的，更不在乎这些，因为它们没有感觉；你不献上，也不受责备，献上了，也不会得到恩惠。} 找不到比这更正确的观点了，\Emph{没有}比这更真实的，而且\Emph{任何人}都可以采纳，哪怕他愚钝不堪，极难\Emph{被说服}。因为谁会如此愚钝，以至于向那些没有知觉的献上、宰杀牺牲，或者认为应该向那些在本质和福乐状态上远离这些的献上这些呢？\par

2. 谁是真神？你们说。要用普通而简单的语言回答你们，我们不知道；因为我们怎能认识那些我们从未见过的呢？我们惯常听你们说，有无数位神，都被算作神明；但如果这些神确实存在，且\Emph{是真实的}神，如\Person{特伦提乌斯}所认为的那样，必然的结果便是他们与自己的名称相符；就是说，他们须是我们全都看到他们应有的样子，\Emph{且配得}以此名称呼；不仅如此——我们长话短说——\Emph{他们须是}宇宙的主、全能的天主自己那样的神，我们提到祂的名时，有足够的认识和理解，称祂为真天主。因为在神性方面，一位神与另一位神毫无差异；在性质不变的情况下，种类相同的事物，其各部分不可能有或多或少之分。既然这是确定的，那么，他们就应当从来不受生，而应是不死不灭的，不求诸于外，不从物质资源中汲取任何世俗的快乐。\par

3. 那么，如果这些事是如此，我们想首先向你们了解这一点——你们举行祭献的原因、理由是什么；\Emph{然后}，这些祭献本身给神祇带来什么益处，又保留给他们什么好处。因为无论做什么都应有一个原因，不应脱离理性，以致迷失于无用的工作，在虚妄而无定的游移中摇摆不定。是天上的神祇靠这些祭品维生，必须供给材料来维持他们各部分的联合吗？有谁会如此无知，或对神是什么如此缺乏认识，竟至认为他们是靠任何养分维持的，认为是给他们的食物使他们能在无尽的永存中生活并延续下去呢？因为凡是依靠外部原因和自身之外事物维持的，必定是有死的，且当赖以生存的任何事物开始缺乏时，便走向灭亡。此外，\Emph{人不可能设想任何人会相信这一点}，因为我们看到，在带到他们祭台前的这些物品中，没有任何东西加到或触及神祇的实体；因为要么是焚香，在炭火上熔化消失，要么只是牺牲的生命被献给神祇，其血被狗舔尽；要么若有什么肉被放在祭台上，同样被火烧毁，\Emph{且}化为灰烬——除非也许神夺取了牺牲的灵魂，或急切地嗅取从燃烧的祭台\Emph{升起的}浓烟和蒸汽，以烧肉散发出的气味为食，那肉仍带着血而湿漉漉的，淌着原先的汁液。然而，若如常言所说，神没有身体，完全不可触摸，那没有身体者怎么可能被属于身体的东西滋养呢？——有死的怎么能支撑不死的，并辅助和赋予生命力于它所不能触及的呢？因此，举行祭献的这一理由似乎并不成立；也没有人可以说，人们持续举行祭献的原因是为了以此滋养神祇，靠着吃这些祭品来维持他们。\par

4. 如果并非如此，难道不是有牺牲者被宰杀祭献于神，并被投在他们的烈火祭台上，为给他们带来某种快乐与喜悦吗？任何人能说服自己相信，神明因刺激感官的快乐而变得温和，渴望肉欲的享受，并像某些卑贱的生物一样，被惬意的感受所影响，因短暂而逝的欢愉而一时心醉神迷吗？因快乐而被胜过的东西，必被其反面——忧愁——所困扰；也不能\Emph{免于}忧苦的焦虑，那因喜悦而颤抖、因欢欣而任性地得意之物。但神明应当免于这两种偏情，如果我们愿意承认他们是永恒的，且免于会死者之软弱的话。而且，一切快乐，宛如身体的某种谄媚，且诉诸那五个众所周知的感官；但若天上的神明感受到了它，他们必定也分有那些身体，藉此才有通往感官的途径，和\Emph{由之}接受快乐的门径。最后，以宰杀无害的生物为乐、双耳常充塞着它们可怜嚎叫的震响、目睹血流成河、生命随血而逝、隐秘器官被剖开，不但连带着粪便的内脏突出，而且那仍因残余生命而跳动的心脏，以及内脏里颤抖悸动的血管，这算什么快乐呢？我们这些半野蛮的人，不——更直言相告更真实更直率的话——我们这些野蛮人，不幸的需要和恶劣的习惯已训练我们以此等物为食，有时也会对它们动恻隐之心；当我们彻底观看并审视此事时，我们自咎自谴，因为我们忽略了约束世人的法律，打破了原初天然联合我们的纽带。有谁会相信，那些\Emph{本是}善良、仁惠、温柔的神明，因牛群的屠杀而喜悦并充满快乐——倘若这些牛在他们的祭台前惨然倒地气绝？那么，正如所见，祭献中便没有快乐的理由，也没有祭献的道理，因为其中并无\Emph{由之而来的}快乐；而即使偶有某种快乐，也已证明它绝不能归于神明。\par

5. 接下来，我们要检验那个常从普通民众口中听到、并\Emph{发现}其深植于大众信仰中的论证——竟认为向天上神明祭献的目的，是为了让他们平息愤怒与偏情，恢复平静安详的宁静，其炽烈精神之怒气得以消除。而如果我们牢记应时时铭记的定义，即一切激动的情感皆非神明所知，那么结论便是：相信神明从不发怒；而且，任何偏情之远离神明，莫过于那种几近野兽与凶猛动物\Emph{之精神}者，它以暴风雨般的情感激动承受者，并将他们带入毁灭的危险。因为凡被任何扰乱所折磨者，显然能感受痛苦，且脆弱；凡已服从痛苦与脆弱者，必定会死；而愤怒折磨并毁灭那些屈服于它的人：因此，凡陷于愤怒情感者，应被称为会死的。然而，我们确知神明应当永不死亡，且应拥有不死不灭的本性；如果这是清楚确定的，愤怒便远远地与神及其境界隔绝。那么，全然无由，希望平息天上神明中那不合其真福境界之物，何宜之有？\par

6. 但是，按照你的意愿，容许神明惯于如此扰乱，且献上祭献、举行神圣仪式以平息它；那么，是在何时应用这些礼仪，或应于何时呈献呢？——是在他们愤怒、被激动之前，还是当他们已被扰动、不悦之时呢？如果我们必须\Emph{在愤怒被激起前以祭献迎候他们}，以防他们发怒，那么你便是将野兽而非神明带到我们面前——习惯向之投掷食物，好让它们疯狂暴怒，转移伤害的欲望，免得它们一旦被激怒，便咆哮着冲破巢穴的藩篱。但若这些祭献是为了满足已然怒火中烧的神明而献上，我不去追问、也不去考虑，那属于神祇的快乐而崇高的心灵伟大，是否会被渺小人类的过犯所扰乱，是否因一个盲目、常在无知云雾中行走的受造者——犯下任何错误、说了任何损害他们尊严的话——而受伤。\par

7. 但我并不要求这般说明，或被告知神明对人类发怒的原因何在，以致他们因受冒犯而需被抚慰。然而，我要问：他们可曾为凡人制定过任何法律？他们可曾定下何为应做的事，何又为不应做的事？他们应追求什么、避免什么；甚或以何种方式希望自己受敬拜，以致若有人行事与他们的命令相左，他们便以愤怒的报应追究，且若被轻蔑，便准备报复那些放肆逾矩者？依我看来，他们从未制定或规定过任何事物，因为他们既未曾被看见，也不可能被清楚分辨究竟是否存在任何神明。那么，若天上神明因任何理由对那些——他们既不屑于在任何时候显示自身之存在，也未给予或强加任何他们希望人尊敬并全然遵守的法律——的人发怒，有何正义可言呢？\par

8. 但是，如我所说，我并非提起此事，而是任由它悄然消逝。我首要问的乃是：若我宰杀一头猪，一位神便会改变心意，放下他的忿怒与狂暴；若我在他眼前、在他的祭台上献上一只小母鸡、一头牛犊，他便忘记我对他所行的过错，全然消除一切不快之感——这有什么道理可言？这一举动中，究竟传递了什么而能消解他的怨恨？又或是一只鹅、一头山羊、一只孔雀有何用处，竟能以其血使发怒的神得到平息？那么，众神竟将他们受到的侮辱视为可以买卖之事吗？正如小男孩为了让他们止住脾气、停止哭闹，人们给他们小麻雀、玩偶、小马、木偶，让他们以此自娱——不死的神明难道就是这样从你们手中接受这些礼物，从而为这些礼物放下怨恨，与那些冒犯他们的人和好吗？然而我本以为，众神——如果相信他们真会被忿怒所动是正确的话——会不索任何代价或报偿，便放下忿怒与怨恨，赦免有罪者的罪。因为这尤其属于神明的特质：宽赦时慷慨大方，赐予恩惠而不求回报。但如果这一点做不到，那么他们固执于忿怒，远比因收受贿赂而被软化要明智得多。因为当人们有了为罪付出代价的希望时，犯罪者的数目便增加了；而当那赦免\Emph{过犯}者的恩惠可以用钱买来时，人们几乎不会犹豫去行恶。\par

9. 那么，假如某头牛，或你喜欢的任何动物，就是那被宰杀以缓解并平息众神怒气的动物，若能发出人的声音，说出这样的话：\Quote{这，\OriginalTerm{朱庇特}{Jupiter}啊，或不论你是什么神，难道是人道或正当的，又或可算为完全公正的吗：当另一个人犯了罪，我却要遭受宰杀，而你竟允准用我的血来向你赔补，尽管我从未伤害过你，从未有意或无意地冒犯你的神性和尊严，如你所知，我是个哑口无言的生物，未曾偏离我本性的纯朴，也非性情轻浮多变？我可曾不够恭敬谨慎地庆祝你的竞赛？我可曾拖出一个舞者以致触怒你的神明？我可曾指着你发过假誓？我可曾亵渎地偷窃你的财产、掠夺你的庙宇？我可曾铲除最神圣的丛林，或以私人住宅环绕圣地而将其玷污亵渎？那么，究竟是什么理由，让别人的罪用我的血来赎偿，让我这无辜的生命为他人的邪恶付出代价，而这邪恶与我毫无干系？难道因为我是低贱的生物，没有理性与智慧，正如那些自称人类，却以残忍使自己沦为野兽的人所宣称的那样？难道不是同一个自然从相同的本源繁衍并塑造了我和他们吗？难道不是同一生命气息支配着他们和我吗？难道我不是和他们一样呼吸、观看，并同他们一样受其他感官影响吗？他们有肝脏、肺腑、心脏、肠子、肚腹；难道我就没有同样多的肢体吗？他们爱自己的幼雏，并交合生育后代；难道我不也是努力繁衍后代，并在后代生出后喜悦于它吗？可是他们有理性，说出清晰的话音；他们怎么知道我所做的是出于自己的理由，怎么知道我所发出的声音是我独特的言语，只有我们同类才懂呢？请问虔敬，究竟是让我被杀戮、让我被宰杀更公正，还是让人因其所作所为得蒙宽恕、免于惩罚更公正？是谁把铁铸成了剑？难道不是人吗？是谁给各族带来灾难，将奴役强加于万民？难道不是人吗？是谁调制成致命的毒剂，并送给父母、兄弟、妻子、朋友？难道不是人吗？是谁找出或设计出数不胜数的邪恶形式，以至于在千万年的编年史中，或\Emph{甚至}在千万日的记录中，都难以尽述？难道不是人吗？这岂不就是残忍、怪异、野蛮吗？朱庇特啊，难道在你看来，为了平息你而使有罪者逍遥法外，就让我被杀、让我被杀戮，不算是非正义和野蛮的事吗？}\par

由此，业已确立，为平息忿怒的众神而献祭，乃是徒劳；因为理性教导我们：众神从不会发怒，他们不愿一个生命因另一个生命而被毁灭、被宰杀，也不愿无辜者的血来抵消对他们犯下的过犯。\par

10. 但也许有人会说：我们向神明奉献祭献和其他礼物，为的是让他们在某种程度上愿意俯允我们的祈祷，赐予我们顺遂，替我们消灾免祸，使我们常享幸福，真能驱除忧伤，并\Emph{消除一切}因意外事故而威胁我们的\Emph{灾祸}。这一点需要深思；人们通常既不愿听闻，也不肯相信如此轻率之言。因为整个学者群体会立刻向我们猛扑过来，他们断言并证明，一切发生的事都是根据\Emph{\OriginalTerm{命运的谕令}{decrees of fate}}而发生的，从而从我们手中夺走那种见解，声称我们的信仰寄托于虚妄之念。他们会说，世间已发生的、正在发生的以及将要发生的一切，都已在往昔被确定不移地安排好了，有其不可更动的原因，借着这些原因，万事相互联结，在过去与未来之间形成了一条不可攻破的\OriginalTerm{不可改变的必然性之链}{chain of unalterable necessity}。倘若每个人将遭遇的祸福已被注定而不可更改，那就已经是确定无疑的了；而如果这是确定不移的，那么神明的种种佑助、憎恶与恩惠便全无用武之地。因为他们既不能为你成就那不可能之事，正如亦不能阻止那必然发生之事；除非他们愿意的话，还能大力贬损你心中所怀的信仰，以至于他们说，甚至神明本身也徒然受你敬拜，你向他们所献的祈祷亦是多余。因为既然他们无力扭转\Emph{事件的}进程，亦无法改变那由\OriginalTerm{命运}{fate}所定之事，那么，你又有何理由、何必去烦扰、灌满那些在最危急时刻你根本无法信赖其援助之耳呢？\par

11. 最后，倘若神明真能驱除忧愁与忧伤，赐予喜乐与欢愉，那么世间为何有如此众多、如此凄惨的人？为何\Emph{会有}如此多的不幸者，在极其卑微的境况中过着以泪洗面的生活？那些每时每刻在祭台上堆积祭献的人，为何不能免于灾难？难道我们未尝看见，\Emph{学者们说}，他们中的一些人成了疾病的居所，眼睛失明，耳朵失聪，双足不能行走，他们活着如同仅存\Emph{躯干}，\Emph{不能使用}双手；他们被火灾、海难和各种灾难吞噬、压垮\Emph{并}毁灭；他们丧失巨额财富后，靠出卖劳力维生，最后沿街行乞；他们被放逐、被剥夺公权，常处忧伤之中，因丧失子女而悲痛欲绝，\Emph{并}受其他种种不幸的折磨，其种类与形式之多，难以历数？但倘若那些曾受恩惠负有义务的神明，真能替那些配得\Emph{此种恩惠}的人驱除并化解这些灾难，这一切便决不会发生了。然而现在，因为这些灾难中并无\Emph{神明干预}的余地，万物皆由\OriginalTerm{不可避免的必然性}{inevitable necessity}所促成，那已预定的世事进程便持续运行，成就那早已确定之事。\par

或者说，如果天上诸神有能力阻止，却任由不幸的人类陷入如此多的困苦和灾难，那么就该说他们是忘恩负义的。但他们或许会对此作出某些重要回答，而不是应该被欺骗、轻浮和轻蔑的耳朵所接受的那种回答。不过，这一点因为需要过于冗长而啰嗦的讨论，我们便不加解释、不加触及地匆匆略过，只满足于指出这一点：如果你们断言，除非先用母山羊和绵羊的血以及放在祭台上的其他物品买通他们，否则他们就不赐福、不挡灾，那么你们就是在赋予你们的神明不光彩的名声。因为，首先，相信神明的能力和天上诸神的超绝卓越竟将其恩惠出售，先收\Emph{代价}，然后才施与\Emph{恩惠；而且}——这更不像话——除非他们得到\Emph{他们所求的}，否则就不帮助任何人，任由最可怜的人遭受可能降临的任何危险，而他们本可阻挡\Emph{这些危险}并施以援手，这是不合适的。如果有两个\OriginalTerm{祭献}{sacrifice}的人，一个是恶棍，很富有，另一个家道微薄，但以其正直良善值得称赞——如果前者宰杀一百头公牛，以及同样多的母羊和羊羔，穷人则烧一点乳香和一些香料的小碎片——那么，如果神明除非先收到报偿就不施加任何恩赐，那么可想而知，他们就会把恩惠给予富者，转而不看那穷人，因为他礼物微薄，并非由于心意，而是由于财力匮乏。因为，当给予者唯利是图且贪图利益时，恩惠必然根据\Emph{购买它所用}礼物的大小而给予，而那更丰厚的报偿和贿赂——\Emph{尽管这很}可耻——流向给予者的人，就会得到有利的裁决。另一方面，假如两个国家交兵对阵，以同等的祭品充实了神明的祭台，并要求将他们的能力和帮助给予己方以对抗对方，那么，同样可想而知，如果神明被报偿说服而效力，他们便会在双方之间左右为难，呆若木鸡，不知所措，因为他们明白，自己一旦接受了祭品，恩惠就已许出了。因为他们要么帮助这一方也帮助那一方，这是不可能的，\Emph{因为那样的话}他们就会自己与自己交战，与自己的恩惠和意愿相争；要么在\Emph{他们助力的}代价已被付出并收受之后，不帮助任何一方，这又是非常邪恶的。因此，这一切恶名都应远远移除于诸神；绝不应说他们能被报偿和贡品收买而赐福消灾，只要他们是真神，配得这一称号。因为要么一切发生之事必然发生，在神那里没有偏私和恩惠的余地；或者，若排除并弃除命运，将服务之恩惠当作商品出售，就不合乎属天尊严，也不合于其恩赐的施与。\par

我认为，我们已经充分表明，祭品及伴随之物是徒然献给不死的神明的，因为他们既不被这些滋养，也不感到任何快乐，也不会放下愤怒和怨恨，从而赐予好运或驱走和避免厄运。现在我们必须考察另一点，即一些人通常断言并应用于仪式形式的说法。因为他们说，这些神圣礼仪是为向天上诸神表示尊敬而设立的，他们所做这些事，是为了向\Emph{他们}表示尊敬，藉此显扬神明的能力。如果他们同样说，他们保持清醒、睡觉、走动、站立、书写和阅读，也是为了向神明表示尊敬，使他们更加荣耀威严，那又如何呢？从牲畜的血和在祭献中准备的其他东西里，能给他们增加什么实质？能赋予和增添什么能力？因为，凡据说是由某人献上，并向一位更伟大的存在表示的尊敬，都是与他者相关的；它由两部分组成：给予者的让步，和接受者尊敬的增添。例如，若有人看见一位以极大权力和权威著称的人，便为他让路、起立、脱帽并从车上跳下，然后躬身用奴性的卑屈和颤抖的激动向他致敬，我看得出这样表现的尊敬意在何处：通过一方的躬身，极大的\Emph{尊敬}便给予了另一方，而那位被下位者的尊敬所抬举、并抬升至自己等级之上的人，便显得伟大。\par

但所有这些我们正谈论的对尊敬的让与和归给，只在人中间遇见，他们天生的软弱和对凌驾于同胞之上的喜爱，教他们以傲慢为乐，并喜欢比别人更受抬举。然而，我要问，在诸神中间，哪里有容纳尊敬的余地？或者，堆积祭品能给他们带来什么更大的抬举？当\Emph{向他们}祭献牲畜时，他们就变得更为可敬、更有能力了吗？因为这事，有增添给他们什么吗？或者，因神性增加，他们开始成为更\Emph{真实}的神明吗？然而我几乎认为，当说一个人尊敬一位神，并通过奉献某些礼物而抬举他时，这近乎侮辱，不，完全是侮辱。因为，如果尊敬会增加并扩大那受尊敬者的伟大，那么，神明便因那从中接受礼物和所献之尊敬的人而变得更大；这样一来，事情就变成：那位被人的尊敬抬举的神乃是下位者，而另一方面，增添神明能力的人则是他的上位者。\par

15. 那么！有人会说：\Quote{难道你认为根本不应给予诸神任何尊崇吗？}如果你向我们提出这样的神——如果他们存在，他们就该是那样的，并且当我们提到那个名字时，我们觉得我们都意指那样的——我们怎能不给予他们甚至最大的尊崇呢？因为我们从那些对我们有特殊权威的命令中学到，甚至要尊敬所有的人，不论他们的等级、不论他们的处境如何？\Quote{请问，这极大的尊崇是什么呢？}\Emph{你问}\Quote{是一种比你们所献上的更符合本分的尊崇，并且是指向一个更有能力的族类。}\Emph{我们回答}\Quote{首先告诉我们，什么是对神相称、正确而光荣的看法，并且不会因为某种可耻的事而变得不合适、值得责备？}\Emph{你说}我们回答：\Quote{这样一种看法：你相信他们既没有任何与人的相似之处，也不寻求任何在他们之外、从外部而来的东西；然后——这一点已经说过多次——他们不会燃起愤怒之火，他们不会纵情于肉欲快乐，他们不会受贿赂而效劳，他们不会受诱惑去伤害我们的敌人，他们不会出卖自己的仁慈和恩惠，他们不会因被加上尊崇而喜悦，如果得不到，他们也不会愤慨和烦恼；而是——这属于神性——他们以自己的能力认识自己，他们不按别人的谄媚来评价自己。}然而，为了使我们看清所说之事的本质，这是哪一种尊崇呢：在神面前捆绑一只阉羊、一只公羊、一头公牛，并在他眼前宰杀它们？这是哪一种尊崇呢：邀请一位神赴血宴，而你看见他与狗一同取食分享？这是哪一种尊崇呢：点燃木柴堆，以烟雾遮蔽天空，并以阴森的黑暗使神像模糊不清？但如果你们觉得应该就其本身考虑这些行为，而不是按你们的偏见来判断，你们将会发现，你们所说的那些祭台，甚至你们奉献给高级神祇的那些华丽的祭台，都是焚烧不幸动物族类的火葬柴堆，是为最不体面的职分而筑起的土墩，并形成将要充满腐败的场所。\par

16. 你们说什么呢，噢，你们——！那么，那种由焚烧的兽皮、骨头、鬃毛、羔羊的毛和家禽的羽毛所发出和散发的恶臭——\Emph{那}是对神明的恩惠和尊崇吗？神明就是这样受尊崇的吗？当你们准备去他们的庙宇时，你们要洗净一切污秽、沐浴净身、完全洁净后才前往。还有什么比这些更污秽、更不幸、更卑贱的呢？如果他们的感官天性如此，竟喜欢如此残忍的东西，并喜爱那些恶臭——这些恶臭随着呼吸被吸入时，甚至连那些献祭者都无法忍受，更不用说娇嫩的鼻子了。但如果你们认为天上的神是因活物的血被献上而受尊崇，那么你们为什么不用骡子、大象和驴子向他们献祭呢？为什么不用狗、熊、狐狸、骆驼、鬣狗和狮子呢？既然你们也把鸟类算作牺牲，那么你们为什么不用兀鹰、雕、鹳、隼、鹰、渡鸦、雀鹰、猫头鹰，以及连同它们一起的蝾螈、水蛇、蝰蛇、狼蛛来献祭呢？因为确实这些动物也有血，也同样被生命的气息所驱动。在前一类祭献中有什么更精巧，在后一类中有什么更不精巧，以致后者不能增添和增加神的尊威呢？因为，\Emph{我的对手说}，\Quote{用那些滋养、维持我们和我们靠之生存的东西来尊崇天上的神，是正确的；这些东西，正是他们出于神圣的仁慈，惠然赐予我们作食物的。}但是，同一些神还赐给你们莳萝、水芹、芜菁、洋葱、欧芹、食用蓟、萝卜、葫芦、芸香、薄荷、罗勒、驱虫草和细香葱，并命令你们把它们用作食物的一部分；那么，你们为什么不把这些也放在祭台上，把喂牛的野墨角兰撒在它们上面，并把带有辛辣味道的洋葱混在它们中间呢？\par

17. 试想，如果狗——为了使事情显得更清晰，必须设想一种情形——试想，我说，如果狗、驴，以及与他们一道的鹡鸰，如果啁啾的燕子，还有猪，也获得了人的某些情感，竟然认为并假设你们是神，并提议向你们献上祭品以表敬意，不是用其他事物和物质，而是用它们按其天性通常赖以滋养和维持的东西——我们请问你们，你们会觉得这是一种尊崇，还是一种最骇人听闻的侮辱，当燕子为你们宰杀并祝圣苍蝇，鹡鸰祝圣蚂蚁；当驴把干草放在你们的祭台上，并浇奠糠秕；当狗放置骨头，并在你们的神龛前焚烧人粪；最后，当猪把它们从骇人的泥沼和污秽的食道中取出的可怕糊状物倾倒在你们面前时？在这种情况下，你们难道不会因你们的伟大被如此轻侮而怒火中烧，并认为被以污秽迎接是一种极恶毒的冤屈吗？但是，你们回答说，你们以公牛的尸体，并通过宰杀其他活物来尊崇众神。这在那一点上与前者不同呢，既然这些祭品，即便目前还不是，但不久也将成为粪土，并会在很短时间后腐烂？最后，停止在你们的祭台上点火，那时你们将确实看到，那被你们用来增光诸神荣耀的公牛肉，被祝圣过的，会因生虫而肿胀起伏，污染败坏空气，并以污浊的臭气侵染邻近的地区。现在，如果众神命令你们将这些转用于自己，按惯常方式取食它们，你们会远远逃开，咒骂这气味，向众神求恕，并发誓再也不向祂们奉献此类祭品。那么，你们的这种行为岂非嘲弄？岂非承认、表明了你们不知道什么是神性，也不明白这个名字的含义和称号应赋予并应用于何种大能？你们以新种类的食物赋予众神新的尊荣吗？你们以滋味和汁液尊崇祂们，只因那些滋养你们的东西使你们喜悦而感激吗？你们相信众神也会蜂拥而至，享受它们的美味，并且，如同狂吠的狗，为了一口食就放下凶猛，频频向递食者摇尾乞怜吗？\par

18. 我们既在谈论被献祭的动物，那么，有什么理由，什么原因，既然众不朽的神——因为，就我们而论，凡被相信是神的都可以是神——本应同心一意，或应有同一的本性、种类和特质，却并非所有祭品都能平息所有神祇，而是依祭献律法，某些神祇以某些动物为祭？因为，容我重复这问题，有什么理由，这个神要以公牛尊崇，另一个要以山羊或绵羊，这个要以乳猪，那个要以未剪毛的羔羊，这个要以处女牛犊，那个要以有角山羊，这个以不孕母牛，而那个以怀胎母猪，这个以白色，那个以暗色祭品，一个以雌性，另一个却相反，以雄性动物为祭？因为如果祭品是为了向众神表示尊崇和敬意而宰杀以献祭，那么用何种动物的生命来偿还此债，平息祂们的忿怒与怨恨，又有什么关系，或有什么差异呢？还是一个神的祭品之血不如另一个的血那样令某位神祇悦纳，而另一个的血却使祂充满欢欣喜乐？或者，如通常所做的那样，某个神祇因某种敬畏和宗教顾虑而禁戒山羊肉，另一神祇厌恶猪肉，而对另一神祇，羊肉却发臭？而这一个避免硬老牛肉以免劳损他那孱弱的胃，而选择嫩口的乳畜以便更迅速地消化？\par

19. 但你们错了，\Emph{我的对手}说，并且陷入了谬误；因为在向女神献雌性祭品，向男神献雄性祭品时，有着一种隐秘且极为奥妙的理由，非大众所能及。我既不探究，也不追问，这些祭献律法教导或包含些什么；但若理性已证明，真理已宣告，在众神之间并无种类之别，且他们不由任何性别区分，那么这一切推理岂非都应被鄙弃，而那些智者——他们听到有人将性别差异归于不朽之神时便忍不住发笑——他们的意见岂非显为正确？我请问每一个人，他自己心中是否相信，并使自己信服，神明的种族是如此区分，即有男有女，并且被造成具有适合繁衍后代的肢体？\par

但是，如果祭献的律法规定，同类性别应献给同类，即雌性祭品献给女神，相反，雄性祭品献给男神，那么颜色又有什么关联，以致将白色的给这些神，暗色的甚至最黑的给那些神，便是正确而合宜的呢？因为，\Emph{我的对手}说，对天上的神祇，以及那些有能力赐予吉兆者，白色所呈现的愉快外观使其悦纳而吉利；相反，对凶神，以及那些居于阴间居所者，暗色则更讨喜爱，且蒙上忧郁之色调。但是，再说，如果这一推理成立，即阴间区域是完全虚幻空洞的名称，而地下并没有\OriginalTerm{普路同的领域}{Plutonian realms}和居所，那么这也必定使你们关于黑色牲口和地下神祇的观念归于无效。因为，如果没有阴间区域，那么必然也没有\OriginalTerm{阴间神祇}{dii Manium}。因为，既然没有区域，又怎能说有谁居住在其中呢？\par

20. 但是让我们如你所愿地同意，既有\OriginalTerm{冥界}{infernal regions}，也有\OriginalTerm{马涅斯}{Manes}，并且有一些神祇居住在其中，对人绝不友善，主宰着不幸；那么，有什么原因，有什么理由，要把黑色的祭品，甚至是最深色的，带到他们的祭坛前呢？因为黑暗的东西适合黑暗，阴郁的东西令相似的生灵喜悦。那么，你们没有看到——我们也可以愚蠢地开开玩笑，就像你们自己那样——祭品的肉并不黑，它们的骨头、牙齿、脂肪、内脏连同脑子和骨头里柔软的骨髓也不黑。但皮毛却是乌黑的，这些动物的鬃毛也是乌黑的。那么，你们是否只拿从祭品身上拔下的羊毛和细鬃毛来献祭给神祇呢？你们是否任凭那些可怜的生灵，或许被剥夺、被剪毛之后，呼吸天上的空气，在牧场上全然无辜地休息呢？但如果你认为黑色和阴沉颜色的东西令冥界神祇喜悦，为什么不注意在所有其他习惯放置在祭品上的东西也应该是黑色、被烟熏过、颜色可怕的呢？如果献香，就把它染黑；腌渍的粗麦粉，以及所有浇奠，毫无例外地染黑。在奶、油、血中倒入煤烟和灰烬，好让这个失去紫色，让其他的变得惨白。可是，如果你在引入一些白色且保持其明亮的东西时毫无顾忌，那么你们自己就废掉了自己的宗教顾虑和推理，因为你们在举行神圣仪式时没有保持任何单一而普遍的规则。\par

21. 但我们也应该从你们那里得知这一点：如果把一只山羊宰杀献给\OriginalTerm{朱庇特}{Jupiter}，而它通常是被祭献给\OriginalTerm{父利伯尔}{father Liber}和\OriginalTerm{墨丘利}{Mercury}的；或者如果把不生育的母牛献给\OriginalTerm{乌恩西亚}{Unxia}，而你们是把它献给\OriginalTerm{普洛塞庇娜}{Proserpine}的；那么，根据什么惯例和规则来确定这其中有什么罪，犯了什么邪恶或罪行呢？因为就献给神的敬拜而言，用哪种动物的头来偿还你们所欠的荣耀，是毫无区别的。\Emph{我的对手}说，把这些混淆起来是不许可的，把仪式的典礼和赎罪的方式搞乱绝不是小罪。请解释理由，我恳求。因为把特定种类的祭品祝圣于某些神祇，并且也应采用某种祈求方式，这是合宜的。那么，再把特定种类的祭品祝圣于某些神祇，并且也应采用某种祈求方式，这其中的理由又是什么呢？因为这合宜性本身应该有其自己的原因，产生并源自某些理由。你们是要谈论古老和习俗吗？如果是这样，你们向我讲述的仅仅是人的意见、一个盲目生灵的发明；而我，当我要求给我提出理由时，我希望听到要么有东西从天而降，要么——这更是问题所需要的——朱庇特与牛的血有什么关系，以至于应把它献祭给他，而不是献给墨丘利或利伯尔。或者，山羊的自然属性是什么，以至于它们又该适合于这些神祇，而不该适合于朱庇特的祭献？动物在众神之间被分配了吗？是否已经达成并同意了某种契约，以致于这位应当避开属于那位的祭品，那位应当停止把属于另一位的血据为己有是正确的？或者，如同嫉妒的男孩们，他们不愿意让别人分享被献给他们的牲口的享用吗？或者，就如传说中风俗迥异的民族所做的那样，同样的东西被一方认为适合食用，却被另一方拒绝作为食物？\par

22. 那么，如果这些事情是虚幻的，得不到任何理性的支持，那么祭献的奉献本身也是无谓的。因为，既然那第二个命题所由以产生的第一个前提被发现完全是虚浮虚幻的，并且没有建立在任何坚实的基础上，那么，随之而来的事情又怎么能有合适的原因呢？他们说，向\OriginalTerm{地母}{mother Earth}献祭一头怀胎多产的母猪；而向\OriginalTerm{童贞密涅瓦}{virgin Minerva}宰杀一头从未被刺棒驱策去尝试任何劳动的小母牛。然而，我们认为，既不该向一位童贞女神献祭一个处女性祭品，以免在那畜牲身上侵犯了童贞，而这正是女神特别受尊敬的原因；也不该出于对大地的多产力的尊敬，而把怀胎和怀孕的祭品献给大地，因为我们全都渴望并希望那多产力总是以不可抑制的生殖力持续下去。因为，如果因为\OriginalTerm{特里同女神}{Tritonian goddess}是一位童贞女，所以向她献祭处女性祭品是合宜的，又如果因为大地是一位母亲，所以应当同样用怀孕的母猪来款待她，那么，\OriginalTerm{阿波罗}{Apollo}也该由音乐家的祭献来受尊崇，因为他是一位音乐家；\OriginalTerm{埃斯库拉庇乌斯}{Aesculapius}，因为他是一位医生，就该由医生的祭献；而\OriginalTerm{伏尔甘}{Vulcan}，因为他是一位工匠，就该由工匠的祭献；又因为墨丘利能言善辩，就该用雄辩和最能言善辩的人向他献祭。但如果说这话是疯狂，或者用温和的话来说，是胡言乱语，那么这就显示了更大的疯狂：因为大地甚至更加多产，就把怀孕的母猪宰杀献给她；因为密涅瓦是纯洁的，童贞无瑕，就把纯洁的小母牛献给她。\par

23. 至于我们听到你们所说的，有些神是好的，相反，另一些神是坏的，而且更倾向于放纵于任性的恶作剧，并且对前一类神奉行惯常的仪式，为的是让他们展示好意，而对后一类神，为的是让他们不伤害你们——这种说法有什么道理，我们承认我们无法理解。因为说神是最仁慈的，具有温和的性情，不但是虔诚和敬神的，而且也是真实的；但说他们是邪恶和凶险的，则绝不该听信，因为那种神圣的力量远非、并且也绝不同于那种为害的性情。但任何能招致灾祸的东西，首先必须看清它是什么，然后应该远远地将其从神的名下移除。\par

那么，假设我们与你们一致认为神明促成好运与灾祸，即便如此，你们也没有理由诱使某些神明赐予你们繁荣，另一方面又用祭品和犒赏哄骗另一些神明不要伤害你们。首先，因为善神不可能行恶，即使他们未受到崇敬的礼拜——因为凡本性温和安详者，远不会行恶或谋划恶事；而恶神则不知遏制其残暴，即使他们被千头牛羊、千座祭台所诱哄。因为苦味不能变为甜味，干燥不能变为湿润，火的热不能变为冷，任何事物都不能将其对立面接纳并转化为自己的本性。因此，你若用手抚摸毒蛇，或爱抚有毒的蝎子，前者会以毒牙攻击你，后者会蜷缩身体，将毒刺扎入你身；你的爱抚毫无用处，因为这两种生灵被激发作恶，并非出于愤怒的刺痛，而是出于其本性的某种特质。因此，靠祭品讨好凶神是徒劳的，因为无论你做还是不做，他们都遵循自己的本性，由天生的律法和某种必然性驱使，去行他们被造而为之的事。此外，这样一来，两类神明都失去了各自的能力，无法保持自己的特性。因为若善神受崇拜是为得他们眷顾，而对恶神也同样祈求，相反地是为了他们不伤害人，那么结论必然是：慈悲的神明若不收礼物就无恩惠，且由善变恶；而恶神若收供奉，则放弃其恶毒秉性，后来变为善；于是便导致如此结果：要么这些神既不慈悲，那些也不凶恶；要么——这是不可能的——两者皆慈悲，两者又皆凶恶。\par

24. 就这样吧；姑且承认这些最不幸的牲畜并非毫无宗教义务就在神明庙宇中被宰杀，也承认按惯例和习俗所做之事具有某种理性依据：然而，如果宰杀公牛献神，并将整只或整块的动物肉焚烧作为祭献，被视为伟大壮观之举，那么，那些与巫术有关的残余物，被大司祭的奥迹重新纳入圣礼秘法之中，并与宗教事务掺杂，又是什么意思呢？我所说的是这些：\OriginalTerm{香肠}{apexaones}、\OriginalTerm{公山羊香肠}{hircioe}、\OriginalTerm{葬礼香肠}{silicernia}、\OriginalTerm{长香肠}{longavi}，这些是香肠的名称和种类，有些灌满山羊血，有些则是剁碎的肝脏。还有\OriginalTerm{油渣}{toe-doe}、\OriginalTerm{肠衣}{uoenioe}、\OriginalTerm{油丁}{offoe}，不是普通人用的那种，而是被称为\OriginalTerm{尾肉块}{offoe penitoe}的那种offoe——其中第一种是将脂肪切成极小块的油渣，如同美味佳肴；第二种是肠衣的延伸，排泄物通过它排出，所有营养汁液已被沥干；而\OriginalTerm{尾肉块}{offa penita}则是兽尾连同一小块肉被切下。还有\OriginalTerm{睾丸}{polimina}、\OriginalTerm{网油}{omenta}、\OriginalTerm{血粉牛尾}{palasea}，或如某些人称\OriginalTerm{血粉牛尾}{plasea}——其中\OriginalTerm{网油}{omentum}是腹内网膜包裹脏器使之各居其位的部分；\OriginalTerm{血粉牛尾}{plasea}是用面粉和血涂抹的牛尾；\OriginalTerm{睾丸}{polimina}则是我们较文雅地称之为\OriginalTerm{子嗣}{proles}的部分，而俗人通常称其为\OriginalTerm{睾丸}{testes}。还有\OriginalTerm{一种粥}{fitilla}、\OriginalTerm{另一种粥}{frumen}、\OriginalTerm{圣饼}{africia}、\OriginalTerm{圣饼}{gratilla}、\OriginalTerm{圣饼}{catumeum}、\OriginalTerm{圣饼}{cumspolium}、\OriginalTerm{圣饼}{cubula}——前两种是粥的名称，但种类和性质不同；后面一连串名称指献祭的糕饼，因为它们并非以同一种方式制作。因为我们不打算提及从公牛后腿取下的\OriginalTerm{臀肉}{caro strebula}，串在铁签上烤的肉片，先在火上加热然后在炽炭上焙烤的肠子，最后也不提那由四种水果混合制成的腌菜。同样，我们不拟提及\OriginalTerm{肠子}{fendicoe}，也就是\OriginalTerm{肠子}{hiroe}，民众说话时通常称之为\OriginalTerm{肠子}{ilia}；同样也不拟提及\OriginalTerm{瘤胃}{oerumnaoe}，那是食道的始端，反刍动物习惯将食物送下然后再反刍上来；也不提\OriginalTerm{香肠或粥}{magmenta}，\OriginalTerm{香肠或粥}{augmina}，以及其他上千种名称晦涩、却让平民百姓更加敬畏的香肠或粥类。\par

25. 因为若人类所做的一切，尤其是在宗教方面，都应有其原因——在人类所做所行的一切中，若无理由便不应有所作为——那么请告诉我们并说明，将这些事物也献给诸神，并在他们的神圣祭台上焚烧的原因是什么？理由是什么？我们在此耽搁，被最迫切地强制等待这一原因，我们停顿，我们站立不动，渴望了解一位神与粥、与糕饼、与以多种方式准备的各式填馅、以及不同的配料有何关系？诸神是否被丰盛的晚宴或午餐所感动，以至为他们设计无数的盛宴是合宜的？他们是否因胃中的厌恶而困扰，且寻求不同的滋味以消除反感，因此在他们面前摆上时而烤的、时而生的、时而半生半熟的肉？但若诸神喜欢接受所有这些你们称为 \OriginalTerm{祭品份}{proesicioe} 的部分，并且若这些以任何愉悦或快乐的感受取悦他们，那么是什么阻止、什么妨碍你们将这些连同整个动物一并一次性地放置在祭台上？是什么原因、是什么理由，使得单独的大腿肉块、食道、尾巴和尾骨块、仅内脏，以及仅薄膜，被带来向他们致敬？天上的诸神是否被各种调味品所感动？在他们用奢侈丰盛的餐食填饱自己之后，他们是否如常做的那样，将这些小块当作甜点，不是为了平息饥饿，而是为了刺激疲惫的味觉，并重新激发起狼吞虎咽的胃口？哦，诸神的伟大奇妙，无人能领会，无受造物能理解！若确实他们的恩惠是用野兽的睾丸和食道买来的，并且若不看到准备好的内脏和购买的 \OriginalTerm{小块祭品}{offoe} 在他们的祭台上焚烧，他们便不放下愤怒和怨恨。\par

26. 关于焚香和酒，我们现在必须说几句，因为这些也与你们的仪式联系在一起并混合使用，且在你们的宗教行为中大量使用。首先，关于你们所使用的香本身，我们特别要问你们：你们从何处或在何时得以熟知并认识它，以至于你们有正当理由认为它要么配得献给诸神，要么最合乎他们的渴望？因为它几乎是新奇的；它开始在这些地区为人所知并传入诸神的神龛，尚无漫长无尽的岁月。因为在英雄时代，据信并宣称，尚不知香为何物，古作者们证明了这一点，在他们的书中找不到关于它的记载；迷信的根源与母亲——\Place{伊特鲁里亚}——也不知晓它的名望和声誉，小圣堂的礼仪可证；在 \Place{阿尔巴} 隆盛的四百余年里，无人献祭时使用过它；甚至 \Person{罗慕洛} 或 \Person{努马}——精于设计新礼仪——也既不知道它的存在，也不知道它的生长，常用来举行常规献祭的神圣谷物表明了这一点。那么，它的使用究竟何时开始被采纳？或者说，何种求新之欲冲击了古老而悠久的习俗，以至于那么多个世代不必要的东西，在仪式中竟占据了首要位置？因为若没有香，宗教崇拜的履行便不完全，且若一定数量的香对于使上天诸神对人温和恩慈是必要的，那么古人便陷入了罪过，不，他们的整个生命充满了罪疚，因为他们漫不经心地忽略了献上那最适于取悦诸神的东西。但是，若在古代无论人还是神都不寻求这香，那么今天同样也就证实了，那古代并不认为必要，近代却无任何理由而欲求的东西，是被徒劳无益地献上的。\par

27. 最后，为了我们总是固守那条被阐明并确定了的规则和定义，即人所做的一切必须有其原因，我们在此也将坚守之，以便向你们追问原因和理由：为何香被放在诸神的神像前的祭台上，且一旦燃烧，他们就被认为变得友好与温和。他们从这一行为中获得什么，或者说有何事触及他们的心智，以使我们应当合理地判断这些物品被妥善耗费，而非无用地徒然消耗？因为正如你们应当说明为何将香献给诸神，同样随之而来的是你们应当表明，诸神有某种理由不轻蔑地拒绝它，甚至更甚，如此喜爱地渴求它。或许有人会说：我们以此向诸神表示敬意。但我们询问的不是你们的感觉，而是诸神的；也不是问你们做了什么，而是问他们对你为求得恩惠而做之事有多看重。然而，虔诚啊，由火的气味产生、由树胶生成的尊敬是什么或有多大？因为，免得你们恰巧不知这是何种香，或它的来源是什么，它乃是从树皮流出的树胶，正如从杏树、樱桃树流出的一样，滴落时凝固。那么，这就尊敬并显扬了天上的尊严吗？或者，若他们的不悦曾在何时被激起，就会在香烟之前消融，并在愤怒被缓和之后，陷入沉睡？那么，你们为什么不不加区分地焚烧任何树的汁液，而不作任何区别呢？因为若诸神以此受尊敬，且不因焚烧 \Place{潘凯亚} 树胶给他们而不悦，那么在你们神圣的祭台上，烟由何物产生，或香烟之云由何种树胶升起，又有何区别呢？\par

28. 或许有人会说，向天上的神祇供奉香是因为它气味芬芳，给鼻子带来快感，而其余的气味令人不快，因其恶臭而被抛弃。那么，众神难道有鼻孔能呼吸吗？他们难道会吸进呼出气流，使各种气味能穿透他们吗？但如果我们允许这一点，我们就让他们臣服于人的条件，并将他们排除在神性界限之外；因为凡是呼吸并吸入空气、再以同样方式呼出的，必定是有死的，因为它是靠空气滋养而维持的。然而，凡是靠空气滋养而维持的，如果你拿走其维持交流的手段，其生命必然被扼杀，其生命原则必然毁灭丧失。因此，众神倘若也呼吸并吸入伴随空气的香气，那么说他们靠从外界获取之物维生，并且如果气孔被堵塞就会灭亡，这并非不确。最后，你们何以知道，如果他们被气味的甜美所吸引，那些令你们愉悦的事物也同样令他们愉悦，并以相似的感觉迷住并影响你们与众神截然不同的本性呢？难道不可能，那些给你们带来快感的事物，对他们来说反而粗粝可憎吗？因为众神的见解不同，实体也非一元，那么用什么方法才能使品质不同的事物，对于触及之物具有相同的感觉和知觉呢？我们岂不是每天看到，即使在土生土长的生物中，同一样东西对不同种群或甜或苦，有些事物对某些生物是致命的，而对另一些则无害，以致于那些以其怡人芬芳迷住某些生物的同样事物，却对另一些生物的身体散发出致命的气息吗？但这原因不在于事物本身，事物不可能同时是致命和有益、甘甜和苦涩的；而是由于每个个体被构造成接受外界印象的方式，才受到影响：他的状况并非由事物影响所导致，而是源自其自身感官的本性，及其与外界事物的联系。然而，这一切远离众神，与他们相距甚远。因为如果确实如智者所信，他们是\OriginalTerm{无形体的}{incorporeal}，也不依赖任何身体力量的卓越，那么气味对他们无效，蒸腾的烟气也不能通过感官触动他们，即便你点上最上等的一千磅香，整个天空被浓密蒸汽的黑暗所笼罩。因为那不具有身体力量和\OriginalTerm{有形体的}{corporeal}实体的，不可能被有形实体触动；而气味是有形体的，正如鼻子触及气味时所显示的：因此，按照理性，它不可能被神感知，因为神没有身体，也没有任何感觉和思想。\par

29. 酒和香一起使用；对此，我们同样要问，为何在香燃烧时要将酒浇在上面。如果做这件事不说明理由，也不道明原由，那么你们的这种行为，现在就不能归于可笑的错误，而更坦白地说，是归于疯狂、愚蠢、盲目。因为，正如已经屡次说过的，一切行为都应有其明显的原因，不应陷入任何隐晦的黑暗。因此，如果你们对自己所做的有信心，就请揭示、指明为何献上那汁液；也就是说，为何将酒倒在祭台上。难道是众神之躯感到极度口渴，需要用水份来调节其干燥吗？他们如同人一般，习惯于饮食结合吗？同样地，在享用糕饼、羹汤以及为尊崇他们而宰杀的牺牲等固体食物之后，他们是否也用频频举杯的葡萄酒灌醉自己，让自己欢愉，好让食物更容易软化、彻底消化呢？我恳求，拿酒给不死的神喝吧；拿出高脚杯、大碗、长柄勺和酒杯；当他们用公牛、豪宴和珍馐填饱肚子时——免得哪块仓促吞下的肉块在穿过胃脏时卡住，赶快，跑上去，拿纯酒给\Person{朱庇特}，那至善至高的神，免得他被噎住。他想要放屁却放不出；除非那障碍消散化解，否则他的呼吸极有可能停止中断，天将失去其统治者而荒废。\par

30. 但\Emph{我的反对者}说，你们毫无理由地侮辱我们，因为我们并非为了这些原因向天上的神浇奠酒，好像我们认为他们会渴、会喝，或因尝到酒的甘甜而喜悦。献酒给他们是为了尊崇他们；为使他们的卓越变得更崇高、更显赫，我们在他们的祭坛上灌奠，并以半熄的余烬发出芬芳，以表我们的敬意。如果你们相信神因接受酒而变得慈祥，或者认为只要往燃烧的炭火上洒落几滴酒就算是对他们的大尊敬，那么还能对神施加什么更大的侮辱呢？我们并非对没有理性或没有一般理解力的人说话：在你们里面也有智慧，也有理解力，你们心里凭自己的判断知道我们说的是真实的。但对于那些完全不愿按事物的本来面目思考、不愿与自己交谈的人，我们又能怎样呢？因为你们做的是你们看见在做的事，而不是你们确信应该做的事，因为在你们那里，没有理性的习俗比基于仔细考察真理而对事物本性的认知更占上风。神与酒有何相干？或者酒里有什么或多么大的能力，以致倾倒之后，神的卓越就增加，他的尊严就被认为得到尊崇？我说，神与酒有何相干？酒与\OriginalTerm{维纳斯}{Venus}的追求紧密相连，削弱一切美德的力量，与端庄和贞洁的得体相敌对——它常常刺激人的心灵，促使人疯狂、发狂，迫使神因狂言秽语而毁掉自己的权威。那么，把酒当作尊荣献上，而你若滥饮，就不知自己在做什么，不知自己在说什么，最终被人辱骂，成了醉汉、奢侈放荡之徒，这难道不是不虔敬，全然亵渎神圣吗？\par

31. 它值得把那些话也引出来，就是在献酒时惯常使用和祈求的话：\Quote{愿神明因我们带来的这酒而受敬拜。}\Person{特雷巴提乌斯}说，\Quote{我们带来的}这句话是为了这个目的加上的，并为此提出，好让所有存放在橱柜和储藏室里的酒，就是从中取出所浇之酒的那些酒，不至于开始成为神圣，而被剥夺人的使用。因此，加了这句话之后，只有带到那个地方的酒才是神圣的，其余的不被祝圣。那么，这是何等样的尊崇，在其中好像给神明强加了一个条件，叫祂不要索求超过所给的？或者，若不用言语限制，那位神会伸张太多欲望，抢夺其恳求者的储藏，这又是何等的贪婪？\Quote{愿神明因我们带来的这酒而受敬拜：}这是一个侮辱，并非尊崇。因为倘若那位神明希望更多，并对所带的感到不满足，那该怎么办！被迫有条件地接受尊崇，难道不能说这位神受了大委屈吗？因为，如果不加上这一限制，所有地窖里的酒都必然变作神圣，那么显然，那位被违反其意愿而加以限制的神是受了侮辱，而你们在献祭时自己违犯了神圣礼仪的义务，因为你们没有提供你们看到神希望赐给自己的那么多酒。\Quote{愿神明因我们带来的这酒而受敬拜：}这无非是说，\Quote{按我所选定的来受敬拜；接受我所规定的尊贵，接受我通过一个严格的约定所决定和限定的那么多崇敬吧？}哦，诸神的崇高，超绝的能力，你们本该以一切礼仪的守则来尊崇敬拜，然而敬拜者却给它加上条件，他以约规和契约来朝拜它，因害怕一句话，就限制它过量饮酒的欲望！\par

32. 但按你们所愿，就算在酒和香中有尊崇吧，就算诸神明的愤怒和不悦能被宰杀祭牲所平息吧：难道众神也被花环、花冠和花朵所感动，也被铜器的叮当声、钹的震动，被手鼓，也被和谐的笛声所感动吗？响板的喀哒声有什么效果，以致神明一听见它们，就认为已蒙崇敬，并在遗忘中放下了他们火热的怨恨之情？或者，正如小孩子们一听到哗啷棒的声音就被吓住，停止了无谓的啼哭，全能的神明也同样被笛子的啸声所抚慰吗？而他们听了钹的乐声就变得温和，怒气就消了吗？你们早晨伴着笛乐唱的那些呼唤有什么意思？天上的神是睡着了吗，以致他们该回到自己的岗位？你们用吉祥的致意将他们交托给那些沉睡，望他们安康，这又是什么意思？他们是从睡眠中被唤醒；而且为了能被睡意所征服，必须听安慰的摇篮曲吗？\Emph{我的反对者}说，今天是众神之母的净化日。难道众神会变脏；而为了去除污秽，那些洗涤他们的人需要水，甚至需要些灰来擦洗吗？明天是\OriginalTerm{朱庇特}{Jupiter}的宴席。想来朱庇特要进餐，必须用盛大的筵席来饱足，因禁食而长期充满对食物的渴望，在惯常的间隔后饥饿着。\OriginalTerm{埃斯库拉庇乌斯}{Aesculapius}的收葡萄节正在举行。众神于是耕葡萄园，聚拢采摘者，为自己的用途榨酒。\OriginalTerm{刻瑞斯}{Ceres}的\OriginalTerm{神宴}{lectisternium}将在下一个伊都斯日，因为众神有卧榻；为了能躺在更柔软的垫子上，当垫子被压下去之后，枕头要被抖松。今天是\OriginalTerm{特露斯}{Tellus}的生日；因为众神是出生的，并有一些节庆日，在那日子确定他们开始呼吸。\par

33. 然而，你们所庆祝的那些称为\OriginalTerm{花神节}{Floralia}和\OriginalTerm{大母神节}{Megalensia}的赛会，以及所有你们愿意视为神圣并当作宗教义务的其余活动，它们有何理由，有何原因，使得它们必须建立、创始并以神灵之名来命名？你们反对者会说：\Quote{神明借此受到尊崇}；如果神明对人犯下的任何过犯有任何记忆，他们便会放下、摆脱它，并向我们重新显出恩慈，恢复友谊。但这又是什么缘故，如果表演荒谬的事情，让无聊的家伙在众人眼前戏耍，就能使他们变得完全平静温和呢？难道\Person{普劳图斯}的\WorkTitle{安菲特律翁}被表演和朗诵，能使\Person{朱庇特}放下他的愤怒吗？或者，如果以舞蹈表现\Person{欧罗巴}、\Person{勒达}、\Person{伽倪墨得斯}或\Person{达那厄}，朱庇特就能抑制住他的冲动吗？\OriginalTerm{大母神}{Magna Mater}看到演员把\Person{阿提斯}的古老故事重新搬演出来，就会变得更加平静温和吗？\Person{维纳斯}看到模仿者也在芭蕾舞中扮演\Person{阿多尼斯}，就会忘记她的不快吗？\Person{索福克勒斯}名为\WorkTitle{特拉基斯少女}的悲剧，或\Person{欧里庇得斯}的\WorkTitle{赫拉克勒斯}上演时，\Person{阿尔喀得斯}的怒气就会消散吗？或者，\Person{弗洛拉}如果在她的节日里看到那些可耻的行为上演，看到妓院被弃置而转向剧场，她就会认为这是向她表示的敬意吗？那么，将最卑鄙下流的事奉献并祝圣给神明，而对此刚正不阿的心灵会嫌恶地转身离去，而你们的法律已判定这些表演者为人所不齿，并视为无耻之徒，这难道不正是贬低神明的尊严吗？其实，神明们以模仿演员为乐；那超越一切的无上卓越——它是任何官能都无法领会的——最乐意地张开耳朵聆听这些戏剧，而它们中大部分他们都知道将自己卷入其中，成为嘲笑的对象；正如事实那样，他们从剃光头的呆子的滑稽、响板声、掌声、可耻的行为言词以及巨大的红色阳具中得到愉悦。更有甚者，如果他们看到一些人让自己柔弱得如同女人一般，一些人无用地大声喊叫，一些人毫无缘由地东奔西跑，一些人在友情未断时用带血的拳击手套彼此打击致残，那些竞相发言时气也不换、鼓腮吹气、大声喊出空洞誓愿的人，神明们是否会在惊羡中向天空伸出双手，被这些奇观所感动，爆发出欢呼，并再次对人类施以恩惠呢？如果这些事情能使神明忘却他们的怨恨，如果喜剧、\OriginalTerm{阿特拉笑剧}{Atellanae fabulae}和哑剧能给他们带来最大的快乐，那么你们为何迟迟不直言，神明们自己也嬉戏、行淫、跳舞、作曲唱淫歌，并且摆臀扭胯呢？因为无论他们自己做出这些事，还是乐于从别人做这些事中得到快乐，这之间有何区别或又有什么关系呢？\par

34. 那么，这些虚妄的观点从何而来，或因何缘由而产生？由此可以清楚地看出，大体上是因为人们无法认识天主是什么，祂的本质、本性、实体、特性是什么；祂是否有形象，还是不受任何身体轮廓的限制；祂有无行动，是一直醒寤，还是有时沉入酣眠；祂是奔跑、端坐、行走，还是免于这些运动和静止。如我所说，既然无法知晓这一切，也无法凭借任何理性的力量来辨别，他们便陷入了这些空想的信念，以至按自己的模样塑造了神明，并将自己所拥有的行动、处境和欲望中的那种本性赋予了它们。但如果他们能认识到自己是毫无价值的受造物，并且自己与一只小蚂蚁之间并无太大区别，那么他们自然会停止以为自己与天上的神明有何共通之处，而将自身卑微渺小限制在恰当的范围之内。但现在，他们看到自己有着面孔、眼睛、头、面颊、耳朵、鼻子，以及我们四肢和肌肉的其余部分，于是便以为神明也是以同样的方式被造的，以为神性具体人形之躯；因为他们觉察到自己会欢喜快乐，又会因极不愉快的事而悲伤，便认为神灵在喜乐场合也欢喜，在不悦之时也沮丧。他们看到自己受赛会的影响，便以为天上神明的性情也能通过享受赛会而得以安抚；因为他们乐于用热水浴来恢复精神，便以为沐浴所带来的洁净也为上天诸神所喜悦。我们人类收获葡萄，他们便以为并相信神明也收获并运进他们的葡萄；我们有诞辰，他们便断言天上的大能也有诞辰。然而，如果他们能把疾病、病患和身体不适归于神明，他们就会毫不犹豫地说，神明们患有脾病、烂眼和疝气，因为他们自己既患有脾病，又常常烂眼，并被巨大的疝气所压垮。\par

35. 来吧：既然讨论已延长并引至这些要点，就让我们将双方的意见摆出来，通过简短的比较，断定是你们对上界神祇的观念更佳，还是我们的想法更可取、远为尊贵和公义，且能赋予神性以应有的尊严。首先，你们宣称，那些你们认为或相信存在、并在所有神庙中为其竖立肖像和雕塑的神祇，乃是从男女生殖的种子，在性交的必要条件下出生和产生的。而我们正好相反：如果他们是真神，拥有此名的权威、能力和尊严，我们则认为他们必须要么是\OriginalTerm{不受生}{unbegotten}——因为这样相信才是虔诚的——要么，若他们出生而有起始，那么至高天主以何种方式创造他们，或自祂允准他们进入其自身神性的永恒存有以来，经历了多少世代，这唯有至高天主知道。你们认为神祇有性别，有些是男神，有些是女神；我们全然否认天界的大能者以性别区分，因为这一区分只赋予了地上的受造物，宇宙的创造者愿意它们互相拥抱和生育，藉着肉欲不断繁衍后代。你们认为他们像人，是按凡人的面容塑造而成；我们认为那些肖像不着边际，因为形貌属于可朽的身体；就算他们真的有形貌，我们也以最大的诚挚和确信起誓，没有人能领会它。你们说他们各有手艺，如同工匠；我们听到你们说这些话时便发笑，因为我们认定并认为诸神不需要职业，而且确定而明显地，这些职业是为济贫而设的。\par

36. 你们说，他们中有的制造纷争，有的降下瘟疫，有的煽动情欲和疯狂，有的甚至主宰战争，以流血为乐；但我们，恰恰相反，断定这些事远离神祇的本性；或者说，若真有任何施与这些灾祸给可悲凡人的神明，我们也主张他们远非神类的本性，且不应被称为这个名号。你们认为神祇会发怒、烦扰，并耽溺臣服于其他心理情感；我们认为这类情绪与他们绝缘，因为这些只适合野蛮的存在以及那些作为凡人而死者。你们认为他们会欣喜、快乐，并与人和好，因为其冒犯的感受被兽血和牺牲的宰杀所平息；我们则认为，在\OriginalTerm{天神}{celestials}里没有嗜血之好，他们也不致冷酷到非待宰杀动物饱饫之后才放下怨恨。你们认为，藉着酒和香，便是向神祇献上尊崇，增其威严；我们却判定，若有人以为神明因烟熏而更可敬，或以为自己献上几滴酒便是以足够的敬畏与尊重向神明祈求，这真是可惊骇可怪异之事。你们相信，藉着铙钹的震响与笙箫之音，藉着赛马和戏剧表演，诸神既欢喜又受感动，且他们先前怀有的愤怒情绪会因而得到满足而平复；我们却认为这不合宜，不仅如此，我们判定这是不可信的：那些在完美高度上超越一切卓越成千上万倍者，竟会对智者发笑的事物感到欢喜快乐，而这些事物除对粗鄙庸俗教养的幼童外，看似并无任何魅力。\par

37. 既然事实如此，我们的意见和你们的意见之间有着如此巨大的差异，那么，究竟是我们不虔敬，还是你们虔敬呢？因为对虔敬与否的判断必须基于双方的观念。因为，一个人若为自己造一个像，当作神明来敬拜，或者宰杀无辜的牲畜，在祝圣的祭台上焚烧，决不可被视为专心于宗教。是观念构成宗教，以及对诸神的正确思考方式，好叫你们不以为他们渴望任何与他们超绝地位相悖的事物，这是显而易见的。因为我们看见所有献给他们的东西都在我们眼前被吞没，除了配得上神祇且最与其名号相称的观念以外，还能说有什么从我们这里达于他们呢？这些才是最可靠的礼物，这些才是真正的祭献；因为粥、香和肉只是喂养吞噬的火焰，与\OriginalTerm{亡者祭礼}{parentalia}非常相合。\par

38. 如果——\Emph{我的论敌}说——不死的神明不能发怒，其本性不会因任何情绪而激动或困扰，那么历史记载、编年史中所记的，受某些烦恼所动，降下瘟疫、不育、歉收和其他灾难给国家与民族的神明，并且又因祭献而得以平息和满足，遂放下燃烧的怒气，将气候与时令转变为更美好的状态，这些又是什么意思呢？大地的轰响、地震，据说是因为竞技会举办得草率疏忽，其性质与细节未被顾及，但后来重新举办，以勤勉的谨慎重演，神明的恐怖便得以平息，他们又恢复了对人的关怀与友谊，这又作何解释？此后多少次——遵照\OriginalTerm{预言者}{seers}的命令和\OriginalTerm{占卜师}{diviners}的回答——献上祭品，从居住在海外的民族那里召请某些神明，为他们建立神龛，将某些画像和雕像安放在更高的柱子上，于是对迫近危险的恐惧被转移，最麻烦的敌人被击败，共和国既因屡次的喜乐胜利，又因获得若干行省而得以扩展！当然，如果神明们鄙视祭献、竞技会和其他崇拜行为，不认为因\OriginalTerm{赎罪祭献}{expiratory offerings}而受尊荣，这一切就不会发生。那么，当这些被献上时，神明们一切的狂怒与愤恨便得以平息，那些充满恐惧的事变成顺遂，显然这一切并非在没有神明意愿的情况下完成，而责备我们献祭给他们，是徒劳的，且表明全然的无知。\par

39. 这样，我们在言谈中触及了问题的核心，就是那问题所系之处，就是讨论的真实且最内在的部分。我们理当搁置迷信的恐惧，放下偏袒，细察这些事到底是如此，还是远非如此，并应与这名号与权能的概念分别开来。因为我们并不否认，编年史家的著作中载有你们引用来反对我们的所有那些事；因为就连我们自己，也按着各人的能力和才分读过，也知道有过记载，说：有一次，在庆祝至尊\Person{朱庇特}的\OriginalTerm{罗马竞技}{ludi circenses}会上，一个主人将他用棍棒殴打过的一个贱奴拖过竞技场的中央，随后又照惯例，处以十字架之刑。此后，当竞技会结束，赛跑刚过不久，一场瘟疫便开始困扰国家；而且每一日都带来比前一日更严重的恶疫，人民成群地死亡，在一次梦中，\Person{朱庇特}对一个因出身卑微而不为人知的乡下人说，他应去见执政官，指出那个跳舞者使他不悦，如果对竞技会献上应有的尊重，并以勤勉的谨慎重新举办，于国家更有裨益。而此人竟完全忽略了这事，要么因为他以为那是一个虚空的梦，即便告诉那些人，他们也不会相信；要么因为他记着自己的微不足道，回避且害怕接近那些如此有权势的人；朱庇特便迁怒于迟延者，降罚使他儿子们死去。后来，当朱庇特以死亡威胁那人自己，除非他去宣告他对那舞者的不悦——这人既已染上瘟疫，连同他自己也正发着瘟疫的热度，在死亡的恐惧中，他遵照邻人的意思，被抬到元老院；当他说出所见异象之后，传染的热病便消失了。于是议决重办竞技会，一方面对表演施以极大的谨慎，人民便恢复了从前的健康。\par

40. 但我们也并不否认知道另一件事，即曾有一时，国家与共和国陷入困境，因一场可怕的瘟疫不断传染人民，将他们夺去；或因敌人强大，且因屡战屡胜，几乎威胁要夺取其自由——遵照\OriginalTerm{预言者}{seers}的命令和建议，从居住在海外的民族中召请了某些神明，并为他们建造了宏伟的庙宇以示尊崇；结果瘟疫的暴力消退了，极频繁的胜利被取得，敌人的势力被击破，帝国的领土得以扩展，无数的行省落入你们的管辖之下。但此事也未能逃出我们的认知：我们见过记载说，当\Place{卡比托利欧山}被雷电击中，山上的许多其他东西，连同立在巨柱上的\Person{朱庇特}像，也被从他的位置抛掷出去。此后，\OriginalTerm{占卜者}{soothsayers}给出回答，说由烈火与杀戮、法律的毁灭和正义的倾覆，尤其由本民族的敌人和一帮不虔的阴谋家，预兆出残酷且极悲惨的灾祸；然而这些事不能避免，不，那些可咒的毒谋不能揭露，除非朱庇特再次牢固地树立在一根更高的柱子上，转向东方，面向升起太阳的光芒。他们的话是可信的，因为，当柱子被竖起，雕像转向太阳之后，秘密便被揭发，那些罪恶昭彰，遭受了惩罚。\par

41. 上述种种，确实有一种神奇的表象——更确切地说，人们相信它们具有这种表象——如果它们传入人们耳中时，正如它们被陈述的那样；我们并不否认，其中确实有某些东西，正如常言所说，摆在最前面，可以震骇人耳，并且因其类似真理而欺骗人。但如果你仔细审视所发生的事，那些人物及其快乐，你就会发现，其中毫无配得上神祇的东西，而且，正如已多次说过的，毫无值得归于此种族的光辉与威严的东西。因为，首先，谁会相信，一位神竟会以漫无目的地奔马为乐，并认为藉此种竞赛被呼召是最愉快的呢？更确切地说，谁会同意，那位\Person{朱庇特}——你们称之为至高神，万物创造者——竟从天庭下来观看骟马竞速，绕场奔跑七圈；而且，虽然他亲自决定了它们不应同等地敏捷，他竟仍然喜欢看它们相互超越，又被超越，有些慌忙中头朝前跌，连人带车翻倒在地，另一些被拖行而致跛，腿骨断裂；他竟认为掺杂琐碎与残酷的蠢行是至高的快乐，而这种快乐，任何人，即使喜好快乐，未受追求庄重与尊严的训练，也会认为是幼稚的，并视作荒谬而摒弃？我再说，谁会相信——反复地重复这个词——那位神圣者，因为一个奴隶被带过竞技场，即将按其应得的受刑和惩罚，竟被激怒，充满了怒火，并准备施行报复，是神呢？因为，如果那奴隶有罪，应受那种惩罚，\Person{朱庇特}为何要感到义愤？当时并没有任何不公正的事发生，不，反而是一个有罪的人正受到应得的惩罚，这是正当的。但如果他无罪，根本不配受罚，那么\Person{朱庇特}自己就是那舞者败坏竞赛的缘由，因为他本可以帮他，却没有给他任何帮助——不，他既要允许自己所不赞同的事，又要向他人索取对他所许可之事的惩罚。那么，他为何还要抱怨，并宣称在那舞者的事上他受了亏待，因为那舞者被带过竞技场中央，去忍受十字架，后背被棍棒与鞭子撕裂？\par

42. 由此能流出何种污染或可憎之物，使竞技场变得不那么洁净，或玷污\Person{朱庇特}，既然他在短短片刻、在数秒之内，目睹了世界各地成千上万的人因各种死法而丧生，遭受各种折磨？\Emph{我的对手}说，他是在竞赛开始庆祝之前被带过去的。如果是出于亵渎的灵和蔑视宗教，我们有理由宽赦\Person{朱庇特}因被轻慢而感到愤慨，以及因他的竞赛没有得到更热切的关注而愤慨。但若是由于错误或意外，那隐秘的过失未被察觉和知晓，\Person{朱庇特}难道不该饶恕人的过犯，宽赦因无知而来的盲目，这岂不才是正当而合宜的吗？然而，这事必须受惩罚。如此之后，还有谁会相信，那位因一个幼稚表演被忽视，便以毁灭一整座城邦来报复和惩罚的神，是神呢？还有谁会相信，他具有任何庄重与尊严，任何恒定的持守，因为他为能迅速重新享受快乐，就将人所呼吸的空气变为有害的毒气，并命瘟疫消灭凡人？如果掌管竞赛的官员过于疏忽，未能查明当天谁被带过了竞技场，由此招致罪责，那么那些不幸的民众做了什么，竟要亲身承受别人过犯的惩罚，且被迫因流行疫病而急促地结束生命呢？更有甚者，那些妇女——她们的软弱使她们不能参与公共事务，那些成年少女，那些小男孩，最后还有那些尚依赖乳母喂养的幼儿——他们做了什么，竟要以同等的严厉受打击，并且在尚未尝到生命的喜乐之前，就感受到死亡的苦涩呢？\par

43. 如果\Person{朱庇特}想要以更隆重的方式重新庆祝他的竞技会；如果他真心想要恢复人们的健康，并使自己引发的灾祸不再蔓延与加剧，那么，他亲自去见执政官本人、去见某位公共司祭、去见\OriginalTerm{大司祭}{pontifex maximus}，或他自己的\OriginalTerm{弗拉门·迪亚利斯}{flamen Dialis}，并在神视中向其透露因那位舞者而出现的竞技会瑕疵与时艰之因，岂不更好？他有什么理由偏要挑选一个惯居乡野、无名无姓、不谙城事、或许连舞者为何物都不知道的人，来宣告他的意愿并求得所愿的满足呢？而如果他确实——假如他是占卜者，他就必定知道——知道这家伙会拒绝服从，那么，改变那人的心意、迫使他甘愿顺从，岂不是比采取更残忍的手段、像强盗一样乱发脾气、毫无理由地迁怒于他人，更合乎神明的本性、更适合一位神？因为，如果这位年老的乡民确因更强的动机而犹豫不决，迟迟不肯遵命，那么他那些不幸的儿女又犯了什么罪，竟使\Person{朱庇特}的忿怒与义愤转向他们，使他们因他人的过错而被夺去生命？又有谁能相信，如此不公、不虔、甚至不顾人情常理的神竟是一位神？在人间，代人受过、将一人的罪行报复于他人身上，都会被视为大罪。但是，我听说，他让那乡民本人也染上了残酷的瘟疫。那么，即便此事看来理当如此，让那位因违命拖延而引发如此盛怒的父亲先感受到惩罚的恐惧，岂不比伤害他的儿女、残害无辜以使他悲伤更妥当、更公平？请问，这等凶残、这等暴虐究竟意谓何物？竟至于在子女夭亡之后，又以自身的危险来吓唬这位父亲！然而，设若他早早就这样做，也就是说，一开始就这样做，那么，不仅无辜的弟兄们不致丧命，神明的恼怒心意也早已昭然。当然，有人会说，当这位乡民履行了宣告神视的责任后，疾病立即离开了，他也立即恢复了健康。但如果他驱除的是他自己吹入那人身上的灾祸，并以虚假的借口自夸，这又有什么值得惊奇的呢？可是，你若细察当时情状，便会发现这解救之中残忍多于仁慈；因为\Person{朱庇特}并未保全那痛失子女、悲不欲生者的生命之乐，反倒让他去品尝孑然一身的孤独与丧亲的剧痛。\par

44. 同样，我们还可以逐一检视其他的叙述，并指出——便如我将要讲述的这个故事一样——在这些叙述及对它们的解说中，关于神明所谈论、所宣扬的，与神明应有的样子相去甚远。我将只补充一两个例子，免得冗赘招厌。你说，在从海外的民族中请来了某些神明，在为他们兴建了庙宇，在祭坛上堆满了祭品之后，瘟疫肆虐的民众便强健康复，攀升的健康气息使瘟疫远遁。我恳请一问：你说的这些神明是哪些？你说是\Person{埃斯库拉庇乌斯}，健康之神，来自\Place{厄庇道鲁斯}，如今定居\Place{台伯河}中央的岛上。若我们对你们的说法严加细究，大可以凭你们自己的权威证明，那位由妇人胎中受孕而生、年复一年抵达生命之期——据你们书中所载，一道霹雳便将他逐出生命与光明的人，绝非神明。但且将这个问题搁置；就依你们的愿望，让\Person{科洛尼斯}之子跻身不朽者之列吧，享上天永恒的福乐。然而，从\Place{厄庇道鲁斯}带来的，除了一条巨蛇，还有什么？如果我们相信\WorkTitle{编年史}，且认可其中确切不移的真理，那么，正如所记载的，再没别的了。那么，我们该怎么说？你们所颂扬的\Person{埃斯库拉庇乌斯}，那位至善可敬、赐予健康、消灾防病、祛除疾厄的神明，竟囿于蛇的形体之内，像蠕虫那般匍匐于地，它们生自污泥；他以颏与胸磨擦土地，蜿蜒盘曲，拖曳前行；而为了能够前进，他以身体前部的力量牵引后半截肢体。\par

45. 而且，正如我们读到他也使用食物，借以维持肉体存在，他有巨大的食道，以便张大嘴吞下寻求的食物；他有肚腹来接收它，并有消化所吃下和吞食的肉的地方，好使血液供给身体，力量得到恢复；他也有排泄通道，借此排出污物，使身体摆脱讨厌的重担。每当他更换地点，准备从一地到另一地时，他并非像神一样隐秘地飞过天上的星辰，瞬间站在需要他出现的地方；相反，就像地上迟钝的动物，他寻求可乘坐的交通工具；他避开海浪；为了安然无恙，他与人们一同上船；那位公共安全之神，竟将自己托付给脆弱的木板和接合的木板。我们不认为你们能证明并显示那条蛇就是\OriginalTerm{埃斯库拉庇乌斯}{Aesculapius}，除非你们愿意提出这样一个托辞，即你们应说神把自己变成了一条蛇，以便他能就自己是谁而欺骗人，或者看看人是什么。但如果你们这样说，你们自己陈述的不一致将显示这种辩护是多么软弱无力。因为如果神躲避被人看见，他就不应该选择以蛇的形象被人看见，因为无论以任何形象，他都不可能是别人，而始终是自己。但另一方面，如果他曾有意让人看见——他本不应拒绝让人们用眼注视他——为什么他不以自己所知的神性权能中的本来面目显现呢？因为这更可取，更好，更配得他的威严，而不是变成一头野兽，变成一个可怕动物的模样，并留下无法决定的争议余地，即他是否是一位真神，或者是不属于崇高神性的某种不同的、相去甚远的东西。\par

46. 但是，\Emph{我的对手}说，如果他不是神，为什么他离开船，爬行到\Place{台伯河}中的岛上之后，立即变成看不见的，不再像从前那样被看见了呢？我们真能知道路上是否有什么东西，在其遮掩下他隐藏了自己，或者地上有什么开口吗？你们自己宣布，说吧，那是什么，或者它应属于哪一类存在，如果你们对某些人物的侍奉本身是确定的。既然情况如此，讨论涉及你们的神，也涉及你们的宗教，那么教导和显示那是什么，是你们的职责，而不是希望听到我们的意见，等待我们的断定。因为对于我们，除了说所发生和被看见的事，就是所有叙述中流传下来的，以及用眼睛观察到的事，还能说什么呢？然而，我们毫不怀疑地说，那是一条体型极为强壮、体长巨大的\OriginalTerm{雌蛇}{colubra}，或者，如果这个名字可鄙，我们说它是一条蛇，我们称它为蛇，或者任何习惯提供给我们的、或语言发展创造的其他名称。因为如果它像蛇一样爬行，不用脚支撑和行走，而是以肚子和胸部贴地；如果它由血肉之质构成，伸展着滑溜的长度；如果它有头和尾，背部覆盖鳞片，散布着各种颜色的斑点；如果它有嘴，布满毒牙，随时准备咬人，那么除了说它出自尘世，尽管体积巨大无朋，尽管在体长和力量之大上超过了\Person{雷古鲁斯}率军攻击所杀的那条蛇之外，我们还能说什么呢？但如果我们想别的，我们就是在颠覆和推翻真理。因此，解释那是什么，它的起源、名称和本性，是你们的职责。因为它怎么可能是神，既然它具有我们提到的那些东西，而神若要为神，并享有这崇高的称号，就不该有这些东西？在它爬行到\Place{台伯河}中的岛上之后，立即无处可见，由此显示它是一位神灵。那么，我们就能知道那里是否有什么东西，在其遮掩下它隐藏了自己，或者地上有什么开口，或者由巨大石块不规则堆砌形成的洞穴和拱顶，它匆匆进入其中，逃避旁观者的目光吗？因为，如果它跃过了河呢？如果它游过了河呢？如果它藏在茂密的森林里了呢？由此推断那条蛇是神，因为它极快地避开了旁观者的眼睛，这是脆弱的推理，因为，根据同样的推理，相反地可以证明它不是神。\par

47. 但是，\Emph{我的对手}说，如果那条蛇不是在场的神灵，为什么在它到来之后，瘟疫的猛烈被制伏，罗马人民的健康得以恢复？我们反过来也要提出这个问题：如果按照命运之书和占卜者的回应，\OriginalTerm{埃斯库拉庇乌斯}{Aesculapius}神奉命被邀请到城中，好使他能使城市免受瘟疫和传染病的感染，安然无恙，并且如你所说，他化为蛇形，没有轻蔑地拒绝邀请而前来——那么，为什么罗马国家如此经常地遭遇这样的灾难，如此经常地一次复一次被撕裂、困扰，因着无数次公民的毁灭而减少成千上万的人呢？因为既然据说神被召唤的目的，是为了彻底驱除所有引发瘟疫的原因，那么国家本应安全，本应一直保持远离瘟疫之风，不受伤害。可是我们却看到，正如前面所说，它一次又一次地因着这些疾病而经历悲惨的季节，它人民的阳刚活力因不小的损失而遭到重创和削弱。那么，\OriginalTerm{埃斯库拉庇乌斯}{Aesculapius}在哪里？那位由可敬神谕所许诺的拯救者在哪里？为什么在为他建造了庙宇，竖立了圣祠之后，他竟允许一个值得他恩惠的国家继续遭瘟疫侵袭，既然他被召唤的目的，是要治愈肆虐的疾病，不让任何令人生畏的此类事物后来悄悄降临到他们身上？\par

48. 但也许有人会说，这样一位\OriginalTerm{神}{god}的关怀之所以被之后和后来的时代所否定，是因为如今人们生活的方式是不敬神且可憎的；相反，它曾给我们的祖先带来帮助，因为他们无可指摘，毫无罪过。现在，如果说在古代所有人都毫无例外是善的，或者后来的时代只产生了恶人而没有别的，那么这话或许还有几分道理可言。但事实上，在各个大的民族、各个国家，甚至在所有城市中，人们的心性、愿望、\OriginalTerm{品格}{mannets}都是混杂的，善人和恶人可以在古代和现代同时存在，那么如果说后来的凡人因为自身的恶而没有得到\OriginalTerm{诸神}{deities}的帮助，就相当愚蠢了。因为如果由于后世恶人的缘故，现代的好人没有得到保护，那么由于古代的恶人，古代的好人也同样不应该得到\OriginalTerm{诸神}{deities}的眷顾。但如果由于古代好人的缘故，古代的恶人也得到了保护，那么后来的时代虽然犯了错，但由于后来的好人的缘故，也应该得到保护。因此，要么是那条\OriginalTerm{蛇}{snake}根本毫无用处，却因为疾病暴力已经减弱衰退时被带到城里而获得了拯救者的名声，要么就得说\OriginalTerm{命运女神}{fates}的赞歌所给的指示完全不是真实的，因为它们所提供的疗方被发现只对一代人有用，而非对所有人都连续有效。\par

49. 但是，\Emph{我的对手}说，同样地，\OriginalTerm{伟大母亲}{Great Mother}按照占卜者的命令从\Place{弗里吉亚的佩西努斯}被请来，也给人民带来了安全与极大的喜乐。因为，一方面，一个长期强敌被赶出了他在\Place{意大利}已占据的位置；另一方面，通过光荣卓越的胜利，城市恢复了古代的荣耀，帝国的疆域大大扩展，无数种族、国家、民族的自由人被剥夺了权利，被套上奴隶的轭，而国内外完成的许多其他事业以不可抗拒的力量确立了罗马民族的声望与尊严。如果历史讲述的是真相，没有在记载事件时掺入虚假，那么据说从弗里吉亚由\Person{阿塔罗斯王}送来的，实际上不过是一块不大的\OriginalTerm{石头}{stone}，可握于一人手中而无压迫感——颜色暗黑——表面不平，有小小的棱角凸起，今天我们都看到它被放在那个形象中代替面孔，粗糙未经雕琢，使那雕像的面容毫无生气。\par

50. 那么，我们该说什么呢？那位著名的迦太基人\Person{汉尼拔}，一个强而有力的敌人，曾让\Place{罗马}的命运在怀疑与不确定中战栗，让它的伟大动摇——他是被一块石头赶出\Place{意大利}的吗？他是被一块石头征服的吗？他是因为一块石头而变得恐惧、怯懦、一反常态的吗？至于罗马再次崛起至权力与王权至尊的高度，难道一点都不是由于智慧、人的力量吗？在恢复昔日卓越的过程中，那么多伟大的将领们，凭借他们的军事技能或对事务的精通，难道没有提供任何帮助吗？是石头给了某些人力量，给了另一些人软弱吗？是它将一些人从成功中推下，又将另一些人看似无望覆灭的命运提升起来了的吗？又有谁会相信，一块从地里取出的、没有感觉的、颜色乌黑、质地暗沉的石头，竟会是\OriginalTerm{诸神之母}{mother of the gods}？或者，还有谁会听从这样的说法——因为这是唯一的选择——任何神祇的大能住在燧石碎片中，在其块体之内，隐藏在其脉络之中？如果那\OriginalTerm{佩西努斯石头}{Pessinuntine stone}中没有神祇，胜利又是如何取得的？我们也许可以说，是由于士兵们的热忱与勇气，由于训练、时机、智慧、理性；我们也可以归因于\OriginalTerm{命运}{fate}，以及轮转不定的\OriginalTerm{运气}{fortune}。但如果局势的改善、成功与胜利的重新获得是藉由石头的帮助，那么，当共和国因如此多、如此庞大的军队遭受屠杀而倾颓，濒临彻底毁灭时，那位\OriginalTerm{弗里吉亚母亲}{Phrygian mother}又在哪里呢？她为何不冲到步步逼近的强敌面前？为何不在这些可怕打击落下之前将其粉碎并击退呢？正是这些打击流尽了所有的血，甚至生命元气都衰竭殆尽。\Emph{我的对手}说，她那时还没被请来，也没被要求施恩。就算如此；但是，一位仁慈的帮助者从不需要被要求，总是主动伸出援手。你说，她无法驱逐敌人并将其赶走，因为当时还与意大利隔着重重大海和陆地。但对于一位神祇，如果真是神，没有任何事物是遥远的，对她而言，大地只是一个点，万物皆因她的旨意而建立。\par

51. 可是，就算神祇确实临在于那块石头中，正如你们要求我们相信的那样：有没有哪个凡人，即使再怎么轻信、再怎么愿意听信任何虚构故事，会认为她在那时是一位女神，或者现在应该那样称呼和命名她，因为她时而想要这些，时而又要求那些，遗弃和轻视自己的崇拜者，离开较卑微的行省，与更强大、更富有的人民结盟，真正喜爱战争，希望置身于战斗、屠杀、死亡和流血之中？如果神祇的特性——如果他们真的是真神，是那些理应按照这个词的意义和神性大能来称呼的——是不做任何邪恶、任何不义之事，不偏不倚地同等恩待所有人，那么会有任何人相信，她拥有神圣的起源，或展现了与神祇相称的仁慈吗？她介入人类的纷争，毁灭一部分人的力量，赐予和向另一些人施恩，剥夺一些人的自由，将另一些人推至权力的顶峰——为了使一个城邦脱颖而出，她生来就是人类的祸害，征服了无罪的世界？\par

\SourceRecord{Translated by Hamilton Bryce and Hugh Campbell.；Ante-Nicene Fathers；Vol. 6.；Edited by Alexander Roberts, James Donaldson, and A. Cleveland Coxe.；Buffalo, NY: Christian Literature Publishing Co.,；1886.；<http://www.newadvent.org/fathers/06317.htm>.}
